\providecommand{\XSp}{\mathcal{X}_{\mathrm{Sp}}}
\providecommand{\dR}{\mathrm{dR}}
\providecommand{\Jinf}{J^{\infty}}
\providecommand{\Tinf}{T^{\infty}}
\providecommand{\PDE}{\operatorname{PDE}}
\providecommand{\Coalg}{\operatorname{Coalg}}
\providecommand{\Map}{\operatorname{Map}}
\providecommand{\Eq}{\operatorname{Eq}}
\providecommand{\fib}{\operatorname{fib}}
\providecommand{\hofib}{\operatorname{hofib}}
\providecommand{\cofib}{\operatorname{cofib}}
\providecommand{\Tot}{\operatorname{Tot}}
\providecommand{\Sec}{\operatorname{Sec}}
\providecommand{\Lag}{\operatorname{Lag}}
\providecommand{\EL}{\operatorname{EL}}
\providecommand{\Helm}{\operatorname{Helm}}
\providecommand{\Crit}{\operatorname{Crit}}
\providecommand{\Prim}{\operatorname{Prim}}
\providecommand{\Hess}{\operatorname{Hess}}
\providecommand{\Jac}{\operatorname{Jac}}
\providecommand{\Rvar}{\mathcal R_{\mathrm{var}}}
\providecommand{\Cvar}{C_{\mathrm{var}}}
\providecommand{\Bvar}{B_{\mathrm{var}}}
\providecommand{\Omegavar}{\Omega_{\mathrm{var}}}
\providecommand{\Omegainf}{\Omega_{\mathrm{inf}}}
\providecommand{\Oadd}{\mathcal O^{\mathrm{add}}}
\providecommand{\DRvar}{\mathrm{DR}_{\mathrm{var}}}
\providecommand{\Vfull}{\mathbb V^{\infty}}
\providecommand{\Vfirst}{\mathbb V^{(1)}}
\providecommand{\Vsecond}{\mathbb V^{(2)}}
\providecommand{\BiFil}{\operatorname{BiFil}^{\mathrm{comp}}_{\ge0}}
\providecommand{\ClVar}{\operatorname{ClVar}}
\providecommand{\EulMod}{\operatorname{Eul}}
\providecommand{\HessDesc}{\operatorname{HessDesc}}
\providecommand{\RealH}{\operatorname{Real}_{H}}
\providecommand{\TransH}{\operatorname{Trans}_{H}}
\providecommand{\Lep}{\operatorname{Lep}}
\providecommand{\Fld}{\operatorname{Fld}}
\providecommand{\Act}{\operatorname{Act}}
\providecommand{\Dens}{\operatorname{Dens}}
\providecommand{\Real}{\operatorname{Real}}
\providecommand{\HelmLift}{\operatorname{HelmLift}}
\providecommand{\LagPrim}{\operatorname{LagPrim}}
\providecommand{\one}{\mathbf 1}
\providecommand{\id}{\operatorname{id}}
\providecommand{\pr}{\operatorname{pr}}

\chapter{Brave-New Variational Geometry and Derived Euler--Lagrange Theory}
\label{chap:brave-new-variational}

\section{Purpose, continuity, and scope}
\label{sec:var-purpose}

Chapters 9 and 10 constructed the square-zero-generated relative de Rham geometry, its effective infinitesimal quotients, parameterized stable coefficient categories, normalized infinitesimal cochains, complete totalization filtrations, coherent infinitesimal differentials, coherently closed cochains, exact-cochain classifiers, primitive spaces, and first-order cotangent linearization in the represented geometric sector. Chapter 11 used the same effective infinitesimal relation to construct the brave-new jet comonad
\[
    \Jinf_{X/S}:\mathcal C_X\longrightarrow\mathcal C_X,
    \qquad
    \mathcal C_X:=(\XSp)_{/X},
\]
the formal-disk monad \(\Tinf_{X/S}\), the adjunction
\[
    \Tinf_{X/S}\dashv\Jinf_{X/S},
\]
and the presentation-independent equivalences
\[
    \PDE^{\mathrm{fi}}_{X/S}
    \simeq
    \Coalg_{\Jinf_{X/S}}(\mathcal C_X)
    \simeq
    (\XSp)_{/Q_{X/S}}.
\]
Thus an intrinsic formally integrable PDE is an object over the effective infinitesimal quotient \(Q_{X/S}\); its jet coaction, formal-disk action, Cartan transport, and all higher descent coherences are outputs of effective descent rather than additional structure.

The purpose of the present chapter is to construct the variational geometry of these intrinsic Cartan objects and, in the represented smooth sector, to integrate differential Lagrangian densities to an actual brave-new action morphism. The classical organizing precedent is Vinogradov's secondary calculus and the variational bicomplex; see \cite{PuglieseSparanoVitagliano2023,Vitolo2000}. Kock's infinitesimal-cochain presentation of differential forms and the scheme-theoretic extension of Breen--Messing provide a complementary first-infinitesimal precedent. The comparison with these classical theories is retained, but it is not the organizing definition of the brave-new construction.

The primary variational object of this chapter is a complete bifiltered object
\[
    \mathfrak{Var}_{X/S}(A;M)
    \in
    \operatorname{Fil}^{\mathrm{comp}}_{\ge0}
    \bigl(
      \operatorname{Fil}^{\mathrm{comp}}_{\ge0}(\mathcal E_Q)
    \bigr),
    \qquad Q:=Q_{X/S},
\]
constructed from the Cartan--vertical double infinitesimal relation. The associated coherent double cochain object and its signed variational bicomplex are derived presentations of this bifiltered object. In particular, the equations
\[
    d_V^2\simeq0,
    \qquad
    d_H^2\simeq0,
    \qquad
    d_Hd_V+d_Vd_H\simeq0
\]
record only the first visible cells of a product-indexed coherent object carrying specified nullhomotopies, interchange homotopies, higher compatible nullhomotopies, and all higher pasting relations supplied by the complete-filtration/coherent-cochain correspondence.

Horizontal Cartan descent is likewise treated as a complete descent resolution rather than as an additional differential imposed after the fact. After horizontal totalization, the secondary form spectrum is canonically
\[
    \Sec^p(A;M)
    \simeq
    \Omegainf^p(\bar A/Q;M),
\]
so the horizontally descended variational differential is the intrinsic coherent infinitesimal differential of the descent object \(\bar A/Q\). The nonzero horizontal bidegrees before totalization retain the filtration and coherent resolution data of Cartan descent; they are not identified with powers or quotients of a classical contact ideal.

The chapter keeps full finite-infinitesimal variation separate from its first split square-zero shadow. The full one-variation coefficient is obtained from the intrinsic first relation simplex, while the one-stage coefficient is the right Kan extension of the actual spectrum-valued first-difference functor on the \(\infty\)-category of structured split square-zero variation diagrams. The raw right Kan extension is proved already to satisfy spectral Nisnevich descent. In the represented sector, universal affine split variations are constructed from the Chapter 2 cotangent complex and are initial in the right-Kan indexing categories, giving a recognition theorem for the intrinsically defined coefficient. No tangent object is inserted into the general definition.

The moduli-theoretic part of the chapter is organized by universal maps before their individual homotopy fibres are taken. Lagrangians, coherently closed full first variations, and first Euler sources form a universal diagram of objects over \(Q\). Helmholtz lifts, Lagrangian realizations, and primitive spaces are the corresponding derived fibres. Likewise, the Euler equation is promoted from a construction attached separately to every Lagrangian to a universal derived critical family over the moduli object of Lagrangians; each \(\Crit(L)\) is obtained from that family by pullback.

The first and mixed second variations are further recognized as the first two terms of a canonical finite-cubical variation system. For every \(k\ge1\), the chapter constructs the \(\infty\)-category of \(k\)-fold split square-zero variations, its total-fibre coefficient \(K_M^{(k)}\), the corresponding right-Kan coefficient \(\mathbb V_M^{(k)}(A)\), its spectral Nisnevich descent and arbitrary base change, its canonical \(\Sigma_k\)-symmetry, and the alternating \(k\)-fold variation of a Lagrangian. This construction uses only finite tensor products of the already constructed split square-zero algebras, finite limits, right Kan extension, and the same Cartesian-centre descent argument used in degree one. No polynomiality, Goodwillie-excision, multilinearity, or Hessian axiom is imposed.

This cubical formulation exposes a genuinely higher feature of the theory. The intrinsic second variation is a function of a complete bi-infinitesimal square and can detect the mixed \(N_1\otimes N_2\) lift even when its two side variations are fixed. Accordingly, the unrestricted first categorical prerequisite for a Hessian is formulated as descent of the intrinsic mixed second variation along the functor from complete bi-variations to their pair of side variations. The derived space of such descents is retained as a homotopy type; the chapter neither truncates it nor assumes that it is inhabited. The on-shell mixed second variation is obtained by pullback to the Euler critical locus. In the represented smooth scalar sector, however, the mixed boundary algebra is constructed as a structured square-zero extension by \(N_1\otimes_BN_2\), smoothness proves the resulting mixed-fill torsor, and criticality canonically kills its Euler-controlled translation. The scalar stabilized cotangent coefficient admits a constructed affine-point linear refinement whose unit underlies \(\Oadd\); that refinement is dualizable, so the descended critical second variation determines a symmetric bilinear Hessian and its adjoint Jacobi operator. No corresponding bilinear statement is inferred for the unrestricted off-shell coefficient.

The action sector remains inside the previously constructed Thom-stable differential coefficient theory. For represented smooth separated finite-presentation
\[
    x:X\longrightarrow S
\]
and \(D\in\mathcal R_S^0\), the chapter constructs a big differential mapping-spectrum coefficient on all affine tests, its big dualizing-density coefficient over \(Q\), and the big action-value coefficient over \(S\). Differential exceptional adjunction gives a canonical equivalence
\[
    \mathcal I^\partial_{x,D}:
    \Pi_qM^\partial_{x,D}
    \xrightarrow{\sim}
    \mathbb A^\partial_{x,D}.
\]
A differential Lagrangian therefore defines a morphism of higher field moduli
\[
    \mathbf S_L:
    \mathbf{Fld}_S(A)
    \longrightarrow
    \Omega_S^\infty\bigl(\mathbb A^\partial_{x,D}[n]\bigr),
\]
whose ordinary action on global fields is only its global-points shadow. Variation commutes with this integration as a natural equivalence on the corresponding variation \(\infty\)-categories; the pointwise first- and second-variation formulas are evaluations of that natural statement.

All higher-categorical strengthening in this chapter is obtained from structures already constructed in Chapters 2 and 9--11, together with the finite affine spectral Nisnevich module-descent theorem proved in the represented Jacobi subsection: effective descent, finite limits in the native \(\infty\)-topos, stable dependent adjunctions, right Beck--Chevalley, complete totalization, coherent filtrations, structured split square-zero extensions, right Kan extensions, and finite Nisnevich descent. The horizontal tower is additionally organized into a derived moduli object of relative Lepage realizations carrying a tautological presymplectic transgression; no preferred integration-by-parts operator or Poincar\'e--Cartan representative is selected. No new axiom of representability, orientation, multiplicativity of the contact filtration, all-degree HKR comparison, admissibility, tangent-generation, multilinearity, or symmetry is introduced. Representability is used only in explicitly marked recognition, Jacobi, local-model, and differential-action sectors.

\begin{convention}[Ambient Cartan geometry]
\label{conv:var-ambient}
Throughout the chapter fix
\[
    x:X\longrightarrow S
\]
in \(\XSp\), put
\[
    Q:=Q_{X/S},
    \qquad
    \epsilon:=\epsilon_{X/S}:X\twoheadrightarrow Q,
\]
and write
\[
    \mathcal C_X:=(\XSp)_{/X},
    \qquad
    \mathcal D_Q:=(\XSp)_{/Q}.
\]
A \emph{Cartan object} means a \(\Jinf_{X/S}\)-coalgebra
\[
    (A,\rho_A),
    \qquad
    \rho_A:A\longrightarrow\Jinf_{X/S}A.
\]
Its corresponding effective-descent object over \(Q\) is denoted
\[
    a:\bar A\longrightarrow Q.
\]
Thus Chapter 11 supplies a canonical equivalence
\[
    A\simeq\epsilon^*\bar A=X\times_Q\bar A.
\]
\end{convention}

\begin{convention}[Arbitrary pullback of Cartan descent data]
\label{conv:var-arbitrary-cartan-pullback}
Let
\[
    c:Q'\longrightarrow Q
\]
be an arbitrary morphism.  Put
\[
    X':=X\times_QQ',
    \qquad
    \epsilon':=\epsilon\times_QQ':X'\twoheadrightarrow Q'.
\]
Because effective epimorphisms are stable under pullback, \(\epsilon'\) is an effective epimorphism.  For a Cartan descent object
\[
    a:\bar A\longrightarrow Q
\]
put
\[
    \bar A':=\bar A\times_QQ',
    \qquad
    A':=(\epsilon')^*\bar A'
       =X'\times_{Q'}\bar A'.
\]
We call \(A'\) the \emph{pulled-back Cartan descent object} over \(Q'\).  Effective descent for \(\epsilon'\) identifies objects over \(Q'\) with coalgebras for the descent comonad
\[
    J_{\epsilon'}:=(\epsilon')^*\Pi_{\epsilon'}
\]
on \((\XSp)_{/X'}\), and equips \(A'\) with its canonical descent coaction.  Its horizontal relation is the \v{C}ech nerve of \(\epsilon'\), with canonical degreewise equivalences
\[
    (X')^{\times_{Q'}(q+1)}
    \simeq
    Q'\times_QX^{\times_Q(q+1)}.
\]
For an arbitrary \(c:Q'\to Q\), no identification of \(J_{\epsilon'}\) with a brave-new jet comonad \(\Jinf_{X'/S'}\) is asserted.  If \(c\) is induced by a represented geometric base change \(S'\to S\), then the Chapter 11 base-change theorem identifies
\[
    Q'\simeq Q_{X'/S'},
    \qquad
    J_{\epsilon'}\simeq\Jinf_{X'/S'},
\]
and \(A'\) is the usual pulled-back Cartan object in the Chapter 11 sense.
\end{convention}

\begin{convention}[Parameterized stable coefficients]
\label{conv:var-stable}
For every \(T\in\XSp\), retain the Chapter 10 notation
\[
    \mathcal E_T:=\operatorname{Sp}((\XSp)_{/T})
\]
for the parameterized stable category, with tensor unit \(\one_T\), fibrewise tensor product, internal Hom, and stable dependent adjoint triple
\[
    \Sigma_f\dashv f^*\dashv\Pi_f
\]
for every \(f:T\to U\). Pullback is exact, strong symmetric monoidal, preserves all small limits and colimits, and satisfies the coherent Beck--Chevalley equivalences constructed in Chapter 10.
\end{convention}

\begin{convention}[Coefficient notation]
\label{conv:var-coefficients}
A general stable coefficient is denoted
\[
    M\in\mathcal E_Q.
\]
Its pullback to an object \(T\to Q\) is denoted \(M_T\). Unless a general coefficient is displayed, the scalar coefficient is
\[
    \Oadd_Q\in\operatorname{CAlg}(\mathcal E_Q)
\]
from Chapter 10. For a relative relation over \(Q\), coefficient-valued normalized infinitesimal cochains are obtained by applying the Chapter 10 normalization construction to the cosimplicial stable object
\[
    [p]\longmapsto\Pi_{v_p}v_p^*M.
\]
No tensor-product identification of these coefficient-valued cochains with cotangent objects is implicit in this notation.
\end{convention}

\begin{convention}[Universes]
\label{conv:var-universes}
Use the universe enlargement of Chapters 9--10. In particular, affine-point categories, all finite local and field-valued cubical variation categories, side-data categories, their comma and functor categories, and the right Kan extensions and mapping spaces below are formed after choosing the same permitted small skeleta. No additional size convention is introduced.
\end{convention}

\subsection{Chapter-level output}
\label{subsec:var-output-preview}

For a Cartan object \((A,\rho_A)\) with descent object \(a:\bar A\to Q\) and a stable coefficient \(M\in\mathcal E_Q\), the chapter constructs the Cartan--vertical double relation
\[
    \Rvar^{p,q}(A)
    =
    X^{\times_Q(q+1)}\times_QR_{\bar A/Q,p},
\]
and its bicosimplicial stable coefficient object
\[
    \Cvar^{p,q}(A;M)
    =
    \Pi_{r_{p,q}}r_{p,q}^*M.
\]
Horizontal effective descent identifies
\[
    \Tot_H\Cvar^{\bullet,\bullet}(A;M)
    \simeq
    C_{\dR}^{\bullet}(\bar A/Q;M),
\]
and therefore
\[
    \DRvar(A;M)
    \simeq
    \mathrm{DR}_{\mathrm{eff}}(\bar A/Q;M).
\]

Before passing to associated graded, the chapter constructs the complete Cartan--vertical bifiltration
\[
    \boxed{
    \mathfrak{Var}_{X/S}(A;M)
    :=
    F^{\bullet,\bullet}\DRvar(A;M)
    }
\]
as an object of
\[
    \operatorname{Fil}^{\mathrm{comp}}_{\ge0}
    \bigl(
      \operatorname{Fil}^{\mathrm{comp}}_{\ge0}(\mathcal E_Q)
    \bigr).
\]
Its iterated associated graded is
\[
    \operatorname{gr}_V^p\operatorname{gr}_H^q\mathfrak{Var}_{X/S}(A;M)
    \simeq
    \Omegavar^{p,q}(A;M)[-p-q],
\]
independently of the order of the two associated-graded constructions. The complete-filtration/coherent-cochain equivalence then produces a coherent double cochain object whose unary structure maps are \(\partial_V,\partial_H\); the signed maps
\[
    d_V:=\partial_V,
    \qquad
    d_H|_{(p,q)}:=(-1)^p\partial_H
\]
form the signed variational bicomplex shadow with the full hierarchy of coherent relations retained.

The assignment \((A,M)\mapsto\mathfrak{Var}_{X/S}(A;M)\) is promoted to an \(\infty\)-functor, contravariant in the descended Cartan geometry and covariant in coefficients. Secondary forms, closed forms, horizontal realization objects, primitive objects, the coherent variational differential, and all their base-change maps are functorial constructions from this primary bifiltered object.

The first variational layer consists of the intrinsic full one-variation coefficient \(\Vfull_M(A)\), the right-Kan one-stage coefficient \(\Vfirst_M(A)\), and the natural restriction
\[
    \lambda^1_{\mathrm{loc}}:
    \Vfull_M(A)\longrightarrow\Vfirst_M(A).
\]
When \(a:\bar A\to Q\) is representable, the cotangent complex constructs universal affine split variations \(u_c\) which are initial in the right-Kan indexing categories; this gives the represented recognition
\[
    \Vfirst_M(A)
    \simeq
    L^{\mathrm{st}}_{\bar A/Q}(M)
\]
without changing the intrinsic definition.

At the moduli level the chapter constructs a universal diagram
\[
    \boxed{
    \mathbf{Lag}^n_{A/Q}(M)
    \longrightarrow
    \mathbf{ClVar}^n_{A/Q}(M)
    \longrightarrow
    \mathbf{Eul}^n_{A/Q}(M),
    }
\]
whose composite is the universal first Euler map. Helmholtz-lift, Lagrangian-realization, and primitive objects are its derived fibres. Pulling the universal first Euler section to the universal Lagrangian family gives a universal derived Euler critical object; every \(\Crit(L)\) is its pullback along the point classifying \(L\).

The complete horizontal filtration is further used to construct, for every Lagrangian and every stage \(q\ge1\), a derived relative Lepage-realization object
\[
    \mathbf{Lep}^{\,q}(L)\longrightarrow Q.
\]
Its points compare the canonical coherently closed full Euler lift with another coherent lift of the same first Euler source through a specified horizontal realization, and the defining homotopy gives the intrinsic relative integration-by-parts factorization.  Over this moduli object there is a tautological boundary realization whose coherent vertical differential factors through the complete horizontal transgression coefficient, producing the universal presymplectic transgression
\[
    \omega_L^{(q)}:
    \one_{\mathbf{Lep}^{\,q}(L)}
    \longrightarrow
    (\pi_L^q)^*\TransH^{2,q}(A;M)[n].
\]
No preferred Lepage representative or integration-by-parts operator is selected.

For every \(k\ge1\), the chapter also constructs a finite-cubical variation coefficient
\[
    \boxed{
    \mathbb V_M^{(k)}(A)
    }
\]
from the total fibre of the \(k\)-cube of stable coefficient sections on the split square-zero cube
\[
    T[N_1,\ldots,N_k].
\]
The construction satisfies centre-pullback Nisnevich descent, arbitrary base change, and the canonical \(\Sigma_k\)-symmetry. A Lagrangian determines its alternating higher variation
\[
    \Delta_L^{(k)}:
    \one_{\bar A}\longrightarrow\mathbb V_M^{(k)}(A)[n].
\]
The cases \(k=1\) and \(k=2\) recover the first square-zero variation and the symmetric mixed second variation. A Hessian is not built into \(\mathbb V_M^{(2)}\): instead the chapter constructs the derived side-descent problem which is a necessary first step toward such an operator. The existing affine example proves that unrestricted off-shell descent is not automatic.  In the represented smooth scalar sector, the chapter constructs an affine-point linear coefficient category \(\mathcal Q_{\bar A}\), a finite-locally-free refinement \(L^{\mathrm{lin}}_{\bar A/Q}\) whose underlying parameterized spectrum is \(L^{\mathrm{st}}_{\bar A/Q}\), and its dual \(T^{\mathrm{lin}}_{\bar A/Q}\). The structured mixed-boundary extension, the smooth mixed-fill torsor, and critical Euler contraction construct the on-shell side descent and the underlying Jacobi linearization
\[
    \Jac_L:
    j_L^*T^{\mathrm{st}}_{\bar A/Q}
    \longrightarrow
    j_L^*L^{\mathrm{st}}_{\bar A/Q}[n].
\]
The critical Hessian is constructed as the symmetric bilinear morphism
\[
    \Hess_L^{\mathrm{crit}}:
    j_L^*T^{\mathrm{lin}}_{\bar A/Q}
    \otimes
    j_L^*T^{\mathrm{lin}}_{\bar A/Q}
    \longrightarrow
    \mathcal O^{\mathrm{lin}}_{C_L}[n]
\]
in \(\mathcal Q_{C_L}\), and is identified with the side descent of the intrinsic on-shell mixed second variation.

In the final represented smooth differential sector, the big dualizing-density coefficient \(M^\partial_{x,D}\), the action-value coefficient \(\mathbb A^\partial_{x,D}\), and exceptional integration produce a universal action morphism on the field-moduli object. For every finite \(k\), the induced \(k\)-fold difference of the action is naturally the differential exceptional integral of the corresponding coefficient-valued cubical difference. The first Euler section represents the \(k=1\) local variation, and critical field families are therefore first-order stationary. The relative Lepage and presymplectic-transgression constructions apply in particular to the differential density coefficient, but no spacetime boundary trace is manufactured without an actual coefficient morphism into the relevant exceptional pullback, and coherent transgression of the stage system requires the compatible family of such morphisms specified in Chapter 10. No converse stationarity theorem, Noether theorem, BRST datum, or gauge-symmetry structure is imposed.

\section{Cartan descent objects}
\label{sec:var-cartan-descent}

\begin{proposition}[Cartan descent equivalence]
\label{prop:var-cartan-descent}
For every Cartan object \((A,\rho_A)\), there is a canonically associated object
\[
    a:\bar A\longrightarrow Q
\]
and a canonical equivalence
\[
    A\simeq X\times_Q\bar A.
\]
The assignment extends to an equivalence of \(\infty\)-categories
\[
    \Coalg_{\Jinf_{X/S}}(\mathcal C_X)
    \xrightarrow{\sim}
    (\XSp)_{/Q}.
\]
For Cartan objects \(A,B\), the induced map on mapping spaces is an equivalence; hence a Cartan morphism is exactly the pullback along \(\epsilon\) of its descended morphism over \(Q\), with all higher homotopies retained.
\end{proposition}

\begin{proof}
This is the effective-descent/Eilenberg--Moore equivalence constructed in Chapter 11 for the effective epimorphism \(\epsilon:X\twoheadrightarrow Q\). Under that equivalence the forgetful functor to \(\mathcal C_X\) is \(\epsilon^*\). Therefore the underlying object corresponding to \(\bar A\to Q\) is \(\epsilon^*\bar A=X\times_Q\bar A\). Full faithfulness identifies the entire mapping spaces, not merely homotopy classes.
\end{proof}

\begin{proposition}[Cofree descent object]
\label{prop:var-cofree-descent}
Let \(Y\to X\) be an \(X\)-bundle. The cofree Cartan object
\[
    K(Y):=(\Jinf Y,\Delta_Y)
\]
corresponds under Proposition~\ref{prop:var-cartan-descent} to
\[
    \Pi_\epsilon Y\longrightarrow Q.
\]
Consequently
\[
    \Jinf Y\simeq\epsilon^*\Pi_\epsilon Y
\]
as an object over \(X\), compatibly with the cofree coaction.
\end{proposition}

\begin{proof}
Chapter 11 identifies the jet comonad with \(\epsilon^*\Pi_\epsilon\) as a comonad. Under the comonadic comparison for \(\epsilon^*\dashv\Pi_\epsilon\), the cofree coalgebra on \(Y\) is the pullback of \(\Pi_\epsilon Y\), with coaction induced by the adjunction unit. This is exactly \(K(Y)\).
\end{proof}

\begin{remark}[The variational base is constructed, not chosen]
\label{rem:var-descent-base}
The object \(\bar A\to Q\) is not an auxiliary quotient or a chosen presentation of the Cartan object. It is the object corresponding to \((A,\rho_A)\) under the equivalence of Proposition~\ref{prop:var-cartan-descent}. Every construction below is functorial in \(\bar A\to Q\), and therefore presentation-independent from the outset.
\end{remark}

\section{Vertical infinitesimal geometry}
\label{sec:var-vertical-geometry}

\begin{definition}[Vertical effective quotient and relation nerve]
\label{def:var-vertical-relation}
Apply the relative de Rham construction of Chapters 9--10 to
\[
    a:\bar A\longrightarrow Q.
\]
Write \(\bar A_{\dR/Q}\) for the relative de Rham object and factor the relative unit by
\[
    \bar A
    \xrightarrow{\epsilon_A}
    G_A:=Q_{\bar A/Q}
    \xrightarrow{\jmath_A}
    \bar A_{\dR/Q},
\]
where \(\epsilon_A\) is an effective epimorphism and \(\jmath_A\) is a monomorphism. For \(p\ge0\), put
\[
    V_{A,p}:=R_{\bar A/Q,p}
    =
    \underbrace{
    \bar A\times_{\bar A_{\dR/Q}}\cdots\times_{\bar A_{\dR/Q}}\bar A
    }_{p+1\text{ factors}}.
\]
By effective-image invariance,
\[
    V_{A,p}
    \simeq
    \underbrace{
    \bar A\times_{G_A}\cdots\times_{G_A}\bar A
    }_{p+1\text{ factors}}.
\]
Write
\[
    v_p:V_{A,p}\longrightarrow Q
\]
for the structural morphism.
\end{definition}

\begin{proposition}[Functoriality and base change]
\label{prop:var-vertical-functoriality}
A Cartan morphism \(A\to B\) induces a morphism of simplicial relation objects
\[
    V_{A,\bullet}\longrightarrow V_{B,\bullet}
\]
over \(Q\), compatibly with the effective quotients and with all higher functoriality coherences. For every \(c:Q'\to Q\), if \(\bar A':=\bar A\times_QQ'\), then
\[
    Q_{\bar A'/Q'}\simeq G_A\times_QQ',
\]
and
\[
    R_{\bar A'/Q',p}
    \simeq
    V_{A,p}\times_QQ'
\]
compatibly with every simplicial operator.
\end{proposition}

\begin{proof}
Under Proposition~\ref{prop:var-cartan-descent}, a Cartan morphism is a morphism of the descended objects over \(Q\). Apply Chapter 10 functoriality of the relative de Rham unit, functorial image factorization, and the effective relation nerve. The base-change assertions are the arbitrary base-change equivalences for relative de Rham objects, effective images, and relation nerves from Chapters 9--10.
\end{proof}

\section{The Cartan--vertical double infinitesimal relation}
\label{sec:var-double-relation}

\begin{definition}[Horizontal Cartan relation]
\label{def:var-horizontal-relation}
For \(q\ge0\), put
\[
    H_q
    :=
    \underbrace{X\times_Q\cdots\times_QX}_{q+1\text{ factors}}.
\]
Then \(H_\bullet\) is the \v{C}ech nerve of the effective epimorphism \(\epsilon:X\twoheadrightarrow Q\).
\end{definition}

\begin{definition}[Double relation]
\label{def:var-double-relation}
For \(p,q\ge0\), define
\[
    \boxed{
    \Rvar^{p,q}(A):=H_q\times_QV_{A,p}.
    }
\]
Write
\[
    e_{p,q}:\Rvar^{p,q}(A)\longrightarrow V_{A,p}
\]
for the projection and
\[
    r_{p,q}:=v_pe_{p,q}:\Rvar^{p,q}(A)\longrightarrow Q
\]
for its structural morphism.
\end{definition}

\begin{theorem}[Bisimplicial variational relation]
\label{thm:var-bisimplicial-relation}
The assignment
\[
    ([p],[q])\longmapsto\Rvar^{p,q}(A)
\]
extends canonically to a functor
\[
    \Rvar^{\bullet,\bullet}(A):
    \Delta^{\mathrm{op}}\times\Delta^{\mathrm{op}}
    \longrightarrow
    (\XSp)_{/Q}.
\]
The first simplicial direction is the vertical relation nerve of \(\bar A\to G_A\), the second is the horizontal \v{C}ech nerve of \(X\to Q\), and the two structures commute with their full higher coherences.
\end{theorem}

\begin{proof}
Every vertical simplicial operator acts only on the ordered \(\bar A\)-factors of \(V_{A,p}\), while every horizontal operator acts only on the ordered \(X\)-factors of \(H_q\). Both actions are induced by functoriality of finite limits over the common base \(Q\). Their composites are therefore the two canonically equivalent pastings of the same finite-limit diagram. The simplicial identities and all higher coherences are inherited from the two relation nerves and the coherent functoriality of pullback.
\end{proof}

\begin{proposition}[Horizontal boundary]
\label{prop:var-horizontal-boundary}
There is a canonical equivalence of simplicial objects over \(\bar A\)
\[
    \Rvar^{0,\bullet}(A)
    \simeq
    \check C_\bullet(A\to\bar A),
\]
where
\[
    A=X\times_Q\bar A\longrightarrow\bar A
\]
is the pullback of \(\epsilon\). Hence this morphism is an effective epimorphism and
\[
    \left|\Rvar^{0,\bullet}(A)\right|\simeq\bar A.
\]
\end{proposition}

\begin{proof}
Since \(V_{A,0}=\bar A\),
\[
    \Rvar^{0,q}(A)=X^{\times_Q(q+1)}\times_Q\bar A,
\]
which is the \(q\)-th term of the \v{C}ech nerve of \(X\times_Q\bar A\to\bar A\). Effective epimorphisms are stable under pullback in an \(\infty\)-topos.
\end{proof}

\begin{proposition}[Vertical boundary]
\label{prop:var-vertical-boundary}
For every \(p\ge0\), there are canonical equivalences
\[
    \Rvar^{p,0}(A)
    \simeq
    R_{A/X,p}
\]
compatible with the complete vertical simplicial structure.
\end{proposition}

\begin{proof}
The morphism \(A\to X\) is the base change of \(\bar A\to Q\) along \(\epsilon:X\to Q\). Apply arbitrary base change of relative de Rham relation nerves:
\[
    R_{A/X,p}
    \simeq
    X\times_QR_{\bar A/Q,p}
    =
    \Rvar^{p,0}(A).
\]
\end{proof}

\begin{remark}[No linearization has occurred]
\label{rem:var-double-nonlinear}
The double relation retains the complete square-zero-generated finite-infinitesimal relation of \(\bar A/Q\) together with the complete effective Cartan descent relation. It is defined for arbitrary objects of the native big spectral Nisnevich \(\infty\)-topos and does not require a tangent bundle, a vector bundle of contact directions, or representability.
\end{remark}

\section{Double stable cochains and horizontal effective descent}
\label{sec:var-double-cochains}

\begin{definition}[Double variational cochains]
\label{def:var-double-cochains}
For \(M\in\mathcal E_Q\), define
\[
    \boxed{
    \Cvar^{p,q}(A;M)
    :=
    \Pi_{r_{p,q}}r_{p,q}^*M
    \in\mathcal E_Q.
    }
\]
For \(M=\Oadd_Q\), write \(\Cvar^{p,q}(A)\).
\end{definition}

\begin{proposition}[Bicosimplicial coefficient object]
\label{prop:var-bicosimplicial-cochains}
The double relation induces a canonical bicosimplicial object
\[
    \Cvar^{\bullet,\bullet}(A;M):
    \Delta\times\Delta\longrightarrow\mathcal E_Q.
\]
It is natural in Cartan morphisms and coefficient morphisms and is compatible with arbitrary base change.
\end{proposition}

\begin{proof}
A simplicial map
\[
    u:\Rvar^{p',q'}(A)\longrightarrow\Rvar^{p,q}(A)
\]
over \(Q\) gives, from the adjunction unit for \(u^*\dashv\Pi_u\), a restriction morphism
\[
    \Pi_{r_{p,q}}r_{p,q}^*M
    \longrightarrow
    \Pi_{r_{p',q'}}r_{p',q'}^*M.
\]
Functoriality of the adjunction unit gives the cosimplicial identities and higher coherences. Naturality is inherited from Theorem~\ref{thm:var-bisimplicial-relation}. Base change follows from right Beck--Chevalley.
\end{proof}

\begin{proposition}[Scalar algebra structure]
\label{prop:var-double-algebra}
For \(M=\Oadd_Q\), every \(\Cvar^{p,q}(A)\) is canonically a commutative algebra object of \(\mathcal E_Q\), all bicosimplicial maps are maps of commutative algebras, and every partial or complete totalization below carries its limit commutative-algebra structure.
\end{proposition}

\begin{proof}
Pullback is strong symmetric monoidal and \(\Pi_f\) is lax symmetric monoidal. Applying this to the commutative algebra \(\Oadd_Q\) gives commutative algebras in every bidegree. Limits in \(\operatorname{CAlg}(\mathcal E_Q)\) are created in \(\mathcal E_Q\).
\end{proof}

\begin{definition}[Coefficient-valued intrinsic de Rham cochains]
\label{def:var-coefficient-dr}
For the relative morphism \(a:\bar A\to Q\), define
\[
    C_{\dR}^p(\bar A/Q;M)
    :=
    \Pi_{v_p}v_p^*M.
\]
The relation nerve makes \(C_{\dR}^{\bullet}(\bar A/Q;M)\) a cosimplicial object of \(\mathcal E_Q\). Define
\[
    \mathrm{DR}_{\mathrm{eff}}(\bar A/Q;M)
    :=
    \Tot_p C_{\dR}^{\bullet}(\bar A/Q;M).
\]
For \(M=\Oadd_Q\), this is the effective intrinsic de Rham algebra of Chapter 10.
\end{definition}

\begin{theorem}[Horizontal effective descent]
\label{thm:var-horizontal-descent-cochains}
There is a canonical equivalence of vertical cosimplicial objects in \(\mathcal E_Q\)
\[
    \boxed{
    \Tot_H\Cvar^{\bullet,\bullet}(A;M)
    \simeq
    C_{\dR}^{\bullet}(\bar A/Q;M).
    }
\]
In vertical degree \(p\), this is the equivalence
\[
    \Tot_q\Cvar^{p,q}(A;M)
    \simeq
    \Pi_{v_p}v_p^*M.
\]
The equivalence is natural in \(A\), in \(M\), and under arbitrary base change.
\end{theorem}

\begin{proof}
Fix \(p\). The simplicial object
\[
    e_{p,\bullet}:\Rvar^{p,\bullet}(A)\longrightarrow V_{A,p}
\]
is the \v{C}ech nerve of
\[
    e_{p,0}:X\times_QV_{A,p}\longrightarrow V_{A,p},
\]
which is the pullback of the effective epimorphism \(\epsilon:X\to Q\). Therefore \(e_{p,0}\) is an effective epimorphism.

Apply Chapter 10 effective descent for parameterized spectra to the object \(v_p^*M\in\mathcal E_{V_{A,p}}\). It gives
\[
    v_p^*M
    \simeq
    \Tot_q
    \Pi_{e_{p,q}}e_{p,q}^*v_p^*M
\]
in \(\mathcal E_{V_{A,p}}\). Applying the exact right adjoint \(\Pi_{v_p}\), which preserves all limits, gives
\[
\begin{aligned}
    \Pi_{v_p}v_p^*M
    &\simeq
    \Tot_q
    \Pi_{v_p}\Pi_{e_{p,q}}e_{p,q}^*v_p^*M\\
    &\simeq
    \Tot_q
    \Pi_{r_{p,q}}r_{p,q}^*M\\
    &=
    \Tot_q\Cvar^{p,q}(A;M).
\end{aligned}
\]
The construction uses only the horizontal \v{C}ech nerve and is natural with respect to every vertical simplicial operator. Hence the degreewise equivalences assemble into an equivalence of vertical cosimplicial objects. The naturality and base-change statements are those of effective descent and right Beck--Chevalley.
\end{proof}

\begin{definition}[Total variational de Rham object]
\label{def:var-total-dr}
Define
\[
    \DRvar(A;M)
    :=
    \Tot_{\Delta\times\Delta}
    \Cvar^{\bullet,\bullet}(A;M).
\]
\end{definition}

\begin{theorem}[Total variational object is intrinsic de Rham descent]
\label{thm:var-total-identification}
There is a canonical equivalence
\[
    \boxed{
    \DRvar(A;M)
    \simeq
    \mathrm{DR}_{\mathrm{eff}}(\bar A/Q;M).
    }
\]
For scalar coefficients this is an equivalence of commutative algebra objects. It is natural in Cartan morphisms, coefficients, and arbitrary base change.
\end{theorem}

\begin{proof}
Limits commute with limits in \(\mathcal E_Q\). Therefore
\[
\begin{aligned}
    \DRvar(A;M)
    &\simeq
    \Tot_p\Tot_q\Cvar^{p,q}(A;M)\\
    &\simeq
    \Tot_p C_{\dR}^p(\bar A/Q;M)
\end{aligned}
\]
by Theorem~\ref{thm:var-horizontal-descent-cochains}. For scalar coefficients, all degreewise equivalences are induced from maps of commutative algebras and totalization is computed in commutative algebras.
\end{proof}

\section{Vertical normalization and the complete bifiltration}
\label{sec:var-normalization}
\label{sec:var-contact-filtration}

\begin{definition}[Vertically normalized horizontal cosimplicial object]
\label{def:var-vertical-normalization}
For every \(p\ge0\), apply Chapter 10 normalization in the vertical cosimplicial direction while retaining the horizontal direction. Define
\[
    \boxed{
    \Bvar^{p,\bullet}(A;M)
    :=
    N_V^p\Cvar^{\bullet,\bullet}(A;M)
    \in
    \operatorname{Fun}(\Delta,\mathcal E_Q).
    }
\]
Explicitly, in horizontal degree \(q\),
\[
    \Bvar^{p,q}(A;M)
    =
    \fib\left(
      \Cvar^{p,q}(A;M)
      \longrightarrow
      M^p_{\sigma,V}(\Cvar^{\bullet,q}(A;M))
    \right),
\]
where \(M^p_{\sigma,V}\) is the vertical codegeneracy matching object.
\end{definition}

\begin{definition}[Doubly normalized variational cochains]
\label{def:var-double-normalization}
Apply horizontal normalization to the horizontal cosimplicial object \(\Bvar^{p,\bullet}\). Define
\[
    \boxed{
    \Omegavar^{p,q}(A;M)
    :=
    N_H^q\Bvar^{p,\bullet}(A;M).
    }
\]
\end{definition}

\begin{proposition}[Commutation of the two normalizations]
\label{prop:var-normalization-commutes}
There is a canonical equivalence
\[
    N_H^qN_V^p\Cvar
    \simeq
    N_V^pN_H^q\Cvar.
\]
The equivalence is natural in the bicosimplicial object and compatible with arbitrary base change.
\end{proposition}

\begin{proof}
Each normalization is the total fibre of a finite cube generated by the relevant codegeneracies. The vertical and horizontal codegeneracy maps commute because \(\Cvar\) is bicosimplicial. Therefore the two iterated normalizations are the two orders of taking the limit of the same finite product cube. Limits commute with limits in the stable category \(\mathcal E_Q\).
\end{proof}

\subsection{The two complete totalization towers}

For a cosimplicial object \(C^\bullet\) in a stable \(\infty\)-category, write
\[
    D^m(C):=\Tot_{\Delta_{\le m}}C,
    \qquad
    D^\infty(C):=\Tot_\Delta C,
\]
and use the convention \(D^{-1}(C)=0\). Chapter 10 associates to this tower the complete decreasing filtration
\[
    F^m(C):=\fib\bigl(D^\infty(C)\to D^{m-1}(C)\bigr).
\]
We apply this construction inside stable functor categories before taking the remaining totalization.

\begin{convention}[Nonpositive extension of the complete filtrations]
\label{conv:var-nonpositive-filtration}
Whenever the complete-filtration/coherent-cochain equivalence of Chapter 10 is applied below, every decreasing filtration is extended to a \(\mathbb Z\)-indexed filtration by making it constant in nonpositive degrees. Write
\[
    p_+:=\max(p,0),
    \qquad
    q_+:=\max(q,0).
\]
Thus
\[
    \boxed{
    \mathfrak F_V^p\Cvar
    :=
    \mathfrak F_V^{p_+}\Cvar
    }
\]
and, after the horizontal filtration is formed,
\[
    \boxed{
    F^{p,q}\DRvar
    :=
    F^{p_+,q_+}\DRvar.
    }
\]
This single convention removes any overlap ambiguity on the nonpositive quadrant and is exactly the constant nonpositive extension required by the Chapter 10 equivalence
\[
    \operatorname{Fil}^{\mathrm{comp}}_{\ge0}
    \simeq
    \operatorname{CCh}_{\ge0}.
\]
\end{convention}

\begin{definition}[Vertical complete filtration before horizontal totalization]
\label{def:var-vertical-filtered-cosimplicial}
Regard
\[
    \Cvar^{\bullet,\bullet}(A;M)
\]
as a vertical cosimplicial object in the stable category
\[
    \operatorname{Fun}(\Delta,\mathcal E_Q)
\]
of horizontal cosimplicial objects. For \(p\ge0\), define
\[
    \mathfrak F_V^p\Cvar
    :=
    \fib\left(
      \Tot_V\Cvar
      \longrightarrow
      \Tot^V_{\Delta_{\le p-1}}\Cvar
    \right)
    \in\operatorname{Fun}(\Delta,\mathcal E_Q).
\]
Thus \(\mathfrak F_V^\bullet\Cvar\) is the complete vertical totalization filtration formed \emph{before} horizontal totalization.
\end{definition}

\begin{definition}[Complete Cartan--vertical bifiltration]
\label{def:var-complete-bifiltration}
For \(p,q\ge0\), define
\[
    \boxed{
    F^{p,q}\DRvar(A;M)
    :=
    \fib\left(
       \Tot_H\mathfrak F_V^p\Cvar
       \longrightarrow
       \Tot^H_{\Delta_{\le q-1}}\mathfrak F_V^p\Cvar
    \right).
    }
\]
The transition maps in the two variables make
\[
    (p,q)\longmapsto F^{p,q}\DRvar(A;M)
\]
a decreasing \(\mathbb N^2\)-filtered object of \(\mathcal E_Q\). We regard this complete bifiltered object, rather than its associated graded, as the primary variational invariant and write
\[
    \boxed{
    \mathfrak{Var}_{X/S}(A;M)
    :=
    F^{\bullet,\bullet}\DRvar(A;M).
    }
\]
\end{definition}

\begin{proposition}[Recovery of the one-directional filtrations]
\label{prop:var-bifiltration-edges}
There are canonical equivalences
\[
    F^{p,0}\DRvar
    \simeq
    \fib\left(
      \DRvar\longrightarrow
      \Tot^V_{\Delta_{\le p-1}}\Tot_H\Cvar
    \right)
    =:F_V^p\DRvar
\]
and
\[
    F^{0,q}\DRvar
    \simeq
    \fib\left(
      \DRvar\longrightarrow
      \Tot^H_{\Delta_{\le q-1}}\Tot_V\Cvar
    \right)
    =:F_H^q\DRvar.
\]
In particular the double filtration restricts on its two axes to the vertical and horizontal complete-totalization filtrations used below.
\end{proposition}

\begin{proof}
For \(q=0\), horizontal totalization over \(\Delta_{\le-1}\) is zero, so
\[
    F^{p,0}\DRvar\simeq\Tot_H\mathfrak F_V^p\Cvar.
\]
Horizontal totalization is a limit and therefore preserves the defining fibre of \(\mathfrak F_V^p\Cvar\), which gives the first formula. For \(p=0\), the vertical filtration stage is \(\mathfrak F_V^0\Cvar\simeq\Tot_V\Cvar\), giving the second formula.
\end{proof}

\begin{proposition}[Symmetry of the bifiltration construction]
\label{prop:var-bifiltration-symmetry}
If one first forms the horizontal complete-totalization filtration in the stable category of vertical cosimplicial objects and then applies the vertical filtration to each horizontal stage, the resulting \(\mathbb N^2\)-filtered object is canonically equivalent to Definition~\ref{def:var-complete-bifiltration}. The equivalence satisfies the identity, composition, and higher pasting coherences inherited from commutation of limits.
\end{proposition}

\begin{proof}
For fixed \(p,q\), both orders of construction identify canonically with the total fibre of the same rectangular diagram
\[
\begin{tikzcd}[column sep=huge]
\Tot_V\Tot_H\Cvar
  \arrow[r]
  \arrow[d]
&
\Tot^V_{\Delta_{\le p-1}}\Tot_H\Cvar
  \arrow[d]
\\
\Tot_V\Tot^H_{\Delta_{\le q-1}}\Cvar
  \arrow[r]
&
\Tot^V_{\Delta_{\le p-1}}
\Tot^H_{\Delta_{\le q-1}}\Cvar .
\end{tikzcd}
\]
Indeed, taking the vertical fibre first and then the horizontal fibre gives Definition~\ref{def:var-complete-bifiltration}, while taking the horizontal fibre first and then the vertical fibre gives the reverse-order construction. Every corner is a limit of the same bicosimplicial diagram and every arrow is induced by restriction of the corresponding indexing category. The canonical Fubini equivalences for limits therefore identify the two total fibres. Naturality of these Fubini equivalences with respect to identity, composition, and restriction of the two partial-totalization indices gives the asserted higher pasting coherences.
\end{proof}

\begin{theorem}[Bifiltration associated graded]
\label{thm:var-bifiltration-graded}
The bifiltration of Definition~\ref{def:var-complete-bifiltration} is complete in each variable and its iterated associated graded satisfies
\[
    \boxed{
    \operatorname{gr}_V^p\operatorname{gr}_H^q
    \DRvar(A;M)
    \simeq
    \Omegavar^{p,q}(A;M)[-p-q].
    }
\]
The reverse order of associated graded is canonically equivalent:
\[
    \operatorname{gr}_H^q\operatorname{gr}_V^p\DRvar
    \simeq
    \operatorname{gr}_V^p\operatorname{gr}_H^q\DRvar.
\]
All comparison maps are natural and commute with arbitrary base change.
\end{theorem}

\begin{proof}
Apply Chapter 10's stable Dold--Kan layer theorem to the vertical cosimplicial object in \(\operatorname{Fun}(\Delta,\mathcal E_Q)\). The \(p\)-th vertical layer is
\[
    \Bvar^{p,\bullet}(A;M)[-p].
\]
The horizontal complete filtration is exact in that stable functor category. Its \(q\)-th associated-graded layer is therefore
\[
    N_H^q\Bvar^{p,\bullet}(A;M)[-p-q]
    =
    \Omegavar^{p,q}(A;M)[-p-q].
\]
Reversing the two directions gives
\[
    N_V^pN_H^q\Cvar[-p-q],
\]
which is canonically equivalent by Proposition~\ref{prop:var-normalization-commutes}. Completeness in each direction is Chapter 10 completeness applied in the corresponding stable functor category.
\end{proof}


\begin{lemma}[Iterated complete filtrations and coherent double cochains]
\label{lem:var-iterated-filtered-cochains}
Let \(\mathcal D\) be a stable \(\infty\)-category admitting all small limits. In particular, every stable functor and complete-filtration category below admits the sequential limits used in the Chapter 10 complete-filtration construction. There is a canonical equivalence
\[
    \operatorname{Fil}^{\mathrm{comp}}_{\ge0}
    \bigl(
      \operatorname{Fil}^{\mathrm{comp}}_{\ge0}(\mathcal D)
    \bigr)
    \simeq
    \operatorname{CCh}_{\ge0}
    \bigl(
      \operatorname{CCh}_{\ge0}(\mathcal D)
    \bigr).
\]
Under the pointed indexing category \(\Xi_{\ge0}^{\mathrm{op}}\) of Chapter 10, an object on the right is equivalently a separately pointed coherent double cochain object
\[
    C:
    \Xi_{\ge0}^{\mathrm{op}}
    \times
    \Xi_{\ge0}^{\mathrm{op}}
    \longrightarrow
    \mathcal D
\]
satisfying
\[
    C(*,q)\simeq0,
    \qquad
    C(p,*)\simeq0
\]
for every \(p,q\ge0\), whose restrictions to the two coordinate directions are coherent nonnegative cochain objects. Interchanging the two factors preserves these separate pointedness conditions and induces the canonical reverse-order comparison, including identity, composition, interchange, and all higher coherent cells.

Applied to Definition~\ref{def:var-complete-bifiltration}, the bidegree-\((p,q)\) term is canonically
\[
    \Omegavar^{p,q}(A;M),
\]
and the two unary structure maps are the connecting morphisms
\[
    \partial_V:
    \Omegavar^{p,q}\longrightarrow\Omegavar^{p+1,q},
    \qquad
    \partial_H:
    \Omegavar^{p,q}\longrightarrow\Omegavar^{p,q+1}.
\]
\end{lemma}

\begin{proof}
Because \(\mathcal D\) admits all small limits, the functor category of \(\mathbb Z\)-indexed filtrations admits all small limits pointwise. The full subcategory
\[
    \operatorname{Fil}^{\mathrm{comp}}_{\ge0}(\mathcal D)
\]
is stable, and the sequential limits entering the Chapter 10 construction are computed pointwise and remain available there. Finite limits, finite colimits, fibres, and the completeness condition are therefore computed in the ambient filtration category. Hence Chapter 10's complete-filtration/coherent-cochain equivalence may first be applied with target
\(\operatorname{Fil}^{\mathrm{comp}}_{\ge0}(\mathcal D)\):
\[
\begin{aligned}
    \operatorname{Fil}^{\mathrm{comp}}_{\ge0}
    \bigl(
      \operatorname{Fil}^{\mathrm{comp}}_{\ge0}(\mathcal D)
    \bigr)
    &\simeq
    \operatorname{CCh}_{\ge0}
    \bigl(
      \operatorname{Fil}^{\mathrm{comp}}_{\ge0}(\mathcal D)
    \bigr).
\end{aligned}
\]
The same equivalence is natural in exact functors preserving the sequential limits involved. Applying it pointwise in the remaining coherent-cochain variable therefore gives
\[
    \operatorname{CCh}_{\ge0}
    \bigl(
      \operatorname{Fil}^{\mathrm{comp}}_{\ge0}(\mathcal D)
    \bigr)
    \simeq
    \operatorname{CCh}_{\ge0}
    \bigl(
      \operatorname{CCh}_{\ge0}(\mathcal D)
    \bigr).
\]
Writing coherent cochains as pointed functors on
\(\Xi_{\ge0}^{\mathrm{op}}\), currying identifies the last category with the full subcategory of
\[
    \operatorname{Fun}\left(
      \Xi_{\ge0}^{\mathrm{op}}\times\Xi_{\ge0}^{\mathrm{op}},
      \mathcal D
    \right)
\]
consisting of the separately pointed functors \(C\) for which
\[
    C(*,q)\simeq0,
    \qquad
    C(p,*)\simeq0
\]
for all \(p,q\ge0\). Indeed, inner pointedness gives \(C(p,*)\simeq0\) for every \(p\), while outer pointedness says that \(C(*,-)\) is the zero coherent cochain object and hence gives \(C(*,q)\simeq0\) for every \(q\); conversely these two conditions curry to an iterated pointed functor. The flip of the two product factors preserves the separately pointed subcategory and gives the reverse-order comparison. Because it is an actual functor of indexing \(\infty\)-categories, it carries all identity, composition, interchange, and higher simplices.

For the bifiltration of Definition~\ref{def:var-complete-bifiltration}, completeness in the horizontal variable is built into its definition. Completeness in the vertical variable follows from
\[
    \lim_p\mathfrak F_V^p\Cvar\simeq0
\]
in the stable category of horizontal cosimplicial objects, since horizontal totalization, partial horizontal totalization, and the fibre defining \(F^{p,q}\) are limits. The vertical associated-graded term is
\[
    \Bvar^{p,\bullet}[-p],
\]
and applying the horizontal associated-graded calculation gives
\[
    N_H^q\Bvar^{p,\bullet}[-p-q]
    =
    \Omegavar^{p,q}[-p-q].
\]
Undoing the grading shifts in the complete-filtration/coherent-cochain equivalence identifies the two first structure maps with the displayed connecting morphisms.
\end{proof}

\subsection{Coherent boundaries and the signed bicomplex}

\begin{construction}[Unsigned coherent boundaries]
\label{con:var-unsigned-boundaries}
The complete vertical filtration in the stable category of horizontal cosimplicial objects gives coherent connecting morphisms
\[
    \partial_V^{p,q}:
    \Omegavar^{p,q}(A;M)
    \longrightarrow
    \Omegavar^{p+1,q}(A;M).
\]
Similarly the complete horizontal filtration gives coherent connecting morphisms
\[
    \partial_H^{p,q}:
    \Omegavar^{p,q}(A;M)
    \longrightarrow
    \Omegavar^{p,q+1}(A;M).
\]
Each family includes the specified nullhomotopies of consecutive composites and all higher coherent cochain relations supplied by the complete-filtration/coherent-cochain correspondence of Chapter 10.
\end{construction}

\begin{proposition}[Mixed coherent commutation]
\label{prop:var-mixed-commutation}
There is a canonical specified homotopy
\[
    \partial_H\partial_V
    \simeq
    \partial_V\partial_H
\]
on every bidegree \(\Omegavar^{p,q}\), compatible with the square-zero nullhomotopies and with all higher coherent relations of the two cochain directions.
\end{proposition}

\begin{proof}
Apply Lemma~\ref{lem:var-iterated-filtered-cochains} to the complete bifiltration \(F^{\bullet,\bullet}\DRvar\). The resulting coherent double cochain object is a single separately pointed functor on
\[
    \Xi_{\ge0}^{\mathrm{op}}
    \times
    \Xi_{\ge0}^{\mathrm{op}},
\]
in the sense of Lemma~\ref{lem:var-iterated-filtered-cochains}.
Its two coordinate unary maps are precisely \(\partial_V\) and \(\partial_H\). The two composites
\[
    \partial_H\partial_V,
    \qquad
    \partial_V\partial_H
\]
are therefore the two boundary paths around the same product square in the indexing \(\infty\)-category. The product functor supplies their canonical interchange homotopy. Because the square-zero nullhomotopies and every higher cochain relation are restrictions of the same product-indexed coherent object, the interchange homotopy is automatically compatible with them and with all higher cells. The reverse-order identification is the flip of the two factors from Lemma~\ref{lem:var-iterated-filtered-cochains}.
\end{proof}

\begin{definition}[Signed brave-new variational bicomplex]
\label{def:var-signed-bicomplex}
Define
\[
    d_V^{p,q}:=\partial_V^{p,q},
\]
and
\[
    d_H^{p,q}:=(-1)^p\partial_H^{p,q}.
\]
Here multiplication by \((-1)^p\) is the canonical sign automorphism of the mapping spectrum in the stable \(\infty\)-category.
\end{definition}

\begin{lemma}[Koszul rephasing of the coherent double cochain]
\label{lem:var-koszul-rephasing}
Let \(C^{p,q}\) be the coherent double cochain object obtained from Lemma~\ref{lem:var-iterated-filtered-cochains}, with unsigned unary maps \(\partial_V,\partial_H\). For a word \(w\) in the letters \(V,H\), beginning in vertical degree \(p\), let
\[
    h(w):=\#\{i\mid w_i=H\}
\]
and
\[
    \operatorname{vh}(w)
    :=
    \#\{(i,j)\mid i<j,\ w_i=V,\ w_j=H\}.
\]
Put
\[
    \kappa_p(w)
    :=
    p\,h(w)+\operatorname{vh}(w)
    \in\mathbb Z/2.
\]
Then the composite of the signed unary maps along \(w\) is canonically
\[
    d_w
    =
    (-1)^{\kappa_p(w)}\partial_w.
\]
For every coherent unsigned comparison
\[
    \partial_w\simeq\partial_{w'}
\]
supplied by the product-indexed double cochain object, multiplication by the central sign automorphism gives a corresponding signed comparison
\[
    d_w
    \simeq
    (-1)^{\kappa_p(w)-\kappa_p(w')}d_{w'}.
\]
These signed comparisons satisfy identity, composition, interchange, and all higher pasting coherences. In particular the adjacent mixed square is rephased by \(-1\).
\end{lemma}

\begin{proof}
Each horizontal unary map applied after a prefix of \(w\) is multiplied by \((-1)^{p+r}\), where \(r\) is the number of vertical letters in that prefix. Multiplying these signs over all horizontal letters gives
\[
    (-1)^{
      p\,h(w)+
      \#\{(i,j)\mid i<j,\ w_i=V,\ w_j=H\}
    }
    =
    (-1)^{\kappa_p(w)},
\]
which proves the composite formula.

Now let an unsigned coherent cell compare two boundary paths \(w,w'\). Multiplication by \(-1\) is a central natural automorphism of every mapping spectrum in a stable \(\infty\)-category, so the unsigned cell may be transported by the scalar
\[
    (-1)^{\kappa_p(w)-\kappa_p(w')}.
\]
For three paths \(w,w',w''\),
\[
    \kappa_p(w)-\kappa_p(w'')
    =
    \bigl(\kappa_p(w)-\kappa_p(w')\bigr)
    +
    \bigl(\kappa_p(w')-\kappa_p(w'')\bigr)
    \pmod2.
\]
Hence the rephasing scalar for a pasted comparison is the product of the rephasing scalars of its pieces. The same identity applies on every face of every higher simplex of the product coherent indexing category. Therefore the transported cells satisfy all identity, composition, interchange, and higher pasting relations.

For the two paths around one mixed square, written in order of application,
\[
    w=VH,
    \qquad
    w'=HV,
\]
one has
\[
    \kappa_p(VH)=p+1,
    \qquad
    \kappa_p(HV)=p,
\]
so their signed comparison carries the factor \(-1\).
\end{proof}

\begin{theorem}[Coherent brave-new variational bicomplex]
\label{thm:var-bicomplex}
The maps \(d_V,d_H\) carry specified coherent nullhomotopies
\[
    d_V^2\simeq0,
    \qquad
    d_H^2\simeq0,
\]
and a specified mixed homotopy
\[
    \boxed{
    d_Hd_V+d_Vd_H\simeq0.
    }
\]
The image of this mixed homotopy in the homotopy category is the ordinary signed anticommutation relation. All higher compatibility data is inherited from the explicit complete bifiltration through the Koszul rephasing of Lemma~\ref{lem:var-koszul-rephasing}. The construction is natural in Cartan morphisms and coefficients and commutes with arbitrary base change.
\end{theorem}

\begin{proof}
For a word containing only vertical arrows, \(\kappa_p(w)=0\), so the coherent vertical nullhomotopies and all their higher relations are unchanged. For two horizontal arrows,
\[
    \kappa_p(HH)=2p=0
    \quad\text{in }\mathbb Z/2,
\]
so the coherent horizontal nullhomotopy and its higher relations are likewise transported without changing their total sign. Thus
\[
    d_V^2\simeq0,
    \qquad
    d_H^2\simeq0
\]
with the complete one-directional coherent structures.

For the mixed square, Proposition~\ref{prop:var-mixed-commutation} supplies the unsigned comparison
\[
    \partial_H\partial_V
    \simeq
    \partial_V\partial_H.
\]
Lemma~\ref{lem:var-koszul-rephasing} rephases the two paths by opposite signs:
\[
\begin{aligned}
    d_Hd_V
    &\simeq
    (-1)^{p+1}\partial_H\partial_V,
    \\
    d_Vd_H
    &\simeq
    (-1)^p\partial_V\partial_H.
\end{aligned}
\]
Hence the transported interchange cell is equivalently a specified nullhomotopy
\[
    d_Hd_V+d_Vd_H\simeq0.
\]
The cocycle calculation in Lemma~\ref{lem:var-koszul-rephasing} proves that this mixed nullhomotopy is compatible with the two square-zero nullhomotopies and with every higher coherent cell. Naturality and base change follow from the corresponding properties of the unsigned product coherent object and the naturality of the central sign automorphisms.
\end{proof}

\subsection{The variational construction as an \(\infty\)-functor}
\label{subsec:var-infinity-functor}

\begin{definition}[Complete bifiltered target]
\label{def:var-bifiltered-target}
Put
\[
    \BiFil(\mathcal E_Q)
    :=
    \operatorname{Fil}^{\mathrm{comp}}_{\ge0}
    \bigl(
      \operatorname{Fil}^{\mathrm{comp}}_{\ge0}(\mathcal E_Q)
    \bigr).
\]
Morphisms are morphisms of the two complete decreasing filtrations, with their full mapping spaces.
\end{definition}

\begin{theorem}[Higher functoriality of the variational object]
\label{thm:var-infinity-functor}
Fix \(X\to S\) and \(Q=Q_{X/S}\). The construction above extends canonically to an \(\infty\)-functor
\[
    \boxed{
    \mathfrak{Var}_{X/S}:
    \bigl((\XSp)_{/Q}\bigr)^{\mathrm{op}}
    \times
    \mathcal E_Q
    \longrightarrow
    \BiFil(\mathcal E_Q),
    }
\]
whose value on \((\bar A\to Q,M)\) is
\[
    \mathfrak{Var}_{X/S}(A;M)
    =
    F^{\bullet,\bullet}\DRvar(A;M).
\]
A morphism \(f:\bar A\to\bar B\) over \(Q\) acts contravariantly by restriction of relation cochains, while a coefficient morphism \(M\to N\) acts covariantly. The functor retains the complete two-variable filtration maps and their entire mapping-space functoriality.

Composing with the iterated complete-filtration/coherent-cochain equivalence gives a functor
\[
    \Phi^{(2)}\mathfrak{Var}_{X/S}:
    \bigl((\XSp)_{/Q}\bigr)^{\mathrm{op}}
    \times
    \mathcal E_Q
    \longrightarrow
    \operatorname{CCh}_{\ge0}
    \bigl(\operatorname{CCh}_{\ge0}(\mathcal E_Q)\bigr),
\]
and the signed brave-new variational bicomplex is obtained from this coherent double cochain functor by the Koszul rephasing of Lemma~\ref{lem:var-koszul-rephasing}.
\end{theorem}

\begin{proof}
A morphism \(f:\bar A\to\bar B\) over \(Q\) induces, by functoriality of the relative de Rham construction, a simplicial map
\[
    V_{A,\bullet}\longrightarrow V_{B,\bullet}.
\]
Taking the product with the fixed horizontal Cartan relation gives a map of bisimplicial relations. Contravariant restriction along this map and the unit maps for the dependent adjunctions give a map of bicosimplicial coefficient objects
\[
    \Cvar^{\bullet,\bullet}(B;M)
    \longrightarrow
    \Cvar^{\bullet,\bullet}(A;M).
\]
Coefficient morphisms act covariantly degreewise. Composition, identities, and higher homotopies are inherited from functoriality of finite limits and of the stable dependent adjunctions.

Partial totalizations, complete totalizations, normalization, fibres, and the two complete filtration constructions are functors between stable \(\infty\)-categories. Applying them degreewise therefore produces the stated functor into \(\BiFil(\mathcal E_Q)\). The coherent double-cochain functor follows from the naturality of the Chapter 10 complete-filtration equivalence. Koszul rephasing is by central sign automorphisms of mapping spectra and is natural in every morphism of the coherent double-cochain object.
\end{proof}

\begin{remark}[The bicomplex is a signed shadow]
\label{rem:var-bicomplex-shadow}
The object \((\Omegavar^{p,q},d_V,d_H)\) is useful notation and recovers the familiar shape of a variational bicomplex, but it is not the primary invariant. The primary object is \(\mathfrak{Var}_{X/S}(A;M)\), equivalently its product-indexed coherent double cochain object. Passing only to the unary signed maps forgets the explicit higher cells that witness square-zero, interchange, and all higher compatibility relations.
\end{remark}

\begin{definition}[Cartan--vertical filtration]
\label{def:var-contact-filtration}
The vertical edge
\[
    F_{\mathcal C}^p\DRvar(A;M)
    :=
    F^{p,0}\DRvar(A;M)
\]
of the complete bifiltration is called the \emph{Cartan--vertical filtration}. The notation \(F_{\mathcal C}\) is retained to facilitate comparison with classical contact geometry. When no confusion can arise we also call it the Cartan-contact filtration, but this is comparison terminology rather than an ideal-theoretic definition. Membership in \(F_{\mathcal C}^p\) is called vertical variational degree at least \(p\).
\end{definition}

\begin{remark}[No exterior-form identification]
\label{rem:var-no-exterior}
The notation \(\Omegavar^{p,q}\) records the two normalization degrees. It does not identify these objects with exterior powers of horizontal and vertical cotangent complexes. In particular the vertical degree-one object retains the higher infinitesimal information which Chapter 10 showed can be killed by first-order cotangent linearization.
\end{remark}

\begin{remark}[Contact filtration is not a contact ideal filtration]
\label{rem:var-contact-not-ideal}
No assertion is made that \(F_{\mathcal C}^p\) is the \(p\)-th power of an ideal. Chapter 10 constructs ordered Alexander--Whitney products on normalized infinitesimal cochains and proves their associative, unital, and Leibniz identities in the homotopy category, but explicitly does not assert the multiplicative filtration required to identify the totalization tower with powers of a contact ideal. The present Cartan--vertical filtration is the vertical edge of the complete bifiltration itself.
\end{remark}

\section{Secondary forms and the full variational operator}
\label{sec:var-secondary-forms}

\begin{definition}[Secondary form spectra]
\label{def:var-secondary}
For every \(p\ge0\), define
\[
    \boxed{
    \Sec^p(A;M)
    :=
    \Tot_H\Bvar^{p,\bullet}(A;M).
    }
\]
Thus horizontal totalization is applied while the horizontal direction is still cosimplicial. For scalar coefficients write \(\Sec^p(A):=\Sec^p(A;\Oadd_Q)\).
\end{definition}

\begin{definition}[Coefficient-valued normalized infinitesimal cochains]
\label{def:var-coefficient-normalized}
For \(p\ge0\), define
\[
    \Omegainf^p(\bar A/Q;M)
    :=
    N^pC_{\dR}^{\bullet}(\bar A/Q;M).
\]
For scalar coefficients this is the Chapter 10 normalized infinitesimal cochain object.
\end{definition}

\begin{theorem}[Secondary forms are intrinsic normalized cochains]
\label{thm:var-secondary-identification}
For every \(p\ge0\), there is a canonical equivalence
\[
    \boxed{
    \Sec^p(A;M)
    \simeq
    \Omegainf^p(\bar A/Q;M).
    }
\]
The equivalence is natural in Cartan morphisms and coefficient morphisms and commutes with arbitrary base change.
\end{theorem}

\begin{proof}
Vertical normalization is a finite limit. Hence it commutes with horizontal totalization:
\[
\begin{aligned}
    \Sec^p(A;M)
    &=
    \Tot_HN_V^p\Cvar\\
    &\simeq
    N_V^p\Tot_H\Cvar.
\end{aligned}
\]
By Theorem~\ref{thm:var-horizontal-descent-cochains}, the horizontal totalization is canonically the vertical cosimplicial object \(C_{\dR}^{\bullet}(\bar A/Q;M)\). Therefore
\[
    \Sec^p(A;M)
    \simeq
    N^pC_{\dR}^{\bullet}(\bar A/Q;M)
    =
    \Omegainf^p(\bar A/Q;M).
\]
All operations used are natural and commute with the base-change equivalences already constructed.
\end{proof}

\begin{remark}[Horizontal degree is a complete descent resolution]
\label{rem:var-horizontal-descent-resolution}
The equivalence
\[
    \Sec^p(A;M)
    \simeq
    \Omegainf^p(\bar A/Q;M)
\]
means that the horizontal direction has been completely descended. Thus the horizontal bidegrees \(q>0\) of \(\Omegavar^{p,q}\) encode the complete Cartan descent resolution, its filtration, and its higher coherence data before totalization; after horizontal totalization they do not define an additional independent differential on \(\Sec^p\). This distinction is essential in the classical comparison: the brave-new horizontal filtration is intrinsic complete descent data and is not identified with a power filtration by a contact ideal.
\end{remark}

\begin{corollary}[Contact associated graded]
\label{cor:var-contact-graded}
There are canonical equivalences
\[
    \boxed{
    \operatorname{gr}_{\mathcal C}^p\DRvar(A;M)
    \simeq
    \Sec^p(A;M)[-p]
    \simeq
    \Omegainf^p(\bar A/Q;M)[-p].
    }
\]
In particular the Cartan--vertical filtration is complete and separated.
\end{corollary}

\begin{proof}
The vertical associated graded of the double totalization is the horizontal totalization of \(\Bvar^{p,\bullet}[-p]\) by the vertical layer theorem. Apply Definition~\ref{def:var-secondary} and Theorem~\ref{thm:var-secondary-identification}.
\end{proof}

\begin{theorem}[Full brave-new variational operator]
\label{thm:var-operator}
The coherent vertical boundary descends through horizontal totalization to canonical maps
\[
    \boxed{
    \delta^{\infty}_{\mathrm{var}}:
    \Sec^p(A;M)
    \longrightarrow
    \Sec^{p+1}(A;M).
    }
\]
These maps form a coherent nonnegative cochain object. In particular, consecutive composites carry specified nullhomotopies
\[
    (\delta^{\infty}_{\mathrm{var}})^2\simeq0
\]
and the full hierarchy of higher coherent relations.
\end{theorem}

\begin{proof}
The vertical complete-filtration boundary
\[
    \partial_V:\Bvar^{p,\bullet}\longrightarrow\Bvar^{p+1,\bullet}
\]
is a morphism in the category of horizontal cosimplicial objects, by the functoriality used in Proposition~\ref{prop:var-mixed-commutation}. Horizontal totalization therefore gives the displayed map. The specified vertical nullhomotopies and higher coherent simplices are natural transformations of horizontal cosimplicial objects and hence totalize with the map.
\end{proof}

\begin{theorem}[Identification with the intrinsic infinitesimal differential]
\label{thm:var-operator-identification}
Under Theorem~\ref{thm:var-secondary-identification}, the full variational operator is canonically identified with the coherent infinitesimal differential of \(\bar A/Q\):
\[
\begin{tikzcd}[column sep=huge]
\Sec^p(A;M)
  \arrow[r,"\delta^{\infty}_{\mathrm{var}}"]
  \arrow[d,"\sim"']
& \Sec^{p+1}(A;M)
  \arrow[d,"\sim"]
\\
\Omegainf^p(\bar A/Q;M)
  \arrow[r,"d_{\mathrm{inf}}"']
& \Omegainf^{p+1}(\bar A/Q;M).
\end{tikzcd}
\]
The comparison identifies the specified nullhomotopies of the square and all higher cochain coherences.
\end{theorem}

\begin{proof}
After horizontal effective descent, the vertical complete totalization tower of the double cochain object is exactly the complete totalization tower of \(C_{\dR}^{\bullet}(\bar A/Q;M)\). Both coherent differentials are the connecting morphisms of that same exact tower. The complete-filtration/coherent-cochain equivalence is functorial on the full coherent indexing category, so it identifies the complete coherent cochain objects rather than only their unary boundary maps.
\end{proof}

\section{Coherently closed forms, exact classifiers, and primitive spaces}
\label{sec:var-closed-primitive}

\begin{definition}[Coherently closed secondary forms]
\label{def:var-closed}
Let
\[
    F_{\mathrm{Tot}}^p(\bar A/Q;M)
\]
denote the complete totalization filtration of \(C_{\dR}^{\bullet}(\bar A/Q;M)\). For \(p\ge0\) and \(n\in\mathbb Z\), define
\[
    \boxed{
    \Sec^{p,\mathrm{cl}}(A;M;n)
    :=
    F_{\mathrm{Tot}}^p(\bar A/Q;M)[n+p].
    }
\]
Through Theorem~\ref{thm:var-total-identification}, this is equivalently the corresponding vertical stage of the variational totalization.
\end{definition}

\begin{proposition}[Underlying secondary form and retained closedness]
\label{prop:var-underlying-closed}
There is a canonical map
\[
    u_p:
    \Sec^{p,\mathrm{cl}}(A;M;n)
    \longrightarrow
    \Sec^p(A;M)[n].
\]
A point of the source is a shifted secondary \(p\)-form together with a specified nullhomotopy of its variational boundary and the complete hierarchy of higher compatibility data. The homotopy fibre of the induced map on zeroth spaces over a chosen form is the space of coherent closedness structures on that form.
\end{proposition}

\begin{proof}
The associated-graded map of the complete filtration is
\[
    F_{\mathrm{Tot}}^p
    \longrightarrow
    \operatorname{gr}^pF_{\mathrm{Tot}}
    \simeq
    \Omegainf^p(\bar A/Q;M)[-p].
\]
Shift by \(n+p\) and use Theorem~\ref{thm:var-secondary-identification}. The interpretation of the fibre is the complete-filtration/coherent-cochain theorem of Chapter 10: a lift through all later filtration stages is precisely the retained closedness hierarchy.
\end{proof}

\begin{construction}[Exact variational classifier]
\label{con:var-exact-classifier}
For \(p\ge1\), the adjacent-stage cofiber sequence gives a canonical connecting morphism
\[
    \boxed{
    \delta^{\mathrm{cl}}_{\mathrm{var}}:
    \Sec^{p-1}(A;M)[n]
    \longrightarrow
    \Sec^{p,\mathrm{cl}}(A;M;n).
    }
\]
Its underlying secondary form under \(u_p\) is \(\delta^{\infty}_{\mathrm{var}}\), equipped with the canonical coherent closedness structure induced by the specified nullhomotopy of \((\delta^{\infty}_{\mathrm{var}})^2\) and all higher relations.
\end{construction}

\begin{definition}[Primitive space]
\label{def:var-primitive}
For a point
\[
    \eta:\one_Q\longrightarrow\Sec^{p,\mathrm{cl}}(A;M;n),
\]
define
\[
    \boxed{
    \Prim_{\mathrm{var}}(\eta)
    :=
    \hofib_{\eta}\left(
      \Map_{\mathcal E_Q}(\one_Q,\Sec^{p-1}(A;M)[n])
      \longrightarrow
      \Map_{\mathcal E_Q}(\one_Q,\Sec^{p,\mathrm{cl}}(A;M;n))
    \right).
    }
\]
\end{definition}

\begin{remark}[Primitives are a space]
\label{rem:var-exactness-space}
The existence of a primitive is the nonemptiness of \(\Prim_{\mathrm{var}}(\eta)\). No truncation is applied: distinct primitives, paths between them, automorphisms, and higher homotopies remain part of the answer.
\end{remark}

\section{Universal moduli maps and their homotopy fibres}
\label{sec:var-moduli}

\begin{definition}[Moduli objects from zeroth spaces]
\label{def:var-moduli}
For \(p\ge0\) and \(n\in\mathbb Z\), define pointed objects of the slice \((\XSp)_{/Q}\) by
\[
    \boxed{
    \boldsymbol\Omega^p_{\mathrm{var},A/Q}(M;n)
    :=
    \Omega_Q^\infty\bigl(\Sec^p(A;M)[n]\bigr),
    }
\]
and
\[
    \boxed{
    \boldsymbol\Omega^{p,\mathrm{cl}}_{\mathrm{var},A/Q}(M;n)
    :=
    \Omega_Q^\infty\bigl(\Sec^{p,\mathrm{cl}}(A;M;n)\bigr).
    }
\]
\end{definition}

\begin{proposition}[Parameterized representing property]
\label{prop:var-moduli-representing}
For every \(t:T\to Q\), there are natural equivalences
\[
\begin{aligned}
    \Map_Q\bigl(T,\boldsymbol\Omega^p_{\mathrm{var},A/Q}(M;n)\bigr)
    &\simeq
    \Omega^\infty
    \operatorname{map}_{\mathcal E_T}
    \bigl(\one_T,t^*\Sec^p(A;M)[n]\bigr),
    \\
    \Map_Q\bigl(T,\boldsymbol\Omega^{p,\mathrm{cl}}_{\mathrm{var},A/Q}(M;n)\bigr)
    &\simeq
    \Omega^\infty
    \operatorname{map}_{\mathcal E_T}
    \bigl(\one_T,t^*\Sec^{p,\mathrm{cl}}(A;M;n)\bigr).
\end{aligned}
\]
The full variational operator, underlying-closed-form maps, and exact classifiers induce the corresponding morphisms of pointed moduli objects by functoriality of \(\Omega_Q^\infty\).
\end{proposition}

\begin{proof}
For \(E\in\mathcal E_Q\), adjunction and the stable dependent-sum formula give
\[
\begin{aligned}
    \Map_Q(T,\Omega_Q^\infty E)
    &\simeq
    \Omega^\infty
    \operatorname{map}_{\mathcal E_Q}(\Sigma_t\one_T,E)\\
    &\simeq
    \Omega^\infty
    \operatorname{map}_{\mathcal E_T}(\one_T,t^*E).
\end{aligned}
\]
Apply this to the displayed coefficient objects. No separate representability theorem and no evaluation of a cochain functor on a generally nonrepresented moduli object is required.
\end{proof}

\begin{construction}[Universal exact-form moduli morphism]
\label{con:var-universal-primitive-map}
Applying \(\Omega_Q^\infty\) to the exact variational classifier gives a canonical morphism of objects over \(Q\)
\[
    \boxed{
    \mathbf d^{\mathrm{cl}}_{\mathrm{var}}:
    \boldsymbol\Omega^{p-1}_{\mathrm{var},A/Q}(M;n)
    \longrightarrow
    \boldsymbol\Omega^{p,\mathrm{cl}}_{\mathrm{var},A/Q}(M;n).
    }
\]
The parameterized primitive object of a chosen coherently closed form is the homotopy fibre of this universal morphism over the corresponding section. Thus the universal moduli map is primary and individual primitive spaces are its fibres.
\end{construction}

\begin{definition}[Parameterized primitive object]
\label{def:var-parameterized-primitive}
Let
\[
    \eta:\one_Q\longrightarrow
    \Sec^{p,\mathrm{cl}}(A;M;n)
\]
be a coherently closed secondary form, and write the corresponding section again as
\[
    \eta:Q\longrightarrow
    \boldsymbol\Omega^{p,\mathrm{cl}}_{\mathrm{var},A/Q}(M;n).
\]
Define the \emph{parameterized primitive object}
\[
    \boxed{
    \mathbf{Prim}_{\mathrm{var}}(\eta)
    :=
    Q
    \times_{\boldsymbol\Omega^{p,\mathrm{cl}}_{\mathrm{var},A/Q}(M;n)}
    \boldsymbol\Omega^{p-1}_{\mathrm{var},A/Q}(M;n),
    }
\]
where the second map is induced by the exact variational classifier
\(\delta^{\mathrm{cl}}_{\mathrm{var}}\).
\end{definition}

\begin{proposition}[Parameterized primitive representing property]
\label{prop:var-parameterized-primitive}
For every \(t:T\to Q\), there is a natural equivalence
\[
    \boxed{
    \Map_Q\bigl(T,\mathbf{Prim}_{\mathrm{var}}(\eta)\bigr)
    \simeq
    \Prim_{\mathrm{var},T}(t^*\eta),
    }
\]
where the right side is the homotopy fibre, over the pulled-back closed form \(t^*\eta\), of
\[
    \Omega^\infty
    \operatorname{map}_{\mathcal E_T}
    \bigl(\one_T,t^*\Sec^{p-1}(A;M)[n]\bigr)
    \longrightarrow
    \Omega^\infty
    \operatorname{map}_{\mathcal E_T}
    \bigl(\one_T,t^*\Sec^{p,\mathrm{cl}}(A;M;n)\bigr).
\]
In particular,
\[
    \Map_Q\bigl(Q,\mathbf{Prim}_{\mathrm{var}}(\eta)\bigr)
    \simeq
    \Prim_{\mathrm{var}}(\eta).
\]
The parameterized primitive object commutes with arbitrary pullback in \(Q\).
\end{proposition}

\begin{proof}
Apply Proposition~\ref{prop:var-moduli-representing} to the homotopy pullback defining
\(\mathbf{Prim}_{\mathrm{var}}(\eta)\). Mapping out of \(T\) preserves that homotopy pullback, and the two mapping objects are exactly the spaces of pulled-back ordinary and coherently closed forms. The final base-change assertion follows because \(\Omega_Q^\infty\), pullback of parameterized spectra, and finite homotopy pullbacks commute with pullback.
\end{proof}

\section{Lagrangian classes and their local full variation}
\label{sec:var-lagrangians}

\begin{proposition}[Degree-zero secondary coefficient]
\label{prop:var-degree-zero}
There is a canonical equivalence
\[
    \boxed{
    \Sec^0(A;M)
    \simeq
    \Pi_a a^*M.
    }
\]
\end{proposition}

\begin{proof}
Normalization in degree zero is the degree-zero cosimplicial object. Therefore
\[
    \Sec^0(A;M)
    \simeq
    \Omegainf^0(\bar A/Q;M)
    =
    C_{\dR}^0(\bar A/Q;M)
    =
    \Pi_a a^*M.
\]
\end{proof}

\begin{definition}[Brave-new Lagrangian class]
\label{def:var-lagrangian}
A shifted Lagrangian class of degree \(n\) with coefficient \(M\) is a morphism
\[
    L:\one_Q\longrightarrow\Sec^0(A;M)[n].
\]
By Proposition~\ref{prop:var-degree-zero} and the adjunction \(a^*\dashv\Pi_a\), it has a canonically adjoint local section
\[
    \boxed{
    \ell_L:\one_{\bar A}\longrightarrow a^*M[n].
    }
\]
\end{definition}

\begin{remark}[Lagrangian versus action]
\label{rem:var-lagrangian-action}
For an arbitrary coefficient \(M\), the section \(\ell_L\) is only a local secondary coefficient on the Cartan descent object and no action functional is implicit. Section~\ref{sec:var-action-functional} later specializes \(M\) to the canonically constructed dualizing-density coefficient \(M^\partial_{x,D}\); in that sector the required morphism into \(x_e^!D\) is constructed from the coefficient itself and the Chapter 10 exceptional trace produces the action.
\end{remark}

\subsection{The local full one-variation coefficient}

Let
\[
    R_1:=R_{\bar A/Q,1}
    =
    \bar A\times_{\bar A_{\dR/Q}}\bar A.
\]
Write
\[
    c,e:R_1\rightrightarrows\bar A
\]
for the first and second vertex maps and
\[
    s:\bar A\longrightarrow R_1
\]
for the diagonal degeneracy. Thus \(cs=\id_{\bar A}=es\).

\begin{definition}[Local full one-variation coefficient]
\label{def:var-full-local-coefficient}
Let
\[
    \varepsilon_c:
    c^*\Pi_c
    \longrightarrow
    \id_{\mathcal E_{R_1}}
\]
be the counit of \(c^*\dashv\Pi_c\).  Define the restriction morphism
\[
    \boxed{
    \operatorname{ev}_s:
    \Pi_ce^*a^*M
    \longrightarrow
    a^*M
    }
\]
to be the composite
\[
\begin{aligned}
    \Pi_ce^*a^*M
    &\simeq
    s^*c^*\Pi_ce^*a^*M\\
    &\xrightarrow{\;s^*(\varepsilon_c)\;}
    s^*e^*a^*M
    \simeq
    a^*M,
\end{aligned}
\]
where \(cs=\id_{\bar A}=es\).  Define
\[
    \boxed{
    \Vfull_M(A)
    :=
    \fib\left(
      \Pi_c e^*a^*M
      \xrightarrow{\operatorname{ev}_s}
      a^*M
    \right)
    \in\mathcal E_{\bar A}.
    }
\]
Thus the displayed map is the type-correct dependent-adjunction form of restriction along the diagonal section.
\end{definition}

\begin{proposition}[Pushforward of the local full coefficient]
\label{prop:var-full-local-pushforward}
There is a canonical equivalence
\[
    \boxed{
    \Pi_a\Vfull_M(A)
    \simeq
    \Sec^1(A;M).
    }
\]
\end{proposition}

\begin{proof}
Since \(ac=ae=v_1\), composition of dependent products gives
\[
    \Pi_a\Pi_c e^*a^*M
    \simeq
    \Pi_{v_1}v_1^*M
    =
    C_{\dR}^1(\bar A/Q;M).
\]
Likewise \(\Pi_aa^*M=C_{\dR}^0(\bar A/Q;M)\). The functor \(\Pi_a\) is exact, hence preserves the defining fibre of \(\Vfull_M(A)\). The resulting fibre is the degree-one normalization
\[
    \fib(C_{\dR}^1\to C_{\dR}^0)
    =
    \Omegainf^1(\bar A/Q;M)
    \simeq
    \Sec^1(A;M).
\]
\end{proof}

\begin{definition}[Local full Euler section]
\label{def:var-full-euler-section}
The full variational derivative
\[
    \delta^{\infty}_{\mathrm{var}}L:
    \one_Q\longrightarrow\Sec^1(A;M)[n]
    \simeq
    \Pi_a\Vfull_M(A)[n]
\]
has, by \(a^*\dashv\Pi_a\), a canonical adjoint
\[
    \boxed{
    \mathsf E_L^{\infty}:
    \one_{\bar A}\longrightarrow\Vfull_M(A)[n].
    }
\]
\end{definition}

\begin{definition}[Full critical locus]
\label{def:var-full-critical}
Let
\[
    \mathbf E_L^{\infty}
    :=
    \Omega_{\bar A}^\infty\bigl(\Vfull_M(A)[n]\bigr)
    \longrightarrow\bar A.
\]
The section \(\mathsf E_L^{\infty}\) and the zero map determine two sections
\[
    e_L^{\infty},0:\bar A\rightrightarrows\mathbf E_L^{\infty}.
\]
Define
\[
    \boxed{
    \Crit^{\infty}(L)
    :=
    \Eq(e_L^{\infty},0)
    }
\]
in \((\XSp)_{/\bar A}\).
\end{definition}

\begin{remark}[Full criticality]
\label{rem:var-full-critical-strong}
The object \(\Crit^{\infty}(L)\) kills the complete local finite-infinitesimal variation, not merely its one-stage square-zero image. Chapter 10's affine-line calculation proves that the full and first one-variation coefficients can differ nontrivially: the first-order linearization has a nonzero kernel on \(\pi_0\). No assertion is made here that the specific kernel class exhibited there is itself the full Euler variation of a Lagrangian, so that calculation alone does not prove that \(\Crit^{\infty}(L)\to\Crit(L)\) fails to be an equivalence for a particular \(L\).
\end{remark}

\section{The constructed one-stage square-zero variation coefficient}
\label{sec:var-square-zero-sector}

\begin{definition}[Split square-zero variation diagram]
\label{def:var-split-sqz-test}
Let \(T=\operatorname{Spec}B\) be an affine spectral scheme, let
\[
    a_T:T\longrightarrow\bar A
\]
be an affine point, and let \(N\in\operatorname{Mod}_B\) be connective. The split square-zero algebra of Chapter 2 gives maps
\[
    B\longrightarrow B\oplus N\longrightarrow B
\]
whose composite is the identity. Write
\[
    T[N]:=\operatorname{Spec}(B\oplus N).
\]
Contravariantly, the projection \(B\oplus N\to B\) gives the structured square-zero immersion
\[
    i_N:T\longrightarrow T[N],
\]
and the split inclusion \(B\to B\oplus N\) gives a retraction
\[
    r_N:T[N]\longrightarrow T,
    \qquad
    r_Ni_N=\id_T.
\]
A \emph{split square-zero variation diagram} centred at \(a_T\) is a map
\[
    \widetilde a:T[N]\longrightarrow\bar A
\]
over \(Q\), together with the specified equality
\[
    \widetilde a i_N\simeq a_T.
\]
The structural map \(T[N]\to Q\) is equivalently \(a\,a_T r_N\), where \(a:\bar A\to Q\).
\end{definition}

\begin{definition}[The variation diagram category]
\label{def:var-sqz-category}
Let
\[
    \mathcal V^{(1)}_{\bar A/Q}
\]
be the \(\infty\)-category whose objects are split square-zero variation diagrams
\[
    v=(T,N,a_T,\widetilde a)
\]
as above.  A morphism
\[
    v'=(T',N',a_{T'},\widetilde a')
    \longrightarrow
    v=(T,N,a_T,\widetilde a)
\]
consists of a morphism of centres
\[
    g:T'\longrightarrow T
\]
over \(\bar A\), a morphism of connective modules
\[
    \alpha:g^*N\longrightarrow N',
\]
and the induced morphism of structured split square-zero extensions
\[
    \widehat g_\alpha:
    T'[N']
    \longrightarrow
    T'[g^*N]
    \simeq
    T[N]\times_TT'
    \longrightarrow
    T[N],
\]
together with a specified equivalence
\[
    \widetilde a'
    \simeq
    \widetilde a\,\widehat g_\alpha
\]
compatible with the central sections and split retractions.  The map
\(\widehat g_\alpha\) is the affine opposite of the morphism of split square-zero algebras induced by the Chapter 2 functoriality
\(g^*N\to N'\).

Write
\[
    \pi_1:\mathcal V^{(1)}_{\bar A/Q}
    \longrightarrow
    \mathcal C^{\mathrm{aff}}_{/\bar A}
\]
for the centre functor.  Pullback of the complete structured variation diagram along a morphism of affine centres gives Cartesian lifts for \(\pi_1\); hence \(\pi_1\) is a Cartesian fibration.
\end{definition}

\begin{definition}[First-difference coefficient on variation diagrams]
\label{def:var-first-difference-functor}
For
\[
    v=(T,N,a_T,\widetilde a)
\]
define
\[
    \boxed{
    K_M^{(1)}(v)
    :=
    \fib\left(
      \Gamma_{T[N]}^{\mathrm{st}}(M_{T[N]})
      \longrightarrow
      \Gamma_T^{\mathrm{st}}(M_T)
    \right).
    }
\]
The arrow is restriction along \(i_N\). A structured morphism of variation diagrams gives a restriction morphism on these fibres; hence
\[
    K_M^{(1)}:
    (\mathcal V^{(1)}_{\bar A/Q})^{\mathrm{op}}
    \longrightarrow
    \operatorname{Sp}
\]
is a functor.
\end{definition}

\begin{definition}[Centre-pullback Nisnevich topology on split variations]
\label{def:var-first-variation-topology}
Equip \(\mathcal V^{(1)}_{\bar A/Q}\) with the Grothendieck topology generated by Cartesian pullback of affine spectral Nisnevich covering families along
\(\pi_1\). Thus, for a variation
\[
    v=(T,N,a_T,\widetilde a)
\]
and an affine spectral Nisnevich covering family
\[
    \{T_i\longrightarrow T\}_{i\in I},
\]
the Cartesian pullbacks \(v_i\to v\) form a covering family. We call this the \emph{centre-pullback Nisnevich topology}.
\end{definition}

\begin{lemma}[Cartesian-centre right-Kan descent]
\label{lem:var-cartesian-centre-kan}
Let \((\mathcal C,\tau)\) be an \(\infty\)-site whose topology is generated by a pullback-stable class \(\mathscr B\) of covering families, and assume that every \(\tau\)-covering sieve contains a covering family belonging to \(\mathscr B\). Let
\[
    \pi:\mathcal V\longrightarrow\mathcal C
\]
be a Cartesian fibration, and give \(\mathcal V\) the Grothendieck topology generated by Cartesian pullbacks of the families in \(\mathscr B\). If
\[
    K:\mathcal V^{\mathrm{op}}\longrightarrow\operatorname{Sp}
\]
is a sheaf for this topology, then
\[
    \operatorname{Ran}_{\pi^{\mathrm{op}}}K
\]
is a \(\tau\)-sheaf on \(\mathcal C\).
\end{lemma}

\begin{proof}
We first prove the space-valued statement. Let
\[
    \pi^*:
    \operatorname{PSh}(\mathcal C;\mathcal S)
    \longrightarrow
    \operatorname{PSh}(\mathcal V;\mathcal S)
\]
be precomposition. If \(R\hookrightarrow h_C\) is a covering sieve in \(\mathcal C\), then
\[
    \pi^*R\longrightarrow\pi^*h_C
\]
is a local equivalence for the centre-pullback topology. Indeed, a point of \(\pi^*h_C(v)\) is a morphism
\[
    f:\pi(v)\longrightarrow C.
\]
The pullback sieve \(f^*R\) covers \(\pi(v)\). By the basis hypothesis it contains a covering family
\[
    \{c_i\longrightarrow\pi(v)\}_{i\in I}
\]
belonging to \(\mathscr B\). Cartesian lifting along \(\pi\) gives a covering family
\[
    \{v_i\longrightarrow v\}_{i\in I}
\]
of \(v\), and on every member of that family the restriction of \(f\) factors through \(R\). Thus
\[
    \pi^*R\longrightarrow\pi^*h_C
\]
is a monomorphism which is locally surjective, hence becomes an equivalence after sheafification.

The \(\tau\)-local equivalences in \(\operatorname{PSh}(\mathcal C;\mathcal S)\) form the strongly saturated class generated by the covering-sieve inclusions \(R\hookrightarrow h_C\). Precomposition preserves colimits and retracts, so the preceding calculation implies that \(\pi^*\) carries the strongly saturated class generated by these inclusions into the strongly saturated class of local equivalences on \(\mathcal V\). Hence \(\pi^*\) carries every local equivalence on \(\mathcal C\) to a local equivalence on \(\mathcal V\). If \(K:\mathcal V^{\mathrm{op}}\to\mathcal S\) is a sheaf and \(P\to P'\) is a local equivalence on \(\mathcal C\), the right-Kan adjunction gives
\[
\begin{aligned}
    \Map\bigl(P',\operatorname{Ran}_{\pi^{\mathrm{op}}}K\bigr)
    &\simeq
    \Map(\pi^*P',K)
    \\
    &\xrightarrow{\sim}
    \Map(\pi^*P,K)
    \\
    &\simeq
    \Map\bigl(P,\operatorname{Ran}_{\pi^{\mathrm{op}}}K\bigr).
\end{aligned}
\]
Hence \(\operatorname{Ran}_{\pi^{\mathrm{op}}}K\) is a space-valued sheaf.

Now let \(K\) be spectrum-valued. For every \(n\in\mathbb Z\), the functor
\[
    \Omega^\infty\Sigma^n:\operatorname{Sp}\to\mathcal S
\]
preserves limits. Since right Kan extension is computed by limits,
\[
    \Omega^\infty\Sigma^n
    \operatorname{Ran}_{\pi^{\mathrm{op}}}K
    \simeq
    \operatorname{Ran}_{\pi^{\mathrm{op}}}
    \bigl(\Omega^\infty\Sigma^nK\bigr).
\]
The spectrum-valued sheaf condition on \(K\) implies that every \(\Omega^\infty\Sigma^nK\) is a space-valued sheaf. By the first part, every displayed right Kan extension is a sheaf on \(\mathcal C\). The family \(\{\Omega^\infty\Sigma^n\}_{n\in\mathbb Z}\) detects equivalences of spectra, so \(\operatorname{Ran}_{\pi^{\mathrm{op}}}K\) is a spectrum-valued sheaf.
\end{proof}

\begin{definition}[Raw first-variation presheaf and its sheafification]
\label{def:var-first-coeff-functor}
Define the spectrum-valued presheaf on affine points of \(\bar A\)
\[
    \Vfirst_{M,\mathrm{pre}}(A)
    :=
    \operatorname{Ran}_{\pi_1^{\mathrm{op}}}K_M^{(1)}.
\]
Thus, at an affine point \(c:T\to\bar A\), its value is
\[
    \Vfirst_{M,\mathrm{pre}}(A)(T,c)
    \simeq
    \lim_{\mathcal J_c}K_M^{(1)},
    \qquad
    \mathcal J_c:=(c^{\mathrm{op}}\downarrow\pi_1^{\mathrm{op}}).
\]
Geometrically, an object of \(\mathcal J_c\) is a variation diagram \(v\) together with a morphism \(\pi_1(v)\to c\) of affine centres; morphisms in \(\mathcal J_c\) reverse the corresponding morphisms of variation diagrams because the right Kan extension is taken along \(\pi_1^{\mathrm{op}}\). Define the \emph{one-stage first-variation coefficient} by spectral Nisnevich sheafification:
\[
    \boxed{
    \Vfirst_M(A)
    :=
    a_{\mathrm{Nis}}^{\operatorname{Sp}}
    \Vfirst_{M,\mathrm{pre}}(A)
    \in
    \operatorname{Shv}_{\mathrm{Nis}}
    (\mathcal C^{\mathrm{aff}}_{/\bar A};\operatorname{Sp})
    \simeq
    \mathcal E_{\bar A}.
    }
\]
\end{definition}

\begin{theorem}[First-difference descent and the raw right-Kan sheaf]
\label{thm:var-first-kan-sheaf}
The functor
\[
    K_M^{(1)}:
    (\mathcal V^{(1)}_{\bar A/Q})^{\mathrm{op}}
    \longrightarrow
    \operatorname{Sp}
\]
satisfies descent for the centre-pullback Nisnevich topology of Definition~\ref{def:var-first-variation-topology}. Moreover the raw right Kan extension
\[
    \Vfirst_{M,\mathrm{pre}}(A)
    =
    \operatorname{Ran}_{\pi_1^{\mathrm{op}}}K_M^{(1)}
\]
is already a spectral Nisnevich sheaf on
\(\mathcal C^{\mathrm{aff}}_{/\bar A}\). Consequently the sheafification unit is an equivalence
\[
    \boxed{
    \Vfirst_{M,\mathrm{pre}}(A)
    \xrightarrow{\sim}
    \Vfirst_M(A).
    }
\]
For an affine centre \(c:T\to\bar A\) and a variation \(w\) equipped with a morphism of its centre to \(c\), the right-Kan limit therefore supplies a canonical evaluation morphism
\[
    \boxed{
    \Vfirst_M(A)(T,c)
    \longrightarrow
    K_M^{(1)}(w).
    }
\]
\end{theorem}

\begin{proof}
Let
\[
    \{T_i\to T\}_{i\in I}
\]
be an affine spectral Nisnevich cover and let \(v_i\) be the Cartesian pullback of
\[
    v=(T,N,a_T,\widetilde a)
\]
to \(T_i\). Write \(N_i\) for the pulled-back module. Functoriality and base change of split square-zero algebras give
\[
    T_i[N_i]\simeq T[N]\times_TT_i,
\]
and the Čech nerve of the family \(T_i[N_i]\to T[N]\) is the pullback of the Čech nerve of \(T_i\to T\). Effective descent for parameterized spectra therefore gives
\[
    \Gamma_{T[N]}^{\mathrm{st}}(M_{T[N]})
    \simeq
    \Tot_i
    \Gamma_{T_i[N_i]}^{\mathrm{st}}(M_{T_i[N_i]})
\]
and
\[
    \Gamma_T^{\mathrm{st}}(M_T)
    \simeq
    \Tot_i
    \Gamma_{T_i}^{\mathrm{st}}(M_{T_i}).
\]
Fibres commute with limits in spectra, so
\[
    K_M^{(1)}(v)
    \simeq
    \Tot_iK_M^{(1)}(v_i).
\]
Thus \(K_M^{(1)}\) is a sheaf for the centre-pullback topology. Lemma~\ref{lem:var-cartesian-centre-kan}, applied to the Cartesian centre functor \(\pi_1\), now shows that
\[
    \operatorname{Ran}_{\pi_1^{\mathrm{op}}}K_M^{(1)}
\]
is already a spectral Nisnevich sheaf. Hence the sheafification unit is an equivalence.

Finally, at \(c:T\to\bar A\), the right-Kan value is
\[
    \lim_{(c^{\mathrm{op}}\to\pi_1^{\mathrm{op}}v)\in(c^{\mathrm{op}}\downarrow\pi_1^{\mathrm{op}})}
    K_M^{(1)}(v).
\]
Every object \(w\to c\) of this indexing category has its canonical limit projection, which gives the asserted evaluation morphism.
\end{proof}

\begin{remark}[No square-zero variation structure is supplied]
\label{rem:var-no-first-datum}
The object \(\Vfirst_M(A)\) is not defined by postulating a tangent bundle or a class of admissible variations. It is the explicit right Kan extension from the complete category of split structured square-zero diagrams, viewed through the equivalent sheafification of Theorem~\ref{thm:var-first-kan-sheaf}. The lift spaces, maps between lifts, automorphisms, and higher homotopies enter the indexing \(\infty\)-category and are not collapsed to connected components.
\end{remark}

\begin{proposition}[A split variation maps canonically to the intrinsic relation]
\label{prop:var-sqz-to-relation}
Every variation diagram
\[
    v=(T,N,a_T,\widetilde a)
\]
determines a canonical morphism
\[
    \boxed{
    \sigma_v:T[N]\longrightarrow R_1
    }
\]
over \(\bar A\times_Q\bar A\), whose two vertex components are
\[
    a_Tr_N:T[N]\longrightarrow\bar A,
    \qquad
    \widetilde a:T[N]\longrightarrow\bar A.
\]
The construction is functorial in \(v\) and compatible with arbitrary base change.
\end{proposition}

\begin{proof}
The two maps \(a_Tr_N\) and \(\widetilde a\) agree after restriction along the structured square-zero immersion \(i_N:T\to T[N]\). By Chapters 9--10, \(i_N\) is inverted by the relative de Rham localization over \(Q\). Hence the two maps have canonically equivalent composites
\[
    T[N]\longrightarrow\bar A_{\dR/Q}.
\]
Equivalently, apply Chapter 10's canonical factorization theorem for a de Rham-equivalent infinitesimal neighbourhood to the diagram
\[
    T\xrightarrow{i_N}T[N]
    \longrightarrow
    \bar A\times_Q\bar A,
\]
where the last map is \((a_Tr_N,\widetilde a)\). The target completion of the diagonal is canonically
\[
    R_1=\bar A\times_{\bar A_{\dR/Q}}\bar A.
\]
The factorization theorem gives \(\sigma_v\), functorially and compatibly with base change.
\end{proof}

\begin{construction}[Restriction of full one-variation to split square-zero tests]
\label{con:var-local-square-zero-restriction}
Fix an affine point \(c:T\to\bar A\). A section of \(\Vfull_M(A)\) over \(T\) is represented by a normalized coefficient-valued function on intrinsic one-simplices centred over \(T\). Let
\[
    v=(T',N',a_{T'},\widetilde a')
\]
be a variation diagram together with a map of centres \(T'\to T\) over \(\bar A\). Restrict first from \(T\) to \(T'\), and then along
\[
    \sigma_v:T'[N']\to R_1.
\]
Because the full one-cochain is normalized, its restriction along the central section \(T'\to T'[N']\) is canonically zero. Hence the restriction lands in
\[
    K_M^{(1)}(v).
\]
These maps form a cone over the right-Kan-extension diagram and therefore induce a natural morphism of presheaves
\[
    \Vfull_M(A)|_{\mathcal C^{\mathrm{aff}}_{/\bar A}}
    \longrightarrow
    \Vfirst_{M,\mathrm{pre}}(A).
\]
Using the equivalence
\[
    \Vfirst_{M,\mathrm{pre}}(A)\xrightarrow{\sim}\Vfirst_M(A)
\]
of Theorem~\ref{thm:var-first-kan-sheaf}, define
\[
    \boxed{
    \lambda^1_{\mathrm{loc}}:
    \Vfull_M(A)
    \longrightarrow
    \Vfirst_M(A).
    }
\]
\end{construction}

\begin{proposition}[Affine-test base change for split variations]
\label{prop:var-first-affine-base-change}
Let \(c:Q'\to Q\), put
\[
    \bar A':=\bar A\times_QQ',
    \qquad
    M':=c^*M,
\]
and let \(t:T\to\bar A'\) be an affine point. Write \(t_0:T\to\bar A\) for its composite. Then composition with \(\bar A'\to\bar A\) induces an equivalence between:
\begin{enumerate}
\item the comma \(\infty\)-category of split square-zero variation diagrams over \(\bar A'/Q'\) whose centres map to \(t\);
\item the comma \(\infty\)-category of split square-zero variation diagrams over \(\bar A/Q\) whose centres map to \(t_0\).
\end{enumerate}
Under this equivalence the first-difference coefficient functors are identified by stable Beck--Chevalley.
\end{proposition}

\begin{proof}
A variation diagram in the second category consists of a split thickening
\[
    T'\xrightarrow{i_N}T'[N]\xrightarrow{r_N}T'
\]
with a centre map \(T'\to T\to\bar A'\to\bar A\) and a map
\[
    \widetilde a:T'[N]\longrightarrow\bar A
\]
over \(Q\) whose restriction to \(T'\) is the centre. The centre map supplies
\[
    T'[N]\xrightarrow{r_N}T'\longrightarrow T\longrightarrow Q'.
\]
By the defining condition that \(\widetilde a\) is over \(Q\), its composite with \(\bar A\to Q\) is the composite of this map to \(Q'\) with \(Q'\to Q\). Hence the pullback universal property of
\[
    \bar A'=\bar A\times_QQ'
\]
gives a unique lift
\[
    T'[N]\longrightarrow\bar A'
\]
through a contractible space, and the lift restricts to the prescribed centre. The same construction on morphisms is inverse to forgetting from \(\bar A'\) to \(\bar A\). This proves the categorical equivalence.

For corresponding variation diagrams, the Cartesian coefficient square and right Beck--Chevalley identify the pullbacks of \(M\) and \(M'\) on \(T'[N]\) and \(T'\). Stable sections preserve the finite fibre defining \(K^{(1)}\), so the first-difference functors agree.
\end{proof}


\begin{lemma}[Affine-point evaluation of parameterized pullback]
\label{lem:var-affine-point-pullback}
Let \(f:Z'\to Z\) be a morphism in \(\XSp\), let \(E\in\mathcal E_Z\), and let \(t:T\to Z'\) be an affine point. Then there is a canonical equivalence
\[
    \boxed{
    (f^*E)(T,t)
    \simeq
    E(T,ft).
    }
\]
It is natural in \(f\), \(E\), and \(t\).
\end{lemma}

\begin{proof}
Under the affine-point presentation, evaluation is represented by the relative suspension spectrum. Hence
\[
\begin{aligned}
    (f^*E)(T,t)
    &\simeq
    \operatorname{map}_{\mathcal E_{Z'}}
    (\Sigma_t\one_T,f^*E)
    \\
    &\simeq
    \operatorname{map}_{\mathcal E_Z}
    (\Sigma_f\Sigma_t\one_T,E)
    \\
    &\simeq
    \operatorname{map}_{\mathcal E_Z}
    (\Sigma_{ft}\one_T,E)
    \\
    &\simeq
    E(T,ft).
\end{aligned}
\]
The comparison \(\Sigma_f\Sigma_t\simeq\Sigma_{ft}\) and the adjunction
\(\Sigma_f\dashv f^*\) carry their canonical composition coherences, giving the stated naturality.
\end{proof}

\begin{theorem}[Naturality and arbitrary base change of the one-stage coefficient]
\label{thm:var-first-local-naturality}
The morphism
\[
    \lambda^1_{\mathrm{loc}}:\Vfull_M(A)\longrightarrow\Vfirst_M(A)
\]
is natural in the Cartan descent object and in coefficient morphisms. For every \(c:Q'\to Q\), put \(\bar A':=\bar A\times_QQ'\), \(M':=c^*M\), and let \(A'\) be the pulled-back Cartan descent object of Convention~\ref{conv:var-arbitrary-cartan-pullback}. There is a canonical equivalence in \(\mathcal E_{\bar A'}\)
\[
    \boxed{
    c_{\bar A}^*\Vfirst_M(A)
    \simeq
    \Vfirst_{M'}(A'),
    }
\]
where \(c_{\bar A}:\bar A'\to\bar A\) is the projection. Under this equivalence the pullback of \(\lambda^1_{\mathrm{loc}}\) is the corresponding local restriction map over \(Q'\). No passage to \(\pi_0\) occurs.
\end{theorem}

\begin{proof}
Naturality in \(A\) and \(M\) follows directly from functoriality of the intrinsic relation, coefficient restriction, right Kan extension, and the canonical equivalence of Theorem~\ref{thm:var-first-kan-sheaf}.

For base change, let \(t:T\to\bar A'\) be an arbitrary affine point and write \(t_0:T\to\bar A\) for its composite. Lemma~\ref{lem:var-affine-point-pullback} identifies the value of the pulled-back coefficient with the value of the original coefficient at \(t_0\):
\[
    \bigl(c_{\bar A}^*\Vfirst_M(A)\bigr)(T,t)
    \simeq
    \Vfirst_M(A)(T,t_0).
\]
By Theorem~\ref{thm:var-first-kan-sheaf}, both sides may be computed by their raw right-Kan formulas. Proposition~\ref{prop:var-first-affine-base-change} identifies the complete comma \(\infty\)-categories computing
\[
    \Vfirst_M(A)(T,t_0)
    \qquad\text{and}\qquad
    \Vfirst_{M'}(A')(T,t)
\]
and identifies the coefficient diagrams under stable Beck--Chevalley. Their limits are therefore canonically equivalent. These equivalences are natural under affine restriction, and the affine-point presentation of \(\mathcal E_{\bar A'}\) detects equivalences, giving
\[
    c_{\bar A}^*\Vfirst_M(A)
    \simeq
    \Vfirst_{M'}(A').
\]
The construction of \(\lambda^1_{\mathrm{loc}}\) is the natural cone over these same right-Kan diagrams, so the pointwise equivalence identifies its pullback with the base-changed cone.
\end{proof}

\begin{definition}[First-variation secondary coefficient]
\label{def:var-first-secondary}
Define
\[
    \boxed{
    \Sec^{1,(1)}(A;M)
    :=
    \Pi_a\Vfirst_M(A)
    \in\mathcal E_Q.
    }
\]
Pushing forward Construction~\ref{con:var-local-square-zero-restriction} and using Proposition~\ref{prop:var-full-local-pushforward} gives
\[
    \boxed{
    \lambda^1_{\mathrm{var}}:
    \Sec^1(A;M)
    \longrightarrow
    \Sec^{1,(1)}(A;M).
    }
\]
\end{definition}

\section{Represented first-order comparison}
\label{sec:var-first-represented}

Assume in this section that
\[
    a:\bar A\longrightarrow Q
\]
is representable. Thus, for every represented test \(T\to Q\), the pullback
\[
    \bar A_T:=T\times_Q\bar A
\]
is represented by a spectral scheme over \(T\). No representability of \(Q\) itself is assumed.

\begin{lemma}[Affine relative split-lift classification]
\label{lem:var-affine-relative-lifts}
Let
\[
    c:T=\operatorname{Spec}C\longrightarrow\bar A
\]
be an affine point, write
\[
    b:=ac:T\longrightarrow Q,
    \qquad
    \bar A_c:=T\times_Q\bar A,
\]
and let
\[
    p_c:\bar A_c\longrightarrow T,
    \qquad
    s_c:T\longrightarrow\bar A_c
\]
be the projection and the section induced by \(c\). For every connective \(C\)-module \(N\), there is a canonical natural equivalence
\[
    \boxed{
    \hofib_{s_c}\left(
      \Map_T(T[N],\bar A_c)
      \longrightarrow
      \Map_T(T,\bar A_c)
    \right)
    \simeq
    \Omega^\infty
    \operatorname{map}_{\operatorname{Mod}_C}
    \left(
      \Gamma(T,s_c^*L_{\bar A_c/T}),
      N
    \right).
    }
\]
Equivalently, regarding \(T[N]\) as an object over \(Q\) through
\[
    T[N]\xrightarrow{r_N}T\xrightarrow{b}Q,
\]
the left side is canonically the space of maps
\[
    \widetilde c:T[N]\longrightarrow\bar A
\]
over \(Q\) equipped with a specified equivalence
\[
    \widetilde c\,i_N\simeq c.
\]
The equivalence is natural under affine restriction of \(c\) and under arbitrary base change of the representable morphism \(a\).
\end{lemma}

\begin{proof}
Because \(a\) is representable, \(\bar A_c\to T\) is a morphism of spectral schemes. The Chapter 2 global relative cotangent complex
\[
    L_{\bar A_c/T}\in\operatorname{QCoh}(\bar A_c)_{\ge0}
\]
represents relative derivations. Pulling the relative derivation universal property back along the section \(s_c\) identifies lifts of \(s_c\) across the split structured square-zero thickening
\[
    i_N:T\longrightarrow T[N]
\]
with the space
\[
    \Omega^\infty
    \operatorname{map}_{\operatorname{QCoh}(T)}
    (s_c^*L_{\bar A_c/T},\widetilde N),
\]
where \(\widetilde N\) is the quasi-coherent module corresponding to \(N\). Since \(T=\operatorname{Spec}C\) is affine,
\[
    \operatorname{QCoh}(T)\simeq\operatorname{Mod}_C
\]
and
\[
    \Gamma(T,s_c^*L_{\bar A_c/T})
\]
is the corresponding connective \(C\)-module, giving the displayed equivalence.

The pullback universal property of
\[
    \bar A_c=T\times_Q\bar A
\]
identifies a map \(T[N]\to\bar A_c\) over \(T\) with a map \(T[N]\to\bar A\) over \(Q\) whose structural map to \(T\) is \(r_N\); the condition on the central section is preserved under this identification. Naturality under affine restriction follows from global cotangent base change and functoriality of the relative derivation correspondence. The same base-change theorem applies after arbitrary pullback of the representable morphism \(a\).
\end{proof}

\begin{construction}[Universal affine first variation]
\label{con:var-universal-first-variation}
For an affine point
\[
    c:T=\operatorname{Spec}C\longrightarrow\bar A,
\]
put
\[
    \boxed{
    N_c
    :=
    \Gamma(T,s_c^*L_{\bar A_c/T}).
    }
\]
Under Lemma~\ref{lem:var-affine-relative-lifts}, the identity
\[
    \id_{N_c}:N_c\longrightarrow N_c
\]
determines a canonical lift
\[
    \widetilde s_c^{\mathrm{univ}}:
    T[N_c]\longrightarrow\bar A_c
\]
over \(T\). Composing with \(\bar A_c\to\bar A\) gives
\[
    \widetilde c_{\mathrm{univ}}:
    T[N_c]\longrightarrow\bar A.
\]
Together with the canonical central section and split retraction this defines
\[
    \boxed{
    u_c:=(T,N_c,c,\widetilde c_{\mathrm{univ}})
    \in
    \mathcal V^{(1)}_{\bar A/Q}.
    }
\]
We call \(u_c\) the \emph{universal affine first variation at \(c\)}.
\end{construction}

\begin{proposition}[Cartesian functoriality of the universal affine variations]
\label{prop:var-universal-first-cartesian}
Let
\[
    g:(T',c')\longrightarrow(T,c)
\]
be a morphism of affine points over \(\bar A\), so \(c'=cg\). Then cotangent base change gives a canonical equivalence
\[
    \boxed{
    N_{c'}\simeq g^*N_c.
    }
\]
Under this equivalence the Cartesian pullback of the complete variation diagram \(u_c\) along \(g\) is canonically equivalent to \(u_{c'}\). Hence
\[
    c\longmapsto u_c
\]
defines a Cartesian section of
\[
    \pi_1:
    \mathcal V^{(1)}_{\bar A/Q}
    \longrightarrow
    \mathcal C^{\mathrm{aff}}_{/\bar A}.
\]
The section is compatible with arbitrary base change of \(a\).
\end{proposition}

\begin{proof}
There is a Cartesian square of spectral schemes
\[
\begin{tikzcd}
\bar A_{c'} \arrow[r] \arrow[d,"p_{c'}"']
& \bar A_c \arrow[d,"p_c"]
\\
T' \arrow[r,"g"']
& T .
\end{tikzcd}
\]
Global cotangent base change gives
\[
    L_{\bar A_{c'}/T'}
    \simeq
    g_{\bar A}^*L_{\bar A_c/T}.
\]
Pulling back along the canonical sections and using
\[
    g_{\bar A}s_{c'}=s_cg
\]
gives
\[
    s_{c'}^*L_{\bar A_{c'}/T'}
    \simeq
    g^*s_c^*L_{\bar A_c/T}.
\]
Since both centres are affine, taking the corresponding modules yields
\[
    N_{c'}\simeq g^*N_c.
\]
The identity class \(\id_{N_c}\) pulls back to the identity class \(\id_{N_{c'}}\). Naturality of Lemma~\ref{lem:var-affine-relative-lifts} therefore identifies the pulled-back universal lift with the lift defining \(u_{c'}\), through a contractible space of compatible choices. The arbitrary-base-change statement is the same calculation after replacing \(a\) by any of its pullbacks.
\end{proof}

\begin{definition}[Coefficient-valued stabilized first-order coefficient]
\label{def:var-coeff-cotangent}
For \(M\in\mathcal E_Q\), define a spectrum-valued presheaf on affine points of \(\bar A\) by
\[
    \boxed{
    L^{\mathrm{st}}_{\bar A/Q,\mathrm{pre}}(M)(T,c)
    :=
    K_M^{(1)}(u_c).
    }
\]
Define its spectral Nisnevich sheafification by
\[
    \boxed{
    L^{\mathrm{st}}_{\bar A/Q}(M)
    :=
    a_{\mathrm{Nis}}^{\operatorname{Sp}}
    L^{\mathrm{st}}_{\bar A/Q,\mathrm{pre}}(M)
    \in
    \mathcal E_{\bar A}.
    }
\]
For scalar coefficients write
\[
    L^{\mathrm{st}}_{\bar A/Q}
    :=
    L^{\mathrm{st}}_{\bar A/Q}(\Oadd_Q).
\]
\end{definition}

\begin{construction}[Cotangent restriction of full one-cochains]
\label{con:var-coeff-cotangent-restriction}
For every affine point \(c:T\to\bar A\), Proposition~\ref{prop:var-sqz-to-relation} applied to \(u_c\) gives
\[
    \sigma_{u_c}:T[N_c]\longrightarrow R_{\bar A/Q,1}.
\]
Restriction of a normalized full one-cochain along \(\sigma_{u_c}\) is canonically zero on the central section and therefore defines a natural morphism
\[
    \Vfull_M(A)(T,c)
    \longrightarrow
    K_M^{(1)}(u_c)
    =
    L^{\mathrm{st}}_{\bar A/Q,\mathrm{pre}}(M)(T,c).
\]
Proposition~\ref{prop:var-universal-first-cartesian} makes these maps natural under affine restriction. Sheafification therefore gives
\[
    \boxed{
    \lambda^1_{\mathrm{cot},M}:
    \Vfull_M(A)
    \longrightarrow
    L^{\mathrm{st}}_{\bar A/Q}(M).
    }
\]
Using Proposition~\ref{prop:var-full-local-pushforward}, define
\[
    \boxed{
    \lambda^1_{\bar A/Q,M}:
    \Sec^1(A;M)
    \longrightarrow
    \Pi_aL^{\mathrm{st}}_{\bar A/Q}(M)
    }
\]
to be \(\Pi_a\lambda^1_{\mathrm{cot},M}\) under
\[
    \Sec^1(A;M)\simeq\Pi_a\Vfull_M(A).
\]
For scalar coefficients write
\[
    \lambda^1_{\bar A/Q}:=
    \lambda^1_{\bar A/Q,\Oadd_Q}.
\]
\end{construction}

\begin{theorem}[Universal factorization through the affine universal variation]
\label{thm:var-sqz-factor-first-neighbourhood}
Let
\[
    v=(T'=\operatorname{Spec}C',N,c',\widetilde a)
\]
be a split square-zero variation diagram and let
\[
    g:T'\longrightarrow T
\]
be a morphism of centres over \(\bar A\) from \(c'\) to \(c\). There is a canonically contractible space of morphisms of complete split variation diagrams
\[
    \boxed{
    v\longrightarrow u_c
    }
\]
over \(g\). Equivalently, the variation \(v\) determines a canonical relative derivation
\[
    \delta_v:
    N_{c'}
    \longrightarrow
    N
\]
and, using
\[
    N_{c'}\simeq g^*N_c,
\]
the corresponding structured map of split thickenings
\[
    \widehat g_{\delta_v}:
    T'[N]\longrightarrow T[N_c].
\]
The intrinsic relation maps satisfy a canonical factorization
\[
\begin{tikzcd}[column sep=huge]
T'[N]
  \arrow[rr,"\sigma_v"]
  \arrow[dr,"\widehat g_{\delta_v}"']
&& R_{\bar A/Q,1}
\\
& T[N_c]
  \arrow[ur,"\sigma_{u_c}"'] &
\end{tikzcd}
\]
over \(\bar A\times_Q\bar A\). The construction is functorial in variation diagrams and compatible with arbitrary base change.
\end{theorem}

\begin{proof}
Because the structural map of \(v\) is
\[
    T'[N]\xrightarrow{r_N}T'\xrightarrow{ac'}Q,
\]
the pullback universal property identifies \(\widetilde a\) with a map
\[
    T'[N]\longrightarrow\bar A_{c'}
\]
over \(T'\) extending the canonical section \(s_{c'}\). Lemma~\ref{lem:var-affine-relative-lifts} therefore associates to \(v\) a relative derivation
\[
    \delta_v:N_{c'}\longrightarrow N.
\]
Proposition~\ref{prop:var-universal-first-cartesian} identifies
\[
    N_{c'}\simeq g^*N_c.
\]

By Definition~\ref{def:var-sqz-category}, a morphism of variations
\[
    v\longrightarrow u_c
\]
over \(g\) is equivalently a module map
\[
    \alpha:g^*N_c\longrightarrow N
\]
together with a specified path between \(\widetilde a\) and the lift obtained from the universal variation \(u_c\) by functoriality of split square-zero algebras along \(\alpha\). Under Lemma~\ref{lem:var-affine-relative-lifts}, this mapping space is
\[
    \left\{
      (\alpha,\gamma)
      \ \middle|\
      \alpha\in
      \Omega^\infty
      \operatorname{map}_{\operatorname{Mod}_{C'}}
      (N_{c'},N),
      \quad
      \gamma:\delta_v\simeq\alpha
    \right\}.
\]
It is the path space with fixed initial point \(\delta_v\), hence is contractible. Taking \(\alpha=\delta_v\) gives the displayed structured map \(\widehat g_{\delta_v}\).

The map \(\sigma_v\) and the composite
\[
    \sigma_{u_c}\widehat g_{\delta_v}
\]
have the same two vertex components
\[
    c'r_N,\widetilde a:T'[N]\rightrightarrows\bar A
\]
through the specified path contained in the morphism of variations. Functoriality and contractible uniqueness of the intrinsic relation factorization of Proposition~\ref{prop:var-sqz-to-relation} therefore give
\[
    \sigma_v
    \simeq
    \sigma_{u_c}\widehat g_{\delta_v}.
\]
The construction uses only cotangent base change, functoriality of split structured square-zero extensions, and the intrinsic relation factorization, so it is functorial and compatible with arbitrary base change.
\end{proof}

\begin{proposition}[Initial object in the right-Kan indexing category]
\label{prop:var-universal-first-initial}
For an affine point \(c:T\to\bar A\), let
\[
    \mathcal J_c
    :=
    (c^{\mathrm{op}}\downarrow\pi_1^{\mathrm{op}})
\]
be the comma \(\infty\)-category computing
\(\operatorname{Ran}_{\pi_1^{\mathrm{op}}}K_M^{(1)}\) at \(c\). The universal variation \(u_c\) of Construction~\ref{con:var-universal-first-variation} determines an initial object of \(\mathcal J_c\).
\end{proposition}

\begin{proof}
An object of \(\mathcal J_c\) is a variation \(v\) together with a morphism from its centre to \(c\). A morphism from the object represented by \(u_c\) to that object in the comma category is, after reversing the \(\mathrm{op}\) in \(\pi_1^{\mathrm{op}}\), precisely a morphism of structured variation diagrams
\[
    v\longrightarrow u_c
\]
over the given centre map. Theorem~\ref{thm:var-sqz-factor-first-neighbourhood} identifies this mapping space with a contractible space. This is the defining initiality condition.
\end{proof}

\begin{theorem}[Represented identification of the one-stage coefficient]
\label{thm:var-first-represented-comparison}
If \(a:\bar A\to Q\) is representable, the presheaf
\[
    L^{\mathrm{st}}_{\bar A/Q,\mathrm{pre}}(M)
\]
is already a spectral Nisnevich sheaf, and there is a canonical equivalence
\[
    \boxed{
    \theta_M:
    L^{\mathrm{st}}_{\bar A/Q}(M)
    \xrightarrow{\sim}
    \Vfirst_M(A).
    }
\]
Under this equivalence,
\[
    \boxed{
    \lambda^1_{\mathrm{loc}}
    \simeq
    \theta_M\lambda^1_{\mathrm{cot},M}.
    }
\]
Consequently
\[
    \boxed{
    \lambda^1_{\mathrm{var}}
    \simeq
    \Pi_a(\theta_M)\lambda^1_{\bar A/Q,M}.
    }
\]
All equivalences are natural under affine restriction, coefficient morphisms, and arbitrary base change.
\end{theorem}

\begin{proof}
Evaluate the raw right Kan extension at an affine point \(c:T\to\bar A\). Proposition~\ref{prop:var-universal-first-initial} gives
\[
    \Vfirst_{M,\mathrm{pre}}(A)(T,c)
    \simeq
    K_M^{(1)}(u_c)
    =
    L^{\mathrm{st}}_{\bar A/Q,\mathrm{pre}}(M)(T,c).
\]
Proposition~\ref{prop:var-universal-first-cartesian} makes these equivalences natural under affine restriction. Hence there is a canonical equivalence of presheaves
\[
    \Vfirst_{M,\mathrm{pre}}(A)
    \simeq
    L^{\mathrm{st}}_{\bar A/Q,\mathrm{pre}}(M).
\]
Theorem~\ref{thm:var-first-kan-sheaf} proves that the left side is already a spectral Nisnevich sheaf. Therefore the right side is also a sheaf, its sheafification unit is an equivalence, and after the defining sheafifications the displayed equivalence is \(\theta_M\).

For an arbitrary variation \(v\) over an affine centre mapping to \(c\), Theorem~\ref{thm:var-sqz-factor-first-neighbourhood} gives
\[
    \sigma_v
    \simeq
    \sigma_{u_c}\widehat g_{\delta_v}.
\]
Thus restriction of a normalized full one-cochain along \(\sigma_v\) is restriction first along \(\sigma_{u_c}\) and then along the structured map classified by \(\delta_v\). The cone defining \(\lambda^1_{\mathrm{loc}}\) is therefore the cone obtained from \(\lambda^1_{\mathrm{cot},M}\) through the initial-object identification. The right-Kan universal property gives
\[
    \lambda^1_{\mathrm{loc}}
    \simeq
    \theta_M\lambda^1_{\mathrm{cot},M}.
\]
Applying \(\Pi_a\) and Proposition~\ref{prop:var-full-local-pushforward} gives the final displayed identity. Arbitrary base-change compatibility also follows from Theorem~\ref{thm:var-first-local-naturality}, together with the affine-point construction above.
\end{proof}

\begin{proposition}[Scalar affine-point cotangent identification]
\label{prop:var-scalar-affine-cotangent}
For scalar coefficients,
\[
    \boxed{
    L^{\mathrm{st}}_{\bar A/Q}(T,c)
    \simeq
    N_c
    =
    \Gamma(T,s_c^*L_{\bar A_c/T})
    }
\]
for every affine point \(c:T=\operatorname{Spec}C\to\bar A\). For every represented affine base test
\[
    t:T\longrightarrow Q,
\]
the pullback of \(L^{\mathrm{st}}_{\bar A/Q}\) to
\[
    \bar A_T:=T\times_Q\bar A
\]
is canonically the Chapter 10 stabilized cotangent coefficient
\[
    L^{\mathrm{st}}_{\bar A_T/T},
\]
and the pullback of \(\lambda^1_{\bar A/Q}\) is the Chapter 10 first-order linearization for \(\bar A_T/T\).
\end{proposition}

\begin{proof}
For \(M=\Oadd_Q\), the scalar base-change equivalence gives
\[
    M_{T[N_c]}\simeq\Oadd_{T[N_c]},
    \qquad
    M_T\simeq\Oadd_T.
\]
Chapter 10's represented scalar-section calculation identifies
\[
    \Gamma_{T[N_c]}^{\mathrm{st}}(\Oadd_{T[N_c]})
    \simeq
    C\oplus N_c
\]
and
\[
    \Gamma_T^{\mathrm{st}}(\Oadd_T)\simeq C.
\]
Therefore
\[
    K_{\Oadd_Q}^{(1)}(u_c)
    \simeq
    \fib(C\oplus N_c\to C)
    \simeq
    N_c.
\]
Theorem~\ref{thm:var-first-represented-comparison} identifies this raw affine-point value with the scalar sheaf \(L^{\mathrm{st}}_{\bar A/Q}\).

Now let \(t:T\to Q\) be represented and affine. The morphism
\[
    \bar A_T\longrightarrow T
\]
is represented by spectral schemes. At every affine point \(c:U\to\bar A_T\), the local construction above is the Chapter 10 first structured diagonal construction for the represented morphism after base change to \(U\), because both are classified by the identity universal derivation of
\[
    \Gamma(U,c^*L_{\bar A_T/T}).
\]
Chapter 10 identifies its stabilized coefficient with that cotangent module and identifies restriction of normalized one-cochains with the first-order cotangent linearization. Affine-point recognition gives the two asserted equivalences on \(\bar A_T\).
\end{proof}

\begin{proposition}[Recovery of the Chapter 10 first diagonal coefficient]
\label{prop:var-universal-first-diagonal-comparison}
Assume in addition that \(Q\) is represented. Then \(\bar A\to Q\) is a morphism of spectral schemes and Chapter 10 constructs
\[
    \bar A
    \xrightarrow{i_1}
    D^{(1)}_{\bar A/Q}
    \xrightarrow{\nu_1}
    \bar A\times_Q\bar A
\]
with first projection
\[
    h_1:D^{(1)}_{\bar A/Q}\longrightarrow\bar A.
\]
For every \(M\in\mathcal E_Q\), there is a canonical equivalence
\[
    \boxed{
    L^{\mathrm{st}}_{\bar A/Q}(M)
    \simeq
    \fib\left(
      \Pi_{h_1}h_1^*a^*M
      \longrightarrow
      a^*M
    \right).
    }
\]
For every affine point \(c:T\to\bar A\), the pullback
\[
    T\times_{\bar A}D^{(1)}_{\bar A/Q}
\]
with its second vertex map to \(\bar A\) is canonically equivalent as a complete split variation diagram to \(u_c\).
\end{proposition}

\begin{proof}
Chapter 10 constructs \(D^{(1)}_{\bar A/Q}\) from the canonical first diagonal class. After base change along \(c:T\to\bar A\), the coefficient is
\[
    \Gamma(T,c^*L_{\bar A/Q})
    \simeq
    \Gamma(T,s_c^*L_{\bar A_c/T})
    =
    N_c
\]
by cotangent base change. Chapter 10's affine calculation proves that the second vertex of this pulled-back structured neighbourhood represents the identity universal derivation of \(N_c\). Construction~\ref{con:var-universal-first-variation} is classified by the same point of the derivation space, and the structured square-zero classification gives a canonical equivalence through a contractible space.

Right Beck--Chevalley and exactness of stable sections then identify the affine-point value of the displayed global fibre with
\[
    K_M^{(1)}(u_c)
    =
    L^{\mathrm{st}}_{\bar A/Q}(M)(T,c).
\]
These equivalences are natural under affine restriction, so affine-point recognition gives the asserted global equivalence.
\end{proof}

\begin{remark}[No tensor formula for general coefficients]
\label{rem:var-no-coeff-tensor-cotangent}
The equivalence of Theorem~\ref{thm:var-first-represented-comparison} does not assert
\[
    L^{\mathrm{st}}_{\bar A/Q}(M)
    \simeq
    L^{\mathrm{st}}_{\bar A/Q}\otimes a^*M.
\]
The general coefficient-valued first-order object is the spectral Nisnevich sheaf represented affine-pointwise by
\[
    (T,c)\longmapsto K_M^{(1)}(u_c).
\]
When \(Q\) is represented, Proposition~\ref{prop:var-universal-first-diagonal-comparison} recovers the corresponding dependent-product fibre formula. The existing dependent-product formalism makes \(\Pi\) lax symmetric monoidal but does not supply the tensor identification required for the displayed tensor formula.
\end{remark}

\section{First variation, Euler sections, and derived critical loci}
\label{sec:var-euler-lagrange}

\begin{remark}[Why the descended vertical differential is variational]
\label{rem:var-secondary-euler-meaning}
Before horizontal totalization, \(\partial_V\) is the coherent vertical infinitesimal boundary of the Cartan--vertical double relation. Horizontal totalization then performs effective Cartan descent through the complete \v{C}ech direction. Thus
\[
    \delta^{\infty}_{\mathrm{var}}
\]
is the vertical infinitesimal differential on already horizontally descended secondary classes. This is the intrinsic passage from raw vertical variation to a secondary variational class modulo complete horizontal descent. No integration-by-parts splitting or source-form projection is used in this definition.  Section~\ref{sec:var-lepage-presymplectic} later constructs the derived moduli of relative integration-by-parts realizations without altering the Euler class. The first square-zero image below is therefore the brave-new Euler--Lagrange class; Theorem~\ref{thm:var-euler-represents-first-variation} proves its universal first-variation property, while Example~\ref{ex:var-classical-empty-equation} records that the analogous classical secondary differential is the ordinary Euler--Lagrange operator in the empty-equation variational bicomplex.
\end{remark}

\begin{definition}[First variational operator]
\label{def:var-first-operator}
Define
\[
    \boxed{
    \delta^{(1)}_{\mathrm{var}}
    :=
    \lambda^1_{\mathrm{var}}\delta^{\infty}_{\mathrm{var}}:
    \Sec^0(A;M)
    \longrightarrow
    \Sec^{1,(1)}(A;M).
    }
\]
\end{definition}

\subsection{Universal variational moduli and the Euler map}
\label{sec:var-universal-variational-moduli}

\begin{definition}[Universal Lagrangian, closed-variation, and Euler-source objects]
\label{def:var-universal-moduli-diagram}
For \(n\in\mathbb Z\), define objects over \(Q\)
\[
    \boxed{
    \mathbf{Lag}^n_{A/Q}(M)
    :=
    \Omega_Q^\infty\bigl(\Sec^0(A;M)[n]\bigr),
    }
\]
\[
    \boxed{
    \mathbf{ClVar}^n_{A/Q}(M)
    :=
    \Omega_Q^\infty\bigl(\Sec^{1,\mathrm{cl}}(A;M;n)\bigr),
    }
\]
and
\[
    \boxed{
    \mathbf{Eul}^n_{A/Q}(M)
    :=
    \Omega_Q^\infty\bigl(\Sec^{1,(1)}(A;M)[n]\bigr).
    }
\]
A degree-\(n\) Lagrangian class, a coherently closed full first variation, and a first Euler source are respectively sections of these three objects over \(Q\).
\end{definition}

\begin{theorem}[Universal variational moduli diagram]
\label{thm:var-universal-moduli-diagram}
The exact classifier and first-order restriction induce a canonical diagram in the slice \((\XSp)_{/Q}\)
\[
    \boxed{
    \mathbf{Lag}^n_{A/Q}(M)
    \xrightarrow{\;\mathbf d^{\mathrm{cl}}_{\mathrm{var}}\;}
    \mathbf{ClVar}^n_{A/Q}(M)
    \xrightarrow{\;\mathbf\lambda^{1,\mathrm{cl}}_{\mathrm{var}}\;}
    \mathbf{Eul}^n_{A/Q}(M).
    }
\]
Here
\[
    \mathbf\lambda^{1,\mathrm{cl}}_{\mathrm{var}}
    :=
    \Omega_Q^\infty\bigl((\lambda^1_{\mathrm{var}}[n])\circ u_1\bigr)
\]
is induced by the stable composite
\[
    \Sec^{1,\mathrm{cl}}(A;M;n)
    \xrightarrow{u_1}
    \Sec^1(A;M)[n]
    \xrightarrow{\lambda^1_{\mathrm{var}}[n]}
    \Sec^{1,(1)}(A;M)[n].
\]
Its composite
\[
    \boxed{
    \mathbf{EL}_{A/Q,M}^n:
    \mathbf{Lag}^n_{A/Q}(M)
    \longrightarrow
    \mathbf{Eul}^n_{A/Q}(M)
    }
\]
is the universal first Euler map. On a section \(L:Q\to\mathbf{Lag}^n_{A/Q}(M)\), it evaluates to the Euler--Lagrange class \(\EL(L)\).

The diagram is natural in Cartan morphisms and coefficient morphisms and commutes with arbitrary pullback in \(Q\).
\end{theorem}

\begin{proof}
Apply \(\Omega_Q^\infty\) to the stable maps
\[
    \Sec^0(A;M)[n]
    \xrightarrow{\delta^{\mathrm{cl}}_{\mathrm{var}}}
    \Sec^{1,\mathrm{cl}}(A;M;n)
    \xrightarrow{\lambda^1_{\mathrm{var}}u_1}
    \Sec^{1,(1)}(A;M)[n].
\]
Their composite is \(\delta^{(1)}_{\mathrm{var}}\) by definition. Naturality and arbitrary base change follow from the corresponding properties of the stable coefficient maps and from pullback compatibility of the relative zeroth-space functor.
\end{proof}

\begin{corollary}[Inverse-problem objects are derived fibres]
\label{cor:var-inverse-problem-fibres-universal}
For a first Euler source \(e:Q\to\mathbf{Eul}^n_{A/Q}(M)\), the parameterized Helmholtz-lift object is the homotopy fibre of
\[
    \mathbf{ClVar}^n_{A/Q}(M)
    \longrightarrow
    \mathbf{Eul}^n_{A/Q}(M)
\]
over \(e\), while the parameterized first-order Lagrangian-realization object is the homotopy fibre of
\[
    \mathbf{EL}_{A/Q,M}^n:
    \mathbf{Lag}^n_{A/Q}(M)
    \longrightarrow
    \mathbf{Eul}^n_{A/Q}(M)
\]
over \(e\). The map from Lagrangian realizations to Helmholtz lifts is induced by the first arrow of the universal diagram, and its fibre over a chosen coherent lift is the parameterized primitive object.
\end{corollary}

\begin{proof}
This is homotopy-fibre pasting for Theorem~\ref{thm:var-universal-moduli-diagram}. It is the internal version of the inverse-problem fibration proved below.
\end{proof}

\begin{definition}[Euler--Lagrange class and local Euler section]
\label{def:var-euler-class}
For a Lagrangian class
\[
    L:\one_Q\to\Sec^0(A;M)[n],
\]
define
\[
    \boxed{
    \EL(L)
    :=
    \delta^{(1)}_{\mathrm{var}}L:
    \one_Q\to\Sec^{1,(1)}(A;M)[n].
    }
\]
Under the equivalence
\[
    \Sec^{1,(1)}(A;M)=\Pi_a\Vfirst_M(A),
\]
the adjoint local section is
\[
    \boxed{
    \mathsf{EL}_L:
    \one_{\bar A}\longrightarrow\Vfirst_M(A)[n].
    }
\]
It is canonically the composite
\[
    \mathsf E_L^{\infty}
    \xrightarrow{\lambda^1_{\mathrm{loc}}}
    \Vfirst_M(A)[n].
\]
\end{definition}

\begin{proposition}[Full versus first variation]
\label{prop:var-full-vs-first}
The diagram
\[
\begin{tikzcd}[column sep=huge]
\one_{\bar A}
  \arrow[r,"\mathsf E_L^{\infty}"]
& \Vfull_M(A)[n]
  \arrow[r,"\lambda^1_{\mathrm{loc}}"]
& \Vfirst_M(A)[n]
\end{tikzcd}
\]
has composite \(\mathsf{EL}_L\). Under the adjunction \(a^*\dashv\Pi_a\), this local factorization is adjoint to the global factorization
\[
\begin{tikzcd}[column sep=huge]
\one_Q
  \arrow[r,"\delta^{\infty}_{\mathrm{var}}L"]
& \Sec^1(A;M)[n]
  \arrow[r,"\lambda^1_{\mathrm{var}}"]
& \Sec^{1,(1)}(A;M)[n].
\end{tikzcd}
\]
No equivalence between the two variation coefficients is asserted. There exists a represented scalar example, namely the Chapter 10 affine-line example, in which the induced first-order map on \(\pi_0\) has a nonzero kernel element.
\end{proposition}

\begin{proof}
The local statement is Definition~\ref{def:var-euler-class}. For a local section
\[
    s:\one_{\bar A}\simeq a^*\one_Q\longrightarrow V,
\]
its adjoint is the composite
\[
    \one_Q
    \xrightarrow{\eta_{\one_Q}}
    \Pi_a a^*\one_Q
    \simeq
    \Pi_a\one_{\bar A}
    \xrightarrow{\Pi_as}
    \Pi_aV,
\]
where \(\eta\) is the unit of \(a^*\dashv\Pi_a\). Naturality of the adjunction identifies the adjoint of \(\lambda^1_{\mathrm{loc}}s\) with \(\Pi_a(\lambda^1_{\mathrm{loc}})\) composed with the adjoint of \(s\). Taking \(s=\mathsf E_L^\infty\), using Proposition~\ref{prop:var-full-local-pushforward} and Definition~\ref{def:var-first-secondary}, gives exactly the displayed global factorization. The final assertion is the represented comparison of Theorem~\ref{thm:var-first-represented-comparison} together with the affine-line calculation of Chapter 10.
\end{proof}

\begin{definition}[First-variation moduli object]
\label{def:var-first-moduli}
Define the pointed object over \(\bar A\)
\[
    \boxed{
    \mathbf E_L^{(1)}
    :=
    \Omega_{\bar A}^\infty\bigl(\Vfirst_M(A)[n]\bigr).
    }
\]
The local Euler section and the zero map induce sections
\[
    e_L^{(1)},0:\bar A\rightrightarrows\mathbf E_L^{(1)}.
\]
\end{definition}

\begin{definition}[Derived Euler--Lagrange critical locus]
\label{def:var-critical-locus}
Define
\[
    \boxed{
    \Crit(L)
    :=
    \Eq(e_L^{(1)},0)
    }
\]
in \((\XSp)_{/\bar A}\). Equivalently, \(\Crit(L)\) is the homotopy zero fibre of the local Euler section.
\end{definition}

\begin{proposition}[Full criticality maps to first criticality]
\label{prop:var-critical-comparison}
The local restriction \(\lambda^1_{\mathrm{loc}}\) induces a canonical morphism over \(\bar A\)
\[
    \boxed{
    \Crit^{\infty}(L)
    \longrightarrow
    \Crit(L).
    }
\]
\end{proposition}

\begin{proof}
Apply \(\Omega_{\bar A}^\infty\) to \(\lambda^1_{\mathrm{loc}}\). It carries the full Euler section to the first Euler section and zero to zero. Functoriality of homotopy pullbacks gives the comparison of zero fibres.
\end{proof}

\begin{theorem}[Euler criticality is intrinsically formally integrable]
\label{thm:var-critical-intrinsic-pde}
The structure morphism
\[
    \Crit(L)\longrightarrow\bar A\longrightarrow Q
\]
makes \(\Crit(L)\) an intrinsic formally integrable PDE:
\[
    \boxed{
    \Crit(L)\in(\XSp)_{/Q}=\PDE^{\mathrm{fi}}_{X/S}.
    }
\]
The same statement holds for \(\Crit^{\infty}(L)\). Their jet-coalgebra structures, formal-disk actions, Cartan transport, and all higher descent coherences are obtained by applying the Chapter 11 effective-descent equivalence; no additional integrability datum is supplied.
\end{theorem}

\begin{proof}
Both critical loci are objects over \(\bar A\), and \(\bar A\) is an object over \(Q\). Therefore their composites to \(Q\) are objects of the slice \((\XSp)_{/Q}\). Chapter 11 defines \(\PDE^{\mathrm{fi}}_{X/S}\) to be this slice and constructs its equivalent jet-coalgebra and formal-disk algebra descriptions by effective descent.
\end{proof}

\subsection{The universal derived Euler critical family}
\label{subsec:var-universal-critical-family}

\begin{construction}[Universal local Euler section]
\label{con:var-universal-euler-section}
Let
\[
    P:=\mathbf{Lag}^n_{A/Q}(M),
    \qquad
    p:P\to Q,
\]
and form
\[
    \bar A_P:=\bar A\times_QP.
\]
Let \(A_P\) denote the pulled-back Cartan descent object of Convention~\ref{conv:var-arbitrary-cartan-pullback} for \(p:P\to Q\); equivalently, it is the descent coalgebra for the pulled-back effective epimorphism \(X\times_QP\twoheadrightarrow P\) corresponding to \(\bar A_P\to P\). By base change,
\[
    p^*\Sec^0(A;M)
    \simeq
    \Sec^0(A_P;p^*M).
\]
The identity map
\[
    \id_P:P\longrightarrow
    P=\Omega_Q^\infty\bigl(\Sec^0(A;M)[n]\bigr)
\]
is, by the parameterized representing property after pullback to \(P\), adjoint to a tautological Lagrangian
\[
    L_{\mathrm{univ}}:
    \one_P\longrightarrow
    \Sec^0(A_P;p^*M)[n].
\]
Applying the first variational operator and adjunction along \(\bar A_P\to P\) gives the universal local Euler section
\[
    \boxed{
    \mathsf{EL}_{\mathrm{univ}}:
    \one_{\bar A_P}
    \longrightarrow
    \Vfirst_{p^*M}(A_P)[n].
    }
\]
\end{construction}

\begin{definition}[Universal derived Euler critical family]
\label{def:var-universal-critical-family}
Put
\[
    \mathbf E_{\mathrm{univ}}^{(1)}
    :=
    \Omega_{\bar A_P}^\infty
    \bigl(\Vfirst_{p^*M}(A_P)[n]\bigr)
    \longrightarrow
    \bar A_P.
\]
The universal local Euler section \(\mathsf{EL}_{\mathrm{univ}}\) and the zero map determine two sections
\[
    e_{\mathrm{univ}}^{(1)},0:
    \bar A_P
    \rightrightarrows
    \mathbf E_{\mathrm{univ}}^{(1)}.
\]
Define
\[
    \boxed{
    \mathbf{Crit}^{(1)}_{A/Q}(M;n)
    :=
    \Eq(e_{\mathrm{univ}}^{(1)},0)
    \longrightarrow
    \bar A_P
    =
    \bar A\times_Q\mathbf{Lag}^n_{A/Q}(M).
    }
\]
This is the universal first Euler critical family.

Likewise, applying the full variational operator to the tautological Lagrangian and adjunction along \(\bar A_P\to P\) gives the universal full Euler section
\[
    \mathsf E_{\mathrm{univ}}^{\infty}:
    \one_{\bar A_P}
    \longrightarrow
    \Vfull_{p^*M}(A_P)[n].
\]
Put
\[
    \mathbf E_{\mathrm{univ}}^{\infty}
    :=
    \Omega_{\bar A_P}^\infty
    \bigl(\Vfull_{p^*M}(A_P)[n]\bigr)
    \longrightarrow
    \bar A_P,
\]
and let
\[
    e_{\mathrm{univ}}^{\infty},0:
    \bar A_P
    \rightrightarrows
    \mathbf E_{\mathrm{univ}}^{\infty}
\]
be the induced sections. Define
\[
    \boxed{
    \mathbf{Crit}^{\infty}_{A/Q}(M;n)
    :=
    \Eq(e_{\mathrm{univ}}^{\infty},0)
    \longrightarrow
    \bar A_P.
    }
\]
\end{definition}

\begin{theorem}[Individual critical loci are pullbacks of the universal family]
\label{thm:var-critical-family-pullback}
For every Lagrangian class
\[
    L:Q\longrightarrow\mathbf{Lag}^n_{A/Q}(M),
\]
there are canonical Cartesian squares
\[
\begin{tikzcd}
\Crit(L) \arrow[r] \arrow[d]
& \mathbf{Crit}^{(1)}_{A/Q}(M;n) \arrow[d]
\\
\bar A \arrow[r, "{(\operatorname{id}_{\bar A},\,L\circ a)}"']
& \bar A\times_Q\mathbf{Lag}^n_{A/Q}(M)
\end{tikzcd}
\]
and similarly for \(\Crit^{\infty}(L)\). Consequently the assignment \(L\mapsto\Crit(L)\) is obtained by base change from one derived object over the higher moduli of Lagrangians.
\end{theorem}

\begin{proof}
The tautological Lagrangian on \(P\) pulls back along \(L:Q\to P\) to the given Lagrangian. Arbitrary base change of \(\Vfirst\), of the first variational operator, and of adjunction identifies the pulled-back universal Euler section with \(\mathsf{EL}_L\). Pullback compatibility of the relative zeroth-space functor therefore identifies the pulled-back section \(e_{\mathrm{univ}}^{(1)}\) with \(e_L^{(1)}\), and identifies zero with zero. Homotopy equalizers commute with pullback in the native \(\infty\)-topos, giving the first square. The full statement is identical with \(\Vfull\) in place of \(\Vfirst\).
\end{proof}

\begin{remark}[Critical loci form a higher family]
\label{rem:var-critical-family-higher}
The universal critical family retains automorphisms of Lagrangians, paths between Lagrangians, and the induced higher identifications among their Euler equations. Thus the passage from a Lagrangian to its Euler PDE is not merely an assignment on connected components; it is a morphism of higher moduli encoded by a single derived family.
\end{remark}

\subsection{The ambient Euler jet morphism for a cofree field object}

\begin{construction}[Constructed ambient Euler jet morphism]
\label{con:var-el-ambient}
Suppose the Cartan object is cofree:
\[
    (A,\rho_A)=K(Y)=(\Jinf Y,\Delta_Y)
\]
for an \(X\)-bundle \(Y\to X\). By Proposition~\ref{prop:var-cofree-descent},
\[
    \bar A=\Pi_\epsilon Y.
\]
Put
\[
    F_L:=\Crit(L)\longrightarrow\Pi_\epsilon Y
\]
over \(Q\), and define
\[
    \mathcal E_L:=\epsilon^*F_L.
\]
Applying \(\epsilon^*\) to \(F_L\to\Pi_\epsilon Y\) and using
\[
    \epsilon^*\Pi_\epsilon Y\simeq\Jinf Y
\]
gives a canonical morphism
\[
    \boxed{
    i_L:\mathcal E_L\longrightarrow\Jinf Y.
    }
\]
The descent object \(F_L\to Q\) equips \(\mathcal E_L\) with its canonical jet coaction, and \(i_L\) is automatically a coalgebra morphism.
\end{construction}

\begin{proposition}[The ambient Euler morphism is derived from descent]
\label{prop:var-el-ambient-cartan}
The map
\[
    i_L:(\mathcal E_L,\rho_L)\longrightarrow K(Y)
\]
is the Cartan morphism corresponding, under the Chapter 11 Eilenberg--Moore equivalence, to the ordinary morphism
\[
    F_L\longrightarrow\Pi_\epsilon Y
\]
over \(Q\). In particular no Cartan compatibility, prolongation equation, or centre-completeness datum is part of the input.
\end{proposition}

\begin{proof}
The equivalence
\[
    (\XSp)_{/Q}
    \simeq
    \Coalg_{\Jinf}(\mathcal C_X)
\]
preserves mapping spaces. Therefore the map of descended objects determines a unique Cartan morphism through a contractible space of choices. Its underlying map is obtained by pulling back along \(\epsilon\), which is precisely Construction~\ref{con:var-el-ambient}.
\end{proof}

\begin{remark}[Cartan prolongation of the constructed ambient morphism]
\label{rem:var-el-prolongation}
Chapter 11 may still be applied to the constructed ambient morphism \(i_L\), giving
\[
    \mathcal E_L^{\infty}
    \simeq
    \Jinf\mathcal E_L
    \times_{\Jinf\Jinf Y}
    \Jinf Y.
\]
This is a construction attached to the chosen ambient jet morphism. It is not what makes the Euler equation intrinsically formally integrable: Theorem~\ref{thm:var-critical-intrinsic-pde} has already done that before any ambient morphism is formed.
\end{remark}

\section{Full Helmholtz theory and coherent closed lifts}
\label{sec:var-helmholtz}

\begin{definition}[Full Helmholtz boundary]
\label{def:var-helmholtz}
For a full secondary one-form
\[
    \eta:\one_Q\longrightarrow\Sec^1(A;M)[n],
\]
define its full Helmholtz boundary by
\[
    \boxed{
    \Helm^{\infty}(\eta)
    :=
    \delta^{\infty}_{\mathrm{var}}\eta:
    \one_Q\longrightarrow\Sec^2(A;M)[n].
    }
\]
\end{definition}

\begin{proposition}[Full Euler forms carry a canonical coherent Helmholtz lift]
\label{prop:var-helmholtz-euler}
For every Lagrangian class \(L\), the full variation
\[
    \eta_L:=\delta^{\infty}_{\mathrm{var}}L
\]
has a canonical lift
\[
    \eta_L^{\mathrm{cl}}
    :=
    \delta^{\mathrm{cl}}_{\mathrm{var}}L
    :
    \one_Q\longrightarrow
    \Sec^{1,\mathrm{cl}}(A;M;n).
\]
Its underlying form is \(\eta_L\), and it includes a specified nullhomotopy
\[
    \Helm^{\infty}(\eta_L)
    =
    (\delta^{\infty}_{\mathrm{var}})^2L
    \simeq0
\]
together with all higher compatibility data.
\end{proposition}

\begin{proof}
Apply Construction~\ref{con:var-exact-classifier} in degree one to \(L\). The Chapter 10 exact-classifier theorem identifies its underlying cochain with the coherent infinitesimal boundary and identifies the retained closedness structure with the specified nullhomotopy of the square and all higher relations. Transport this statement through Theorem~\ref{thm:var-operator-identification}.
\end{proof}

\begin{definition}[Helmholtz lift space of a first Euler source]
\label{def:var-helmholtz-lift}
Let
\[
    e:\one_Q\longrightarrow\Sec^{1,(1)}(A;M)[n]
\]
be a first-order source. Define its \emph{space of coherent full Helmholtz lifts} by
\[
    \boxed{
    \HelmLift(e)
    :=
    \hofib_e\left(
      \Map_{\mathcal E_Q}
      (\one_Q,\Sec^{1,\mathrm{cl}}(A;M;n))
      \xrightarrow{\lambda^1_{\mathrm{var}}u_1}
      \Map_{\mathcal E_Q}
      (\one_Q,\Sec^{1,(1)}(A;M)[n])
    \right).
    }
\]
A point consists of a coherently closed full finite-infinitesimal one-form \(\eta\), its complete closedness hierarchy, and a specified equivalence
\[
    \lambda^1_{\mathrm{var}}u_1(\eta)\simeq e.
\]
\end{definition}

\begin{definition}[Parameterized Helmholtz-lift object]
\label{def:var-parameterized-helmholtz-lift}
Use the universal moduli notation
\[
    \mathbf E^{(1)}_{A/Q}(M;n)
    :=
    \mathbf{Eul}^n_{A/Q}(M),
    \qquad
    \mathbf C^{1,\mathrm{cl}}_{A/Q}(M;n)
    :=
    \mathbf{ClVar}^n_{A/Q}(M).
\]
Regard a first-order source
\(
 e:\one_Q\to\Sec^{1,(1)}(A;M)[n]
\)
as the corresponding section
\(e:Q\to\mathbf E^{(1)}_{A/Q}(M;n)\). Define
\[
    \boxed{
    \mathbf{HelmLift}(e)
    :=
    Q
    \times_{\mathbf E^{(1)}_{A/Q}(M;n)}
    \mathbf C^{1,\mathrm{cl}}_{A/Q}(M;n),
    }
\]
where the second map is induced by \(\lambda^1_{\mathrm{var}}u_1\).
\end{definition}

\begin{proposition}[Parameterized Helmholtz representing property]
\label{prop:var-parameterized-helmholtz-lift}
For every \(t:T\to Q\),
\[
    \Map_Q\bigl(T,\mathbf{HelmLift}(e)\bigr)
\]
is canonically the space of coherently closed full lifts of the pulled-back first source \(t^*e\). In particular
\[
    \boxed{
    \Map_Q\bigl(Q,\mathbf{HelmLift}(e)\bigr)
    \simeq
    \HelmLift(e).
    }
\]
The object \(\mathbf{HelmLift}(e)\) commutes with arbitrary pullback in \(Q\).
\end{proposition}

\begin{proof}
This is the mapping-space interpretation of the defining homotopy pullback, using Proposition~\ref{prop:var-moduli-representing}. Pullback preserves the coefficient objects, zeroth spaces, and finite homotopy pullbacks.
\end{proof}

\begin{theorem}[Every Euler source has its constructed Helmholtz lift]
\label{thm:var-euler-helmholtz-lift}
For every Lagrangian \(L\), the exact classifier \(\eta_L^{\mathrm{cl}}\) of Proposition~\ref{prop:var-helmholtz-euler} determines canonically a point
\[
    \boxed{
    \eta_L^{\mathrm{cl}}
    \in
    \HelmLift(\EL(L)).
    }
\]
\end{theorem}

\begin{proof}
The underlying form of \(\eta_L^{\mathrm{cl}}\) is \(\delta^{\infty}_{\mathrm{var}}L\). Applying \(\lambda^1_{\mathrm{var}}\) gives, by definition,
\[
    \lambda^1_{\mathrm{var}}\delta^{\infty}_{\mathrm{var}}L
    =
    \EL(L).
\]
The equality is a canonical comparison in the stable \(\infty\)-category, so it determines the required point of the homotopy fibre.
\end{proof}

\section{The derived inverse problem for first Euler sources}
\label{sec:var-inverse-problem}

\begin{definition}[First-order Lagrangian realization space]
\label{def:var-first-lagrangian-primitive}
For
\[
    e:\one_Q\longrightarrow\Sec^{1,(1)}(A;M)[n],
\]
define
\[
    \boxed{
    \LagPrim^{(1)}(e)
    :=
    \hofib_e\left(
      \Map_{\mathcal E_Q}(\one_Q,\Sec^0(A;M)[n])
      \xrightarrow{\delta^{(1)}_{\mathrm{var}}}
      \Map_{\mathcal E_Q}(\one_Q,\Sec^{1,(1)}(A;M)[n])
    \right).
    }
\]
A point is a Lagrangian class together with a specified equivalence between its first variation and \(e\).
\end{definition}

\begin{definition}[Parameterized first-order Lagrangian realization object]
\label{def:var-parameterized-first-lagrangian-primitive}
For a first source \(e\), define
\[
    \boxed{
    \mathbf{LagPrim}^{(1)}(e)
    :=
    Q
    \times_{\mathbf E^{(1)}_{A/Q}(M;n)}
    \Omega_Q^\infty\bigl(\Sec^0(A;M)[n]\bigr),
    }
\]
where the second map is induced by \(\delta^{(1)}_{\mathrm{var}}\).
\end{definition}

\begin{theorem}[Parameterized derived inverse-problem fibration]
\label{thm:var-parameterized-inverse-problem}
The exact variational classifier induces a canonical morphism over \(Q\)
\[
    \boxed{
    \mathbf{LagPrim}^{(1)}(e)
    \longrightarrow
    \mathbf{HelmLift}(e).
    }
\]
For a parameterized coherent full lift
\[
    \eta:T\longrightarrow\mathbf{HelmLift}(e),
\]
the homotopy fibre after pullback to \(T\) is canonically the parameterized primitive object
\[
    \boxed{
    \mathbf{Prim}_{\mathrm{var},T}(\eta).
    }
\]
All these objects and the fibre identification commute with arbitrary pullback.
\end{theorem}

\begin{proof}
The factorization
\[
    \Sec^0(A;M)[n]
    \xrightarrow{\delta^{\mathrm{cl}}_{\mathrm{var}}}
    \Sec^{1,\mathrm{cl}}(A;M;n)
    \xrightarrow{\lambda^1_{\mathrm{var}}u_1}
    \Sec^{1,(1)}(A;M)[n]
\]
induces the corresponding factorization on zeroth-space objects over \(Q\). Taking the two homotopy pullbacks along the section \(e\) gives the displayed map. Homotopy-fibre pasting in the slice \((\XSp)_{/Q}\) identifies its fibre over \(\eta\) with the pullback defining \(\mathbf{Prim}_{\mathrm{var},T}(\eta)\). Pullback in an \(\infty\)-topos preserves finite limits, proving the final assertion.
\end{proof}

\begin{theorem}[Derived inverse-problem fibration]
\label{thm:var-inverse-problem}
There is a canonical map of spaces
\[
    \boxed{
    \LagPrim^{(1)}(e)
    \longrightarrow
    \HelmLift(e)
    }
\]
which sends a Lagrangian realization \(L\) to its constructed coherent full lift \(\delta^{\mathrm{cl}}_{\mathrm{var}}L\). For a point
\[
    \eta\in\HelmLift(e),
\]
the homotopy fibre of this map over \(\eta\) is canonically
\[
    \boxed{
    \Prim_{\mathrm{var}}(\eta).
    }
\]
Equivalently, \(\LagPrim^{(1)}(e)\) is the total space of the family of complete primitive spaces parameterized by \(\HelmLift(e)\).
\end{theorem}

\begin{proof}
Take global \(Q\)-points of Theorem~\ref{thm:var-parameterized-inverse-problem}. Propositions~\ref{prop:var-parameterized-helmholtz-lift} and \ref{prop:var-parameterized-primitive} identify the resulting spaces with \(\HelmLift(e)\) and \(\Prim_{\mathrm{var}}(\eta)\), while the defining pullback of \(\mathbf{LagPrim}^{(1)}(e)\) gives \(\LagPrim^{(1)}(e)\). Since mapping out of the terminal object preserves limits, the internal fibre square becomes the asserted homotopy-fibre square of spaces.
\end{proof}

\begin{remark}[Two inverse problems are retained]
\label{rem:var-two-inverse-problems}
For a chosen full coherently closed one-form \(\eta\), the primitive space \(\Prim_{\mathrm{var}}(\eta)\) is the full finite-infinitesimal inverse problem. For a first Euler source \(e\), Theorem~\ref{thm:var-inverse-problem} retains the additional moduli of coherent full lifts before taking primitives. The distinction is forced by the possible nontrivial kernel of \(\lambda^1_{\mathrm{var}}\).
\end{remark}

\section{Horizontal realization and complete transgression}
\label{sec:var-horizontal-transgression}

The horizontal direction of the double relation is a complete effective descent direction. The correct intrinsic replacement for choosing a raw horizontal representative is therefore the space of lifts through the complete horizontal totalization tower.

\begin{definition}[Horizontal totalization filtration on secondary forms]
\label{def:var-horizontal-filtration}
For fixed \(p\ge0\) and \(q\ge0\), put
\[
    \Sec_H^{p,<q}(A;M)
    :=
    \Tot_{\Delta_{\le q-1}}
    \Bvar^{p,\bullet}(A;M),
\]
with \(\Sec_H^{p,<0}=0\), and define
\[
    \boxed{
    F_H^q\Sec^p(A;M)
    :=
    \fib\left(
      \Sec^p(A;M)
      \longrightarrow
      \Sec_H^{p,<q}(A;M)
    \right).
    }
\]
Thus \(F_H^0\Sec^p\simeq\Sec^p\) and \(F_H^{q+1}\to F_H^q\) is the complete decreasing horizontal totalization filtration.
\end{definition}

\begin{proposition}[Horizontal associated graded]
\label{prop:var-horizontal-graded}
The filtration of Definition~\ref{def:var-horizontal-filtration} is complete and satisfies
\[
    \boxed{
    \operatorname{gr}_H^qF_H^{\bullet}\Sec^p(A;M)
    \simeq
    \Omegavar^{p,q}(A;M)[-q].
    }
\]
The vertical differential preserves the entire horizontal filtration and induces the vertical coherent boundary on every associated-graded bidegree.
\end{proposition}

\begin{proof}
Apply Chapter 10's stable Dold--Kan layer theorem to the horizontal cosimplicial object \(\Bvar^{p,\bullet}\). Its degree-\(q\) normalization is \(\Omegavar^{p,q}\), so the layer is \(\Omegavar^{p,q}[-q]\). The vertical boundary is a morphism of horizontal cosimplicial objects before totalization, hence induces a morphism of the complete horizontal towers and their associated graded.
\end{proof}

\begin{construction}[Universal horizontal-realization morphism]
\label{con:var-universal-horizontal-realization}
For every \(p,q,n\), applying \(\Omega_Q^\infty\) to the filtration inclusion gives a canonical morphism
\[
    \boxed{
    \mathbf{RealH}^{p,q}_{A/Q}(M;n):
    \Omega_Q^\infty\bigl(F_H^q\Sec^p(A;M)[n]\bigr)
    \longrightarrow
    \boldsymbol\Omega^p_{\mathrm{var},A/Q}(M;n).
    }
\]
The horizontal realization space of a particular secondary form is the homotopy fibre of this universal map over that form.
\end{construction}

\begin{definition}[Horizontal realization space]
\label{def:var-horizontal-realization}
Let
\[
    \eta:\one_Q\longrightarrow\Sec^p(A;M)[n]
\]
be a secondary form. Define its space of horizontal \(q\)-realizations by
\[
    \boxed{
    \RealH^q(\eta)
    :=
    \hofib_{\eta}\left(
      \Map_{\mathcal E_Q}
      (\one_Q,F_H^q\Sec^p(A;M)[n])
      \longrightarrow
      \Map_{\mathcal E_Q}
      (\one_Q,\Sec^p(A;M)[n])
    \right).
    }
\]
A point is a lift of the secondary form into the \(q\)-th horizontal filtration tail. Equivalently, it is a coherent trivialization of the image of \(\eta\) in the partial horizontal totalization \(\Sec_H^{p,<q}(A;M)[n]\), together with all compatibility data retained by that stage of the tower.
\end{definition}

\begin{definition}[Parameterized horizontal realization object]
\label{def:var-parameterized-horizontal-realization}
Regard
\[
    \eta:\one_Q\longrightarrow\Sec^p(A;M)[n]
\]
as the corresponding section
\[
    \eta:Q\longrightarrow
    \Omega_Q^\infty\bigl(\Sec^p(A;M)[n]\bigr).
\]
Define
\[
    \boxed{
    \mathbf{RealH}^{q}(\eta)
    :=
    Q
    \times_{\Omega_Q^\infty\Sec^p(A;M)[n]}
    \Omega_Q^\infty F_H^q\Sec^p(A;M)[n].
    }
\]
Then for every \(t:T\to Q\),
\[
    \Map_Q\bigl(T,\mathbf{RealH}^{q}(\eta)\bigr)
\]
is the horizontal realization space of \(t^*\eta\), and in particular
\[
    \Map_Q\bigl(Q,\mathbf{RealH}^{q}(\eta)\bigr)
    \simeq
    \RealH^q(\eta).
\]
The parameterized realization object commutes with arbitrary pullback.
\end{definition}

\begin{definition}[Complete horizontal transgression coefficient]
\label{def:var-horizontal-transgression}
Define
\[
    \boxed{
    \TransH^{p,q}(A;M)
    :=
    \fib\left(
      F_H^q\Sec^p(A;M)
      \longrightarrow
      \Sec^p(A;M)
    \right).
    }
\]
Since \(F_H^q\Sec^p\) is itself the fibre of
\[
    \Sec^p\longrightarrow\Sec_H^{p,<q},
\]
stability gives a canonical equivalence
\[
    \boxed{
    \TransH^{p,q}(A;M)
    \simeq
    \Omega\,\Sec_H^{p,<q}(A;M).
    }
\]
For a test \(t:T\to Q\), a section of \(t^*F_H^q\Sec^p\) together with a specified nullhomotopy of its image in \(t^*\Sec^p\) determines canonically a point of
\[
    \Omega^\infty
    \Gamma_T^{\mathrm{st}}
    \bigl(t^*\TransH^{p,q}(A;M)\bigr).
\]
Such a point is called a \emph{complete horizontal transgression}.
\end{definition}

\begin{proposition}[Differences of horizontal realizations]
\label{prop:var-horizontal-differences}
Let \(\eta\) be a secondary form and let
\[
    \widetilde\eta_0,\widetilde\eta_1
    \in
    \RealH^q(\eta)
\]
be two horizontal realizations. Their difference, together with the specified path identifying their images with the same point \(\eta\), determines canonically a point of the complete horizontal transgression space associated with \(\TransH^{p,q}(A;M)\). Paths and higher homotopies between realizations map to the corresponding higher homotopies of transgressions.
\end{proposition}

\begin{proof}
The source and target of
\[
    F_H^q\Sec^p\longrightarrow\Sec^p
\]
are spectra. Subtract the two lifted sections in the source. Their images in \(\Sec^p\) are identified through the paths contained in the two points of \(\RealH^q(\eta)\), producing a specified nullhomotopy of the difference. By the universal property of the homotopy fibre, the pair determines a point of \(\TransH^{p,q}\). Functoriality of fibres gives the higher statement.
\end{proof}

\begin{proposition}[Vertical variation of horizontal realizations]
\label{prop:var-vertical-horizontal-realization}
The vertical differential induces morphisms of complete horizontal towers
\[
    d_V:
    F_H^q\Sec^p(A;M)
    \longrightarrow
    F_H^q\Sec^{p+1}(A;M)
\]
and therefore maps
\[
    \RealH^q(\eta)
    \longrightarrow
    \RealH^q(\delta^{\infty}_{\mathrm{var}}\eta)
\]
for every secondary form \(\eta\). It also induces
\[
    d_V^{\mathrm{tr}}:
    \TransH^{p,q}(A;M)
    \longrightarrow
    \TransH^{p+1,q}(A;M).
\]
The induced maps retain the coherent nullhomotopy of \(d_V^2\) and all higher relations.
\end{proposition}

\begin{proof}
The first statement is Proposition~\ref{prop:var-horizontal-graded} before passing to associated graded. The two squares
\[
\begin{tikzcd}
F_H^q\Sec^p \arrow[r,"d_V"] \arrow[d]
& F_H^q\Sec^{p+1} \arrow[d]
\\
\Sec^p \arrow[r,"\delta^{\infty}_{\mathrm{var}}"']
& \Sec^{p+1}
\end{tikzcd}
\]
commute through the natural coherence of the horizontal tower. Functoriality of homotopy fibres gives the maps on realization and transgression spaces. The vertical coherent-cochain structure is a natural transformation of the complete horizontal towers, so its higher relations pass to the fibres.
\end{proof}


\begin{remark}[No preferred Lepage representative]
\label{rem:var-no-canonical-lepage}
The horizontal realization tower does not select a Poincar\'e--Cartan form, a horizontal homotopy operator, or a source-form splitting.  This remains essential: the complete horizontal direction is effective Cartan descent, and the present chapter does not identify it with a chosen complex of exterior horizontal forms.  What can be constructed intrinsically is the derived moduli object of \emph{relative Lepage realizations}: one compares the canonical coherently closed full Euler lift with another coherently closed lift of the same first Euler source and asks that their difference admit a specified lift into a chosen stage of the horizontal totalization filtration.  The construction below retains the complete space of such comparison data and does not select one of its points.
\end{remark}

\begin{remark}[Presymplectic data is relative to Lepage realization]
\label{rem:var-presymplectic-interface}
A classical presymplectic current is obtained by vertically differentiating a boundary potential.  In the brave-new theory the corresponding unconditional object is therefore not a distinguished current on \(Q\), but a universal transgression over the derived moduli of Lepage realizations.  A chosen point of that moduli object produces a particular presymplectic transgression; changing the point is measured by the already constructed complete horizontal transgression coefficient.  No preferred representative is inserted.
\end{remark}

\section{Derived Lepage realizations and universal presymplectic transgression}
\label{sec:var-lepage-presymplectic}

The preceding horizontal tower already contains the complete obstruction and ambiguity theory for Cartan boundary data.  This section packages that information relative to the Euler source of a fixed Lagrangian.  The construction is deliberately relative: it compares coherently closed full lifts of the same first Euler source.  Consequently it does not identify the full brave-new one-form coefficient with a classical source-form complex and does not assert that a preferred Poincar\'e--Cartan representative exists.

\subsection{The relative Lepage realization object}

\begin{construction}[Universal difference of full Euler lifts]
\label{con:var-lepage-difference}
Let
\[
    L:\one_Q\longrightarrow\Sec^0(A;M)[n]
\]
be a Lagrangian, put
\[
    e_L:=\EL(L):
    \one_Q\longrightarrow\Sec^{1,(1)}(A;M)[n],
\]
and let
\[
    \eta_L^{\mathrm{cl}}
    =
    \delta_{\mathrm{var}}^{\mathrm{cl}}L:
    \one_Q\longrightarrow
    \Sec^{1,\mathrm{cl}}(A;M;n)
\]
be its canonical coherently closed full Euler lift.  Write
\[
    H_L:=\mathbf{HelmLift}(e_L)
    \xrightarrow{h_L}Q
\]
for the parameterized Helmholtz-lift object.

By the defining homotopy pullback of \(H_L\), there is a tautological coherently closed lift
\[
    \eta_{\mathrm{taut}}:
    \one_{H_L}
    \longrightarrow
    h_L^*\Sec^{1,\mathrm{cl}}(A;M;n)
\]
whose first-order image is \(h_L^*e_L\).  Pulling back \(\eta_L^{\mathrm{cl}}\) to \(H_L\), applying the underlying-form map \(u_1\), and subtracting in the stable category gives the canonical section
\[
    \boxed{
    \kappa_L:
    \one_{H_L}
    \longrightarrow
    h_L^*\Sec^1(A;M)[n],
    \qquad
    \kappa_L
    :=
    h_L^*u_1(\eta_L^{\mathrm{cl}})
    -
    u_1(\eta_{\mathrm{taut}}).
    }
\]
Both summands have the same image under \(h_L^*\lambda^1_{\mathrm{var}}\), so \(\kappa_L\) carries the canonical nullhomotopy of its first-order image.
\end{construction}

\begin{definition}[Stage-\(q\) relative Lepage realization object]
\label{def:var-lepage-object}
For \(q\ge1\), regard \(\kappa_L\) as the corresponding morphism
\[
    H_L\longrightarrow
    \Omega_Q^\infty\bigl(\Sec^1(A;M)[n]\bigr)
\]
over \(Q\).  Define
\[
    \boxed{
    \mathbf{Lep}^{\,q}(L)
    :=
    H_L
    \times_{\Omega_Q^\infty\Sec^1(A;M)[n]}
    \Omega_Q^\infty F_H^q\Sec^1(A;M)[n].
    }
\]
The second map is the universal horizontal-realization morphism of Construction~\ref{con:var-universal-horizontal-realization}.  Write
\[
    \pi_L^q:\mathbf{Lep}^{\,q}(L)\longrightarrow H_L\longrightarrow Q
\]
for the structural morphism.
\end{definition}

\begin{proposition}[Meaning and representing property of relative Lepage realizations]
\label{prop:var-lepage-representing}
Let \(t:T\to Q\).  A \(T\)-point of \(\mathbf{Lep}^{\,q}(L)\) consists of
\begin{enumerate}
\item a coherently closed full lift
\[
    \eta:
    \one_T\longrightarrow
    t^*\Sec^{1,\mathrm{cl}}(A;M;n)
\]
of the pulled-back first Euler source \(t^*e_L\);
\item a lift
\[
    \Theta:
    \one_T\longrightarrow
    t^*F_H^q\Sec^1(A;M)[n]
\]
whose image in \(t^*\Sec^1(A;M)[n]\) is equipped with a specified equivalence
\[
    \Theta\longmapsto
    t^*u_1(\eta_L^{\mathrm{cl}})-u_1(\eta);
\]
\item all higher compatibility data belonging to those two homotopy fibres.
\end{enumerate}
Equivalently,
\[
    \Map_Q\bigl(T,\mathbf{Lep}^{\,q}(L)\bigr)
\]
is the total space, over the \(T\)-family of coherent Helmholtz lifts of \(t^*e_L\), of the horizontal \(q\)-realization spaces of their differences from the canonical full Euler lift.

In particular define the global relative Lepage-realization space by
\[
    \boxed{
    \Lep^q(L)
    :=
    \Map_Q\bigl(Q,\mathbf{Lep}^{\,q}(L)\bigr).
    }
\]
The canonical lift \(\eta_L^{\mathrm{cl}}\), together with the zero horizontal realization of the zero difference, gives a distinguished basepoint
\[
    \ell_{L,0}^{\,q}:Q\longrightarrow\mathbf{Lep}^{\,q}(L).
\]
This basepoint is only the tautological zero comparison; it is not asserted to be a classical source-form or Poincar\'e--Cartan normalization.
\end{proposition}

\begin{proof}
Apply the mapping-space interpretation of the two defining homotopy pullbacks: first that of \(\mathbf{HelmLift}(e_L)\), then that of Definition~\ref{def:var-lepage-object}.  The displayed description is exactly the universal property of the resulting iterated homotopy pullback.  For the basepoint, take the point \(\eta_L^{\mathrm{cl}}\in\HelmLift(e_L)\) of Theorem~\ref{thm:var-euler-helmholtz-lift}.  Its difference from itself is the zero section of \(\Sec^1[n]\), and the zero section of \(F_H^q\Sec^1[n]\) is a lift of that zero section.
\end{proof}


\begin{construction}[Tautological complete Lepage boundary realization]
\label{con:var-tautological-lepage}
By Definition~\ref{def:var-lepage-object}, over \(\mathbf{Lep}^{\,q}(L)\) there is a tautological section
\[
    \boxed{
    \Theta_L^{(q)}:
    \one_{\mathbf{Lep}^{\,q}(L)}
    \longrightarrow
    (\pi_L^q)^*F_H^q\Sec^1(A;M)[n].
    }
\]
Its image in \((\pi_L^q)^*\Sec^1(A;M)[n]\) is equipped with the tautological equivalence
\[
    \Theta_L^{(q)}
    \longmapsto
    (\pi_L^q)^*u_1(\eta_L^{\mathrm{cl}})
    -
    u_1(\eta_{\mathrm{taut}}).
\]
We call \(\Theta_L^{(q)}\) the \emph{universal complete Lepage boundary realization of stage \(q\)}.  The terminology refers to its universal property in the horizontal descent tower; no exterior-form identification is implicit.
\end{construction}


\begin{theorem}[Relative brave-new integration by parts]
\label{thm:var-relative-integration-by-parts}
Let
\[
    \iota_H^q:
    F_H^q\Sec^1(A;M)
    \longrightarrow
    \Sec^1(A;M)
\]
be the horizontal filtration inclusion.  Over \(\mathbf{Lep}^{\,q}(L)\), write
\[
    E_L^{(q)}
    :=
    u_1(\eta_{\mathrm{taut}})
\]
for the tautological coherently closed full representative of the first Euler source.  Then the defining homotopy of the Lepage pullback is equivalently a canonical coherent factorization
\[
    \boxed{
    (\pi_L^q)^*\delta_{\mathrm{var}}^\infty L
    \simeq
    E_L^{(q)}
    +
    \iota_H^q\Theta_L^{(q)}.
    }
\]
Moreover
\[
    (\pi_L^q)^*(\lambda^1_{\mathrm{var}})\,E_L^{(q)}
    \simeq
    (\pi_L^q)^*\EL(L),
\]
and the horizontal term carries the specified nullhomotopy of its first-order image.  Thus a point of \(\mathbf{Lep}^{\,q}(L)\) is a derived \emph{relative integration-by-parts realization}: it consists of a coherent full Euler representative and a complete horizontal boundary realization comparing it with the canonical full first variation.

For \(q=1\) this is the one-stage brave-new analogue of a decomposition customarily written
\[
    d_VL=E_L+d_H\Theta,
\]
but the displayed boxed identity, rather than that classical notation, is the intrinsic statement.
\end{theorem}

\begin{proof}
By Construction~\ref{con:var-lepage-difference}, the section classified by the first map in Definition~\ref{def:var-lepage-object} is
\[
    (\pi_L^q)^*\delta_{\mathrm{var}}^\infty L
    -
    E_L^{(q)}.
\]
The homotopy pullback defining \(\mathbf{Lep}^{\,q}(L)\) supplies a lift of this difference through \(F_H^q\Sec^1[n]\) and a specified equivalence between its image under \(\iota_H^q\) and that difference.  Rearranging in the stable category gives the boxed formula.  The first-order statement is the tautological property of \(\eta_{\mathrm{taut}}\in\mathbf{HelmLift}(e_L)\).  Subtracting the two first-order images gives the retained nullhomotopy on the horizontal term.
\end{proof}


\subsection{Universal presymplectic transgression}

\begin{theorem}[Universal presymplectic transgression]
\label{thm:var-universal-presymplectic}
For every \(q\ge1\), vertical variation of the tautological boundary realization, together with the coherent closedness structures of the two full Euler lifts, determines canonically a section
\[
    \boxed{
    \omega_L^{(q)}:
    \one_{\mathbf{Lep}^{\,q}(L)}
    \longrightarrow
    (\pi_L^q)^*\TransH^{2,q}(A;M)[n].
    }
\]
It is natural in \(L\), in coefficient morphisms, in Cartan morphisms, and under arbitrary pullback.  We call \(\omega_L^{(q)}\) the \emph{universal stage-\(q\) presymplectic transgression}.
\end{theorem}

\begin{proof}
Apply Proposition~\ref{prop:var-vertical-horizontal-realization} to the tautological section
\[
    \Theta_L^{(q)}:
    \one\to
    (\pi_L^q)^*F_H^q\Sec^1[n].
\]
This gives
\[
    d_V\Theta_L^{(q)}:
    \one\longrightarrow
    (\pi_L^q)^*F_H^q\Sec^2[n].
\]
Its image in \((\pi_L^q)^*\Sec^2[n]\) is the vertical differential of
\[
    (\pi_L^q)^*u_1(\eta_L^{\mathrm{cl}})
    -
    u_1(\eta_{\mathrm{taut}}).
\]
Both \(\eta_L^{\mathrm{cl}}\) and \(\eta_{\mathrm{taut}}\) are points of the coherently closed object \(\Sec^{1,\mathrm{cl}}\).  Their retained closedness structures therefore supply a specified nullhomotopy of that difference in \(\Sec^2[n]\), together with all higher compatibility cells.  The pair
\[
    \bigl(d_V\Theta_L^{(q)},\text{ specified nullhomotopy}\bigr)
\]
is consequently a section of the homotopy fibre
\[
    \fib\left(
      F_H^q\Sec^2\longrightarrow\Sec^2
    \right)
    =
    \TransH^{2,q}.
\]
Functoriality is inherited from the parameterized Helmholtz object, the horizontal realization tower, and the vertical coherent differential.
\end{proof}

\begin{theorem}[Complete presymplectic recurrence]
\label{thm:var-presymplectic-recurrence}
The section \(\omega_L^{(q)}\) is the canonical complete-horizontal-transgression witness for the equality of the two coherent vertical boundaries
\[
    d_Vu_1(\eta_L^{\mathrm{cl}})
    \simeq
    d_Vu_1(\eta_{\mathrm{taut}})
\]
after transport through the stage-\(q\) Lepage realization.  Equivalently, its image under the fibre inclusion
\[
    \TransH^{2,q}\longrightarrow F_H^q\Sec^2
\]
is \(d_V\Theta_L^{(q)}\), while its composite to \(\Sec^2\) is canonically nullhomotopic by the difference of the two retained Helmholtz nullhomotopies.

On every associated-graded horizontal layer this recurrence is the image of the mixed coherent relation
\[
    d_Hd_V+d_Vd_H\simeq0
\]
of Theorem~\ref{thm:var-bicomplex}.  Thus whenever a represented comparison identifies a chosen horizontal realization with an ordinary one-step boundary potential, the unary signed shadow is the classical presymplectic recurrence.
\end{theorem}

\begin{proof}
The first statement is the construction in the proof of Theorem~\ref{thm:var-universal-presymplectic}.  Passing to an associated-graded horizontal layer commutes with the vertical differential by Proposition~\ref{prop:var-horizontal-graded}.  The comparison of the two orders of boundary is therefore the associated-graded image of the mixed coherent cell of Theorem~\ref{thm:var-bicomplex}, with the Koszul sign already incorporated in the definition of \(d_H\).
\end{proof}

\begin{proposition}[Vertical closedness]
\label{prop:var-presymplectic-vertical-closed}
The universal presymplectic transgression carries a canonical coherent vertical closedness:
\[
    \boxed{
    d_V^{\mathrm{tr}}\omega_L^{(q)}
    \simeq0.
    }
\]
The nullhomotopy includes the full hierarchy of higher compatibility data inherited from the coherent relation \(d_V^2\simeq0\).
\end{proposition}

\begin{proof}
The map
\[
    d_V^{\mathrm{tr}}:
    \TransH^{2,q}\longrightarrow\TransH^{3,q}
\]
is induced functorially from the vertical coherent differential on the two horizontal towers.  Apply it to the fibre representative of Theorem~\ref{thm:var-universal-presymplectic}.  The resulting section is represented by \(d_V^2\Theta_L^{(q)}\) together with the vertically differentiated nullhomotopy.  The specified coherent nullhomotopy of \(d_V^2\) and its higher relations identify this section canonically with zero.
\end{proof}

\begin{proposition}[Change of Lepage realization]
\label{prop:var-presymplectic-change}
Fix a coherently closed full lift
\[
    \eta\in\HelmLift(e_L)
\]
and let two points of the fibre of
\[
    \mathbf{Lep}^{\,q}(L)\longrightarrow\mathbf{HelmLift}(e_L)
\]
over \(\eta\) be represented by horizontal realizations
\[
    \Theta_0,\Theta_1.
\]
Their difference determines a complete horizontal transgression
\[
    \tau_{01}\in
    \Omega^\infty\Gamma_Q^{\mathrm{st}}\TransH^{1,q}(A;M)[n]
\]
as in Proposition~\ref{prop:var-horizontal-differences}, and the corresponding presymplectic transgressions satisfy
\[
    \boxed{
    \omega_1-\omega_0
    \simeq
    d_V^{\mathrm{tr}}\tau_{01}.
    }
\]
All paths and higher homotopies between the two Lepage realizations are carried to the corresponding higher homotopies between the presymplectic transgressions.
\end{proposition}

\begin{proof}
The two horizontal realizations have the same image in \(\Sec^1[n]\), so Proposition~\ref{prop:var-horizontal-differences} produces \(\tau_{01}\).  The construction of Theorem~\ref{thm:var-universal-presymplectic} is functorial in the horizontal realization and is obtained by applying the vertical map of Proposition~\ref{prop:var-vertical-horizontal-realization}.  Subtraction in the stable category therefore gives
\[
    \omega_1-\omega_0
    =
    d_V^{\mathrm{tr}}(\Theta_1-\Theta_0),
\]
which is the displayed identity after the universal fibre factorization defining \(\tau_{01}\).
\end{proof}


\subsection{Universal Lepage and presymplectic families over Lagrangian moduli}

\begin{construction}[Universal relative Lepage family]
\label{con:var-universal-lepage-family}
Let
\[
    P:=\mathbf{Lag}^n_{A/Q}(M),
    \qquad
    p:P\longrightarrow Q
\]
and let
\[
    L_{\mathrm{univ}}:
    \one_P\longrightarrow
    \Sec^0(A_P;p^*M)[n]
\]
be the tautological Lagrangian of Construction~\ref{con:var-universal-euler-section}.  Apply Definition~\ref{def:var-lepage-object} after base change to \(P\) and define
\[
    \boxed{
    \mathbf{Lep}^{\,q}_{\mathrm{univ}}
    :=
    \mathbf{Lep}^{\,q}(L_{\mathrm{univ}})
    \longrightarrow P.
    }
\]
It carries the universal boundary realization
\[
    \Theta_{\mathrm{univ}}^{(q)}
\]
and universal presymplectic transgression
\[
    \omega_{\mathrm{univ}}^{(q)}.
\]
\end{construction}

\begin{theorem}[Individual Lepage and presymplectic objects are pullbacks]
\label{thm:var-lepage-family-pullback}
For every Lagrangian section
\[
    L:Q\longrightarrow P,
\]
pullback of Construction~\ref{con:var-universal-lepage-family} along \(L\) gives canonical equivalences
\[
    \boxed{
    L^*\mathbf{Lep}^{\,q}_{\mathrm{univ}}
    \simeq
    \mathbf{Lep}^{\,q}(L),
    }
\]
and identifies
\[
    L^*\Theta_{\mathrm{univ}}^{(q)}
    \simeq
    \Theta_L^{(q)},
    \qquad
    L^*\omega_{\mathrm{univ}}^{(q)}
    \simeq
    \omega_L^{(q)}.
\]
Thus relative Lepage realizations and presymplectic transgressions form one higher family over the moduli object of Lagrangians.
\end{theorem}

\begin{proof}
The universal Lagrangian pulls back to \(L\) by Construction~\ref{con:var-universal-euler-section}.  Arbitrary base change identifies its Euler source, canonical closed full lift, parameterized Helmholtz object, horizontal filtration stages, and transgression coefficients with the corresponding objects for \(L\).  Definition~\ref{def:var-lepage-object} is a finite homotopy pullback of those objects, so it commutes with pullback.  The tautological sections and Theorem~\ref{thm:var-universal-presymplectic} are functorial in the same diagram, proving the remaining identities.
\end{proof}

\begin{remark}[What has and has not been constructed]
\label{rem:var-lepage-scope}
The object \(\mathbf{Lep}^{\,q}(L)\) is a derived moduli object of relative horizontal realization data.  Its basepoint is the tautological comparison of the canonical full Euler lift with itself, while nontrivial points compare that lift with other coherent full lifts through the horizontal filtration.  The construction is therefore noncommittal about classical source normalization.  It does not assert a preferred integration-by-parts operator, an all-degree HKR collapse, or an identification of \(F_H^\bullet\) with a complex of exterior horizontal forms.  A later represented comparison may single out geometrically meaningful subspaces or points; the universal presymplectic transgression above is already available over the entire derived realization object.
\end{remark}


\section{Higher cubical variations and multilinearization descent}
\label{sec:var-second-variation}
\label{sec:var-higher-variations}

The first split square-zero variation and the mixed bi-square-zero variation are the first two cases of one finite-cubical construction. This section constructs that system uniformly. No polynomial or multilinear approximation is imposed: higher variation is cubical before any possible multilinear descent.

\begin{definition}[Split \(k\)-cube of square-zero directions]
\label{def:var-k-cube-test}
Let \(T=\operatorname{Spec}B\), let \(k\ge1\), and let
\[
    N_1,\ldots,N_k\in\operatorname{Mod}_B
\]
be connective. For a subset \(I\subseteq[k]:=\{1,\ldots,k\}\), define
\[
    B[N_I]
    :=
    \bigotimes_{i\in I}{}_B(B\oplus N_i),
    \qquad
    B[N_\varnothing]:=B,
\]
and
\[
    T[N_I]:=\operatorname{Spec}B[N_I].
\]
For \(I\subseteq J\), augmentation in the factors \(J\setminus I\) gives a morphism of algebras
\[
    B[N_J]\longrightarrow B[N_I]
\]
and hence a face inclusion
\[
    j_{I,J}:T[N_I]\longrightarrow T[N_J].
\]
Unit maps give compatible retractions
\[
    p_{J,I}:T[N_J]\longrightarrow T[N_I].
\]
The full vertex is denoted
\[
    T[N_1,\ldots,N_k]:=T[N_{[k]}].
\]
\end{definition}

\begin{proposition}[Underlying module of a split \(k\)-cube]
\label{prop:var-k-cube-module}
There is a canonical equivalence of \(B\)-modules
\[
    \boxed{
    B[N_1,\ldots,N_k]
    \simeq
    \bigoplus_{I\subseteq[k]}
    \bigotimes_{i\in I}N_i,
    }
\]
where the empty tensor product is \(B\). Multiplication is induced by disjoint union of index sets and is coherently zero whenever a factor \(N_i\) is repeated. The symmetric monoidal coherence gives a canonical \(\Sigma_k\)-action permuting the labelled directions.
\end{proposition}

\begin{proof}
Expand the finite tensor product
\[
    \bigotimes_{i=1}^k(B\oplus N_i)
\]
using distributivity of tensor product over finite direct sums in \(\operatorname{Mod}_B\). Choosing either \(B\) or \(N_i\) in each factor gives one summand for each subset \(I\subseteq[k]\). Since the augmentation ideal in each individual factor is square-zero, multiplication vanishes whenever the same index occurs twice. Permutation of tensor factors gives the coherent symmetric-group action.
\end{proof}

\begin{definition}[Structured \(k\)-variation diagram]
\label{def:var-k-variation}
A \emph{split structured \(k\)-variation} centred at an affine point \(a_T:T\to\bar A\) is a map
\[
    \widetilde a_{[k]}:
    T[N_1,\ldots,N_k]
    \longrightarrow
    \bar A
\]
over \(Q\), where the structural map to \(Q\) is explicitly
\[
    T[N_1,\ldots,N_k]
    \xrightarrow{r_{[k]}}
    T
    \xrightarrow{a_T}
    \bar A
    \xrightarrow{a}
    Q,
\]
together with a specified identification of its restriction to the central vertex \(T\) with \(a_T\). For every \(I\subseteq[k]\), write
\[
    \widetilde a_I:
    T[N_I]\longrightarrow\bar A
\]
for the restriction along \(j_{I,[k]}\).

Let
\[
    \mathcal V^{(k)}_{\bar A/Q}
\]
be the \(\infty\)-category of complete structured \(k\)-variation diagrams. A morphism consists of a morphism of affine centres \(g:T'\to T\), module maps
\[
    \alpha_i:g^*N_i\longrightarrow N_i',
    \qquad 1\le i\le k,
\]
the induced morphism of the complete split \(k\)-cubes, and a specified compatible equivalence of the maps to \(\bar A\). Write
\[
    \pi_k:
    \mathcal V^{(k)}_{\bar A/Q}
    \longrightarrow
    \mathcal C^{\mathrm{aff}}_{/\bar A}
\]
for the centre functor. Cartesian pullback of the entire cube along a morphism of centres gives Cartesian lifts, so \(\pi_k\) is a Cartesian fibration. For \(k=1\), this construction is canonically the category \(\mathcal V^{(1)}_{\bar A/Q}\) and centre functor \(\pi_1\) of Definition~\ref{def:var-sqz-category}; we henceforth use that identification rather than introducing a second first-variation category.
\end{definition}

\begin{definition}[Centre-pullback Nisnevich topology on \(k\)-variations]
\label{def:var-k-variation-topology}
Equip \(\mathcal V^{(k)}_{\bar A/Q}\) with the Grothendieck topology generated by Cartesian pullback of affine spectral Nisnevich covering families along \(\pi_k\).
\end{definition}

\begin{definition}[The \(k\)-th cubical difference coefficient]
\label{def:var-k-difference-functor}
For \(v\in\mathcal V^{(k)}_{\bar A/Q}\), let
\[
    \Gamma_v(I)
    :=
    \Gamma_{T[N_I]}^{\mathrm{st}}(M_{T[N_I]}),
    \qquad I\subseteq[k].
\]
Restriction along the face inclusions makes \(I\mapsto\Gamma_v(I)\) a \(k\)-cube of spectra with arrows from larger subsets to smaller subsets. Define its total fibre by
\[
    \boxed{
    K_M^{(k)}(v)
    :=
    \operatorname{TotFib}_{I\subseteq[k]}\Gamma_v(I)
    :=
    \fib\left(
      \Gamma_v([k])
      \longrightarrow
      \lim_{I\subsetneq[k]}\Gamma_v(I)
    \right).
    }
\]
Structured morphisms of \(k\)-variations induce restriction maps on the entire cube and hence on its total fibre, giving a functor
\[
    K_M^{(k)}:
    (\mathcal V^{(k)}_{\bar A/Q})^{\mathrm{op}}
    \longrightarrow
    \operatorname{Sp}.
\]
For \(k=1\), this is \(K_M^{(1)}\). For \(k=2\), it is canonically
\[
    K_M^{(2)}(v_{12})
    \simeq
    \fib\left(
      \Gamma_{T[N_1,N_2]}^{\mathrm{st}}(M)
      \longrightarrow
      \Gamma_{T[N_1]}^{\mathrm{st}}(M)
      \times_{\Gamma_T^{\mathrm{st}}(M)}
      \Gamma_{T[N_2]}^{\mathrm{st}}(M)
    \right).
\]
\end{definition}

\begin{lemma}[Alternating projector for a split finite cube]
\label{lem:var-split-cube-projector}
Let \(\mathcal D\) be a stable \(\infty\)-category and let
\[
    X:\mathcal P([k])^{\mathrm{op}}\longrightarrow\mathcal D
\]
be a finite \(k\)-cube. For \(I\subseteq J\), write
\[
    \rho_{J,I}:X_J\longrightarrow X_I
\]
for the restriction map. Suppose that the cube is equipped with coherent sections
\[
    s_{I,J}:X_I\longrightarrow X_J,
    \qquad
    \rho_{J,I}s_{I,J}\simeq\id_{X_I},
\]
compatible with chains of inclusions and with the cubical structure. For
\[
    \widehat i:=[k]\setminus\{i\},
\]
put
\[
    P_i:=s_{\widehat i,[k]}\rho_{[k],\widehat i}
    \in\operatorname{End}_{\mathcal D}(X_{[k]}).
\]
Then the \(P_i\) are coherently commuting idempotents, and the endomorphism
\[
    D_{[k]}
    :=
    (\id-P_k)\circ\cdots\circ(\id-P_1)
\]
is independent, through the canonical coherent comparisons, of the chosen ordering of the factors. There is a canonical natural factorization
\[
    \boxed{
    X_{[k]}
    \xrightarrow{\;\overline D_{[k]}\;}
    \operatorname{TotFib}(X)
    \xrightarrow{\;\iota\;}
    X_{[k]}
    }
\]
such that
\[
    \iota\overline D_{[k]}\simeq D_{[k]},
    \qquad
    \overline D_{[k]}\iota\simeq\id_{\operatorname{TotFib}(X)}.
\]
Moreover
\[
    \boxed{
    D_{[k]}
    \simeq
    \sum_{I\subseteq[k]}
    (-1)^{k-|I|}
    s_{I,[k]}\rho_{[k],I}.
    }
\]
The factorization is natural in coherently split cubes and is compatible with reindexing by bijections of the finite direction set.
\end{lemma}

\begin{proof}
We argue by induction on \(k\). For \(k=1\), write the split arrow as
\[
    X_1\xrightarrow{\rho}X_0,
    \qquad
    s:X_0\longrightarrow X_1,
    \qquad
    \rho s\simeq\id_{X_0}.
\]
Because \(\rho\) has the specified section \(s\), the fibre sequence splits canonically as
\[
    \fib(\rho)\oplus X_0
    \xrightarrow{\;\sim\;}
    X_1,
\]
where the two summand maps are the fibre inclusion \(\iota\) and \(s\). Let
\[
    \overline D_1:X_1\longrightarrow\fib(\rho)
\]
be the first component of the inverse equivalence. Then
\[
    \iota\overline D_1\simeq\id-s\rho,
    \qquad
    \overline D_1\iota\simeq\id_{\fib(\rho)}.
\]
This proves the claim in degree one.

Assume the statement for \(k-1\). Separate the last coordinate and define two split \((k-1)\)-cubes
\[
    X^1_J:=X_{J\cup\{k\}},
    \qquad
    X^0_J:=X_J,
    \qquad
    J\subseteq[k-1].
\]
Restriction in the last direction gives a morphism of split cubes
\[
    \rho_k:X^1\longrightarrow X^0
\]
with coherent section induced by the given cubical sections. Fubini for finite limits gives a canonical equivalence
\[
    \operatorname{TotFib}(X)
    \simeq
    \fib\left(
      \operatorname{TotFib}(X^1)
      \longrightarrow
      \operatorname{TotFib}(X^0)
    \right).
\]
By the induction hypothesis, the alternating projector in the first \(k-1\) directions factors naturally through \(\operatorname{TotFib}(X^1)\). Naturality with respect to \(\rho_k\) and its section makes the induced arrow
\[
    \operatorname{TotFib}(X^1)
    \longrightarrow
    \operatorname{TotFib}(X^0)
\]
a split arrow. Applying the degree-one construction to this arrow and composing with the inductive projector gives
\[
    \overline D_{[k]}:
    X_{[k]}\longrightarrow\operatorname{TotFib}(X).
\]
Its composite with the total-fibre inclusion is
\[
    (\id-P_k)\circ\cdots\circ(\id-P_1),
\]
which is \(D_{[k]}\), and the same induction gives
\[
    \overline D_{[k]}\iota\simeq\id.
\]

For distinct \(i,j\), cubical compatibility of the restrictions and sections gives
\[
    P_iP_j\simeq P_jP_i,
\]
with the full higher compatibility inherited from the split cube. Hence the factors \(\id-P_i\) commute coherently. Expanding their product in the additive mapping spectrum gives
\[
    D_{[k]}
    \simeq
    \sum_{J\subseteq[k]}
    (-1)^{|J|}
    \prod_{j\in J}P_j.
\]
If \(I=[k]\setminus J\), compatibility of the cubical splittings identifies
\[
    \prod_{j\in J}P_j
    \simeq
    s_{I,[k]}\rho_{[k],I}.
\]
Since \(|J|=k-|I|\), this is the displayed alternating formula. Every step is functorial in maps of split cubes. A bijection of finite direction sets only relabels the coherently commuting factors, proving the reindexing statement.
\end{proof}

\begin{definition}[Higher cubical variation coefficient]
\label{def:var-k-coefficient}
Define
\[
    \mathbb V^{(k)}_{M,\mathrm{pre}}(A)
    :=
    \operatorname{Ran}_{\pi_k^{\mathrm{op}}}K_M^{(k)}
\]
on affine points of \(\bar A\), and define
\[
    \boxed{
    \mathbb V^{(k)}_M(A)
    :=
    a_{\mathrm{Nis}}^{\operatorname{Sp}}
    \mathbb V^{(k)}_{M,\mathrm{pre}}(A)
    \in\mathcal E_{\bar A}.
    }
\]
For \(k=1\) this notation agrees canonically with \(\Vfirst_M(A)\); for \(k=2\) write also \(\Vsecond_M(A)\).
\end{definition}

\begin{theorem}[Cubical-difference descent and the raw right-Kan sheaf]
\label{thm:var-k-kan-sheaf}
For every finite \(k\ge1\), the functor \(K_M^{(k)}\) satisfies descent for the centre-pullback Nisnevich topology, and the raw right Kan extension
\[
    \mathbb V^{(k)}_{M,\mathrm{pre}}(A)
\]
is already a spectral Nisnevich sheaf. Hence the sheafification unit is a canonical equivalence
\[
    \boxed{
    \mathbb V^{(k)}_{M,\mathrm{pre}}(A)
    \xrightarrow{\sim}
    \mathbb V^{(k)}_M(A).
    }
\]
At an affine centre \(c:T\to\bar A\), the right-Kan value is the limit over
\[
    \mathcal J_c^{(k)}
    :=
    (c^{\mathrm{op}}\downarrow\pi_k^{\mathrm{op}}),
\]
and every structured \(k\)-variation over a centre mapping to \(c\) receives the corresponding canonical evaluation morphism from that limit.
\end{theorem}

\begin{proof}
Let \(\{T_i\to T\}\) be an affine spectral Nisnevich cover and let \(v_i\) be the Cartesian pullback of a \(k\)-variation \(v\). For every subset \(I\subseteq[k]\), split square-zero base change gives
\[
    T_i[N_{I,i}]
    \simeq
    T[N_I]\times_TT_i.
\]
Effective descent for parameterized spectra therefore identifies \(\Gamma_v(I)\) with the totalization of the corresponding section spectra on the Čech nerve. The total fibre of a finite cube is a finite limit, and finite limits commute with the descent totalization. Hence
\[
    K_M^{(k)}(v)
    \simeq
    \Tot_iK_M^{(k)}(v_i).
\]
Thus \(K_M^{(k)}\) is local for the centre-pullback topology. Lemma~\ref{lem:var-cartesian-centre-kan} applied to \(\pi_k\) proves that the raw right Kan extension is already local. The final evaluation statement is the canonical projection from the right-Kan limit.
\end{proof}

\begin{proposition}[Finite-set reindexing of cubical variations]
\label{prop:var-finite-set-reindexing}
Let \(I\) be a nonempty finite set. Replacing \([k]\) by \(I\) in Definitions~\ref{def:var-k-cube-test}--\ref{def:var-k-coefficient} gives the \(\infty\)-category
\[
    \mathcal V^{(I)}_{\bar A/Q},
\]
the total-fibre coefficient
\[
    K_M^{(I)},
\]
and the right-Kan coefficient
\[
    \mathbb V_M^{(I)}(A).
\]
For every bijection
\[
    \sigma:I\xrightarrow{\sim}I',
\]
relabeling the square-zero directions gives a canonical equivalence
\[
    \sigma_*:
    \mathcal V^{(I)}_{\bar A/Q}
    \xrightarrow{\sim}
    \mathcal V^{(I')}_{\bar A/Q}
\]
over the affine centre category, together with a canonical identification of coefficient functors and hence a canonical equivalence
\[
    \mathbb V_M^{(I)}(A)
    \xrightarrow{\sim}
    \mathbb V_M^{(I')}(A).
\]
These equivalences satisfy the identity and composition coherences for finite bijections and are compatible with coefficient morphisms. For \(I=[k]\), the automorphism group of \(I\) is \(\Sigma_k\), and this reindexing action is the action of Theorem~\ref{thm:var-k-symmetry}.
\end{proposition}

\begin{proof}
For a family \((N_i)_{i\in I}\), the split parameter algebra is the finite tensor product
\[
    B[N_I]
    :=
    \bigotimes_{i\in I}{}_B(B\oplus N_i),
\]
and its face cube is indexed by the power set of \(I\). A bijection \(\sigma:I\to I'\) relabels the tensor factors and the power-set cube. Symmetric monoidal coherence gives the resulting equivalence of parameter algebras and all face maps, coherently under identity and composition of bijections. The same relabeling therefore gives the displayed equivalence of structured variation categories over the unchanged centre.

The section-spectrum cube is carried to the relabeled section-spectrum cube. Since total fibre is functorial in equivalences of finite cubes, this identifies \(K_M^{(I)}\) with the pullback of \(K_M^{(I')}\). Right Kan extension along the centre functors transports this identification to the right-Kan coefficients. Coefficient morphisms commute termwise with the construction, so all stated coherences are inherited from symmetric monoidal reindexing, functoriality of finite limits, and right Kan extension.
\end{proof}

\begin{definition}[Evaluation of a higher variation section]
\label{def:var-k-evaluation}
Let
\[
    e:\one_{\bar A}\longrightarrow\mathbb V_M^{(k)}(A)[n]
\]
be a section and let \(w\in\mathcal V^{(k)}_{\bar A/Q}\) be a structured \(k\)-variation whose centre maps to an affine point \(c:T\to\bar A\). Restriction of \(e\) to \(T\), followed by the canonical projection from the right-Kan limit of Theorem~\ref{thm:var-k-kan-sheaf}, defines
\[
    \boxed{
    \langle e,w\rangle
    \in
    \Omega^\infty K_M^{(k)}(w)[n].
    }
\]
This evaluation is natural in structured morphisms of \(k\)-variations. For \(k=1\) it agrees with Definition~\ref{def:var-euler-evaluation-on-variation}.
\end{definition}

\begin{theorem}[Arbitrary base change of higher cubical variation coefficients]
\label{thm:var-k-base-change}
Let \(c:Q'\to Q\), put \(\bar A':=\bar A\times_QQ'\) and \(M':=c^*M\), and let \(A'\) be the pulled-back Cartan descent object of Convention~\ref{conv:var-arbitrary-cartan-pullback}. Then for every finite \(k\ge1\) there is a canonical equivalence
\[
    \boxed{
    c_{\bar A}^*\mathbb V_M^{(k)}(A)
    \simeq
    \mathbb V_{M'}^{(k)}(A').
    }
\]
It is coherent in composition of base changes and in coefficient morphisms, and is natural under the finite-set reindexing of Proposition~\ref{prop:var-finite-set-reindexing}.
\end{theorem}

\begin{proof}
Let \(t:T\to\bar A'\) be affine and let \(t_0:T\to\bar A\) be its composite. A \(k\)-variation over \(\bar A/Q\) whose centre maps to \(t_0\) carries the structural map
\[
    T'[N_1,\ldots,N_k]
    \xrightarrow{r_{[k]}}
    T'\longrightarrow T\longrightarrow Q'.
\]
Its given map to \(\bar A\) has the same composite to \(Q\). The pullback universal property of \(\bar A'=\bar A\times_QQ'\) therefore lifts the complete \(k\)-variation uniquely through a contractible space. This construction on objects, morphisms, and higher homotopies gives an equivalence of the comma \(\infty\)-categories \(\mathcal J_t^{(k)}\) and \(\mathcal J_{t_0}^{(k)}\). Right Beck--Chevalley identifies every vertex of the coefficient cube and hence its total fibre. Taking the right-Kan limits and using affine-point recognition gives the displayed equivalence. Every step is unchanged by relabeling a finite direction set through a bijection, so the base-change equivalence is natural under Proposition~\ref{prop:var-finite-set-reindexing}.
\end{proof}

\begin{theorem}[Symmetric-group action]
\label{thm:var-k-symmetry}
Permutation of the labelled square-zero directions gives an action of \(\Sigma_k\) on \(\mathcal V^{(k)}_{\bar A/Q}\) over the affine centre category. The symmetric monoidal coherence identifies the coefficient functor with each of its permuted pullbacks. Consequently there is a canonical coherent \(\Sigma_k\)-action
\[
    \boxed{
    \Sigma_k\curvearrowright\mathbb V_M^{(k)}(A).
    }
\]
\end{theorem}

\begin{proof}
Permuting the factors in
\[
    \bigotimes_{i=1}^k(B\oplus N_i)
\]
gives the required equivalence of complete cubes and is coherent for multiplication in \(\Sigma_k\) by the symmetric monoidal coherence. The total-fibre construction is functorial in equivalences of cubes, and right Kan extension transports the resulting action to \(\mathbb V_M^{(k)}(A)\).
\end{proof}


\begin{construction}[Alternating higher variation of a Lagrangian]
\label{con:var-k-variation-lagrangian}
Let
\[
    \ell_L:\one_{\bar A}\longrightarrow a^*M[n]
\]
be the local section of a Lagrangian. For a structured \(k\)-variation \(v\), the coefficient cube
\[
    I\longmapsto\Gamma_v(I)
\]
is coherently split: for \(I\subseteq J\), restriction along \(j_{I,J}\) is split by pullback along the retraction \(p_{J,I}\). Let
\[
    \overline D_v:
    \Gamma_v([k])
    \longrightarrow
    K_M^{(k)}(v)
\]
be the canonical alternating projector of Lemma~\ref{lem:var-split-cube-projector}. Define
\[
    \boxed{
    \Delta_v^{(k)}L
    :=
    \overline D_v\bigl(\widetilde a_{[k]}^*\ell_L\bigr)
    \in
    \Omega^\infty K_M^{(k)}(v)[n].
    }
\]
If \(\ell_I\) denotes the pullback of \(\ell_L\) along \(\widetilde a_I:T[N_I]\to\bar A\), further pulled to the full vertex along \(p_{[k],I}\), then under the total-fibre inclusion the preceding point has full-vertex representative
\[
    \boxed{
    \sum_{I\subseteq[k]}
    (-1)^{k-|I|}\ell_I.
    }
\]
These points are natural in structured \(k\)-variations and therefore assemble, by the right-Kan universal property, to a section
\[
    \boxed{
    \Delta_L^{(k)}:
    \one_{\bar A}
    \longrightarrow
    \mathbb V_M^{(k)}(A)[n].
    }
\]
The section carries the canonical \(\Sigma_k\)-equivariance induced by Theorem~\ref{thm:var-k-symmetry}.
\end{construction}

\begin{proof}
The face maps and retractions of Definition~\ref{def:var-k-cube-test} give coherent splittings of the section-spectrum cube because
\[
    p_{J,I}j_{I,J}=\id_{T[N_I]}
\]
and hence
\[
    j_{I,J}^*p_{J,I}^*\simeq\id.
\]
Lemma~\ref{lem:var-split-cube-projector} therefore gives the canonical map \(\overline D_v\) into the total fibre and identifies its composite with the total-fibre inclusion with the displayed alternating sum. Naturality in structured \(k\)-variations follows from naturality of the split-cube projector and of coefficient pullback. The finite-set reindexing compatibility of the same lemma, together with Theorem~\ref{thm:var-k-symmetry}, gives the stated \(\Sigma_k\)-equivariance.
\end{proof}

\begin{proposition}[The first cubical variation is the Euler source]
\label{prop:var-k1-euler}
Under the canonical identification
\[
    \mathbb V_M^{(1)}(A)\simeq\Vfirst_M(A),
\]
the section \(\Delta_L^{(1)}\) is canonically the local Euler--Lagrange section
\[
    \mathsf{EL}_L:\one_{\bar A}\to\Vfirst_M(A)[n].
\]
\end{proposition}

\begin{proof}
For \(k=1\), the alternating difference is
\[
    \widetilde a^*\ell_L-r_N^*a_T^*\ell_L.
\]
By the Chapter 10 degree-zero normalized differential, this is the restriction of
\[
    e^*\ell_L-c^*\ell_L
\]
along the canonical map from the split square-zero variation to the intrinsic first relation simplex. This is exactly the construction of \(\lambda^1_{\mathrm{loc}}\mathsf E_L^\infty=\mathsf{EL}_L\).
\end{proof}

\subsection{The intrinsic mixed second variation}

For \(k=2\), write
\[
    B[N_1,N_2]
    =
    (B\oplus N_1)\otimes_B(B\oplus N_2),
\]
so that
\[
    B[N_1,N_2]
    \simeq
    B\oplus N_1\oplus N_2\oplus(N_1\otimes_BN_2)
\]
as a \(B\)-module. The complete two-cube is the Cartesian square
\[
\begin{tikzcd}
T[N_1,N_2] \arrow[r] \arrow[d]
& T[N_2] \arrow[d]
\\
T[N_1] \arrow[r]
& T.
\end{tikzcd}
\]
The second variation coefficient is
\[
    \Vsecond_M(A)
    =
    \mathbb V_M^{(2)}(A),
\]
and the alternating higher variation is
\[
    \boxed{
    \Delta_L^{(2),\mathrm{mix}}
    :=
    \Delta_L^{(2)}:
    \one_{\bar A}\to\Vsecond_M(A)[n].
    }
\]
It carries the canonical transposition symmetry from \(\Sigma_2\).

\begin{proposition}[The mixed second variation is not determined by its two sides]
\label{prop:var-mixed-not-hessian}
Even on discrete classical input, the mixed second variation need not depend only on the two side variations. Let \(k\) be a field, take
\[
    X=S:=\operatorname{Spec}Hk,
    \qquad
    Q:=Q_{X/S}\simeq S=\operatorname{Spec}Hk,
\]
and put
\[
    \bar A:=\operatorname{Spec}H(k[x]),
    \qquad
    M:=\Oadd_Q,
    \qquad
    n:=0.
\]
Take the affine centre
\[
    T=\operatorname{Spec}B=Q,
    \qquad
    B=Hk,
    \qquad
    N_1=N_2=Hk.
\]
Then
\[
    B[N_1,N_2]
    \simeq
    H\!\left(
      k[\varepsilon_1,\varepsilon_2]/
      (\varepsilon_1^2,\varepsilon_2^2)
    \right).
\]
Let \(L\) be the degree-zero scalar Lagrangian adjoint to the local coordinate section \(x\). A bi-variation with centre \(x_0\in k\) may be given by the map corresponding to
\[
    x\longmapsto
    x_0+a\varepsilon_1+b\varepsilon_2+c\varepsilon_1\varepsilon_2.
\]
Its two side restrictions depend only on \(a\) and \(b\), whereas
\[
    \boxed{
    \Delta_L^{(2),\mathrm{mix}}
    =
    c\varepsilon_1\varepsilon_2.
    }
\]
Hence complete bi-variations with identical side variations can have different intrinsic mixed second variations.
\end{proposition}

\begin{proof}
The two side restrictions set respectively \(\varepsilon_2=0\) and \(\varepsilon_1=0\), so they forget \(c\). The four corner values are
\[
\begin{aligned}
\ell_{12}&=x_0+a\varepsilon_1+b\varepsilon_2+c\varepsilon_1\varepsilon_2,\\
\ell_1&=x_0+a\varepsilon_1,\\
\ell_2&=x_0+b\varepsilon_2,\\
\ell_0&=x_0,
\end{aligned}
\]
and their alternating sum is \(c\varepsilon_1\varepsilon_2\). In particular, the choices \(c=0\) and \(c=1\) have identical side variations but give distinct elements of
\[
    \pi_0K_M^{(2)}
    \simeq
    k\varepsilon_1\varepsilon_2.
\]
\end{proof}

\subsection{Hessian as a derived descent problem}

\begin{definition}[Side-data category and side-restriction functor]
\label{def:var-hessian-side-category}
Let
\[
    \mathcal P^{(2)}_{\bar A/Q}
    :=
    \mathcal V^{(1)}_{\bar A/Q}
    \times_{\mathcal C^{\mathrm{aff}}_{/\bar A}}
    \mathcal V^{(1)}_{\bar A/Q}.
\]
An object is a pair of first split variations with the same affine centre. In particular it remembers the centre, the two connective modules \(N_1,N_2\), and the two side maps to \(\bar A\), but no mixed lift on \(T[N_1,N_2]\). Restriction of a complete bi-variation to its two one-dimensional sides defines a functor
\[
    \boxed{
    \partial_2:
    \mathcal V^{(2)}_{\bar A/Q}
    \longrightarrow
    \mathcal P^{(2)}_{\bar A/Q}.
    }
\]
\end{definition}

\begin{definition}[Mixed coefficient on side data]
\label{def:var-hessian-side-coefficient}
For a side pair
\[
    p=(T,N_1,N_2,\widetilde a_1,\widetilde a_2)
    \in\mathcal P^{(2)}_{\bar A/Q},
\]
define
\[
    \boxed{
    K^{(2)}_{M,\partial}(p)
    :=
    \fib\left(
      \Gamma_{T[N_1,N_2]}^{\mathrm{st}}(M)
      \longrightarrow
      \Gamma_{T[N_1]}^{\mathrm{st}}(M)
      \times_{\Gamma_T^{\mathrm{st}}(M)}
      \Gamma_{T[N_2]}^{\mathrm{st}}(M)
    \right).
    }
\]
The structural map of every vertex to \(Q\) is induced from the common centre, so this coefficient is defined without choosing a mixed lift to \(\bar A\). Structured morphisms of side pairs make
\[
    K^{(2)}_{M,\partial}:
    (\mathcal P^{(2)}_{\bar A/Q})^{\mathrm{op}}
    \longrightarrow
    \operatorname{Sp}
\]
a functor. There is a canonical identification
\[
    \boxed{
    \partial_2^*K^{(2)}_{M,\partial}
    \simeq
    K_M^{(2)}.
    }
\]
\end{definition}

\begin{proof}
The total-fibre coefficient uses only the common affine centre, the two modules, their split square-zero parameter cube, and the pullback of \(M\) from \(Q\). It does not use the mixed map \(\widetilde a_{12}:T[N_1,N_2]\to\bar A\). Hence the same four-corner spectrum is functorially attached to the side pair, and pullback along \(\partial_2\) recovers exactly the coefficient already denoted \(K_M^{(2)}\).
\end{proof}

Write \(\mathbb S\) for the constant sphere-spectrum functor on either of the variation indexing \(\infty\)-categories appearing below.

\begin{lemma}[Global sections of a right-Kan cubical coefficient]
\label{lem:var-global-sections-right-kan}
Let
\[
    \mathcal C:=\mathcal C^{\mathrm{aff}}_{/\bar A},
\]
let \(\mathbb S_{\mathcal C}\) be the constant sphere-spectrum presheaf on \(\mathcal C\), and let \(k\ge1\). There is a canonical natural equivalence
\[
\begin{aligned}
    \Map_{\mathcal E_{\bar A}}
    \bigl(
      \one_{\bar A},
      \mathbb V_M^{(k)}(A)[n]
    \bigr)
    \simeq
    \Map_{\operatorname{Fun}((\mathcal V^{(k)}_{\bar A/Q})^{\mathrm{op}},\operatorname{Sp})}
    \bigl(
      \mathbb S,
      K_M^{(k)}[n]
    \bigr).
\end{aligned}
\]
Under this equivalence, a section of \(\mathbb V_M^{(k)}(A)[n]\) is exactly the corresponding coherent natural section of the raw cubical coefficient functor.
\end{lemma}

\begin{proof}
By the Chapter 10 affine-point presentation,
\[
    \mathcal E_{\bar A}
    \simeq
    \operatorname{Shv}_{\mathrm{Nis}}
    (\mathcal C;\operatorname{Sp}),
\]
and the tensor unit corresponds to the spectral Nisnevich sheafification
\[
    a_{\mathrm{Nis}}^{\operatorname{Sp}}\mathbb S_{\mathcal C}
\]
of the constant sphere presheaf. By Theorem~\ref{thm:var-k-kan-sheaf}, the raw right Kan extension
\[
    \operatorname{Ran}_{\pi_k^{\mathrm{op}}}K_M^{(k)}
\]
is already a spectral Nisnevich sheaf and represents \(\mathbb V_M^{(k)}(A)\). Therefore the sheafification adjunction gives
\[
\begin{aligned}
    \Map_{\mathcal E_{\bar A}}
    \bigl(
      \one_{\bar A},
      \mathbb V_M^{(k)}(A)[n]
    \bigr)
    &\simeq
    \Map_{\operatorname{Fun}(\mathcal C^{\mathrm{op}},\operatorname{Sp})}
    \left(
      \mathbb S_{\mathcal C},
      \operatorname{Ran}_{\pi_k^{\mathrm{op}}}K_M^{(k)}[n]
    \right).
\end{aligned}
\]
The right-Kan adjunction identifies the last mapping space with
\[
    \Map_{\operatorname{Fun}((\mathcal V^{(k)}_{\bar A/Q})^{\mathrm{op}},\operatorname{Sp})}
    \left(
      (\pi_k^{\mathrm{op}})^*\mathbb S_{\mathcal C},
      K_M^{(k)}[n]
    \right).
\]
The pullback of the constant sphere presheaf is the constant sphere functor \(\mathbb S\), which gives the displayed equivalence.
\end{proof}

\begin{lemma}[Restriction of cubical right-Kan coefficients]
\label{lem:var-restricted-cubical-right-kan}
Let
\[
    j:C\longrightarrow\bar A
\]
be any morphism in \(\XSp\), and let \(k\ge1\). Define
\[
    \mathcal V_C^{(k)}
    :=
    \mathcal V^{(k)}_{\bar A/Q}
    \times_{\mathcal C^{\mathrm{aff}}_{/\bar A}}
    \mathcal C^{\mathrm{aff}}_{/C},
\]
write
\[
    \pi_{k,C}:
    \mathcal V_C^{(k)}
    \longrightarrow
    \mathcal C^{\mathrm{aff}}_{/C}
\]
for the restricted centre functor, and let
\[
    K_{M,C}^{(k)}
    :=
    K_M^{(k)}\big|_{(\mathcal V_C^{(k)})^{\mathrm{op}}}.
\]
Then there is a canonical equivalence in \(\mathcal E_C\)
\[
    \boxed{
    j^*\mathbb V_M^{(k)}(A)
    \simeq
    \operatorname{Ran}_{\pi_{k,C}^{\mathrm{op}}}
    K_{M,C}^{(k)}.
    }
\]
The raw right Kan extension on the right is already a spectral Nisnevich sheaf on
\(\mathcal C^{\mathrm{aff}}_{/C}\).  The equivalence is natural in \(j\), in coefficient morphisms, and in finite-set reindexing of the cubical directions.
\end{lemma}

\begin{proof}
Let
\[
    c:T\longrightarrow C
\]
be an affine point.  Lemma~\ref{lem:var-affine-point-pullback} and Theorem~\ref{thm:var-k-kan-sheaf} give
\[
\begin{aligned}
    \bigl(j^*\mathbb V_M^{(k)}(A)\bigr)(T,c)
    &\simeq
    \mathbb V_M^{(k)}(A)(T,jc)\\
    &\simeq
    \lim_{
      ((jc)^{\mathrm{op}}\downarrow\pi_k^{\mathrm{op}})
    }
    K_M^{(k)}.
\end{aligned}
\]
On the other hand, the value of the displayed restricted right Kan extension at \(c\) is
\[
    \lim_{
      (c^{\mathrm{op}}\downarrow\pi_{k,C}^{\mathrm{op}})
    }
    K_{M,C}^{(k)}.
\]

There is a canonical equivalence of the two comma \(\infty\)-categories
\[
    \boxed{
    (c^{\mathrm{op}}\downarrow\pi_{k,C}^{\mathrm{op}})
    \simeq
    ((jc)^{\mathrm{op}}\downarrow\pi_k^{\mathrm{op}}).
    }
\]
Indeed, an object of the category on the right is a complete structured \(k\)-variation \(v\) with affine centre
\[
    d:T'\longrightarrow\bar A
\]
together with a morphism
\[
    g:T'\longrightarrow T
\]
of affine centres over \(\bar A\), so that
\[
    d\simeq jc\,g.
\]
The composite
\[
    c\,g:T'\longrightarrow C
\]
therefore lifts the centre \(d\) through \(j\), and together with the same map \(g\) defines an object of the category on the left.  Conversely, forgetting the specified lift of the centre through \(C\) gives an object of the category on the right.  In the presence of the map \(g\), the lift through \(C\) is canonically \(cg\); hence these two constructions are inverse through a contractible space, on objects, morphisms, and all higher simplices.

Under this comma-category equivalence the two coefficient diagrams are literally the same restricted cubical-difference diagram.  Their limits are therefore canonically equivalent.  These equivalences are natural under affine restriction in \(C\), so they identify the underlying affine-point presheaves.

The left side is a parameterized spectrum on \(C\), hence its affine-point presheaf is a spectral Nisnevich sheaf.  Therefore the raw right Kan extension on the right is already a sheaf.  Affine-point recognition gives the displayed equivalence in \(\mathcal E_C\).  Naturality in \(j\), coefficients, and finite-set reindexing follows from the same comma-category construction.
\end{proof}

\begin{corollary}[Restricted global sections of cubical right-Kan coefficients]
\label{cor:var-restricted-global-sections-right-kan}
With the notation of Lemma~\ref{lem:var-restricted-cubical-right-kan}, there is a canonical natural equivalence
\[
\begin{aligned}
    &
    \Map_{\mathcal E_C}
    \left(
      \one_C,
      j^*\mathbb V_M^{(k)}(A)[n]
    \right)
    \\
    &\qquad\simeq
    \Map_{
      \operatorname{Fun}(
        (\mathcal V_C^{(k)})^{\mathrm{op}},
        \operatorname{Sp}
      )
    }
    \left(
      \mathbb S,
      K_{M,C}^{(k)}[n]
    \right).
\end{aligned}
\]
Under this equivalence, a parameterized section on \(C\) is exactly its coherent natural section on the restricted complete \(k\)-variation category.
\end{corollary}

\begin{proof}
By Lemma~\ref{lem:var-restricted-cubical-right-kan}, the target is represented by the raw right Kan extension and that right Kan extension is already a spectral Nisnevich sheaf.  The spectral Nisnevich sheafification adjunction therefore identifies maps from the tensor unit with maps from the constant sphere presheaf into that raw right Kan extension.  The right-Kan adjunction then gives
\[
\begin{aligned}
    &
    \Map_{\operatorname{Fun}(
      (\mathcal C^{\mathrm{aff}}_{/C})^{\mathrm{op}},
      \operatorname{Sp})}
    \left(
      \mathbb S_{\mathcal C^{\mathrm{aff}}_{/C}},
      \operatorname{Ran}_{\pi_{k,C}^{\mathrm{op}}}
      K_{M,C}^{(k)}[n]
    \right)
    \\
    &\qquad\simeq
    \Map_{
      \operatorname{Fun}(
        (\mathcal V_C^{(k)})^{\mathrm{op}},
        \operatorname{Sp}
      )
    }
    \left(
      \mathbb S,
      K_{M,C}^{(k)}[n]
    \right),
\end{aligned}
\]
because pullback of the constant sphere presheaf is again constant.  This is the claimed equivalence.
\end{proof}

\begin{definition}[Derived Hessian-descent space]
\label{def:var-hessian-descent-space}
Under Lemma~\ref{lem:var-global-sections-right-kan}, the mixed second variation is equivalently the natural section
\[
    \Delta_L^{(2),\mathrm{mix}}
    \in
    \Map_{\operatorname{Fun}((\mathcal V^{(2)})^{\mathrm{op}},\operatorname{Sp})}
    \bigl(\mathbb S,K_M^{(2)}[n]\bigr).
\]
Pullback along \(\partial_2\), followed by the canonical identification of Definition~\ref{def:var-hessian-side-coefficient}, gives
\[
\begin{aligned}
    \partial_2^*:
    &\Map_{\operatorname{Fun}((\mathcal P^{(2)})^{\mathrm{op}},\operatorname{Sp})}
    \bigl(\mathbb S,K^{(2)}_{M,\partial}[n]\bigr)
    \\
    &\longrightarrow
    \Map_{\operatorname{Fun}((\mathcal V^{(2)})^{\mathrm{op}},\operatorname{Sp})}
    \bigl(\mathbb S,K_M^{(2)}[n]\bigr).
\end{aligned}
\]
Define
\[
    \boxed{
    \HessDesc(L)
    :=
    \hofib_{\Delta_L^{(2),\mathrm{mix}}}(\partial_2^*).
    }
\]
\end{definition}

\begin{proposition}[Meaning of Hessian descent]
\label{prop:var-hessian-descent-meaning}
A point of \(\HessDesc(L)\) is a coherent section of the mixed total-fibre coefficient depending only on the two side variations, together with a specified equivalence between its pullback to complete bi-variations and the intrinsic mixed second variation. Thus dependence only on two first variation directions is additional descent data, not part of the definition of second variation. A point of \(\HessDesc(L)\) is only the first categorical prerequisite for a Hessian; it does not by itself assert additivity or bilinearity in the two directions.

The affine example of Proposition~\ref{prop:var-mixed-not-hessian} proves that this descent need not exist: two complete bi-variations with the same side-data object have the same target spectrum \(K^{(2)}_{M,\partial}\) but the mixed second variation takes the distinct values determined by their different mixed coefficients.
\end{proposition}

\begin{proof}
The mapping space in the source classifies coherent natural sections of a functor defined entirely on pairs of side variations. Pullback along \(\partial_2\) evaluates such a section on every complete bi-variation through its two sides. The homotopy fibre over \(\Delta_L^{(2),\mathrm{mix}}\) therefore classifies precisely the desired factorization together with all comparison homotopies.

In Proposition~\ref{prop:var-mixed-not-hessian}, fixing \(x_0,a,b\) fixes the side pair, the modules, and hence the spectrum \(K^{(2)}_{M,\partial}\), while changing \(c\) changes the value \(c\varepsilon_1\varepsilon_2\). A section pulled back from the side-data category has one value on that fixed side pair, so it cannot reproduce two different intrinsic mixed values. Hence the descent space is empty for that unrestricted family of examples.
\end{proof}

\begin{definition}[On-shell mixed second variation]
\label{def:var-jacobi}
Let
\[
    j_L:\Crit(L)\longrightarrow\bar A
\]
be the first Euler critical locus. Define
\[
    \boxed{
    \Vsecond_{M,L}
    :=
    j_L^*\Vsecond_M(A),
    \qquad
    \Delta_{L,\mathrm{crit}}^{(2),\mathrm{mix}}
    :=
    j_L^*\Delta_L^{(2),\mathrm{mix}}.
    }
\]
The section \(\Delta_{L,\mathrm{crit}}^{(2),\mathrm{mix}}\) is the \emph{on-shell mixed second variation}. It inherits the canonical transposition symmetry. At this stage no Jacobi operator is inferred from the pullback alone.  Subsection~\ref{subsec:var-represented-jacobi} constructs the required critical side descent and Jacobi operator in the represented smooth scalar sector; outside that sector the present definition remains only the intrinsic on-shell mixed variation.
\end{definition}

\begin{remark}[The coherent square of the variational differential is different]
\label{rem:var-hessian-not-dsquare}
The nullhomotopy
\[
    (\delta^{\infty}_{\mathrm{var}})^2L\simeq0
\]
is the coherent Helmholtz relation in the normalized de Rham direction. The cubical second variation is the four-term cross-effect on a complete split two-cube. The former is an antisymmetrized cochain relation; the latter retains the independent mixed \(N_1\otimes_BN_2\) lift. Their distinction is structural and persists before any possible Hessian descent.
\end{remark}


\subsection{Represented critical Hessian and Jacobi linearization}
\label{subsec:var-represented-jacobi}

The unrestricted mixed second variation remains genuinely cubical.  The present subsection does not change that fact.  It isolates the representably smooth scalar sector in which the first Euler coefficient is a finite locally free cotangent coefficient and the Chapter 2 structured square-zero lifting calculus controls the mixed \(N_1\otimes N_2\) layer.  The mixed boundary extension, its classifying derivation, the smooth fill space, the linear coefficient category used for duality, and the critical descent are all constructed below.  Criticality kills precisely the Euler-controlled dependence on the mixed lift.  Bilinearity and the Jacobi operator are then consequences of the constructed affine-point linear refinement rather than additional axioms.

\begin{construction}[Affine-point linear coefficient category]
\label{con:var-linear-coefficient-category}
Let \(Z\in\XSp\).  Define
\[
    \boxed{
    \mathcal Q_Z
    :=
    \operatorname{QCoh}(Z).
    }
\]
Chapter 1 defines this category for every spectral prestack by affine-point right Kan extension of the module coefficient system.  Equivalently, if \(\mathcal C^{\mathrm{aff}}_{/Z}\) is the affine-point category, then \(\mathcal Q_Z\) is the \(\infty\)-category of Cartesian sections of
\[
    c:T=\operatorname{Spec}B\longmapsto\operatorname{Mod}_B,
\]
whose transition functor along
\[
    g:(T'=\operatorname{Spec}B',c')\longrightarrow(T=\operatorname{Spec}B,c)
\]
is the exact strong symmetric monoidal extension of scalars
\[
    g^*:
    \operatorname{Mod}_B
    \longrightarrow
    \operatorname{Mod}_{B'},
    \qquad
    P\longmapsto B'\otimes_BP.
\]
Concretely, an object \(P\in\mathcal Q_Z\) consists of modules
\[
    P_c\in\operatorname{Mod}_B
\]
for all affine points, together with specified coherent equivalences
\[
    B'\otimes_BP_c
    \xrightarrow{\sim}
    P_{c'}
\]
for every morphism of affine points, satisfying the identity, composition, and all higher compatibility relations.  The pointwise tensor products over the coordinate rings make \(\mathcal Q_Z\) a stable symmetric monoidal \(\infty\)-category.  Its tensor unit is denoted
\[
    \mathcal O_Z^{\mathrm{lin}},
    \qquad
    (\mathcal O_Z^{\mathrm{lin}})_c:=B.
\]

For \(P\in\mathcal Q_Z\), forget the module structures and use the unit maps
\[
    P_c\longrightarrow B'\otimes_BP_c\xrightarrow{\sim}P_{c'}
\]
as affine restriction maps.  This gives a spectrum-valued presheaf
\[
    U_Z^{\mathrm{pre}}(P)
    \in
    \operatorname{Fun}\bigl((\mathcal C^{\mathrm{aff}}_{/Z})^{\mathrm{op}},\operatorname{Sp}\bigr).
\]
Lemma~\ref{lem:var-affine-nis-module-descent} implies that this presheaf is already a spectral Nisnevich sheaf.  Define
\[
    \boxed{
    U_Z(P)
    :=
    U_Z^{\mathrm{pre}}(P)
    \in
    \mathcal E_Z.
    }
\]
Equivalently, the spectral Nisnevich sheafification unit
\[
    U_Z^{\mathrm{pre}}(P)
    \longrightarrow
    a_{\mathrm{Nis}}^{\operatorname{Sp}}U_Z^{\mathrm{pre}}(P)
\]
is an equivalence.  The functor \(U_Z\) is exact.  The unit family is the affine-point restriction of the Chapter 10 scalar coefficient, so
\[
    \boxed{
    U_Z(\mathcal O_Z^{\mathrm{lin}})
    \simeq
    \Oadd_Z.
    }
\]

For a morphism \(f:Z'\to Z\), composition of affine points with \(f\) gives a pullback functor
\[
    f^*:\mathcal Q_Z\longrightarrow\mathcal Q_{Z'},
\]
coherently compatible with identities and composition, and there is a canonical equivalence
\[
    \boxed{
    U_{Z'}f^*\simeq f^*U_Z.
    }
\]
\end{construction}

\begin{proof}
The identification \(\mathcal Q_Z=\operatorname{QCoh}(Z)\) and its Cartesian affine-point description are the Chapter 1 right-Kan definition of quasi-coherent modules on a spectral prestack.  Stability, the symmetric monoidal structure, and exact strong symmetric monoidality of pullback are therefore inherited from that coefficient system.

Let \(P\in\mathcal Q_Z\), let
\[
    c:T=\operatorname{Spec}B\longrightarrow Z
\]
be an affine point, and let
\[
    \{T_i=\operatorname{Spec}B_i\longrightarrow T\}_{i\in I}
\]
be a finite affine spectral Nisnevich cover.  On the iterated intersection indexed by \((i_0,\ldots,i_r)\), Cartesianity gives
\[
    P_{c_{i_0\cdots i_r}}
    \simeq
    \bigl(B_{i_0}\otimes_B\cdots\otimes_BB_{i_r}\bigr)\otimes_BP_c.
\]
The objectwise assertion of Lemma~\ref{lem:var-affine-nis-module-descent} therefore gives
\[
    U_Z^{\mathrm{pre}}(P)(T,c)
    \xrightarrow{\sim}
    \Tot_{[r]\in\Delta}
    \prod_{(i_0,\ldots,i_r)\in I^{r+1}}
    U_Z^{\mathrm{pre}}(P)(T_{i_0\cdots i_r},c_{i_0\cdots i_r}).
\]
The finite affine basis theorem of Chapter 9 says that these families generate the affine spectral Nisnevich topology, so \(U_Z^{\mathrm{pre}}(P)\) is already a spectral Nisnevich sheaf.

A fibre or cofiber sequence in \(\mathcal Q_Z\) is computed affine-pointwise in the module categories, and forgetting to spectra preserves finite limits and finite colimits.  Hence \(U_Z\) is exact without any further sheafification argument.

At an affine point \(c:T=\operatorname{Spec}B\to Z\), the unit family has value \(B\) with the ordinary scalar restriction maps.  Chapter 10 identifies these affine-point values, including their multiplicative naturality, with the restriction of \(\Oadd_Z\), which proves
\[
    U_Z(\mathcal O_Z^{\mathrm{lin}})\simeq\Oadd_Z.
\]

Finally, let \(f:Z'\to Z\) and let \(t:T\to Z'\) be an affine point.  By the affine-point definitions,
\[
\begin{aligned}
    \bigl(U_{Z'}f^*P\bigr)(T,t)
    &\simeq
    (f^*P)_t\\
    &\simeq
    P_{ft}\\
    &\simeq
    U_Z(P)(T,ft)\\
    &\simeq
    \bigl(f^*U_Z(P)\bigr)(T,t).
\end{aligned}
\]
These equivalences are natural under affine restriction and carry the identity, composition, and pullback-pasting coherences of quasi-coherent pullback.  Affine-point recognition in \(\mathcal E_{Z'}\) gives the asserted equivalence \(U_{Z'}f^*\simeq f^*U_Z\).
\end{proof}

\begin{lemma}[Finite affine spectral Nisnevich descent for modules]
\label{lem:var-affine-nis-module-descent}
Let
\[
    \{T_i=\operatorname{Spec}B_i\longrightarrow
      T=\operatorname{Spec}B\}_{i\in I}
\]
be a finite affine spectral Nisnevich covering family, and let \(T_\bullet\to T\) be its \v{C}ech nerve.  Then extension of scalars induces a canonical symmetric monoidal equivalence
\[
    \boxed{
    \operatorname{Mod}_B
    \xrightarrow{\sim}
    \Tot_{[r]\in\Delta}
    \prod_{(i_0,\ldots,i_r)\in I^{r+1}}
    \operatorname{Mod}_{
      B_{i_0}\otimes_B\cdots\otimes_BB_{i_r}
    }.
    }
\]
Equivalently,
\[
    \operatorname{QCoh}(T)
    \xrightarrow{\sim}
    \Tot_{[r]\in\Delta}\operatorname{QCoh}(T_r).
\]
Moreover, for every \(P\in\operatorname{Mod}_B\), forgetting the module structures gives a canonical equivalence of spectra, where \(U\) denotes the underlying-spectrum functor in each affine degree:
\[
    \boxed{
    U(P)
    \xrightarrow{\sim}
    \Tot_{[r]\in\Delta}
    \prod_{(i_0,\ldots,i_r)\in I^{r+1}}
    U\!\left(
      P\otimes_B
      B_{i_0}\otimes_B\cdots\otimes_BB_{i_r}
    \right).
    }
\]
The equivalences are natural under refinement and affine base change.
\end{lemma}

\begin{proof}
Every member \(T_i\to T\) of the covering family is spectral \(\acute{e}\)tale, and the family is jointly surjective on the underlying topological space.  Hence it is a covering family for the spectral \(\acute{e}\)tale topology.  The spectral \(\acute{e}\)tale module-descent theorem says that the affine coefficient system
\[
    \operatorname{Spec}A\longmapsto\operatorname{Mod}_A
\]
is a stack of symmetric monoidal stable \(\infty\)-categories for this topology.  Applied to the present \v{C}ech nerve, it gives a canonical symmetric monoidal equivalence
\[
    \operatorname{Mod}_B
    \xrightarrow{\sim}
    \Tot_{[r]\in\Delta}\operatorname{QCoh}(T_r).
\]

Because the covering family is finite, every term of the \v{C}ech nerve is the finite coproduct
\[
    T_r
    \simeq
    \coprod_{(i_0,\ldots,i_r)\in I^{r+1}}
    \operatorname{Spec}
    \bigl(
      B_{i_0}\otimes_B\cdots\otimes_BB_{i_r}
    \bigr).
\]
A finite coproduct of affine spectral schemes is affine with coordinate algebra the corresponding finite product, and the Chapter 1 decomposition of modules over a finite product gives
\[
    \operatorname{QCoh}(T_r)
    \simeq
    \prod_{(i_0,\ldots,i_r)\in I^{r+1}}
    \operatorname{Mod}_{B_{i_0}\otimes_B\cdots\otimes_BB_{i_r}}.
\]
Substitution gives the displayed categorical equivalence.

It remains to record the objectwise statement used below.  Let \(P\in\operatorname{Mod}_B\), and let \(K\in\operatorname{Sp}\).  Mapping spectra in a limit of stable \(\infty\)-categories are computed by the corresponding limit, so the descent equivalence gives
\[
\begin{aligned}
    \operatorname{map}_{\operatorname{Sp}}(K,U(P))
    &\simeq
    \operatorname{map}_{\operatorname{Mod}_B}(B\otimes K,P)
    \\
    &\simeq
    \Tot_{[r]\in\Delta}
    \prod_{(i_0,\ldots,i_r)\in I^{r+1}}
    \operatorname{map}_{\operatorname{Mod}_{B_{i_0}\otimes_B\cdots\otimes_BB_{i_r}}}
    \left(
      (B_{i_0}\otimes_B\cdots\otimes_BB_{i_r})\otimes K,
      P\otimes_BB_{i_0}\otimes_B\cdots\otimes_BB_{i_r}
    \right)
    \\
    &\simeq
    \operatorname{map}_{\operatorname{Sp}}
    \left(
      K,
      \Tot_{[r]\in\Delta}
      \prod_{(i_0,\ldots,i_r)\in I^{r+1}}
      U\!\left(
        P\otimes_BB_{i_0}\otimes_B\cdots\otimes_BB_{i_r}
      \right)
    \right).
\end{aligned}
\]
This equivalence is natural in \(K\).  Stable Yoneda therefore gives the displayed equivalence of underlying spectra.  Naturality under refinement and affine base change is inherited from functoriality of spectral \(\acute{e}\)tale descent and extension of scalars.
\end{proof}

\begin{proposition}[Effective descent for affine-point linear coefficients]
\label{prop:var-linear-coefficient-effective-descent}
Let
\[
    e:Y\twoheadrightarrow Z
\]
be an effective epimorphism in \(\XSp\), with \v{C}ech nerve
\[
    Y_\bullet,\qquad
    Y_r:=\underbrace{Y\times_Z\cdots\times_ZY}_{r+1\text{ factors}}.
\]
Then restriction induces a canonical symmetric monoidal equivalence
\[
    \boxed{
    \mathcal Q_Z
    \xrightarrow{\sim}
    \Tot_{[r]\in\Delta}\mathcal Q_{Y_r}.
    }
\]
The equivalence is natural in morphisms of effective-epimorphism diagrams and is compatible with pullback.
\end{proposition}

\begin{proof}
Restriction along the maps \(Y_r\to Z\) gives the comparison functor.  We construct its inverse directly on affine points.

Let
\[
    P_\bullet
    \in
    \Tot_{[r]\in\Delta}\mathcal Q_{Y_r}
\]
be a coherent descent datum, and let
\[
    c:T=\operatorname{Spec}B\longrightarrow Z
\]
be an affine point.  Pulling back \(e\) gives an effective epimorphism
\[
    Y_T:=T\times_ZY\twoheadrightarrow T.
\]
Effective epimorphisms of spectral Nisnevich sheaves are locally surjective on affine tests.  Hence the identity section of \(T\) lifts to \(Y_T\) over a Nisnevich covering sieve of \(T\).  By the Chapter 9 finite affine basis theorem, choose inside that sieve a finite affine spectral Nisnevich covering family
\[
    \{T_i=\operatorname{Spec}B_i\longrightarrow T\}_{i\in I}
\]
and compatible lifts
\[
    \widetilde c_i:T_i\longrightarrow Y
\]
of \(c|_{T_i}\).

The degree-zero object \(P_0\in\mathcal Q_Y\) supplies modules
\[
    P_i:=(P_0)_{\widetilde c_i}\in\operatorname{Mod}_{B_i}.
\]
On an iterated intersection
\[
    T_{i_0\cdots i_r}
    :=
    T_{i_0}\times_T\cdots\times_TT_{i_r},
\]
the tuple of lifted points determines a canonical map
\[
    T_{i_0\cdots i_r}
    \longrightarrow
    Y_r.
\]
Evaluation of the degree-\(r\) part of \(P_\bullet\) at this affine point, together with the cosimplicial compatibility of \(P_\bullet\), gives the complete coherent \v{C}ech descent datum for the modules \(P_i\).  Lemma~\ref{lem:var-affine-nis-module-descent} therefore descends them to a module
\[
    P_c\in\operatorname{Mod}_B
\]
through a contractible space of choices.

The construction is independent of the chosen lifting cover.  Indeed, two finite affine lifting covers admit a common finite affine Nisnevich refinement.  After restriction to that refinement the two descended modules induce the same module descent datum, and full faithfulness in Lemma~\ref{lem:var-affine-nis-module-descent} identifies them through a contractible space.

Now let
\[
    g:(T'=\operatorname{Spec}B',c')
    \longrightarrow
    (T=\operatorname{Spec}B,c)
\]
be a morphism of affine points of \(Z\).  Pulling a chosen lifting cover of \(T\) back along \(g\) gives a finite affine Nisnevich lifting cover of \(T'\).  Both
\[
    B'\otimes_BP_c
    \qquad\text{and}\qquad
    P_{c'}
\]
descend the same pulled-back module datum.  Full faithfulness of affine Nisnevich module descent gives a canonical equivalence
\[
    \boxed{
    B'\otimes_BP_c
    \xrightarrow{\sim}
    P_{c'}.
    }
\]
Uniqueness in descent supplies the identity, composition, and all higher compatibility relations.  Thus
\[
    c\longmapsto P_c
\]
defines an object \(P\in\mathcal Q_Z\).

Restricting this reconstructed object to \(Y_\bullet\) recovers \(P_\bullet\), because this is true after every affine pullback and affine Nisnevich lifting cover.  Conversely, beginning with an object of \(\mathcal Q_Z\), restricting it to \(Y_\bullet\), and applying the construction above recovers its module at every affine point of \(Z\).  The two constructions are therefore inverse equivalences.

Tensor products and units are computed pointwise and affine Nisnevich module descent is symmetric monoidal.  Hence the equivalence is symmetric monoidal.  Naturality and pullback compatibility follow by pulling back the lifting covers and using the uniqueness clauses in the same descent theorem.
\end{proof}

\begin{convention}[Represented smooth scalar Jacobi sector]
\label{conv:var-jacobi-sector}
For the remainder of this subsection take
\[
    M=\Oadd_Q
\]
and suppose that
\[
    a:\bar A\longrightarrow Q
\]
is representable and \emph{representably smooth}: for every represented affine test
\[
    t:T\longrightarrow Q,
\]
the pullback
\[
    \bar A_T:=T\times_Q\bar A\longrightarrow T
\]
is a smooth spectral morphism in the sense of Chapter 3.  This condition is preserved by arbitrary pullback in \(Q\).  For an affine point \(c:T=\operatorname{Spec}B\to\bar A\), put
\[
    N_c
    :=
    \Gamma(T,s_c^*L_{\bar A_c/T}).
\]
Proposition~\ref{prop:var-universal-first-cartesian} supplies coherent base-change equivalences
\[
    B'\otimes_BN_c
    \xrightarrow{\sim}
    N_{c'}
\]
under every morphism of affine points.  Hence the family \(c\mapsto N_c\) defines a canonical object
\[
    \boxed{
    L^{\mathrm{lin}}_{\bar A/Q}
    \in
    \mathcal Q_{\bar A}.
    }
\]
By Proposition~\ref{prop:var-scalar-affine-cotangent} and affine-point recognition,
\[
    \boxed{
    U_{\bar A}\bigl(L^{\mathrm{lin}}_{\bar A/Q}\bigr)
    \simeq
    L^{\mathrm{st}}_{\bar A/Q}
    \simeq
    \Vfirst(A).
    }
\]

Because \(a\) is representably smooth, every \(N_c\) is finite locally free over \(B\).  The dual modules \(N_c^\vee\), their evaluation and coevaluation maps, and their triangle homotopies commute with extension of scalars.  They therefore form the dual of \(L^{\mathrm{lin}}_{\bar A/Q}\) in \(\mathcal Q_{\bar A}\).  Write
\[
    \boxed{
    T^{\mathrm{lin}}_{\bar A/Q}
    :=
    \bigl(L^{\mathrm{lin}}_{\bar A/Q}\bigr)^\vee
    }
\]
and define its underlying parameterized tangent coefficient by
\[
    \boxed{
    T^{\mathrm{st}}_{\bar A/Q}
    :=
    U_{\bar A}\bigl(T^{\mathrm{lin}}_{\bar A/Q}\bigr).
    }
\]
Thus
\[
    T^{\mathrm{lin}}_{\bar A/Q}(T,c)
    \simeq
    N_c^\vee.
\]
No tangent object is used outside this represented smooth recognition sector.
\end{convention}

\begin{convention}[Module-valued scalar cubical total fibres]
\label{conv:var-scalar-cubical-module-fibre}
Let
\[
    v=(T=\operatorname{Spec}B,N_1,\ldots,N_k,\widetilde a_{[k]})
    \in
    \mathcal V^{(k)}_{\bar A/Q}.
\]
For scalar coefficients, the Chapter 10 represented scalar-section equivalence identifies every vertex of the complete section-spectrum cube with the commutative \(B\)-algebra
\[
    B[N_I].
\]
All cubical restriction maps are morphisms of \(B\)-algebras and hence \(B\)-linear.  We therefore regard the scalar coefficient cube canonically as a cube in \(\operatorname{Mod}_B\) and define
\[
    K_{\Oadd,B}^{(k)}(v)
    :=
    \operatorname{TotFib}_{\operatorname{Mod}_B}
    \bigl(I\longmapsto B[N_I]\bigr).
\]
The forgetful functor
\[
    \operatorname{Mod}_B\longrightarrow\operatorname{Sp}
\]
creates limits, so there is a canonical equivalence of underlying spectra
\[
    U\bigl(K_{\Oadd,B}^{(k)}(v)\bigr)
    \simeq
    K_{\Oadd}^{(k)}(v)
\]
with the coefficient of Definition~\ref{def:var-k-difference-functor}.  We use the notation
\[
    K_{\Oadd}^{(k)}(v)
\]
for this canonically refined \(B\)-module whenever scalar coefficients are in force.  The same convention applies to the side coefficient \(K^{(2)}_{\Oadd,\partial}\).
\end{convention}

\begin{proposition}[Scalar total fibre of a split cubical variation]
\label{prop:var-scalar-cubical-total-fibre}
Let
\[
    v=(T=\operatorname{Spec}B,N_1,\ldots,N_k,\widetilde a_{[k]})
    \in
    \mathcal V^{(k)}_{\bar A/Q}
\]
be a finite structured cubical variation.  For scalar coefficients there is a canonical natural equivalence of \(B\)-modules
\[
    \boxed{
    K_{\Oadd}^{(k)}(v)
    \simeq
    N_1\otimes_B\cdots\otimes_BN_k.
    }
\]
It is compatible with structured morphisms of cubical variations and with finite-set reindexing.  In particular, for side data
\[
    p=(T,N_1,N_2,\widetilde a_1,\widetilde a_2)
    \in\mathcal P^{(2)}_{\bar A/Q},
\]
there is a canonical equivalence
\[
    \boxed{
    K^{(2)}_{\Oadd,\partial}(p)
    \simeq
    N_1\otimes_BN_2.
    }
\]
\end{proposition}

\begin{proof}
For every \(I\subseteq[k]\), the represented scalar-section calculation of Chapter 10 gives
\[
    \Gamma_{T[N_I]}^{\mathrm{st}}(\Oadd_{T[N_I]})
    \simeq
    B[N_I].
\]
Proposition~\ref{prop:var-k-cube-module} identifies the resulting cube, as a cube of \(B\)-modules, with
\[
    I
    \longmapsto
    \bigoplus_{J\subseteq I}N_J,
    \qquad
    N_J:=\bigotimes_{j\in J}N_j,
\]
where \(N_\varnothing=B\).  The restriction from \(I\) to \(I'\subseteq I\) is the identity on \(N_J\) when \(J\subseteq I'\), and is zero on that summand otherwise.

Finite total fibres commute with finite direct sums in \(\operatorname{Mod}_B\).  Fix \(J\subsetneq[k]\) and choose \(i\notin J\).  In the \(i\)-th cubical direction the \(N_J\)-summand is carried by an identity map, so its total fibre is zero.  The only summand without such an identity direction is
\[
    N_{[k]}=N_1\otimes_B\cdots\otimes_BN_k,
\]
which occurs only at the full vertex and whose total fibre is itself.  This proves the formula.  All decompositions and restriction maps are natural under module maps and finite-set reindexing, proving the compatibility statements.
\end{proof}

\begin{lemma}[Connecting morphisms construct the required structured extension]
\label{lem:var-connecting-structured-extension}
Let \(A\) be a connective commutative ring spectrum and let
\[
    J\longrightarrow P\xrightarrow{\rho}N
    \xrightarrow{\partial_\rho}J[1]
\]
be a fibre sequence of connective \(A\)-modules.  Regard \(J\) as a module over the split square-zero algebra \(A\oplus N\) through the augmentation \(A\oplus N\to A\).  Then the connecting morphism \(\partial_\rho\) canonically determines a structured square-zero extension
\[
    \boxed{
    A\oplus P
    \longrightarrow
    A\oplus N
    }
\]
by \(J\), whose classifying derivation is represented by the morphism
\[
    L_{(A\oplus N)/A}
    \longrightarrow
    J[1]
\]
corresponding to \(\partial_\rho:N\to J[1]\).  The specified Chapter 2 pullback model is canonical and the construction is natural in morphisms of such fibre sequences.
\end{lemma}

\begin{proof}
Since \(J\) is an \((A\oplus N)\)-module through the augmentation, an \(A\)-linear derivation
\[
    A\oplus N
    \longrightarrow
    (A\oplus N)\oplus J[1]
\]
is equivalently an \(A\)-linear morphism \(N\to J[1]\).  Let \(d_{\partial_\rho}\) be the derivation represented by the connecting morphism and \(d_0\) the zero derivation.  Chapter 2 constructs the structured square-zero extension
\[
    (A\oplus N)^{\partial_\rho}
    :=
    (A\oplus N)
    \times_{(A\oplus N)\oplus J[1]}
    (A\oplus N),
\]
formed using these two derivations.

The split square-zero functor
\[
    \operatorname{Mod}_A
    \longrightarrow
    \operatorname{CAlg}(\operatorname{Sp})_{A//A},
    \qquad
    K\longmapsto A\oplus K,
\]
preserves limits.  Therefore
\[
\begin{aligned}
    (A\oplus N)^{\partial_\rho}
    &\simeq
    A\oplus
    \left(
      N\times_{N\oplus J[1]}N
    \right)\\
    &\simeq
    A\oplus\fib(\partial_\rho)\\
    &\simeq
    A\oplus P.
\end{aligned}
\]
This equivalence is over \(A\oplus N\) and identifies the fibre with the given module \(J\), so it is exactly the specified pullback identification required by the Chapter 2 structured square-zero classification.  No recognition from the underlying morphism is used.  Naturality follows from naturality of fibre sequences, connecting morphisms, split square-zero algebras, and pullbacks.
\end{proof}

\begin{lemma}[The mixed layer is a constructed structured square-zero extension of the boundary square]
\label{lem:var-mixed-boundary-square-zero}
Let \(T=\operatorname{Spec}B\) and \(N_1,N_2\in\operatorname{Mod}_B\) be connective.  Put
\[
    B_{\partial}[N_1,N_2]
    :=
    (B\oplus N_1)\times_B(B\oplus N_2)
    \simeq
    B\oplus N_1\oplus N_2.
\]
There is a canonical classifying morphism
\[
    \boxed{
    \eta_{12}:
    L_{B_{\partial}[N_1,N_2]/B}
    \longrightarrow
    (N_1\otimes_BN_2)[1]
    }
\]
and a canonical specified pullback identification exhibiting
\[
    \boxed{
    B[N_1,N_2]
    \longrightarrow
    B_{\partial}[N_1,N_2]
    }
\]
as the structured square-zero extension classified by \(\eta_{12}\).  Thus
\[
    T_{\partial}[N_1,N_2]
    :=
    \operatorname{Spec}B_{\partial}[N_1,N_2]
    \longrightarrow
    T[N_1,N_2]
\]
is a constructed structured square-zero immersion with coefficient \(N_1\otimes_BN_2\).  The construction is natural in \(B,N_1,N_2\), compatible with base change, and carried to itself by the canonical transposition comparison \(N_1\leftrightarrow N_2\).

A pair of side variations with common centre determines a map
\[
    T_{\partial}[N_1,N_2]\longrightarrow\bar A,
\]
and complete mixed lifts are precisely extensions of that map across the displayed structured square-zero immersion.
\end{lemma}

\begin{proof}
Put
\[
    A_1:=B\oplus N_1,
    \qquad
    P:=A_1\otimes_BN_2.
\]
Base change of the split square-zero algebra \(B\oplus N_2\) along \(B\to A_1\) gives
\[
    B[N_1,N_2]
    =
    A_1\otimes_B(B\oplus N_2)
    \simeq
    A_1\oplus P
\]
as augmented \(A_1\)-algebras.  Regarding \(N_2\) as an \(A_1\)-module through \(A_1\to B\), the boundary algebra is canonically
\[
    B_{\partial}[N_1,N_2]
    \simeq
    A_1\oplus N_2.
\]
Under these identifications the boundary restriction is induced by
\[
    \rho:
    P=A_1\otimes_BN_2
    \longrightarrow
    N_2
\]
from the augmentation \(A_1\to B\).

Using \(A_1\simeq B\oplus N_1\) as a \(B\)-module gives
\[
    P\simeq N_2\oplus(N_1\otimes_BN_2),
\]
with \(\rho\) projection to \(N_2\).  Hence
\[
    J:=\fib(\rho)
    \simeq
    N_1\otimes_BN_2.
\]
Multiplication by the augmentation ideal \(N_1\subset A_1\) on \(J\) repeats the first square-zero direction and is coherently zero, so the \(A_1\)-action on \(J\) factors through \(B\).  Lemma~\ref{lem:var-connecting-structured-extension}, applied to
\[
    J\longrightarrow P\xrightarrow{\rho}N_2
    \xrightarrow{\partial_\rho}J[1],
\]
constructs an \(A_1\)-relative classifying morphism
\[
    \eta_{12}^{A_1}:
    L_{B_{\partial}[N_1,N_2]/A_1}
    \longrightarrow
    J[1]
\]
and a specified equivalence
\[
    B[N_1,N_2]
    \simeq
    B_{\partial}[N_1,N_2]
    \times_{B_{\partial}[N_1,N_2]\oplus J[1]}
    B_{\partial}[N_1,N_2]
\]
using its associated derivation and the zero derivation.

The \(A_1\)-relative derivation is in particular \(B\)-linear.  Cotangent transitivity therefore represents it by
\[
    L_{B_{\partial}[N_1,N_2]/B}
    \longrightarrow
    L_{B_{\partial}[N_1,N_2]/A_1}
    \xrightarrow{\eta_{12}^{A_1}}
    J[1],
\]
which is \(\eta_{12}\).  Restriction from \(A_1\)-algebras to \(B\)-algebras preserves the displayed pullback, so the same model supplies the required \(B\)-relative structured extension.  Naturality and base change follow from functoriality of split square-zero base change, fibre sequences, connecting morphisms, cotangent transitivity, and the structured-extension construction.  Interchanging the two factors gives the same construction through symmetric monoidal coherence.

Finally, the boundary scheme is the pushout of the two side schemes over their common centre.  Hence a pair of side maps gives the displayed boundary map, and a complete bi-variation is exactly an extension of it to the full vertex.
\end{proof}

\begin{lemma}[Fixed-extension lifting for represented spectral geometry]
\label{lem:var-fixed-extension-lifting}
Let \(Y\to T\) be a represented morphism of spectral schemes, let
\[
    i:U\longrightarrow U'
\]
be a structured square-zero immersion over \(T\) with connective coefficient \(J\in\operatorname{QCoh}(U)\) and classifying morphism
\[
    \eta_i:L_{U/T}\longrightarrow J[1],
\]
and let \(f:U\to Y\) be a \(T\)-morphism.  Then
\[
    \operatorname{Lift}_T(f;i,Y)
    :=
    \hofib_f\left(
      \Map_T(U',Y)
      \longrightarrow
      \Map_T(U,Y)
    \right)
\]
is canonically the space of nullhomotopies of
\[
    \boxed{
    o(f;i):
    f^*L_{Y/T}
    \longrightarrow
    L_{U/T}
    \xrightarrow{\eta_i}
    J[1].
    }
\]
Consequently the lift space is nonempty exactly when \([o(f;i)]=0\), and, when nonempty, is naturally a torsor under
\[
    \boxed{
    \Omega^\infty
    \operatorname{map}_{\operatorname{QCoh}(U)}
    (f^*L_{Y/T},J).
    }
\]
The equivalence is natural in morphisms of structured extensions and in Cartesian base change.
\end{lemma}

\begin{proof}
A structured square-zero immersion induces a homeomorphism on underlying topological spaces, so the small Zariski \(\infty\)-topoi of \(U\) and \(U'\) identify.  On that fixed topos the Chapter 2 global structured extension is the specified pullback
\[
    \mathcal O_{U'}
    \simeq
    \mathcal O_U
    \times_{\mathcal O_U\oplus J[1]}
    \mathcal O_U
\]
formed from the derivation classified by \(\eta_i\) and the zero derivation.

A lift \(U'\to Y\) over \(T\) restricting to \(f\) is equivalently a lift of the morphism of commutative structure sheaves defining \(f\) through this fixed structured pullback.  Applying the corresponding mapping-space functor gives a homotopy pullback of mapping spaces.  Taking the fibre over \(f\) identifies the lift space with the space of homotopies between the induced derivation and zero.  Cotangent functoriality identifies that derivation with the displayed composite \(o(f;i)\).

When the homotopy fibre is nonempty, translation by any point identifies it as a torsor under the homotopy fibre over zero.  Exactness of the mapping-spectrum functor identifies that fibre with
\[
    \Omega^\infty
    \operatorname{map}_{\operatorname{QCoh}(U)}(f^*L_{Y/T},J).
\]
All constructions occur in the functorial fixed-topos pullback square, which gives the naturality and base-change assertions.
\end{proof}

\begin{lemma}[Square-zero translation formula for shifted scalar sections]
\label{lem:var-square-zero-translation}
Let
\[
    y:Y\longrightarrow T
\]
be a represented morphism of spectral schemes.  Let
\[
    i:U\longrightarrow U'
\]
be a structured square-zero immersion over \(T\) with connective coefficient
\[
    J\in\operatorname{QCoh}(U),
\]
and let
\[
    f:U\longrightarrow Y
\]
be a \(T\)-morphism.  Suppose the lift space
\[
    \operatorname{Lift}_T(f;i,Y)
\]
is nonempty.  Let
\[
    \ell:
    \one_Y
    \longrightarrow
    \Oadd_Y[n]
\]
be a shifted scalar section.

For two lifts
\[
    \widetilde f_0,\widetilde f_1
    \in
    \operatorname{Lift}_T(f;i,Y),
\]
write
\[
    \delta_{01}:
    f^*L_{Y/T}
    \longrightarrow
    J
\]
for their oriented difference under the torsor identification of Lemma~\ref{lem:var-fixed-extension-lifting}; thus \(\widetilde f_1\) is obtained from \(\widetilde f_0\) by translation through \(\delta_{01}\).  The two pulled-back scalar sections have the same restriction to \(U\).  Their difference therefore determines a point
\[
    \widetilde f_1^*\ell-\widetilde f_0^*\ell
    \in
    \Omega^\infty
    \operatorname{map}_{\operatorname{QCoh}(U)}
    (\mathcal O_U,J[n]).
\]
Let
\[
    d\ell_f:
    \mathcal O_U
    \longrightarrow
    f^*L_{Y/T}[n]
\]
be the pullback along \(f\) of the sheaf-theoretic universal differential of \(\ell\).  Then
\[
    \boxed{
    \widetilde f_1^*\ell-\widetilde f_0^*\ell
    \simeq
    \delta_{01}[n]\circ d\ell_f.
    }
\]

More strongly, after choosing either lift as origin, the translation formula is induced by the morphism of spectra
\[
    \boxed{
    \operatorname{map}_{\operatorname{QCoh}(U)}
    (f^*L_{Y/T},J)
    \xrightarrow{\;(-)[n]\circ d\ell_f\;}
    \operatorname{map}_{\operatorname{QCoh}(U)}
    (\mathcal O_U,J[n]).
    }
\]
Hence the formula is natural in morphisms of structured square-zero extensions, in \(f\), in shifted scalar sections, and under Cartesian base change, and it retains paths, loops, and all higher simplices of the lift torsor.
\end{lemma}

\begin{proof}
The assertion is Zariski-local on \(U\).  Since a structured square-zero immersion induces a homeomorphism on underlying topological spaces, \(U\) and \(U'\) have canonically equivalent small spectral Zariski sites.  We may therefore choose compatible affine opens and reduce to
\[
    T=\operatorname{Spec}R,
    \qquad
    Y=\operatorname{Spec}A,
    \qquad
    U=\operatorname{Spec}C,
    \qquad
    U'=\operatorname{Spec}C',
\]
where
\[
    q:C'\longrightarrow C
\]
is a structured square-zero extension by the connective \(C\)-module \(J\), and
\[
    \phi:A\longrightarrow C
\]
represents \(f\).  Let
\[
    \widetilde\phi_0,\widetilde\phi_1:
    A\longrightarrow C'
\]
represent the two lifts.

Form the pullback
\[
    D:=A\times_C C'
    \longrightarrow
    A.
\]
Base change of the specified structured extension makes \(D\to A\) a structured square-zero extension by \(J\), regarded as an \(A\)-module through \(\phi\).  Each lift \(\widetilde\phi_i\) is equivalently a section
\[
    s_i:A\longrightarrow D.
\]
The chosen section \(s_0\) gives a specified nullhomotopy of the classifying derivation of \(D\to A\).  Transporting the structured pullback model along this nullhomotopy identifies \(D\to A\), through a contractible space of compatible choices, with the split structured square-zero extension
\[
    A\oplus J\longrightarrow A
\]
in such a way that
\[
    s_0(a)=(a,0).
\]
Under the same splitting the second section has the form
\[
    s_1(a)=(a,\partial_{01}(a))
\]
for a unique \(R\)-linear derivation
\[
    \partial_{01}:A\longrightarrow J.
\]

The cotangent classifier of \(\partial_{01}\) is the adjoint of the torsor difference
\[
    \delta_{01}:
    L_{A/R}\otimes_A C
    \longrightarrow
    J.
\]
By the universal property of \(L_{A/R}\),
\[
    \boxed{
    \partial_{01}
    \simeq
    \left(
      A\xrightarrow{d_{A/R}}L_{A/R}
      \longrightarrow
      L_{A/R}\otimes_A C
      \xrightarrow{\delta_{01}}
      J
    \right).
    }
\]  Equivalently, the difference of the two original lifts as maps of underlying spectra lands in
\[
    \fib(C'\to C)\simeq J
\]
and is exactly the derivation classified by \(\delta_{01}\).

A degree-\(n\) scalar section on \(Y=\operatorname{Spec}A\) is a point
\[
    \ell:\mathbb S\longrightarrow A[n]
\]
of the underlying spectrum.  Therefore
\[
\begin{aligned}
    \widetilde\phi_1^*\ell-\widetilde\phi_0^*\ell
    &=
    (\widetilde\phi_1-\widetilde\phi_0)[n]\circ\ell\\
    &\simeq
    \delta_{01}[n]\circ d\ell_f.
\end{aligned}
\]
This is an equality induced by composition in mapping spectra, so it applies without restriction on \(n\) and is compatible with every simplex of the derivation torsor.

On overlaps the affine identities agree by functoriality of structured square-zero pullback, cotangent base change, and derivations.  Quasi-coherent Zariski descent therefore glues them to the displayed global equivalence.  The same functoriality gives the asserted naturality and Cartesian base-change compatibility.
\end{proof}

\begin{lemma}[Smooth fillability and the mixed-lift torsor]
\label{lem:var-critical-mixed-fillability}
Let
\[
    c:T=\operatorname{Spec}B\longrightarrow\bar A
\]
be an affine centre and let
\[
    p=(T,N_1,N_2,\widetilde a_1,\widetilde a_2)
    \in\mathcal P^{(2)}_{\bar A/Q}
\]
be side data.  In the represented smooth scalar sector, the space \(\operatorname{Fill}(p)\) of complete mixed lifts is nonempty and carries a canonical torsor structure under
\[
    \boxed{
    \Omega^\infty
    \operatorname{map}_{\operatorname{Mod}_B}
    \left(
      N_c,
      N_1\otimes_BN_2
    \right).
    }
\]
The torsor and its difference map are natural under affine restriction, morphisms of the two side modules, and represented smooth base change.
\end{lemma}

\begin{proof}
Put
\[
    J:=N_1\otimes_BN_2,
    \qquad
    U:=T_{\partial}[N_1,N_2],
    \qquad
    U':=T[N_1,N_2].
\]
By Lemma~\ref{lem:var-mixed-boundary-square-zero}, \(i:U\to U'\) is a constructed structured square-zero immersion with coefficient \(J\).  The two side variations give a boundary map
\[
    f_{\partial}:U\longrightarrow\bar A_c:=T\times_Q\bar A
\]
over \(T\), and \(\operatorname{Fill}(p)=\operatorname{Lift}_T(f_{\partial};i,\bar A_c)\).

Since \(a\) is representably smooth, \(\bar A_c\to T\) is smooth and
\[
    P:=f_{\partial}^*L_{\bar A_c/T}
\]
is finite locally free on the affine spectral scheme \(U\).  Lemma~\ref{lem:var-fixed-extension-lifting} places the obstruction in
\[
    \pi_0
    \operatorname{map}_{\operatorname{QCoh}(U)}(P,J[1]).
\]
Because \(P\) is finite locally free, \(\operatorname{map}(P,J)\) is a retract of a finite product of copies of the connective spectrum of sections of \(J\).  Hence
\[
    \pi_0\operatorname{map}(P,J[1])
    =
    \pi_{-1}\operatorname{map}(P,J)
    =0.
\]
The obstruction therefore vanishes, proving fillability.

The same lifting lemma identifies the fill torsor with
\[
    \Omega^\infty
    \operatorname{map}_{\operatorname{QCoh}(U)}(P,J).
\]
The boundary algebra acts on \(J\) through the central augmentation to \(B\).  Restriction of scalars and tensor-Hom adjunction therefore give
\[
\begin{aligned}
    \operatorname{map}_{\operatorname{QCoh}(U)}(P,J)
    &\simeq
    \operatorname{map}_{\operatorname{Mod}_B}
    \left(
      B\otimes_{B_{\partial}[N_1,N_2]}\Gamma(U,P),
      J
    \right)\\
    &\simeq
    \operatorname{map}_{\operatorname{Mod}_B}(N_c,J),
\end{aligned}
\]
where the final equivalence is cotangent base change along the central section.  Naturality follows from the naturality of the two preceding lemmas and cotangent base change.
\end{proof}

\begin{definition}[Critical side-data categories]
\label{def:var-critical-side-categories}
Let
\[
    j_L:C_L:=\Crit(L)\longrightarrow\bar A
\]
be the first Euler critical locus.  Define
\[
    \mathcal V^{(2)}_{L,\mathrm{crit}}
    :=
    \mathcal V^{(2)}_{\bar A/Q}
    \times_{\mathcal C^{\mathrm{aff}}_{/\bar A}}
    \mathcal C^{\mathrm{aff}}_{/C_L},
\]
\[
    \mathcal P^{(2)}_{L,\mathrm{crit}}
    :=
    \mathcal P^{(2)}_{\bar A/Q}
    \times_{\mathcal C^{\mathrm{aff}}_{/\bar A}}
    \mathcal C^{\mathrm{aff}}_{/C_L},
\]
and let
\[
    \partial_{2,L}:
    \mathcal V^{(2)}_{L,\mathrm{crit}}
    \longrightarrow
    \mathcal P^{(2)}_{L,\mathrm{crit}}
\]
be the restricted side-restriction functor.
\end{definition}

\begin{construction}[Natural-section presentation of the on-shell mixed variation]
\label{con:var-critical-mixed-natural-section}
Apply Corollary~\ref{cor:var-restricted-global-sections-right-kan} with
\[
    C=C_L,
    \qquad
    j=j_L,
    \qquad
    k=2,
    \qquad
    M=\Oadd_Q.
\]
Under the resulting canonical equivalence
\[
\begin{aligned}
    &
    \Map_{\mathcal E_{C_L}}
    \left(
      \one_{C_L},
      j_L^*\Vsecond_{\Oadd_Q}(A)[n]
    \right)
    \\
    &\qquad\simeq
    \Map_{
      \operatorname{Fun}(
        (\mathcal V^{(2)}_{L,\mathrm{crit}})^{\mathrm{op}},
        \operatorname{Sp}
      )
    }
    \left(
      \mathbb S,
      K_{\Oadd}^{(2)}[n]
    \right),
\end{aligned}
\]
define
\[
    \boxed{
    \widetilde\Delta_{L,\mathrm{crit}}^{(2),\mathrm{mix}}
    }
\]
to be the coherent natural section corresponding to the parameterized on-shell section
\[
    \Delta_{L,\mathrm{crit}}^{(2),\mathrm{mix}}
    =
    j_L^*\Delta_L^{(2),\mathrm{mix}}.
\]
Thus the tilde records only the change from the parameterized-spectrum presentation to the restricted right-Kan natural-section presentation; it introduces no new variation datum.
\end{construction}

\begin{proposition}[Critical mixed-lift translation is Euler contraction]
\label{prop:var-critical-mixed-lift-independence}
Let
\[
    p=(T=\operatorname{Spec}B,N_1,N_2,\widetilde a_1,\widetilde a_2)
    \in\mathcal P^{(2)}_{L,\mathrm{crit}}
\]
be critical side data and put
\[
    J:=N_1\otimes_BN_2.
\]
For two fills
\[
    v_0,v_1\in\operatorname{Fill}(p),
\]
let
\[
    \delta(v_0,v_1):N_c\longrightarrow J
\]
be their oriented difference under the torsor of Lemma~\ref{lem:var-critical-mixed-fillability}.  Then
\[
    \boxed{
    \left\langle
      \widetilde\Delta_{L,\mathrm{crit}}^{(2),\mathrm{mix}},
      v_1
    \right\rangle
    -
    \left\langle
      \widetilde\Delta_{L,\mathrm{crit}}^{(2),\mathrm{mix}},
      v_0
    \right\rangle
    \simeq
    \delta(v_0,v_1)\circ\mathsf{EL}_{L,c}
    }
\]
in
\[
    \Omega^\infty J[n].
\]
More strongly, the translation formula is induced by the morphism of spectra
\[
    \boxed{
    \operatorname{map}_{\operatorname{Mod}_B}(N_c,J)
    \xrightarrow{\;(-)[n]\circ\mathsf{EL}_{L,c}\;}
    J[n].
    }
\]

Because \(c\) factors through the derived Euler critical locus, its critical lift supplies a specified nullhomotopy
\[
    h_c:\mathsf{EL}_{L,c}\simeq0.
\]
Precomposition with \(h_c\) gives a specified nullhomotopy of the entire displayed morphism of spectra.  Consequently evaluation of the intrinsic on-shell mixed second variation is coherently constant on the complete fill torsor, including all paths, loops, and higher simplices.
\end{proposition}

\begin{proof}
Put
\[
    U:=T_{\partial}[N_1,N_2],
    \qquad
    U':=T[N_1,N_2].
\]
By Lemma~\ref{lem:var-mixed-boundary-square-zero},
\[
    i:U\longrightarrow U'
\]
is the constructed structured square-zero immersion with coefficient \(J\).  The common side data determines the boundary map
\[
    f_{\partial}:U\longrightarrow
    \bar A_c:=T\times_Q\bar A.
\]
The two fills are two extensions of this same map to \(U'\).

In the four-corner alternating difference defining
\(\Delta_L^{(2),\mathrm{mix}}\), the centre term and the two side terms are identical for \(v_0\) and \(v_1\).  Subtraction therefore leaves only the difference of the two full-vertex pullbacks of the local scalar Lagrangian:
\[
\begin{aligned}
    &
    \left\langle
      \widetilde\Delta_{L,\mathrm{crit}}^{(2),\mathrm{mix}},
      v_1
    \right\rangle
    -
    \left\langle
      \widetilde\Delta_{L,\mathrm{crit}}^{(2),\mathrm{mix}},
      v_0
    \right\rangle
    \\
    &\qquad\simeq
    \widetilde f_1^*\ell_L-\widetilde f_0^*\ell_L.
\end{aligned}
\]
Proposition~\ref{prop:var-scalar-cubical-total-fibre} identifies the scalar mixed total fibre with \(J\).

Apply Lemma~\ref{lem:var-square-zero-translation} to the represented morphism
\[
    \bar A_c\longrightarrow T
\]
and the fixed structured extension \(U\to U'\).  It identifies the preceding difference with contraction of the difference derivation
\[
    f_{\partial}^*L_{\bar A_c/T}
    \longrightarrow
    J
\]
against the first-order differential of \(\ell_L\).

The \(B_{\partial}[N_1,N_2]\)-module structure on \(J\) factors through the central augmentation
\[
    B_{\partial}[N_1,N_2]\longrightarrow B.
\]
The tensor--Hom adjunction and cotangent base change used in the proof of Lemma~\ref{lem:var-critical-mixed-fillability} identify
\[
    \operatorname{map}_{\operatorname{QCoh}(U)}
    \bigl(f_{\partial}^*L_{\bar A_c/T},J\bigr)
    \simeq
    \operatorname{map}_{\operatorname{Mod}_B}(N_c,J),
\]
and carry the lift difference to \(\delta(v_0,v_1)\).

Under the same central restriction, Propositions~\ref{prop:var-scalar-affine-cotangent} and \ref{prop:var-k1-euler} identify the first-order differential of the local scalar Lagrangian with
\[
    \mathsf{EL}_{L,c}:
    \one_T\longrightarrow N_c[n].
\]
The translation formula of Lemma~\ref{lem:var-square-zero-translation} therefore becomes exactly
\[
    \delta(v_0,v_1)\circ\mathsf{EL}_{L,c}.
\]

Finally, the chosen factorization of \(c\) through
\[
    C_L=\Crit(L)
\]
contains the specified path
\[
    h_c:\mathsf{EL}_{L,c}\simeq0.
\]
Composition with \(h_c\) nullhomotopes the complete map of spectra
\[
    \operatorname{map}_{\operatorname{Mod}_B}(N_c,J)
    \longrightarrow
    J[n].
\]
Since this is a homotopy of spectrum maps, every simplex of the fill torsor is simultaneously trivialized, proving the coherent constancy.
\end{proof}

\begin{construction}[The represented critical Hessian section]
\label{con:var-critical-hessian-section}
For critical side data
\[
    p=(T=\operatorname{Spec}B,N_1,N_2,\widetilde a_1,\widetilde a_2),
\]
put
\[
    J:=N_1\otimes_BN_2,
    \qquad
    G_p:=
    \Omega^\infty
    \operatorname{map}_{\operatorname{Mod}_B}(N_c,J).
\]
Lemma~\ref{lem:var-critical-mixed-fillability} identifies
\[
    \operatorname{Fill}(p)
\]
as a nonempty \(G_p\)-torsor.  Evaluate the natural-section presentation
\(\widetilde\Delta_{L,\mathrm{crit}}^{(2),\mathrm{mix}}\) of Construction~\ref{con:var-critical-mixed-natural-section} on complete fills of \(p\).  Using
\[
    \partial_{2,L}^*K^{(2)}_{\Oadd,\partial}
    \simeq
    K_{\Oadd}^{(2)},
\]
this gives
\[
    E_p:
    \operatorname{Fill}(p)
    \longrightarrow
    \Omega^\infty K^{(2)}_{\Oadd,\partial}(p)[n].
\]
Proposition~\ref{prop:var-critical-mixed-lift-independence} gives the coherent translation formula
\[
    E_p(v+\delta)-E_p(v)
    \simeq
    \delta\circ\mathsf{EL}_{L,c}.
\]
The critical nullhomotopy makes the right side coherently zero as a map of spectra.

The morphism
\[
    \operatorname{Fill}(p)\longrightarrow *
\]
is an effective epimorphism of spaces.  Its \v{C}ech nerve is the torsor bar construction
\[
    \operatorname{Fill}(p)\times G_p^{\times\bullet}.
\]
The coherent action-invariance just constructed supplies a descent datum for \(E_p\) on this complete \v{C}ech nerve.  Effective descent therefore determines, through a contractible space of choices, a unique point
\[
    \boxed{
    H_L^{\mathrm{crit}}(p)
    \in
    \Omega^\infty K^{(2)}_{\Oadd,\partial}(p)[n]
    }
\]
whose pullback to every fill is its intrinsic on-shell mixed second variation.

Naturality of the structured mixed extension, the fill torsor, the translation formula, and the critical nullhomotopy makes
\(p\mapsto H_L^{\mathrm{crit}}(p)\) a coherent natural section
\[
    \boxed{
    H_L^{\mathrm{crit}}
    \in
    \Map_{\operatorname{Fun}((\mathcal P^{(2)}_{L,\mathrm{crit}})^{\mathrm{op}},\operatorname{Sp})}
    \bigl(
      \mathbb S,
      K^{(2)}_{\Oadd,\partial}[n]
    \bigr).
    }
\]
No fill is selected or retained as part of the construction.
\end{construction}

\begin{definition}[Critical represented Hessian-descent space]
\label{def:var-critical-hessian-descent}
Pullback along \(\partial_{2,L}\) gives
\[
\begin{aligned}
    \partial_{2,L}^*:
    &\Map_{\operatorname{Fun}((\mathcal P^{(2)}_{L,\mathrm{crit}})^{\mathrm{op}},\operatorname{Sp})}
    \bigl(\mathbb S,K^{(2)}_{\Oadd,\partial}[n]\bigr)
    \\
    &\longrightarrow
    \Map_{\operatorname{Fun}((\mathcal V^{(2)}_{L,\mathrm{crit}})^{\mathrm{op}},\operatorname{Sp})}
    \bigl(\mathbb S,K^{(2)}_{\Oadd}[n]\bigr).
\end{aligned}
\]
Define
\[
    \boxed{
    \HessDesc^{\mathrm{rep}}_{\mathrm{crit}}(L)
    :=
    \hofib_{\widetilde\Delta_{L,\mathrm{crit}}^{(2),\mathrm{mix}}}
    \bigl(\partial_{2,L}^*\bigr).
    }
\]
\end{definition}

\begin{theorem}[Represented on-shell Hessian descent]
\label{thm:var-critical-hessian-descent}
The section \(H_L^{\mathrm{crit}}\) pulls back to the intrinsic on-shell mixed second variation and hence determines a canonical point
\[
    \boxed{
    h_L^{\mathrm{crit}}
    \in
    \HessDesc^{\mathrm{rep}}_{\mathrm{crit}}(L).
    }
\]
Thus the off-shell mixed-lift obstruction is killed on the Euler critical locus in the represented smooth scalar sector by the constructed Euler-contraction descent.
\end{theorem}

\begin{proof}
For a complete critical bi-variation \(v\), the side object \(p=\partial_{2,L}(v)\) has \(v\) as a fill.  The descent defining \(H_L^{\mathrm{crit}}(p)\) therefore pulls back along that fill to
\[
    \left\langle
      \widetilde\Delta_{L,\mathrm{crit}}^{(2),\mathrm{mix}},v
    \right\rangle.
\]
Proposition~\ref{prop:var-critical-mixed-lift-independence} supplies the coherent comparison with every other fill.  Naturality in complete bi-variations is precisely
\[
    \partial_{2,L}^*H_L^{\mathrm{crit}}
    \simeq
    \widetilde\Delta_{L,\mathrm{crit}}^{(2),\mathrm{mix}},
\]
which determines the displayed point of the homotopy fibre.
\end{proof}

\begin{proposition}[Universal critical Hessian tensor]
\label{prop:var-critical-hessian-tensor}
For an affine critical point
\[
    c:T=\operatorname{Spec}B\longrightarrow C_L,
\]
let \(u_c\) be the universal affine first variation.  Evaluating \(H_L^{\mathrm{crit}}\) on the universal side pair \((u_c,u_c)\) and applying Proposition~\ref{prop:var-scalar-cubical-total-fibre} gives
\[
    \widetilde h_{L,c}^{\mathrm{crit}}
    \in
    \Omega^\infty(N_c\otimes_BN_c)[n].
\]
These points commute with affine restriction and assemble to a canonical morphism in \(\mathcal Q_{C_L}\)
\[
    \boxed{
    \widetilde h_L^{\mathrm{crit}}:
    \mathcal O_{C_L}^{\mathrm{lin}}
    \longrightarrow
    j_L^*L^{\mathrm{lin}}_{\bar A/Q}
    \otimes
    j_L^*L^{\mathrm{lin}}_{\bar A/Q}[n].
    }
\]
If two side variations at \(c\) correspond to derivations
\[
    \alpha_i:N_c\longrightarrow N_i,
    \qquad i=1,2,
\]
then
\[
    \boxed{
    H_L^{\mathrm{crit}}(v_1,v_2)
    \simeq
    (\alpha_1\otimes_B\alpha_2)
    \widetilde h_{L,c}^{\mathrm{crit}}.
    }
\]
The tensor \(\widetilde h_L^{\mathrm{crit}}\) is invariant under the symmetry exchanging its two \(L^{\mathrm{lin}}\)-factors.
\end{proposition}

\begin{proof}
The side pair \((u_c,u_c)\) is an object of \(\mathcal P^{(2)}_{L,\mathrm{crit}}\).  The scalar total-fibre calculation identifies its coefficient with \(N_c\otimes_BN_c\).  Proposition~\ref{prop:var-universal-first-cartesian} makes \(c\mapsto u_c\) Cartesian under affine restriction, while Construction~\ref{con:var-critical-hessian-section} is natural in side data.  Hence the resulting maps
\[
    B\longrightarrow N_c\otimes_BN_c[n]
\]
form a Cartesian family and therefore define the displayed morphism in \(\mathcal Q_{C_L}\).

For an arbitrary side variation \(v_i\), Theorem~\ref{thm:var-sqz-factor-first-neighbourhood} gives the canonically contractible space of morphisms \(v_i\to u_c\) classified by \(\alpha_i\).  Naturality carries the universal value through
\[
    N_c\otimes_BN_c
    \xrightarrow{\alpha_1\otimes\alpha_2}
    N_1\otimes_BN_2,
\]
which proves the evaluation formula.  Transposition fixes the universal side pair, and the intrinsic mixed second variation is \(\Sigma_2\)-equivariant.  Since critical descent is canonical, the descended universal tensor is invariant under the corresponding symmetry.
\end{proof}

\begin{theorem}[Bilinearity and symmetry of the represented critical Hessian]
\label{thm:var-critical-hessian-bilinear}
Duality between \(L^{\mathrm{lin}}_{\bar A/Q}\) and \(T^{\mathrm{lin}}_{\bar A/Q}\) carries the universal tensor of Proposition~\ref{prop:var-critical-hessian-tensor} to a canonical morphism in \(\mathcal Q_{C_L}\)
\[
    \boxed{
    \Hess_L^{\mathrm{crit}}:
    j_L^*T^{\mathrm{lin}}_{\bar A/Q}
    \otimes
    j_L^*T^{\mathrm{lin}}_{\bar A/Q}
    \longrightarrow
    \mathcal O_{C_L}^{\mathrm{lin}}[n].
    }
\]
At an affine critical point \(c:T=\operatorname{Spec}B\to C_L\), this is
\[
    \boxed{
    \Hess_{L,c}^{\mathrm{crit}}:
    N_c^\vee\otimes_BN_c^\vee
    \longrightarrow
    B[n].
    }
\]
For side derivations \(\alpha_i:N_c\to N_i\), the corresponding mixed second variation is obtained from the universal tensor by \(\alpha_1\otimes\alpha_2\).  Hence it is additive and \(B\)-linear in each derivation variable, with no multilinearity axiom imposed.  The Hessian is symmetric, with any Koszul signs supplied by the stable symmetric monoidal braiding.
\end{theorem}

\begin{proof}
In the symmetric monoidal category \(\mathcal Q_{C_L}\), dualizability gives
\[
\begin{aligned}
    \operatorname{map}_{\mathcal Q_{C_L}}
    \left(
      \mathcal O_{C_L}^{\mathrm{lin}},
      j_L^*L^{\mathrm{lin}}
      \otimes
      j_L^*L^{\mathrm{lin}}[n]
    \right)
    \simeq{}&
    \operatorname{map}_{\mathcal Q_{C_L}}
    \left(
      j_L^*T^{\mathrm{lin}}
      \otimes
      j_L^*T^{\mathrm{lin}},
      \mathcal O_{C_L}^{\mathrm{lin}}[n]
    \right).
\end{aligned}
\]
The image of \(\widetilde h_L^{\mathrm{crit}}\) is the displayed Hessian.  Evaluating in \(\mathcal Q_{C_L}\) at an affine point gives the ordinary duality between the finite locally free modules \(N_c\) and \(N_c^\vee\), proving the affine formula.

The formula of Proposition~\ref{prop:var-critical-hessian-tensor} expresses every side evaluation through \(\alpha_1\otimes_B\alpha_2\).  Tensor product of modules is additive and \(B\)-linear in each variable, which proves bilinearity as a consequence of the constructed universal tensor.  The symmetry of that tensor dualizes to the symmetry of \(\Hess_L^{\mathrm{crit}}\).
\end{proof}

\begin{definition}[Jacobi operator and Jacobi coefficient]
\label{def:var-jacobi-coefficient}
Adjunction for the dual pair in \(\mathcal Q_{C_L}\) turns the critical Hessian into a morphism
\[
    \boxed{
    \Jac_L^{\mathrm{lin}}:
    j_L^*T^{\mathrm{lin}}_{\bar A/Q}
    \longrightarrow
    j_L^*L^{\mathrm{lin}}_{\bar A/Q}[n].
    }
\]
Applying the exact functor \(U_{C_L}\) defines the parameterized \emph{Jacobi operator}
\[
    \boxed{
    \Jac_L:
    j_L^*T^{\mathrm{st}}_{\bar A/Q}
    \longrightarrow
    j_L^*L^{\mathrm{st}}_{\bar A/Q}[n].
    }
\]
Define
\[
    \mathbb J_L:=\fib(\Jac_L^{\mathrm{lin}})\in\mathcal Q_{C_L},
    \qquad
    \boxed{
    \mathcal J_L
    :=
    U_{C_L}(\mathbb J_L)
    \simeq
    \fib(\Jac_L)
    \in
    \mathcal E_{C_L}.
    }
\]
Its relative zeroth-space object is
\[
    \boxed{
    \mathbf{Jac}(L)
    :=
    \Omega_{C_L}^\infty\mathcal J_L
    \longrightarrow
    C_L.
    }
\]
A point consists of a critical point, a represented first-order tangent variation, and a specified nullhomotopy of its linearized Euler source.
\end{definition}

\begin{proposition}[Affine description of the Jacobi linearization]
\label{prop:var-jacobi-affine}
At an affine critical point \(c:T=\operatorname{Spec}B\to C_L\), the affine-point linear Jacobi morphism is
\[
    \boxed{
    \Jac^{\mathrm{lin}}_{L,c}:
    N_c^\vee
    \longrightarrow
    N_c[n]
    }
\]
adjoint to the bilinear critical second variation.  Equivalently, it is the represented first-order linearization of the brave-new Euler source at the critical point.  The parameterized Jacobi operator \(\Jac_L\) of Definition~\ref{def:var-jacobi-coefficient} is the spectral Nisnevich sheafification of this coherently base-change-compatible affine-point family.
\end{proposition}

\begin{proof}
The first description is Definition~\ref{def:var-jacobi-coefficient} evaluated in the affine-point linear category under the finite-locally-free duality of Convention~\ref{conv:var-jacobi-sector}.  For the second, total fibre of a split two-cube may be taken iteratively in either direction.  Thus the alternating mixed second difference is the first difference, in one direction, of the first difference in the other.  Proposition~\ref{prop:var-k1-euler} identifies the first cubical difference of the Lagrangian with the Euler source.  Theorem~\ref{thm:var-critical-hessian-descent} removes the mixed-lift dependence along the critical locus.  Partial evaluation of the descended Hessian in one first-order direction is therefore the first-order difference of the Euler source in the other direction, and its adjoint is exactly \(\Jac^{\mathrm{lin}}_{L,c}\).  Naturality under affine restriction gives the compatible presheaf map whose sheafification is \(\Jac_L\).
\end{proof}

\begin{proposition}[Jacobi variations are intrinsically formally integrable]
\label{prop:var-jacobi-formal-integrability}
The composite
\[
    \mathbf{Jac}(L)\longrightarrow C_L\longrightarrow Q
\]
makes the Jacobi-variation object an intrinsic formally integrable linearized equation object.  Under the Chapter 11 equivalence it therefore inherits its jet coaction, formal-disk action, Cartan transport, and all higher descent coherences without any additional prolongation datum.
\end{proposition}

\begin{proof}
The object \(\mathcal J_L\) lies in \(\mathcal E_{C_L}\), so \(\mathbf{Jac}(L)\) is an object of the slice over \(C_L\), and \(C_L\) is already an object over \(Q\).  Hence the composite is an object of \((\XSp)_{/Q}\).  Apply the Chapter 11 effective-descent identification exactly as in Theorem~\ref{thm:var-critical-intrinsic-pde}.
\end{proof}

\begin{remark}[Off shell, on shell, and Jacobi]
\label{rem:var-hessian-jacobi-scope}
There are now three distinct levels.  Off shell, \(\Delta_L^{(2),\mathrm{mix}}\) is intrinsically cubical and its unrestricted side descent is the derived problem \(\HessDesc(L)\), which can be empty.  On the Euler critical locus, the mixed second variation is defined without further structure.  In the represented smooth scalar sector, the affine-point linear category, mixed boundary structured extension, fixed-extension lifting theorem, smooth fillability, Euler-contraction identity, and effective descent construct the side-dependent critical second variation.  Its universal side value is a tensor in \(\mathcal Q_{C_L}\), whose dual is the symmetric bilinear critical Hessian and whose partial adjoint gives the underlying Jacobi operator in \(\mathcal E_{C_L}\).  No corresponding bilinear statement is inferred outside that sector and no fillability, multilinearity, symmetry, Hessian, or Jacobi axiom is inserted.
\end{remark}

\section{Functoriality and arbitrary base change}
\label{sec:var-base-change}

\begin{proposition}[Affine-test base change for bi-variations]
\label{prop:var-second-affine-base-change}
Let \(c:Q'\to Q\), put
\[
    \bar A':=\bar A\times_QQ',
    \qquad
    M':=c^*M,
\]
and let \(t:T\to\bar A'\) be an affine point. Composition with \(\bar A'\to\bar A\) induces an equivalence between the comma \(\infty\)-categories of bi-square-zero variation diagrams over \(\bar A'/Q'\) and over \(\bar A/Q\) whose centres map to \(T\). Under this equivalence the mixed second-difference coefficient functors are identified by stable Beck--Chevalley.
\end{proposition}

\begin{proof}
The argument is the two-direction analogue of Proposition~\ref{prop:var-first-affine-base-change}. Given a bi-variation over \(\bar A/Q\) whose centre factors through \(T\to\bar A'\), the common retraction
\[
    r_{12}:T'[N_1,N_2]\longrightarrow T'
\]
and the centre map \(T'\to T\to Q'\) provide the required map from the thickened test to \(Q'\). Its composite to \(Q\) agrees with the structural map of the given bi-variation. The pullback universal property of \(\bar A'=\bar A\times_QQ'\) therefore lifts the complete four-corner diagram uniquely through a contractible space. The construction is inverse to forgetting to \(\bar A\) and is functorial on morphisms. Beck--Chevalley identifies all four coefficient section spectra, and exactness identifies their total fibre \(K_M^{(2)}\).
\end{proof}

\begin{theorem}[Base change of the variational calculus]
\label{thm:var-base-change}
Let \(c:Q'\to Q\), put
\[
    \bar A':=\bar A\times_QQ',
\]
let \(A'\) be the pulled-back Cartan descent object of Convention~\ref{conv:var-arbitrary-cartan-pullback}, and put \(M':=c^*M\). Then every \emph{internal} object and morphism in the complete variational calculus is canonically compatible with pullback. In particular there are coherent base-change equivalences for:
\begin{enumerate}
\item the vertical effective quotient and relation nerve;
\item the double relation \(\Rvar^{p,q}\) and double cochains \(\Cvar^{p,q}\);
\item the two complete filtered objects in the stable functor categories, the complete bifiltration \(F^{p,q}\DRvar\), \(\Bvar^{p,q}\), and \(\Omegavar^{p,q}\);
\item the unsigned and signed coherent maps \(\partial_H,\partial_V,d_H,d_V\);
\item \(\DRvar\), \(\Sec^p\), \(\Sec^{p,\mathrm{cl}}\), \(\delta^{\infty}_{\mathrm{var}}\), exact classifiers, the complete horizontal and Cartan--vertical filtrations, and the transgression coefficients;
\item the parameterized primitive objects \(\mathbf{Prim}_{\mathrm{var}}\), parameterized horizontal realization objects \(\mathbf{RealH}^q\), the relative Lepage-realization objects \(\mathbf{Lep}^{\,q}(L)\), their tautological boundary realizations, and the universal presymplectic transgressions \(\omega_L^{(q)}\);
\item the local full coefficient \(\Vfull\), the right-Kan first-variation coefficient \(\Vfirst\), and \(\lambda^1_{\mathrm{loc}}\) and \(\lambda^1_{\mathrm{var}}\);
\item Lagrangian local sections, full and first Euler sections, and both critical loci;
\item the parameterized Helmholtz-lift objects \(\mathbf{HelmLift}\), parameterized first-order Lagrangian-realization objects \(\mathbf{LagPrim}^{(1)}\), and their internal inverse-problem fibre square;
\item for every finite \(k\ge1\), the right-Kan cubical coefficient \(\mathbb V^{(k)}\), the alternating higher variation \(\Delta_L^{(k)}\), and its \(\Sigma_k\)-symmetry; in degree two, the on-shell mixed second variation.  In the represented smooth scalar sector the affine-point linear coefficient category, the finite-locally-free cotangent and tangent refinements, the structured mixed-boundary extension, the mixed-fill torsor, Jacobi linearization, represented critical Hessian, its canonical critical side descent, and the derived Jacobi-variation object are compatible with arbitrary pullback in the representably smooth sector.
\end{enumerate}
All equivalences satisfy identity, composition, horizontal and vertical pasting, and the higher pullback coherences inherited from Chapters 9--10, effective descent, and the pointwise Kan-extension constructions; in the geometric-base-change sector the Chapter 11 identifications agree with these pullbacks.

The external variation categories
\[
    \mathcal V^{(k)}_{\bar A'/Q'}
    \longrightarrow
    \mathcal V^{(k)}_{\bar A/Q}
\]
and, in degree two, the external side-data categories
\[
    \mathcal P^{(2)}_{\bar A'/Q'}
    \longrightarrow
    \mathcal P^{(2)}_{\bar A/Q}
\]
carry the canonical base-change restriction functors obtained by composition with \(\bar A'\to\bar A\). For every affine point \(t:T\to\bar A'\), these restriction functors induce equivalences on the comma \(\infty\)-categories over \(t\) and its composite \(t_0:T\to\bar A\). In degree two they commute coherently with the side-restriction functors and identify the pulled-back side coefficient diagrams by stable Beck--Chevalley. Consequently they induce a canonical restriction map
\[
    \HessDesc(L)
    \longrightarrow
    \HessDesc(c^*L),
\]
but no arbitrary base-change equivalence of Hessian-descent spaces is asserted.

For the global point spaces
\[
    \Prim_{\mathrm{var}}(\eta),\qquad
    \HelmLift(e),\qquad
    \LagPrim^{(1)}(e),\qquad
    \RealH^q(\eta),\qquad
    \Lep^q(L),
\]
arbitrary pullback induces the canonical restriction maps obtained by taking points of the preceding parameterized base-change equivalences; no arbitrary equivalence of global point spaces is asserted. If \(c\) is an effective epimorphism, these five global spaces are recovered from the corresponding parameterized objects by totalization over its Čech nerve.
\end{theorem}

\begin{proof}
Chapters 9--10 give arbitrary base change of relative de Rham objects, effective images, relation nerves, parameterized spectra, stable dependent adjunctions, and complete totalization filtrations. For the horizontal direction, Convention~\ref{conv:var-arbitrary-cartan-pullback} replaces any appeal to an arbitrary jet-comonad base-change theorem: the pulled-back effective epimorphism \(\epsilon':X\times_QQ'\twoheadrightarrow Q'\) has \v{C}ech nerve canonically \(Q'\times_QH_\bullet\). Hence the horizontal relation, the vertical relation, and every finite pullback defining the double relation base-change canonically. Right Beck--Chevalley identifies the double coefficient objects. When \(Q'\to Q\) comes from an actual represented geometric base change \(S'\to S\), Chapter 11 further identifies this descent comonad with the genuine jet comonad \(\Jinf_{X'/S'}\), but that identification is not used for the arbitrary-\(Q'\) statement.

Parameterized pullback preserves all limits and colimits. It therefore commutes with the complete-filtered constructions made in the stable functor categories, with partial and complete totalizations, matching fibres, iterated filtration fibres, associated graded objects, and the connecting maps that define the coherent boundaries. This proves (1)--(5). The parameterized primitive and horizontal-realization objects are finite homotopy pullbacks of zeroth-space objects built from these coefficients.  The relative Lepage object is an iterated finite homotopy pullback of the parameterized Helmholtz object and the horizontal realization tower, while its tautological boundary realization and presymplectic transgression are functorial sections obtained from the same pullback-stable coefficient maps.  This proves (6).

The local full coefficient is a finite fibre of dependent products, so it base-changes by Beck--Chevalley. For \(\Vfirst\), Lemma~\ref{lem:var-affine-point-pullback} reduces evaluation of a pulled-back parameterized spectrum at an affine point of \(\bar A'\) to evaluation of the original coefficient at the composite affine point of \(\bar A\). Proposition~\ref{prop:var-first-affine-base-change} identifies the complete comma categories and coefficient diagrams computing the raw right-Kan values, and Theorem~\ref{thm:var-first-kan-sheaf} identifies the raw and sheafified coefficients. This proves (7). Item (8) follows from adjunction, zeroth spaces, and finite homotopy limits; item (9) follows from the same finite-limit calculation applied to the parameterized lift-moduli objects.

Theorem~\ref{thm:var-k-base-change}, together with Lemma~\ref{lem:var-affine-point-pullback}, gives the pointwise comparison for every finite cubical right-Kan coefficient. Theorem~\ref{thm:var-k-kan-sheaf} identifies the raw and sheafified coefficients, while the alternating differences and symmetric-group actions are preserved termwise. This proves the internal cubical assertions of (10).  In the representably smooth scalar sector, pullback reindexes the affine-point linear coefficient systems, cotangent base change preserves the Cartesian finite-locally-free cotangent family and its pointwise dual, and Lemmas~\ref{lem:var-mixed-boundary-square-zero} and \ref{lem:var-fixed-extension-lifting} commute with arbitrary pullback; representable smoothness is preserved by base change.  The obstruction-vanishing argument, Euler-contraction identity, effective side descent, Hessian tensor, Jacobi linearization, and Jacobi-variation object are therefore transported by the same pullback comparison.

Composition with \(\bar A'\to\bar A\) gives the displayed restriction functors on the external variation categories. For every affine centre of \(\bar A'\), the comma-category equivalence is the construction in the proof of Theorem~\ref{thm:var-k-base-change}. In degree two the same construction applied to pairs of side variations gives the side-data restriction functor, and restriction of a complete bi-variation to its two sides commutes with base change through the canonical pullback comparison. Stable Beck--Chevalley identifies the corresponding side coefficient diagrams. Precomposition on coherent-section mapping spaces therefore gives the natural restriction map on \(\HessDesc(L)\); no arbitrary equivalence of Hessian-descent spaces is asserted.

Finally, taking global points is a right adjoint and therefore turns the internal pullback comparisons into natural restriction maps, but it does not turn an arbitrary base change into an equivalence of global section spaces. If \(c\) is an effective epimorphism, effective descent identifies each internal moduli object with the totalization of its Čech restrictions; applying global points gives the stated descent result for the primitive, Helmholtz, Lagrangian-realization, horizontal-realization, and relative Lepage spaces. No separate Čech-reconstruction theorem for \(\HessDesc(L)\) is asserted here.
\end{proof}

\section{Formal-\'etale invariance and spectral Nisnevich descent}
\label{sec:var-nisnevich}

\begin{theorem}[Formal-\'etale locality]
\label{thm:var-formal-etale}
Let \(u:U\to X\) be formally \(\acute{e}\)tale in the relative brave-new sense of Chapters 9--11. Put
\[
    Q_U:=Q_{U/S},
    \qquad
    Q:=Q_{X/S},
    \qquad
    c_u:Q_U\longrightarrow Q,
\]
and, for the Cartan descent object \(a:\bar A\to Q\), put
\[
    \bar A_U:=\bar A\times_QQ_U.
\]
The Chapter 11 formal-\'etale comparison identifies the horizontal Cartan relation over \(U/S\) with the pullback along \(c_u\) of the relation over \(X/S\), while Chapter 10 identifies the corresponding vertical relative de Rham geometry of \(\bar A_U/Q_U\) with the same pullback of the vertical geometry of \(\bar A/Q\). Under these identifications, every internal object of the complete variational calculus on the restricted Cartan object is the pullback along \(c_u\) of the corresponding object on \(X/S\). This includes the complete bifiltration, secondary forms, full variational differential, \(\Vfull\), \(\Vfirst\), every finite cubical coefficient \(\mathbb V^{(k)}\), Euler sections and critical loci, the universal variational moduli diagram and universal critical family, parameterized Helmholtz-lift, primitive, first-order Lagrangian-realization, horizontal-realization, and relative Lepage-realization objects, their tautological boundary realizations and universal presymplectic transgressions, transgression coefficients, the higher alternating variations, their symmetric-group actions, and the on-shell mixed second variation.  Whenever the represented smooth scalar Jacobi sector is preserved by the restriction, the represented Jacobi linearization, critical Hessian, critical Hessian descent, and Jacobi-variation object are transported by the same base-change comparison. The external side-data diagram in degree two is carried to the restricted side-data diagram by the base-change restriction functors of Theorem~\ref{thm:var-base-change}. The associated global lift and realization spaces, and the derived Hessian-descent space, carry the induced restriction maps; they are not asserted to be equivalent for a general formally \(\acute{e}\)tale restriction.
\end{theorem}

\begin{proof}
Chapter 11 supplies the formal-\'etale Cartesianity equivalences
\[
    U\simeq X\times_QQ_U,
    \qquad
    R_{U/S,\bullet}\simeq Q_U\times_QR_{X/S,\bullet}
\]
for the horizontal relation and its effective quotient. Chapter 10 supplies
\[
    R_{\bar A_U/Q_U,\bullet}
    \simeq
    Q_U\times_QR_{\bar A/Q,\bullet}
\]
for the vertical relative de Rham relation. Hence, degreewise,
\[
\begin{aligned}
    \Rvar^{p,q}(A_U)
    &=
    U^{\times_{Q_U}(q+1)}
    \times_{Q_U}
    R_{\bar A_U/Q_U,p}
    \\
    &\simeq
    Q_U\times_Q
    \left(
      X^{\times_Q(q+1)}
      \times_Q
      R_{\bar A/Q,p}
    \right)
    \\
    &=
    Q_U\times_Q\Rvar^{p,q}(A),
\end{aligned}
\]
compatibly with both simplicial directions. Right Beck--Chevalley therefore identifies the double cochains after pullback, and Theorem~\ref{thm:var-base-change} transports every internal stable coefficient construction, including the pointwise right-Kan descriptions of \(\Vfirst\) and of every \(\mathbb V^{(k)}\). Its external base-change statement transports the side-data diagram and induces the restriction map on \(\HessDesc\). No separate assertion that a nonaffine pullback of an affine variation centre remains affine is used: the first-order comparison is made on the new affine tests as in Proposition~\ref{prop:var-first-affine-base-change}, and the higher cubical comparison is Theorem~\ref{thm:var-k-base-change}.
\end{proof}

\begin{theorem}[Spectral Nisnevich descent]
\label{thm:var-nisnevich-descent}
Let \(X\) be represented and let \(U\to X\) be the coproduct of a finite spectral Nisnevich covering family, with \v{C}ech nerve \(U_\bullet\to X\). Let
\[
    Q_U:=Q_{U/S}\longrightarrow Q:=Q_{X/S}
\]
be the induced effective epimorphism of Chapter 10--11, and let \(Q_{U_\bullet}\) be its \v{C}ech nerve. For a Cartan descent object \(\bar A\to Q\), put
\[
    \bar A_U:=\bar A\times_QQ_U.
\]
Then \(\bar A_U\to\bar A\) is an effective epimorphism and its \v{C}ech nerve is canonically \(\bar A\times_QQ_{U_\bullet}\). The complete brave-new variational package is recovered by effective descent from its restrictions to this nerve. In particular this holds for
\[
    \Rvar^{p,q},\quad
    \Cvar^{p,q},\quad
    F^{p,q}\DRvar,\quad
    \Bvar^{p,q},\quad
    \Omegavar^{p,q},\quad
    \Sec^p,\quad
    \Sec^{p,\mathrm{cl}},\quad
    \delta^{\infty}_{\mathrm{var}},
\]
the complete horizontal and Cartan--vertical filtrations, \(\Vfull\), \(\Vfirst\), every \(\mathbb V^{(k)}\), full and first Euler sections, individual and universal critical families, the universal variational moduli diagram, the parameterized Helmholtz-lift, primitive, first-order Lagrangian-realization, horizontal-realization, and relative Lepage-realization objects, their tautological boundary realizations and universal presymplectic transgressions, all alternating higher variations and their symmetric-group actions, and the on-shell mixed second variation.  In the represented smooth scalar sector the same statement applies to the represented Jacobi and critical-Hessian package. Taking global points recovers the corresponding global Helmholtz, primitive, inverse-problem, horizontal-realization, and relative Lepage spaces as the totalizations of their restrictions to the Čech nerve.
\end{theorem}

\begin{proof}
Chapter 10 proves that \(Q_U\to Q\) is an effective epimorphism and that \(Q_{U_\bullet}\) is its \v{C}ech nerve. Pullback along \(\bar A\to Q\) therefore shows that \(\bar A_U\to\bar A\) is an effective epimorphism with the stated nerve.

Theorem~\ref{thm:var-base-change} identifies every coefficient object on a member of that nerve with the pullback of the corresponding global coefficient. Effective descent for parameterized spectra now gives
\[
    \mathcal E_{\bar A}
    \simeq
    \Tot_k\mathcal E_{\bar A\times_QQ_{U_k}}.
\]
Consequently the objects \(\Vfull\), \(\Vfirst\), every finite cubical coefficient \(\mathbb V^{(k)}\), the secondary coefficients, the filtered stages, the universal and fibrewise moduli objects, and all maps between them are recovered by totalization from their restrictions. The relation objects themselves satisfy descent in the ambient \(\infty\)-topos. Zeroth spaces and finite homotopy pullbacks preserve these descent limits. Applying global points to the descended parameterized Helmholtz, primitive, inverse-problem, horizontal-realization, and relative Lepage objects gives the asserted totalization formulas for their global spaces.  The universal presymplectic transgression descends as a section because it is constructed functorially from the descended vertical differential and the defining homotopy fibre.

It remains to justify the affine-point linear objects used in the represented smooth scalar Hessian sector, since they live in \(\mathcal Q\) rather than merely in \(\mathcal E\).  The pulled-back critical locus
\[
    C_{L,U}
    :=
    C_L\times_QQ_U
    \longrightarrow
    C_L
\]
is an effective epimorphism, and its \v{C}ech nerve is
\[
    C_L\times_QQ_{U_\bullet}.
\]
Proposition~\ref{prop:var-linear-coefficient-effective-descent} therefore gives
\[
    \mathcal Q_{C_L}
    \simeq
    \Tot_k
    \mathcal Q_{C_L\times_QQ_{U_k}}.
\]
The same proposition applied over \(\bar A\) recovers
\(L^{\mathrm{lin}}_{\bar A/Q}\) and
\(T^{\mathrm{lin}}_{\bar A/Q}\) from their restrictions.  The universal critical tensor, the symmetric bilinear Hessian, and the affine-point linear Jacobi morphism are morphisms between these descended \(\mathcal Q\)-objects, so they and their fibres descend in the displayed totalization.  Applying the exact underlying functor \(U\), compatibly with pullback, recovers the parameterized Jacobi operator and Jacobi-variation coefficient.  Hence the represented Jacobi and critical-Hessian package claimed in the theorem is also recovered by effective descent.  Thus the entire package, including its higher coherent structure, descends from the spectral Nisnevich cover.
\end{proof}

\section{Smooth local models}
\label{sec:var-local-models}

\begin{theorem}[Horizontal local formal-translation model]
\label{thm:var-horizontal-local-translation}
Assume \(X\to S\) is represented by a smooth spectral morphism. After the Zariski shrinking and spectral \(\acute{e}\)tale chart
\[
    q:U\longrightarrow\mathbb A_V^d
\]
constructed in Chapter 11, the horizontal relation nerve over \(U/V\) is identified with the action-groupoid nerve of the intrinsic affine formal disk
\[
    D_V^d
    :=
    V\times_{(\mathbb A_V^d)_{\dR/V}}\mathbb A_V^d.
\]
Consequently the horizontal direction of the restricted double variational relation has the corresponding formal-translation model.
\end{theorem}

\begin{proof}
This is exactly the Chapter 11 local affine formal-disk action theorem for the represented smooth morphism \(U\to V\), followed by product with the restricted vertical relation over the common base.
\end{proof}

\begin{theorem}[Vertical local formal-translation model]
\label{thm:var-vertical-local-translation}
Let \(t:T\to Q\) be a represented local test such that
\[
    \bar A_T:=\bar A\times_QT\longrightarrow T
\]
is represented and smooth. After the local Zariski shrinking and spectral \(\acute{e}\)tale chart supplied by Chapter 11 for this represented smooth morphism, the vertical relation nerve \(R_{\bar A_T/T,\bullet}\) is identified with the corresponding formal-translation action-groupoid nerve. If the hypotheses of Theorem~\ref{thm:var-horizontal-local-translation} hold simultaneously, the restricted double relation is the product of the two commuting formal actions.
\end{theorem}

\begin{proof}
Apply the Chapter 11 formal-translation theorem to the represented smooth morphism \(\bar A_T\to T\). The last assertion follows from Definition~\ref{def:var-double-relation} and the commutation of the two simplicial directions.
\end{proof}

\begin{theorem}[Tubular-completion comparison in the doubly represented smooth sector]
\label{thm:var-local-tubular}
On a local test where both the horizontal morphism and the vertical morphism are represented and smooth, the Chapter 10 tubular-completion models identify every fixed bidegree of \(\Rvar\) with the product of the corresponding intrinsic de Rham completions of the local horizontal and vertical diagonal linear models. A single chosen degree-one tubular comparison in each represented smooth direction therefore gives simultaneous degreewise equivalences of the full intrinsic formal relation objects in every bidegree. No additional compatibility of the chosen tubular models with every simplicial composition map is asserted. The comparison does not identify the normalized cochains with exterior powers and does not imply an HKR collapse.
\end{theorem}

\begin{proof}
Apply the Chapter 10 local tubular-completion theorem separately to the two represented smooth relation nerves and take their product over the local base. The non-HKR statement follows from the explicit Chapter 10 example in which a nonzero normalized one-cochain is killed by first-order cotangent linearization.
\end{proof}

\section{First-order, reduced-classical, and secondary-calculus comparison}
\label{sec:var-classical-comparison}

\begin{proposition}[First-order scalar vertical comparison]
\label{prop:var-first-contact-coefficient}
Assume
\[
    a:\bar A\longrightarrow Q
\]
is representable.  For scalar coefficients, the map
\[
    \lambda^1_{\bar A/Q}:
    \Sec^1(A)
    \simeq
    \Omegainf^1(\bar A/Q)
    \longrightarrow
    \Pi_aL^{\mathrm{st}}_{\bar A/Q}
\]
is characterized by represented affine base change as follows.

For every represented affine base test
\[
    t:T=\operatorname{Spec}R\longrightarrow Q,
\]
put
\[
    \bar A_T:=T\times_Q\bar A,
    \qquad
    a_T:\bar A_T\longrightarrow T.
\]
Then the pullback of \(\lambda^1_{\bar A/Q}\) is the Chapter 10 first-order linearization
\[
    \Omegainf^1(\bar A_T/T)
    \longrightarrow
    \Pi_{a_T}L^{\mathrm{st}}_{\bar A_T/T}.
\]
Consequently the composite
\[
    \boxed{
    \Sec^0(A)
    \xrightarrow{\delta^\infty_{\mathrm{var}}}
    \Sec^1(A)
    \xrightarrow{\lambda^1_{\bar A/Q}}
    \Pi_aL^{\mathrm{st}}_{\bar A/Q}
    }
\]
pulls back over \(T\) to the stabilized universal relative derivation
\[
    \Gamma(\bar A_T,\mathcal O_{\bar A_T})
    \longrightarrow
    \Gamma(\bar A_T,L_{\bar A_T/T}).
\]

More precisely, for every affine open immersion
\[
    j:U=\operatorname{Spec}B
    \hookrightarrow
    \bar A_T,
\]
the restriction of this map is the universal spectral derivation
\[
    \boxed{
    d_{B/R}:
    B\longrightarrow L_{B/R}.
    }
\]
For an arbitrary affine point
\[
    c:U=\operatorname{Spec}B
    \longrightarrow
    \bar A_T,
\]
one instead has the canonical affine-point identification
\[
    \boxed{
    L^{\mathrm{st}}_{\bar A_T/T}(U,c)
    \simeq
    \Gamma\bigl(U,c^*L_{\bar A_T/T}\bigr).
    }
\]
Thus, for a shifted scalar section \(\ell\), its represented first-order variation at \(c\) is the section
\[
    d\ell_c:
    \mathcal O_U
    \longrightarrow
    c^*L_{\bar A_T/T}[n]
\]
obtained by pulling back the sheaf-theoretic universal differential of \(\ell\).

If \(c\) is spectral \(\acute{e}\)tale---in particular, if \(c\) is an affine open immersion---then cotangent transitivity gives
\[
    c^*L_{\bar A_T/T}
    \xrightarrow{\sim}
    L_{U/T},
\]
and therefore
\[
    \Gamma(U,c^*L_{\bar A_T/T})
    \simeq
    L_{B/R}.
\]
In this \(\acute{e}\)tale case, on discrete smooth classical input, taking \(\pi_0\) recovers the ordinary K\"ahler universal derivation
\[
    B\longrightarrow\Omega^1_{B/R}.
\]
\end{proposition}

\begin{proof}
Theorem~\ref{thm:var-first-represented-comparison} identifies the one-stage coefficient with the affine-pointwise stabilized first-order coefficient and identifies \(\lambda^1_{\mathrm{var}}\) with the corresponding cotangent restriction. Proposition~\ref{prop:var-scalar-affine-cotangent} identifies its pullback along every represented affine \(t:T\to Q\) with the Chapter 10 stabilized cotangent coefficient and first-order linearization of
\[
    \bar A_T/T.
\]
Theorem~\ref{thm:var-operator-identification} identifies the pullback of \(\delta^\infty_{\mathrm{var}}\) in degree zero with the Chapter 10 coherent infinitesimal differential of the same represented morphism. Chapter 10 proves that the composite of that degree-zero infinitesimal differential with first-order linearization is the stabilized universal derivation
\[
    \Gamma(\bar A_T,\mathcal O_{\bar A_T})
    \longrightarrow
    \Gamma(\bar A_T,L_{\bar A_T/T}).
\]
After restriction to an affine open \(U=\operatorname{Spec}B\subseteq\bar A_T\), its affine calculation is
\[
    d_{B/R}:B\longrightarrow L_{B/R}.
\]

For an arbitrary affine point \(c:U\to\bar A_T\), Proposition~\ref{prop:var-scalar-affine-cotangent} and the Chapter 10 affine-point description give
\[
    L^{\mathrm{st}}_{\bar A_T/T}(U,c)
    \simeq
    \Gamma(U,c^*L_{\bar A_T/T}),
\]
and naturality of the global cotangent derivation under pullback gives the displayed \(d\ell_c\).

Finally, cotangent transitivity for
\[
    U\xrightarrow{c}\bar A_T\longrightarrow T
\]
gives a cofiber sequence
\[
    c^*L_{\bar A_T/T}
    \longrightarrow
    L_{U/T}
    \longrightarrow
    L_{U/\bar A_T}.
\]
If \(c\) is spectral \(\acute{e}\)tale, the final term vanishes, yielding the asserted equivalence.  On discrete smooth input the degree-zero comparison for the spectral cotangent complex identifies
\[
    \pi_0L_{B/R}
    \simeq
    \Omega^1_{B/R}.
\]
\end{proof}

\begin{definition}[Reduced-classical double relation in the comparison sector]
\label{def:var-reduced-classical-double}
Assume that the Chapter 9 reduced-classical comparison maps identify the horizontal relative de Rham target of \(X/S\) with its reduced-classical target and identify the vertical relative de Rham target of \(\bar A/Q\) with its reduced-classical target. By uniqueness of effective-image factorization, these equivalences identify the corresponding effective quotients. Write
\[
    Q^{\mathrm{red,cl}}\simeq Q
\]
for the identified horizontal effective quotient,
\[
    H^{\mathrm{red,cl}}_\bullet
\]
for the corresponding reduced-classical horizontal relation nerve, and
\[
    V^{\mathrm{red,cl}}_{A,\bullet}
\]
for the reduced-classical vertical relation nerve, all transported over the identified base. Define
\[
    \boxed{
    \mathcal R_{\mathrm{var,red,cl}}^{p,q}(A)
    :=
    H_q^{\mathrm{red,cl}}
    \times_{Q^{\mathrm{red,cl}}}
    V_{A,p}^{\mathrm{red,cl}}.
    }
\]
\end{definition}

\begin{theorem}[Reduced-classical comparison of the double relation]
\label{thm:var-reduced-classical}
Under the hypotheses of Definition~\ref{def:var-reduced-classical-double}, there is a canonical equivalence of bicosimplicial relation objects
\[
    \boxed{
    \Rvar^{\bullet,\bullet}(A)
    \simeq
    \mathcal R_{\mathrm{var,red,cl}}^{\bullet,\bullet}(A).
    }
\]
It induces the corresponding comparison on double stable cochains, the explicit complete bifiltration, and all partial and complete totalizations. Under represented first-order linearization, the degree-one scalar vertical comparison is the classical universal derivation of Proposition~\ref{prop:var-first-contact-coefficient}.
\end{theorem}

\begin{proof}
Chapter 11 constructs the reduced-classical comparison for the horizontal jet relation and proves it to be an equivalence when the relevant Chapter 9 relative de Rham comparison is an equivalence. Chapters 9--10 give the corresponding vertical relative de Rham comparison. By the hypothesis, the effective-image identifications transport both relation nerves over the same identified quotient base. Taking their product gives the displayed bicosimplicial equivalence. Right Beck--Chevalley transports the coefficient objects, and the complete bifiltration is formed only from limits, fibres, and partial totalizations, so the comparison propagates through every stage. The first-order statement is Proposition~\ref{prop:var-first-contact-coefficient}.
\end{proof}

\begin{remark}[The classical comparison is first-order and structural]
\label{rem:var-classical-contact}
The full brave-new normalized cochains retain finite-infinitesimal information beyond exterior cotangent forms, as Chapter 10's affine-line calculation shows.  The present chapter constructs the represented scalar degree-one first-order comparison of Proposition~\ref{prop:var-first-contact-coefficient} and, under the stated reduced-classical hypotheses, the relation-level comparison of Theorem~\ref{thm:var-reduced-classical}.  It does not construct an all-degree exterior realization
\[
    \Omegainf^p(\bar A/Q)
    \longrightarrow
    \bigwedge^pL_{\bar A/Q},
\]
does not identify the Cartan-contact filtration \(F_{\mathcal C}^\bullet\) with powers of a contact ideal, and does not identify the full brave-new bifiltration with the classical variational bicomplex.  The comparison below therefore records the classical secondary-calculus architecture and the already proved degree-one point of contact, rather than a theorem of equivalence between the two filtered complexes.
\end{remark}

\subsection{Comparison with classical secondary calculus}
\label{subsec:var-classical-secondary-special-case}

The classical construction closest to the present architecture is Vinogradov's secondary calculus on a diffiety; see \cite{PuglieseSparanoVitagliano2023,Vitolo2000}.  Let
\[
    (\mathscr E,\mathscr C)
\]
be the infinite prolongation of a formally integrable classical PDE, with Cartan distribution \(\mathscr C\).  Write
\[
    \mathscr C\Omega^\bullet(\mathscr E)
    \subset
    \Omega^\bullet(\mathscr E)
\]
for the differential ideal of contact forms.  Its powers give the contact filtration
\[
    \Omega^\bullet(\mathscr E)
    \supset
    \mathscr C\Omega^\bullet(\mathscr E)
    \supset
    \mathscr C^2\Omega^\bullet(\mathscr E)
    \supset\cdots,
\]
and the associated \(\mathscr C\)-spectral sequence has
\[
    \mathscr C E_0^{p,q}
    =
    \frac{
      \mathscr C^p\Omega^{p+q}(\mathscr E)
    }{
      \mathscr C^{p+1}\Omega^{p+q}(\mathscr E)
    }.
\]
The \(d_0\)-differential is the horizontal differential and
\[
    \mathscr C E_1^{p,q}
    =
    H^q
    \bigl(
      \mathscr C E_0^{p,\bullet},
      d_0
    \bigr)
\]
is the corresponding horizontal cohomology with contact/transverse \(p\)-form coefficients, interpreted in secondary calculus as classical secondary \(p\)-forms in horizontal degree \(q\).  The next differential
\[
    d_1:
    \mathscr C E_1^{p,q}
    \longrightarrow
    \mathscr C E_1^{p+1,q}
\]
is the secondary de Rham differential.

The infinitesimal-simplex presentation of ordinary forms gives an independent first-infinitesimal comparison model.  In the synthetic construction of Kock and its scheme-theoretic extension by Breen--Messing, ordinary differential forms are represented by normalized alternating functions on suitable first infinitesimal simplices, and the exterior differential is induced by the simplicial boundary.  This identifies the classical exterior complex with a normalized first-infinitesimal simplicial model in that classical setting.  It does not, by itself, construct a morphism from the full square-zero-generated brave-new normalized cochains in every degree to the classical exterior complex, nor does it identify the brave-new Cartan-contact filtration with the powers of the classical contact ideal.

Accordingly the classical sequence
\[
    \mathscr C E_1^{0,\bullet}
    \xrightarrow{d_1}
    \mathscr C E_1^{1,\bullet}
    \xrightarrow{d_1}
    \mathscr C E_1^{2,\bullet}
    \longrightarrow\cdots
\]
is a structural benchmark for
\[
    \Sec^0(A;M)
    \xrightarrow{\delta^\infty_{\mathrm{var}}}
    \Sec^1(A;M)
    \xrightarrow{\delta^\infty_{\mathrm{var}}}
    \Sec^2(A;M)
    \longrightarrow\cdots:
\]
both organize a transverse or vertical differential after horizontal Cartan descent.  The present chapter proves only the degree-zero-to-one first-order compatibility of Proposition~\ref{prop:var-first-contact-coefficient} and the conditional reduced-classical relation comparison of Theorem~\ref{thm:var-reduced-classical}; no equivalence of the two complete filtered theories is used.

\begin{example}[The empty equation and the ordinary variational bicomplex]
\label{ex:var-classical-empty-equation}
Let \(\pi:E\to M\) be a smooth classical fibre bundle and take the empty equation
\[
    \mathscr E=J^\infty\pi
\]
with its canonical Cartan distribution.  The preceding contact filtration underlies the ordinary variational bicomplex.  In top horizontal degree the secondary differential
\[
    d_1:
    \mathscr C E_1^{0,\dim M}
    \longrightarrow
    \mathscr C E_1^{1,\dim M}
\]
is the classical Euler--Lagrange operator, and the next map
\[
    d_1:
    \mathscr C E_1^{1,\dim M}
    \longrightarrow
    \mathscr C E_1^{2,\dim M}
\]
is the classical Helmholtz operator.  The inverse problem asks for Lagrangian primitives of Euler sources satisfying the Helmholtz condition in this secondary complex.  This is the classical benchmark for the brave-new sequence of secondary forms, Euler sections, coherent Helmholtz lifts, and primitive spaces constructed above; the example does not identify the classical contact filtration with \(F_{\mathcal C}^\bullet\).
\end{example}

\begin{remark}[Compatibility with the proved degree-one comparison]
The represented scalar comparison of Proposition~\ref{prop:var-first-contact-coefficient} is the constructed degree-zero-to-one point of contact with the classical picture: after first-order restriction, the brave-new infinitesimal differential becomes the cotangent universal derivation, and on discrete smooth input its \(\pi_0\)-shadow is the ordinary K\"ahler derivation.  No extension of this statement to an equivalence of the full normalized complexes or filtrations is asserted here.  The full brave-new theory retains the higher finite-infinitesimal classes discarded by first-order linearization.
\end{remark}


\begin{remark}[Classical interpretation of the relative Lepage object]
\label{rem:var-classical-lepage-shadow}
In the ordinary variational bicomplex, a one-step Lepage decomposition
\[
    d_VL=E_L+d_H\Theta
\]
has the same formal shape as Theorem~\ref{thm:var-relative-integration-by-parts}: an Euler representative of the secondary first variation is compared with the raw first variation by a horizontally trivialized difference, and vertical differentiation of that trivialization produces the presymplectic current.  The present chapter has constructed the brave-new moduli object of such relative realization data and its universal transgression, but the limited classical comparison established above does not identify the complete brave-new horizontal tower with the exterior horizontal complex.  Therefore no unconditional map from every classical Lepage representative to \(\mathbf{Lep}^{\,q}(L)\), and no all-degree HKR equivalence, is asserted here.  The classical formula is the comparison benchmark for the universal property already constructed.
\end{remark}

\begin{proposition}[First-order classical shadow of the represented Jacobi operator]
\label{prop:var-classical-jacobi-shadow}
On discrete smooth classical input in the represented smooth scalar Jacobi sector, the affine \(\pi_0\)-shadow of
\[
    \Jac^{\mathrm{lin}}_{L,c}:
    N_c^\vee
    \longrightarrow
    N_c[n]
\]
is the ordinary first-order linearization, in K\"ahler-cotangent coordinates, of the affine \(\pi_0\)-shadow of the brave-new first Euler source along the critical locus.  The affine \(\pi_0\)-shadow of
\[
    \Hess_L^{\mathrm{crit}}:
    j_L^*T^{\mathrm{lin}}_{\bar A/Q}
    \otimes
    j_L^*T^{\mathrm{lin}}_{\bar A/Q}
    \longrightarrow
    \mathcal O_{C_L}^{\mathrm{lin}}[n]
\]
is the symmetric bilinear pairing adjoint to that first-order linearization.  The parameterized \(\Jac_L\) is the spectral Nisnevich sheafification of these compatible affine linearizations.  No identification with the Jacobi operator or Hessian of the classical variational bicomplex is asserted beyond a sector in which an additional comparison theorem identifies the brave-new Euler source with the classical Euler--Lagrange source.
\end{proposition}

\begin{proof}
Proposition~\ref{prop:var-first-contact-coefficient} identifies the represented scalar first-order infinitesimal differential, after taking \(\pi_0\) on discrete smooth input, with the ordinary K\"ahler universal derivation.  Proposition~\ref{prop:var-jacobi-affine} identifies the affine linear Jacobi morphism with the first-order linearization of the brave-new Euler source, while Theorem~\ref{thm:var-critical-hessian-bilinear} identifies its adjoint with the constructed symmetric critical pairing.  Sheafification gives the parameterized operator without strengthening this affine first-order classical comparison.  The chapter has not identified the full brave-new normalized complex or its secondary differential with the classical variational bicomplex, so no stronger classical Euler--Lagrange/Jacobi identification follows from the established comparison.
\end{proof}

\begin{remark}[Derived critical loci]
\label{rem:var-derived-critical-literature}
The critical loci are homotopy zero fibres in the native \(\infty\)-topos. They retain derived intersection data, automorphisms, and higher equation homotopies rather than replacing the equation by an image subobject; compare the general derived-critical-locus perspective in \cite{Grataloup2020}. No shifted-symplectic structure is imported into the construction.
\end{remark}

\section{Differential densities, the brave-new action functional, and its variation}
\label{sec:var-action-functional}

Assume throughout this section that
\[
    x:X\longrightarrow S
\]
is represented, smooth, separated, and of finite presentation and that \(S\) lies in the standing differential-motivic scope of Chapters 9--10. Write
\[
    q:Q=Q_{X/S}\longrightarrow S
\]
for the structural morphism, so \(x=q\epsilon\). For every represented differential base \(Z\), write
\[
    \one_Z^\partial\in\mathcal R_Z^0
\]
for the tensor unit of the Thom-stable differential category; this is distinct from the tensor unit \(\one_Z\in\mathcal E_Z\). Write
\[
    x_{e!}:\mathcal R_X^0\rightleftarrows\mathcal R_S^0:x_e^!
\]
for the Thom-stable differential exceptional adjunction. Thus
\[
    x_e^!D=\Sigma_{T_x,0}x^{*,0}D
\]
for \(D\in\mathcal R_S^0\), and Chapter 10 supplies the differential Umkehr equivalence
\[
    \mathcal Z^\partial_{x,D}:
    \operatorname{map}_{\mathcal R_X^0}
    (K,x_e^!D)
    \xrightarrow{\sim}
    \operatorname{map}_{\mathcal R_S^0}
    (x_{e!}K,D),
\]
whose value on a morphism is postcomposition with the exceptional trace
\[
    \operatorname{tr}^\partial_{x,D}:
    x_{e!}x_e^!D\longrightarrow D.
\]
No BNV cycle-compatibility hypothesis is required for this primary exceptional integration.

\subsection{Big differential mapping-spectrum coefficients}

\begin{definition}[Big differential mapping-spectrum coefficient]
\label{def:var-big-differential-map-coefficient}
Let \(Z\) be a represented qcqs spectral scheme in the standing differential-motivic range and let
\[
    E,F\in\mathcal R_Z^0.
\]
On the affine-point site \(\mathcal C^{\mathrm{aff}}_{/Z}\), define
\[
    \boxed{
    \underline{\operatorname{map}}^{\partial}_{Z}(E,F)(U,u)
    :=
    \operatorname{map}_{\mathcal R_U^0}
    \bigl(u^{*,0}E,u^{*,0}F\bigr)
    }
\]
for every affine point \(u:U\to Z\). Restriction along a morphism of affine points is induced by differential pullback.
\end{definition}

\begin{theorem}[Descent, base change, and global sections of the big differential mapping coefficient]
\label{thm:var-big-differential-map-coefficient}
The presheaf of Definition~\ref{def:var-big-differential-map-coefficient} is a spectral Nisnevich sheaf and hence determines an object
\[
    \underline{\operatorname{map}}^{\partial}_{Z}(E,F)
    \in
    \mathcal E_Z.
\]
For every represented morphism \(g:Z'\to Z\) between qcqs bases in the standing differential base range there is a canonical coherent equivalence
\[
    \boxed{
    g^*\underline{\operatorname{map}}^{\partial}_{Z}(E,F)
    \simeq
    \underline{\operatorname{map}}^{\partial}_{Z'}
    (g^{*,0}E,g^{*,0}F).
    }
\]
Moreover
\[
    \boxed{
    \Gamma_Z^{\mathrm{st}}
    \underline{\operatorname{map}}^{\partial}_{Z}(E,F)
    \simeq
    \operatorname{map}_{\mathcal R_Z^0}(E,F).
    }
\]
All equivalences are natural in \(E,F\) and satisfy the pullback-pasting coherences inherited from the differential coefficient system.
\end{theorem}

\begin{proof}
Let \(U_\bullet\to U\) be the Čech nerve of a finite affine spectral Nisnevich cover. Chapter 9 proves finite Nisnevich descent for the Thom-stable differential categories:
\[
    \mathcal R_U^0
    \simeq
    \Tot_{[k]\in\Delta}\mathcal R_{U_k}^0.
\]
Under this equivalence the two objects \(u^{*,0}E\) and \(u^{*,0}F\) are the descent objects determined by their restrictions. Mapping spectra in a limit of stable \(\infty\)-categories are computed as the corresponding limits of mapping spectra. Hence
\[
\begin{aligned}
    \operatorname{map}_{\mathcal R_U^0}
    (u^{*,0}E,u^{*,0}F)
    &\simeq
    \Tot_k
    \operatorname{map}_{\mathcal R_{U_k}^0}
    (u_k^{*,0}E,u_k^{*,0}F).
\end{aligned}
\]
Thus the affine-point presheaf satisfies spectral Nisnevich descent.

For base change, let \(t:T\to Z'\) be an affine point. Lemma~\ref{lem:var-affine-point-pullback} identifies evaluation of the left side with the original sheaf at the composite \(gt:T\to Z\), namely
\[
    \operatorname{map}_{\mathcal R_T^0}
    ((gt)^{*,0}E,(gt)^{*,0}F).
\]
Functoriality of differential pullback identifies this with
\[
    \operatorname{map}_{\mathcal R_T^0}
    (t^{*,0}g^{*,0}E,t^{*,0}g^{*,0}F),
\]
which is the value of the right side. Affine-point recognition gives the displayed equivalence and its coherences.

For the global-section equivalence, let \(\operatorname{QCOp}(Z)\) be the site of quasi-compact open spectral subschemes of \(Z\) with finite Zariski covering families, and let
\[
    \operatorname{AffOp}(Z)
    \subset
    \operatorname{QCOp}(Z)
\]
be a chosen small affine-open basis. Because \(Z\) is quasi-separated, every finite intersection of quasi-compact opens is quasi-compact, and every quasi-compact open admits a finite affine Zariski cover. Thus \(\operatorname{AffOp}(Z)\) is a basis for this finite-cover site.

For \(j:U\hookrightarrow Z\) in \(\operatorname{QCOp}(Z)\), define
\[
    \mathcal F(U)
    :=
    \Gamma_U^{\mathrm{st}}
    \left(
      j^*\underline{\operatorname{map}}^{\partial}_{Z}(E,F)
    \right)
\]
and
\[
    \mathcal G(U)
    :=
    \operatorname{map}_{\mathcal R_U^0}
    (j^{*,0}E,j^{*,0}F).
\]
The functor \(\mathcal F\) satisfies descent for finite Zariski covers by effective descent for parameterized spectra: under the descent equivalence the tensor unit and the pulled-back coefficient form compatible objects, and mapping spectra in a limit of stable \(\infty\)-categories are computed as limits. The functor \(\mathcal G\) satisfies the same descent because Chapter 9 identifies \(\mathcal R_U^0\) with the limit of the Thom-stable differential categories on every finite Nisnevich, hence finite Zariski, Čech nerve.

If \(U\) is affine, the identity affine point represents the tensor unit in \(\mathcal E_U\). The base-change equivalence proved above, stable Yoneda, and the defining affine-point formula therefore give a canonical equivalence
\[
    \mathcal F(U)
    \simeq
    \mathcal G(U).
\]
These affine equivalences are natural under restriction in the affine basis. Restriction from finite-Zariski sheaves on \(\operatorname{QCOp}(Z)\) to the affine basis is fully faithful, so the affine equivalences extend uniquely to a natural equivalence
\[
    \mathcal F(U)
    \simeq
    \mathcal G(U)
\]
for every quasi-compact open \(U\subseteq Z\). Taking \(U=Z\) gives
\[
    \Gamma_Z^{\mathrm{st}}
    \underline{\operatorname{map}}^{\partial}_{Z}(E,F)
    \simeq
    \operatorname{map}_{\mathcal R_Z^0}(E,F).
\]
\end{proof}

\subsection{The big dualizing-density coefficient}

\begin{definition}[Big dualizing density coefficient]
\label{def:var-big-density-coefficient}
For \(D\in\mathcal R_S^0\), define
\[
    \boxed{
    \mathscr D_{x,D}^{\mathrm{big}}
    :=
    \underline{\operatorname{map}}^{\partial}_{X}
    (\one_X^\partial,x_e^!D)
    \in
    \mathcal E_X.
    }
\]
\end{definition}

\begin{definition}[Dualizing-density coefficient over \(Q\)]
\label{def:var-dualizing-density-Q}
Define
\[
    \boxed{
    M^\partial_{x,D}
    :=
    \Pi_\epsilon\mathscr D_{x,D}^{\mathrm{big}}
    \in\mathcal E_Q.
    }
\]
This is the \emph{brave-new differential dualizing-density coefficient} associated with \((x,D)\).
\end{definition}

\begin{theorem}[Density realization and arbitrary represented base change]
\label{thm:var-density-realization}
There is a canonical equivalence
\[
    \boxed{
    \Dens_{x,D}:
    \Gamma_Q^{\mathrm{st}}(M^\partial_{x,D})
    \xrightarrow{\sim}
    \operatorname{map}_{\mathcal R_X^0}
    (\one_X^\partial,x_e^!D).
    }
\]
For every represented \(t:T\to S\) in the standing differential base range, put
\[
    X_T:=X\times_ST,
    \qquad
    Q_T:=Q\times_ST,
    \qquad
    D_T:=t^{*,0}D.
\]
Arbitrary base change of the effective quotient gives a canonical equivalence
\[
    \boxed{
    Q_T\simeq Q_{X_T/T},
    }
\]
which is used throughout the remaining parameterized action constructions. Then there are canonical coherent base-change equivalences
\[
    \boxed{
    t_X^*\mathscr D_{x,D}^{\mathrm{big}}
    \simeq
    \mathscr D_{x_T,D_T}^{\mathrm{big}},
    \qquad
    t_Q^*M^\partial_{x,D}
    \simeq
    M^\partial_{x_T,D_T}.
    }
\]
Consequently
\[
    \boxed{
    \Gamma_{Q_T}^{\mathrm{st}}
    \bigl(t_Q^*M^\partial_{x,D}\bigr)
    \simeq
    \operatorname{map}_{\mathcal R_{X_T}^0}
    (\one_{X_T}^\partial,x_{T,e}^!D_T).
    }
\]
\end{theorem}

\begin{proof}
The adjunction \(\epsilon^*\dashv\Pi_\epsilon\) and Theorem~\ref{thm:var-big-differential-map-coefficient} give
\[
\begin{aligned}
    \Gamma_Q^{\mathrm{st}}(M^\partial_{x,D})
    &\simeq
    \Gamma_X^{\mathrm{st}}
    \mathscr D_{x,D}^{\mathrm{big}}
    \\
    &\simeq
    \operatorname{map}_{\mathcal R_X^0}
    (\one_X^\partial,x_e^!D).
\end{aligned}
\]

For represented base change, Theorem~\ref{thm:var-big-differential-map-coefficient} gives
\[
    t_X^*\mathscr D_{x,D}^{\mathrm{big}}
    \simeq
    \underline{\operatorname{map}}^{\partial}_{X_T}
    (\one_{X_T}^\partial,t_X^{*,0}x_e^!D).
\]
The explicit formula
\[
    x_e^!D=\Sigma_{T_x,0}x^{*,0}D
\]
and cotangent and differential Thom base change identify
\[
    t_X^{*,0}x_e^!D
    \simeq
    x_{T,e}^!D_T.
\]
This proves the first base-change equivalence. Right Beck--Chevalley for the Cartesian square of \(\epsilon\) gives the second. Applying the first part of the theorem after base change gives the final section-spectrum equivalence.
\end{proof}

\begin{proposition}[Comparison with the Chapter 10 smooth-site density coefficient]
\label{prop:var-density-smooth-comparison}
Write
\[
    \mathscr U_{x,D}:=(j_{\mathrm{Th},X}x_e^!D)_0
    \in
    \operatorname{Shv}_{\mathrm{Nis}}
    (\operatorname{Sm}^{\mathrm{fp}}_{/X};\operatorname{Sp}).
\]
The restriction of \(\mathscr D_{x,D}^{\mathrm{big}}\) to represented smooth finite-presentation tests, meaning the smooth-site sheaf
\[
    (u:U\to X)
    \longmapsto
    \Gamma_U^{\mathrm{st}}
    \bigl(u^*\mathscr D_{x,D}^{\mathrm{big}}\bigr),
\]
is canonically equivalent to \(\mathscr U_{x,D}\).
\end{proposition}

\begin{proof}
Let \(u:U\to X\) be represented, smooth, and of finite presentation. The base-change equivalence of Theorem~\ref{thm:var-big-differential-map-coefficient}, followed by its global-section equivalence, gives
\[
\begin{aligned}
    \Gamma_U^{\mathrm{st}}
    \bigl(u^*\mathscr D_{x,D}^{\mathrm{big}}\bigr)
    &\simeq
    \operatorname{map}_{\mathcal R_U^0}
    (\one_U^\partial,u^{*,0}x_e^!D)
    \\
    &\simeq
    \operatorname{map}_{\mathcal R_X^0}
    (u_\sharp^0\one_U^\partial,x_e^!D),
\end{aligned}
\]
where the second equivalence is the smooth adjunction.

Since \(\Real_U^0\) is strong symmetric monoidal,
\[
    \one_U^\partial
    \simeq
    \Real_U^0\Sigma_+^\infty h_U.
\]
Chapter 10 smooth-direct-image exchange therefore identifies
\[
    u_\sharp^0\one_U^\partial
    \simeq
    \Real_X^0\Sigma_+^\infty h_U,
\]
where on the right \(h_U\) is the represented smooth \(X\)-test. Stabilized evaluation then gives
\[
    \operatorname{map}_{\mathcal R_X^0}
    (\Real_X^0\Sigma_+^\infty h_U,x_e^!D)
    \simeq
    (j_{\mathrm{Th},X}x_e^!D)_0(U).
\]
These equivalences are natural under smooth restriction and therefore assemble to the asserted comparison of smooth-site sheaves.
\end{proof}

\subsection{The big exceptional-integration coefficient}

\begin{definition}[Big action-value coefficient]
\label{def:var-big-action-coefficient}
Define
\[
    \boxed{
    \mathbb A^\partial_{x,D}
    :=
    \underline{\operatorname{map}}^{\partial}_{S}
    (x_{e!}\one_X^\partial,D)
    \in
    \mathcal E_S.
    }
\]
For every represented affine point \(t:T\to S\), differential exceptional base change identifies its value with
\[
    \boxed{
    \mathbb A^\partial_{x,D}(T,t)
    \simeq
    \operatorname{map}_{\mathcal R_T^0}
    (x_{T,e!}\one_{X_T}^\partial,D_T).
    }
\]
\end{definition}

\begin{theorem}[Big exceptional integration is an equivalence]
\label{thm:var-big-exceptional-integration}
There is a canonical equivalence in \(\mathcal E_S\)
\[
    \boxed{
    \mathcal I^\partial_{x,D}:
    \Pi_qM^\partial_{x,D}
    \xrightarrow{\sim}
    \mathbb A^\partial_{x,D}.
    }
\]
At every represented affine point \(t:T\to S\), its value is the Chapter 10 differential Umkehr equivalence
\[
    \operatorname{map}_{\mathcal R_{X_T}^0}
    (\one_{X_T}^\partial,x_{T,e}^!D_T)
    \xrightarrow{\sim}
    \operatorname{map}_{\mathcal R_T^0}
    (x_{T,e!}\one_{X_T}^\partial,D_T).
\]
In particular
\[
    \boxed{
    \Gamma_S^{\mathrm{st}}(\mathbb A^\partial_{x,D})
    \simeq
    \operatorname{map}_{\mathcal R_S^0}
    (x_{e!}\one_X^\partial,D),
    }
\]
and on global sections \(\mathcal I^\partial_{x,D}\) is exactly the differential Umkehr adjunction.
\end{theorem}

\begin{proof}
Let \(t:T\to S\) be an affine point. Right Beck--Chevalley for \(q:Q\to S\), followed by Theorem~\ref{thm:var-density-realization}, gives
\[
\begin{aligned}
    (\Pi_qM^\partial_{x,D})(T,t)
    &\simeq
    \Gamma_{Q_T}^{\mathrm{st}}
    (t_Q^*M^\partial_{x,D})
    \\
    &\simeq
    \operatorname{map}_{\mathcal R_{X_T}^0}
    (\one_{X_T}^\partial,x_{T,e}^!D_T).
\end{aligned}
\]
The differential exceptional adjunction for \(x_T\) identifies this spectrum with
\[
    \operatorname{map}_{\mathcal R_T^0}
    (x_{T,e!}\one_{X_T}^\partial,D_T),
\]
which is the affine value of \(\mathbb A^\partial_{x,D}\). Differential exceptional base change and composition identify these equivalences under every morphism of affine points. Affine-point recognition therefore assembles them to the displayed equivalence in \(\mathcal E_S\). The global statement is Theorem~\ref{thm:var-big-differential-map-coefficient} at \(Z=S\), and the description of the global map is the counit description of the adjunction.
\end{proof}

\subsection{Cartan field moduli and universal Lagrangian evaluation}

\begin{definition}[Cartan field moduli object]
\label{def:var-field-moduli}
For a Cartan object \((A,\rho_A)\) with descent object \(a:\bar A\to Q\), define
\[
    \boxed{
    \mathbf{Fld}_S(A)
    :=
    \Pi_q\bar A
    \in
    (\XSp)_{/S}.
    }
\]
Its space of global Cartan fields is
\[
    \boxed{
    \Fld(A)
    :=
    \Map_S(S,\mathbf{Fld}_S(A))
    \simeq
    \Map_Q(Q,\bar A).
    }
\]
\end{definition}

\begin{proposition}[Parameterized field interpretation and cofree fields]
\label{prop:var-field-moduli-representing}
For every \(t:T\to S\),
\[
    \boxed{
    \Map_S(T,\mathbf{Fld}_S(A))
    \simeq
    \Map_{Q_T}(Q_T,\bar A_T).
    }
\]
If \(A=K(Y)=\Jinf Y\) is cofree, then
\[
    \boxed{
    \mathbf{Fld}_S(K(Y))
    \simeq
    \Pi_xY,
    }
\]
and consequently
\[
    \Map_S(T,\mathbf{Fld}_S(K(Y)))
    \simeq
    \Map_{X_T}(X_T,Y_T).
\]
\end{proposition}

\begin{proof}
The first equivalence is the adjunction \(q^*\dashv\Pi_q\).  For the cofree object, Proposition~\ref{prop:var-cofree-descent} gives \(\bar A\simeq\Pi_\epsilon Y\). Composition of dependent products gives
\[
    \Pi_q\bar A
    \simeq
    \Pi_q\Pi_\epsilon Y
    \simeq
    \Pi_xY.
\]
Apply the same adjunction after pullback to \(T\).
\end{proof}

\begin{construction}[Universal coefficient-valued Lagrangian evaluation]
\label{con:var-universal-lagrangian-evaluation}
Let \(M\in\mathcal E_Q\) and let
\[
    L:\one_Q\longrightarrow\Sec^0(A;M)[n]
\]
be a Lagrangian with adjoint local section
\[
    \ell_L:\one_{\bar A}\longrightarrow a^*M[n].
\]
Put
\[
    F:=\mathbf{Fld}_S(A),
    \qquad
    f:F\to S.
\]
The counit of \(q^*\dashv\Pi_q\) gives the universal field evaluation
\[
    \operatorname{ev}:
    Q\times_SF\longrightarrow\bar A
\]
over \(Q\). Write
\[
    p:Q\times_SF\to F,
    \qquad
    r:Q\times_SF\to Q
\]
for the projections. Pulling back \(\ell_L\) gives
\[
    \one_{Q\times_SF}
    \longrightarrow
    r^*M[n].
\]
Adjunction along \(p\), followed by right Beck--Chevalley
\[
    \Pi_pr^*M
    \simeq
    f^*\Pi_qM,
\]
gives a section
\[
    \one_F
    \longrightarrow
    f^*\Pi_qM[n].
\]
Equivalently, define the \emph{universal coefficient-valued Lagrangian-evaluation morphism}
\[
    \boxed{
    \mathbf L_L:
    \mathbf{Fld}_S(A)
    \longrightarrow
    \Omega_S^\infty
    \bigl(\Pi_qM[n]\bigr).
    }
\]
\end{construction}

\begin{proposition}[Evaluation on a parameterized field]
\label{prop:var-universal-lagrangian-field-evaluation}
Let \(t:T\to S\) and let
\[
    \phi_T:Q_T\to\bar A_T
\]
be the field corresponding to \(c:T\to\mathbf{Fld}_S(A)\).  The pullback of \(\mathbf L_L\) along \(c\) is the section
\[
    \boxed{
    \phi_T^*\ell_L:
    \one_{Q_T}
    \longrightarrow
    t_Q^*M[n].
    }
\]
\end{proposition}

\begin{proof}
This is the counit characterization of the universal evaluation map.  Pulling
\[
    \operatorname{ev}:Q\times_SF\to\bar A
\]
back along \(c:T\to F\) gives precisely the section \(\phi_T:Q_T\to\bar A_T\).  The adjunction \(p^*\dashv\Pi_p\) and right Beck--Chevalley are natural under that pullback, so the resulting coefficient section is \(\phi_T^*\ell_L\).
\end{proof}

\begin{definition}[Differential Lagrangian]
\label{def:var-differential-lagrangian}
Fix \(D\in\mathcal R_S^0\). A \emph{degree-\(n\) differential Lagrangian with target \(D\)} on the Cartan object \(A\) is a Lagrangian class with coefficient \(M^\partial_{x,D}\):
\[
    \boxed{
    L:\one_Q
    \longrightarrow
    \Sec^0(A;M^\partial_{x,D})[n].
    }
\]
For such an \(L\), Construction~\ref{con:var-universal-lagrangian-evaluation} gives
\[
    \mathbf L_L:
    \mathbf{Fld}_S(A)
    \longrightarrow
    \Omega_S^\infty
    \bigl(\Pi_qM^\partial_{x,D}[n]\bigr).
\]
\end{definition}


\begin{proposition}[Differential-density specialization of the Lepage package]
\label{prop:var-differential-lepage-specialization}
Let \(L\) be a differential Lagrangian with coefficient \(M^\partial_{x,D}\).  The constructions of Section~\ref{sec:var-lepage-presymplectic} specialize canonically to parameterized objects
\[
    \mathbf{Lep}^{\,q}_\partial(L)
    :=
    \mathbf{Lep}^{\,q}(L)
    \longrightarrow Q
\]
and universal sections
\[
    \Theta_{L,\partial}^{(q)}:
    \one_{\mathbf{Lep}^{\,q}_\partial(L)}
    \longrightarrow
    (\pi_L^q)^*
    F_H^q\Sec^1(A;M^\partial_{x,D})[n],
\]
\[
    \boxed{
    \omega_{L,\partial}^{(q)}:
    \one_{\mathbf{Lep}^{\,q}_\partial(L)}
    \longrightarrow
    (\pi_L^q)^*
    \TransH^{2,q}(A;M^\partial_{x,D})[n].
    }
\]
They satisfy the same complete recurrence, coherent vertical closedness, and change-of-realization formula.

The differential exceptional adjunction by itself does not turn \(\omega_{L,\partial}^{(q)}\) into a spacetime-boundary integral.  For a boundary correspondence
\[
    Z\xrightarrow{e}X,
    \qquad
    Z\xrightarrow{\pi}B,
\]
Chapter 10 transgression requires an actual constructed coefficient morphism
\[
    \gamma:e^{*,0}D_X\longrightarrow\pi_e^!D_B.
\]
To transgress a coherent de Rham or horizontal-stage system while retaining normalized cochains, coherently closed cochains, the coherent differential, and its higher relations, one requires a coherent family
\[
    \gamma_q:
    e^{*,0}D_X^q
    \longrightarrow
    \pi_e^!D_B^q
\]
compatible with the stage transitions and with the adjacent-stage exact diagrams.  These coefficient morphisms are not constructed by the exceptional adjunction alone, and no spacetime-boundary transgression is asserted here without them.  Whenever such morphisms have been constructed from additional geometry or coefficients, Chapter 10 differential transgression applies functorially to the displayed coefficient-valued current and preserves the corresponding coherent differential relations.
\end{proposition}

\begin{proof}
The first statement is substitution of \(M^\partial_{x,D}\) into the finite-limit and horizontal-totalization constructions of Section~\ref{sec:var-lepage-presymplectic}.  For the boundary statement, Chapter 10 defines differential transgression by first supplying a coefficient morphism
\[
    e^{*,0}D_X\longrightarrow\pi_e^!D_B
\]
and then applying \(\pi_{e!}\) followed by the exceptional counit.  The exceptional adjunction constructs the counit but not the preceding coefficient morphism.  Chapter 10's coherent-stage transgression theorem further shows that preservation of normalized and coherently closed cochains and of the coherent differential is obtained precisely when the coefficient morphisms form a stagewise family compatible with transitions and adjacent-stage exact diagrams.  This is the stated interface; no additional boundary structure is inserted into the present variational package.
\end{proof}

\subsection{The brave-new action as a big moduli morphism}

\begin{definition}[Big brave-new differential action]
\label{def:var-action-functional}
Let \(L\) be a differential Lagrangian. Define
\[
    \boxed{
    \mathbf S_L:
    \mathbf{Fld}_S(A)
    \longrightarrow
    \Omega_S^\infty
    \bigl(\mathbb A^\partial_{x,D}[n]\bigr)
    }
\]
by
\[
    \mathbf S_L
    :=
    \Omega_S^\infty(\mathcal I^\partial_{x,D}[n])
    \circ
    \mathbf L_L.
\]
Its map on global \(S\)-points is denoted
\[
    \boxed{
    \mathcal S_L:
    \Fld(A)
    \longrightarrow
    \Act^n_{x,D},
    }
\]
where
\[
    \Act^n_{x,D}
    :=
    \Omega^\infty
    \operatorname{map}_{\mathcal R_S^0}
    (x_{e!}\one_X^\partial,D[n]).
\]
\end{definition}

\begin{proposition}[The global action is the exceptional trace]
\label{prop:var-global-action-trace}
For a global Cartan field \(\phi:Q\to\bar A\), Theorem~\ref{thm:var-density-realization} gives a differential density
\[
    \mathcal L_L(\phi):
    \one_X^\partial\longrightarrow x_e^!D[n].
\]
The value of the global action is
\[
    \boxed{
    \mathcal S_L(\phi)
    =
    \left(
    x_{e!}\one_X^\partial
    \xrightarrow{x_{e!}\mathcal L_L(\phi)}
    x_{e!}x_e^!D[n]
    \xrightarrow{\operatorname{tr}^\partial_{x,D}}
    D[n]
    \right).
    }
\]
\end{proposition}

\begin{proof}
By Theorem~\ref{thm:var-big-exceptional-integration}, global sections of \(\mathcal I^\partial_{x,D}\) are exactly the Chapter 10 mapping-spectrum Umkehr equivalence.  Apply it to the global section corresponding under Theorem~\ref{thm:var-density-realization} to \(\mathcal L_L(\phi)\).  The counit description of the Umkehr equivalence gives the displayed trace formula.
\end{proof}

\begin{remark}[Why the canonical action is not forced to be scalar-valued]
\label{rem:var-action-not-scalar}
The canonical orientation-free action value is a morphism
\[
    x_{e!}\one_X^\partial\longrightarrow D[n].
\]
No canonical map \(\one_S^\partial\to x_{e!}\one_X^\partial\) has been constructed for a general smooth separated finite-presentation \(x\). Therefore the chapter does not silently replace the action target by
\[
    \Omega^\infty\operatorname{map}_{\mathcal R_S^0}(\one_S^\partial,D[n]).
\]
If later geometry supplies a differential fundamental or compact-support class
\[
    [X]^\partial:\one_S^\partial\longrightarrow x_{e!}\one_X^\partial,
\]
then precomposition produces the scalarized action
\[
    \one_S^\partial\xrightarrow{[X]^\partial}x_{e!}\one_X^\partial
    \xrightarrow{\mathcal S_L(\phi)}D[n].
\]
No such class is part of the present variational package.
\end{remark}

\begin{proposition}[Represented parameter families recover the differential action]
\label{prop:var-action-smooth-parameter}
Let \(t:T\to S\) be represented in the standing differential base range, and let
\[
    \phi_T:Q_T\longrightarrow\bar A_T
\]
be a \(T\)-family of Cartan fields, corresponding to \(c:T\to\mathbf{Fld}_S(A)\). Under the canonical representing equivalence
\[
\begin{aligned}
    \Map_S\left(
      T,
      \Omega_S^\infty\bigl(\mathbb A^\partial_{x,D}[n]\bigr)
    \right)
    &\simeq
    \Omega^\infty
    \operatorname{map}_{\mathcal R_T^0}
    \left(
      x_{T,e!}\one_{X_T}^\partial,
      D_T[n]
    \right),
\end{aligned}
\]
the composite
\[
    T
    \xrightarrow{c}
    \mathbf{Fld}_S(A)
    \xrightarrow{\mathbf S_L}
    \Omega_S^\infty\bigl(\mathbb A^\partial_{x,D}[n]\bigr)
\]
corresponds canonically to the differential exceptional action
\[
    \boxed{
    \mathcal S_{L,T}(\phi_T):
    x_{T,e!}\one_{X_T}^\partial
    \longrightarrow
    D_T[n].
    }
\]
These identifications are natural under further represented base change. In particular they apply to affine split square-zero and bi-square-zero parameter thickenings.
\end{proposition}

\begin{proof}
The parameterized zeroth-space representing property gives
\[
\begin{aligned}
    \Map_S\left(
      T,
      \Omega_S^\infty\bigl(\mathbb A^\partial_{x,D}[n]\bigr)
    \right)
    &\simeq
    \Omega^\infty
    \Gamma_T^{\mathrm{st}}
    \bigl(t^*\mathbb A^\partial_{x,D}[n]\bigr).
\end{aligned}
\]
By Theorem~\ref{thm:var-big-differential-map-coefficient}, differential exceptional base change, and the global-section equivalence for the big mapping coefficient, the right side is canonically
\[
    \Omega^\infty
    \operatorname{map}_{\mathcal R_T^0}
    \left(
      x_{T,e!}\one_{X_T}^\partial,
      D_T[n]
    \right),
\]
which proves the displayed representing equivalence.

Proposition~\ref{prop:var-universal-lagrangian-field-evaluation} identifies the value of \(\mathbf L_L\) on the field \(c\) with the coefficient section
\[
    \phi_T^*\ell_L:
    \one_{Q_T}\longrightarrow t_Q^*M^\partial_{x,D}[n].
\]
The arbitrary represented base-change equivalence of Theorem~\ref{thm:var-density-realization} realizes this section as a differential density
\[
    \mathcal L_L(\phi_T):
    \one_{X_T}^\partial\longrightarrow x_{T,e}^!D_T[n].
\]
The base change of \(\mathcal I^\partial_{x,D}\) is, by its affine-point construction and differential exceptional Beck--Chevalley, the differential Umkehr equivalence for \(x_T\). Therefore the image of \(\mathcal L_L(\phi_T)\) is the composite
\[
    x_{T,e!}\one_{X_T}^\partial
    \xrightarrow{x_{T,e!}\mathcal L_L(\phi_T)}
    x_{T,e!}x_{T,e}^!D_T[n]
    \xrightarrow{\operatorname{tr}^\partial_{x_T,D_T}}
    D_T[n].
\]
By Definition~\ref{def:var-action-functional}, this is precisely the morphism corresponding to \(\mathbf S_Lc\) under the representing equivalence above. The base-change naturality is inherited from differential exceptional Beck--Chevalley and the base-change coherence of the big differential mapping coefficients.
\end{proof}

\subsection{Split square-zero variations of the field moduli object}

\begin{definition}[Split square-zero field variation]
\label{def:var-field-variation-family}
Let
\[
    F:=\mathbf{Fld}_S(A).
\]
A \emph{split square-zero variation of a \(T\)-field} consists of an affine spectral scheme \(t:T=\operatorname{Spec}B\to S\), a connective \(N\in\operatorname{Mod}_B\), a centre
\[
    c:T\longrightarrow F,
\]
and a lift
\[
    \widetilde c:T[N]\longrightarrow F
\]
over \(S\) together with a specified equivalence
\[
    \widetilde c\,i_N\simeq c.
\]
Under Proposition~\ref{prop:var-field-moduli-representing}, this is equivalently a field
\[
    \phi_T:Q_T\to\bar A_T
\]
and a field
\[
    \widetilde\phi:
    Q_{T[N]}\to\bar A_{T[N]},
\]
where
\[
    Q_T:=Q\times_ST,
    \qquad
    Q_{T[N]}:=Q\times_ST[N].
\]
The restriction of \(\widetilde\phi\) along \(Q_T\to Q_{T[N]}\) is \(\phi_T\).
\end{definition}

\begin{definition}[First difference of a big coefficient-valued field map]
\label{def:var-big-first-difference}
Let \(E\in\mathcal E_S\), let
\[
    \sigma:F\longrightarrow\Omega_S^\infty E[n],
\]
and let \(v=(T,N,c,\widetilde c)\) be a split square-zero field variation. Define
\[
    K_E^{(1)}(v)
    :=
    \fib\left(
      \Gamma_{T[N]}^{\mathrm{st}}(E_{T[N]})
      \longrightarrow
      \Gamma_T^{\mathrm{st}}(E_T)
    \right).
\]
A coefficient morphism \(f:E\to E'\) induces by functoriality of fibres a canonical morphism
\[
    K^{(1)}(f)(v):
    K_E^{(1)}(v)\longrightarrow K_{E'}^{(1)}(v).
\]
The two pulled-back values satisfy
\[
    i_N^*\widetilde c^*\sigma
    \simeq
    c^*\sigma
    \simeq
    i_N^*r_N^*c^*\sigma.
\]
Their difference therefore determines a canonical point
\[
    \boxed{
    \Delta_v\sigma
    :=
    \widetilde c^*\sigma-r_N^*c^*\sigma
    \in
    \Omega^\infty K_E^{(1)}(v)[n].
    }
\]
\end{definition}

\begin{definition}[Evaluation of a first Euler source on a local split variation]
\label{def:var-euler-evaluation-on-variation}
Let
\[
    e:\one_{\bar A}\longrightarrow\Vfirst_M(A)[n]
\]
be a first-order source and let
\[
    w=(U,N_U,a_U,\widetilde a_U)
\]
be a split square-zero variation diagram over \(\bar A/Q\). By Theorem~\ref{thm:var-first-kan-sheaf}, restriction of \(e\) to the affine centre followed by the canonical right-Kan limit projection at \(w\) gives a canonical point
\[
    \boxed{
    \langle e,w\rangle
    \in
    \Omega^\infty K_M^{(1)}(w)[n].
    }
\]
\end{definition}

\begin{lemma}[Affine-point coproduct is an effective epimorphism]
\label{lem:var-affine-point-effective-epi}
Let \(Y\to S\) be an object of \((\XSp)_{/S}\), and form the coproduct over the chosen small affine-point skeleton
\[
    e_Y:
    \coprod_{(U,u)\in\operatorname{AffPt}_S(Y)}
    U
    \longrightarrow
    Y.
\]
Then \(e_Y\) is an effective epimorphism.
\end{lemma}

\begin{proof}
Factor \(e_Y\) by its effective-image factorization
\[
    \coprod_{(U,u)}U
    \twoheadrightarrow
    I
    \lhook\joinrel\longrightarrow
    Y.
\]
Every affine-point map \(u:U\to Y\) factors through \(I\) by construction of the image. Since \(I\to Y\) is a monomorphism, these factorizations are compatible, through contractible spaces of choices, with every morphism in the affine-point indexing \(\infty\)-category. Chapter 9 affine density in slices gives
\[
    \operatorname*{colim}_{(U,u)\in\operatorname{AffPt}_S(Y)}
    U
    \xrightarrow{\sim}
    Y,
\]
so the compatible factorizations induce a section
\[
    Y\longrightarrow I
\]
of \(I\to Y\). A monomorphism admitting a section is an equivalence in an \(\infty\)-topos. Hence \(I\simeq Y\), and \(e_Y\) is an effective epimorphism.
\end{proof}

\begin{construction}[Local variations induced by a field variation]
\label{con:var-induced-local-field-variations}
Let \(v=(T,N,c,\widetilde c)\) correspond to \((\phi_T,\widetilde\phi)\) as in Definition~\ref{def:var-field-variation-family}. For every affine point
\[
    u:U\longrightarrow Q_T,
\]
put
\[
    U[N_u]
    :=
    U\times_TT[N].
\]
Since \(T[N]\to T\) is affine split square-zero, \(U[N_u]\) is affine and \(U\to U[N_u]\) is its pullback split square-zero thickening. Restricting \(\phi_T\) and \(\widetilde\phi\) gives a split square-zero variation diagram
\[
    \boxed{
    v_u\in\mathcal V^{(1)}_{\bar A/Q}
    }
\]
centred at
\[
    \phi_Tu:U\longrightarrow\bar A.
\]
The construction is functorial in affine restriction \(U'\to U\).
\end{construction}

\begin{definition}[Parameterized first-difference coefficient of a field variation]
\label{def:var-parameterized-first-field-difference}
Let \(M\in\mathcal E_Q\) and let
\[
    v=(T,N,c,\widetilde c)
\]
be a split square-zero variation of the field-moduli object with structural map \(t:T\to S\). Write
\[
    M_T:=t_Q^*M\in\mathcal E_{Q_T},
\]
and let
\[
    r_Q:Q_{T[N]}\longrightarrow Q_T,
    \qquad
    i_Q:Q_T\longrightarrow Q_{T[N]}
\]
be the split retraction and central section induced from \(r_N\) and \(i_N\). Define
\[
    \boxed{
    \mathscr K^{(1)}_{M,v}
    :=
    \fib\left(
      \Pi_{r_Q}r_Q^*M_T
      \longrightarrow
      M_T
    \right)
    \in\mathcal E_{Q_T},
    }
\]
where the displayed arrow is restriction along \(i_Q\), using
\[
    \Pi_{r_Q}\Pi_{i_Q}i_Q^*r_Q^*M_T
    \simeq
    \Pi_{r_Qi_Q}M_T
    \simeq
    M_T.
\]
\end{definition}

\begin{lemma}[First-difference Beck--Chevalley and local evaluation]
\label{lem:var-first-difference-beck-chevalley}
Let \(M\in\mathcal E_Q\) and let \(v=(T,N,c,\widetilde c)\) be a split square-zero variation of the field-moduli object. There is a canonical equivalence
\[
    \boxed{
    \Gamma_{Q_T}^{\mathrm{st}}
    \mathscr K^{(1)}_{M,v}
    \simeq
    K_{\Pi_qM}^{(1)}(v)
    \simeq
    \fib\left(
      \Gamma_{Q_{T[N]}}^{\mathrm{st}}(M)
      \longrightarrow
      \Gamma_{Q_T}^{\mathrm{st}}(M)
    \right).
    }
\]
For every affine point \(u:U\to Q_T\), there is a canonical equivalence
\[
    \boxed{
    \mathscr K^{(1)}_{M,v}(U,u)
    \simeq
    K_M^{(1)}(v_u).
    }
\]
Writing \(\phi_T:Q_T\to\bar A\) also for the composite of the descended \(T\)-field with \(\bar A_T\to\bar A\), the right-Kan evaluation morphisms of Theorem~\ref{thm:var-first-kan-sheaf} assemble to a canonical morphism
\[
    \boxed{
    \operatorname{ev}^{(1)}_v:
    \phi_T^*\Vfirst_M(A)
    \longrightarrow
    \mathscr K^{(1)}_{M,v}.
    }
\]
All three constructions are natural in the complete variation diagram and in \(M\).
\end{lemma}

\begin{proof}
Stable sections and the adjunction \(r_Q^*\dashv\Pi_{r_Q}\) give
\[
\begin{aligned}
    \Gamma_{Q_T}^{\mathrm{st}}
    \left(
      \Pi_{r_Q}r_Q^*M_T
    \right)
    &\simeq
    \Gamma_{Q_{T[N]}}^{\mathrm{st}}
    (r_Q^*M_T).
\end{aligned}
\]
Since stable sections preserve fibres, the definition of \(\mathscr K^{(1)}_{M,v}\) therefore gives
\[
    \Gamma_{Q_T}^{\mathrm{st}}
    \mathscr K^{(1)}_{M,v}
    \simeq
    \fib\left(
      \Gamma_{Q_{T[N]}}^{\mathrm{st}}(M)
      \longrightarrow
      \Gamma_{Q_T}^{\mathrm{st}}(M)
    \right).
\]
Right Beck--Chevalley for the Cartesian squares over \(T[N]\) and \(T\) identifies the pullbacks of \(\Pi_qM\) with the corresponding dependent products from \(Q_{T[N]}\) and \(Q_T\). Stable sections of these dependent products are the two displayed section spectra, and naturality identifies the restriction maps. This gives the first boxed equivalence with \(K_{\Pi_qM}^{(1)}(v)\).

Now let \(u:U\to Q_T\) be affine and put
\[
    U[N_u]:=U\times_{Q_T}Q_{T[N]}
    \simeq
    U\times_TT[N].
\]
Right Beck--Chevalley for the square obtained by pulling back \(r_Q\) along \(u\) gives
\[
    u^*\mathscr K^{(1)}_{M,v}
    \simeq
    \fib\left(
      \Pi_{r_u}r_u^*M_U
      \longrightarrow
      M_U
    \right),
\]
where \(r_u:U[N_u]\to U\) is the pulled-back split retraction. Taking stable sections over the affine centre \(U\) identifies its value with
\[
    \fib\left(
      \Gamma_{U[N_u]}^{\mathrm{st}}(M)
      \longrightarrow
      \Gamma_U^{\mathrm{st}}(M)
    \right)
    =
    K_M^{(1)}(v_u).
\]

Finally, affine-point evaluation of the pullback coefficient gives
\[
    \bigl(\phi_T^*\Vfirst_M(A)\bigr)(U,u)
    \simeq
    \Vfirst_M(A)(U,\phi_Tu).
\]
The induced local variation \(v_u\) is an object of the right-Kan indexing category over the centre \(\phi_Tu\), so Theorem~\ref{thm:var-first-kan-sheaf} supplies a canonical map
\[
    \Vfirst_M(A)(U,\phi_Tu)
    \longrightarrow
    K_M^{(1)}(v_u)
    \simeq
    \mathscr K^{(1)}_{M,v}(U,u).
\]
Construction~\ref{con:var-induced-local-field-variations} is functorial under affine restriction, and the right-Kan projections are the structure maps of the universal cone. Hence these affine-point maps are natural. The Chapter 10 affine-point presentation of \(\mathcal E_{Q_T}\) therefore assembles them uniquely to \(\operatorname{ev}^{(1)}_v\).
\end{proof}

\begin{theorem}[Euler--Lagrange represents the square-zero first variation]
\label{thm:var-euler-represents-first-variation}
Let \(L\) be any Lagrangian with coefficient \(M\), let
\[
    \mathbf L_L:
    \mathbf{Fld}_S(A)
    \to
    \Omega_S^\infty(\Pi_qM[n])
\]
be its universal coefficient-valued field evaluation, and let
\[
    v=(T,N,c,\widetilde c)
\]
be a split square-zero field variation. Under the equivalence
\[
    K_{\Pi_qM}^{(1)}(v)
    \simeq
    \Gamma_{Q_T}^{\mathrm{st}}
    \mathscr K^{(1)}_{M,v}
\]
of Lemma~\ref{lem:var-first-difference-beck-chevalley}, the first difference
\[
    \Delta_v\mathbf L_L
\]
is the global section of \(\mathscr K^{(1)}_{M,v}[n]\) represented by the composite
\[
    \boxed{
    \one_{Q_T}
    \xrightarrow{\;\phi_T^*\mathsf{EL}_L\;}
    \phi_T^*\Vfirst_M(A)[n]
    \xrightarrow{\;\operatorname{ev}^{(1)}_v[n]\;}
    \mathscr K^{(1)}_{M,v}[n].
    }
\]
Equivalently, under right Beck--Chevalley it is the first difference
\[
    \widetilde\phi^*\ell_L
    -
    r_Q^*\phi_T^*\ell_L
\]
in the pulled-back coefficient on \(Q_{T[N]}\). For every affine point \(u:U\to Q_T\), its restriction to the induced local variation \(v_u\) is canonically
\[
    \boxed{
    \left.\Delta_v\mathbf L_L\right|_{v_u}
    \simeq
    \langle\mathsf{EL}_L,v_u\rangle.
    }
\]
Hence the first Euler--Lagrange section is exactly the universal split-square-zero first variation of the Lagrangian field-evaluation morphism.
\end{theorem}

\begin{proof}
Degree-zero normalization identifies the coherent infinitesimal boundary of the local section with the ordered two-vertex difference
\[
    d_{\mathrm{inf}}\ell_L
    =
    e^*\ell_L-c^*\ell_L,
\]
with the sign fixed by the Chapter 10 \v{C}ech-face convention. The local full Euler section is the adjoint of this normalized one-cochain. Construction~\ref{con:var-local-square-zero-restriction} defines \(\lambda^1_{\mathrm{loc}}\) by restricting that cochain along
\[
    \sigma_{v_u}:U[N_u]\to R_{\bar A/Q,1}.
\]
Its two vertex components are the centre field pulled back along the retraction and the varied field. Therefore
\[
    \langle\mathsf{EL}_L,v_u\rangle
    =
    \widetilde\phi^*\ell_L
    -
    r_Q^*\phi_T^*\ell_L
\]
after restriction to \(U[N_u]\).

On the other hand, Definition~\ref{def:var-big-first-difference} computes the first difference of the universal evaluation morphism by the same subtraction before affine restriction. Lemma~\ref{lem:var-first-difference-beck-chevalley} identifies this first difference with a global section of \(\mathscr K^{(1)}_{M,v}[n]\). The composite
\[
    \one_{Q_T}
    \xrightarrow{\phi_T^*\mathsf{EL}_L}
    \phi_T^*\Vfirst_M(A)[n]
    \xrightarrow{\operatorname{ev}^{(1)}_v[n]}
    \mathscr K^{(1)}_{M,v}[n]
\]
is a second global section of the same coefficient. At every affine point \(u:U\to Q_T\), the defining affine-point value of \(\operatorname{ev}^{(1)}_v\) identifies the second section with
\[
    \langle\mathsf{EL}_L,v_u\rangle,
\]
while the preceding calculation identifies the first-difference section with the same point.

Under the Chapter 10 affine-point presentation, morphisms in \(\mathcal E_{Q_T}\) are natural transformations of the corresponding affine-point spectral Nisnevich sheaves. The two sections therefore agree because their values agree naturally at every affine point. This proves the global boxed identification, and the displayed local formula is its affine-point evaluation.
\end{proof}

\subsection{First variation of the big action}

\begin{definition}[The \(\infty\)-category of first field variations]
\label{def:var-field-first-variation-category}
For \(F=\mathbf{Fld}_S(A)\), let \(\mathcal V_F^{(1)}\) be the \(\infty\)-category whose objects are split square-zero field variations \((T,N,c,\widetilde c)\) of Definition~\ref{def:var-field-variation-family} and whose morphisms are the corresponding structured morphisms of centres, modules, split thickenings, and maps to \(F\). For \(E\in\mathcal E_S\), the assignment
\[
    v\longmapsto K_E^{(1)}(v)
\]
defines a functor
\[
    K_E^{(1)}:(\mathcal V_F^{(1)})^{\mathrm{op}}\to\operatorname{Sp}.
\]
A morphism \(f:E\to E'\) induces a natural transformation \(K^{(1)}(f):K_E^{(1)}\to K_{E'}^{(1)}\).
\end{definition}

\begin{proposition}[First difference is natural on the variation \(\infty\)-category]
\label{prop:var-first-difference-natural-category}
For a morphism \(\sigma:F\to\Omega_S^\infty E[n]\), the points \(\Delta_v\sigma\) assemble into a natural section of the functor \(K_E^{(1)}[n]\) on \((\mathcal V_F^{(1)})^{\mathrm{op}}\). For every coefficient morphism \(f:E\to E'\), there is a canonical equality of natural sections
\[
    \boxed{
    \Delta^{(1)}(\Omega_S^\infty(f[n])\sigma)
    \simeq
    K^{(1)}(f)\bigl(\Delta^{(1)}\sigma\bigr).
    }
\]
\end{proposition}

\begin{proof}
Every first difference is defined by subtraction of two pullbacks in a stable category. Structured morphisms of field variations carry the two pullbacks and their central identification to the corresponding data for the target variation, so the differences are natural. A coefficient morphism preserves subtraction, the zero object, fibres, and all paths because it is a morphism in a stable \(\infty\)-category, giving the displayed natural identity.
\end{proof}

\begin{theorem}[Brave-new first-variation formula]
\label{thm:var-action-first-variation}
Let \(L\) be a differential Lagrangian. On the entire variation \(\infty\)-category \(\mathcal V_F^{(1)}\), the coefficient equivalence \(\mathcal I^\partial_{x,D}\) induces an equivalence of first-difference functors carrying the natural first-difference section of \(\mathbf L_L\) to that of \(\mathbf S_L\). Equivalently, for every split square-zero field variation
\[
    v=(T,N,c,\widetilde c),
\]
one has
\[
    \boxed{
    \Delta_v\mathbf S_L
    \simeq
    K^{(1)}(\mathcal I^\partial_{x,D})(v)
    \bigl(
      \Delta_v\mathbf L_L
    \bigr).
    }
\]
The map
\[
    K^{(1)}(\mathcal I^\partial_{x,D})(v)
\]
is an equivalence. For every affine \(u:U\to Q_T\), Theorem~\ref{thm:var-euler-represents-first-variation} identifies the local coefficient of \(\Delta_v\mathbf L_L\) with
\[
    \langle\mathsf{EL}_L,v_u\rangle.
\]
Thus the split-square-zero first variation of the big action is literally differential exceptional integration, on the represented affine square-zero parameter \(T[N]\), of the coefficient-valued first difference whose affine local evaluations are represented by the Euler--Lagrange source.
\end{theorem}

\begin{proof}
By Definition~\ref{def:var-action-functional},
\[
    \mathbf S_L
    =
    \Omega_S^\infty(\mathcal I^\partial_{x,D}[n])
    \circ
    \mathbf L_L.
\]
The first-difference construction is a finite fibre in a stable category. Hence every morphism of parameterized spectra induces a morphism of first-difference fibres and preserves subtraction, zero objects, paths, and all higher homotopies. Applying it to the coefficient equivalence \(\mathcal I^\partial_{x,D}\) gives the displayed identity and shows that
\(K^{(1)}(\mathcal I^\partial_{x,D})(v)\) is an equivalence.

The affine description is Theorem~\ref{thm:var-euler-represents-first-variation}. Proposition~\ref{prop:var-action-smooth-parameter} applies to the represented affine bases \(T\) and \(T[N]\), so the two coefficient values and their restriction map are exactly the differential exceptional adjunction data for \(x_T\) and \(x_{T[N]}\). Consequently the displayed identity is the differential exceptional integral of the first variation on the actual square-zero parameter.
\end{proof}

\begin{definition}[Critical-field moduli object]
\label{def:var-critical-field-moduli}
Let
\[
    \Crit(L)\longrightarrow\bar A\longrightarrow Q
\]
be the first Euler critical object. Define
\[
    \boxed{
    \mathbf{CritFld}_S(L)
    :=
    \Pi_q\Crit(L)
    \longrightarrow
    \Pi_q\bar A
    =
    \mathbf{Fld}_S(A).
    }
\]
A \(T\)-point is exactly a \(T\)-family of Cartan fields whose descended section
\[
    Q_T\to\bar A_T
\]
admits a lift to \(\Crit(L)_T\).
\end{definition}

\begin{theorem}[Euler critical fields are stationary]
\label{thm:var-critical-stationary}
Let \(v=(T,N,c,\widetilde c)\) be a split square-zero variation of a \(T\)-field. If the centre
\[
    c:T\longrightarrow\mathbf{Fld}_S(A)
\]
factors through
\[
    \mathbf{CritFld}_S(L)\longrightarrow\mathbf{Fld}_S(A),
\]
then
\[
    \boxed{
    \Delta_v\mathbf S_L\simeq0.
    }
\]
The nullhomotopy is canonical relative to the chosen critical lift and retains all higher homotopies in the variation diagram. No converse is asserted because the chapter does not prove that split square-zero variations of the global field-moduli object generate every affine-local split variation of the descended field which is detected by the right-Kan Euler coefficient.
\end{theorem}

\begin{proof}
A lift of the centre through \(\mathbf{CritFld}_S(L)\) is, by dependent adjunction, a lift of the descended field
\[
    \phi_T:Q_T\to\bar A_T
\]
through \(\Crit(L)_T\). By definition of the homotopy zero fibre, the chosen lift supplies a specified nullhomotopy
\[
    \phi_T^*\mathsf{EL}_L\simeq0
\]
as a section of \(\phi_T^*\Vfirst_M(A)[n]\). Applying the canonical evaluation morphism
\[
    \operatorname{ev}^{(1)}_v:
    \phi_T^*\Vfirst_M(A)
    \longrightarrow
    \mathscr K^{(1)}_{M,v}
\]
of Lemma~\ref{lem:var-first-difference-beck-chevalley} gives a specified nullhomotopy of the corresponding section of \(\mathscr K^{(1)}_{M,v}[n]\). By Theorem~\ref{thm:var-euler-represents-first-variation}, that section is exactly \(\Delta_v\mathbf L_L\). Hence
\[
    \Delta_v\mathbf L_L\simeq0.
\]
Applying Theorem~\ref{thm:var-action-first-variation} carries this nullhomotopy through the coefficient equivalence \(\mathcal I^\partial_{x,D}\) and gives
\[
    \Delta_v\mathbf S_L\simeq0.
\]
Every map used is functorial on the complete variation \(\infty\)-category, so the nullhomotopy retains the higher homotopies specified by the critical lift and by the variation diagram.
\end{proof}

\subsection{Higher cubical variation of the action}
\label{subsec:var-action-higher-variation}

\begin{definition}[Finite cubical variation of a field]
\label{def:var-field-k-variation}
Let \(F=\mathbf{Fld}_S(A)\) and let \(k\ge1\). A \emph{split \(k\)-cubical field variation} consists of an affine centre
\[
    c:T=\operatorname{Spec}B\longrightarrow F,
\]
connective modules \(N_1,\ldots,N_k\), and a lift
\[
    \widetilde c_{[k]}:
    T[N_1,\ldots,N_k]
    \longrightarrow
    F
\]
over \(S\), together with the specified identification on the common centre. Let \(\mathcal V_F^{(k)}\) be the \(\infty\)-category of complete such diagrams and their structured morphisms. For \(I\subseteq[k]\), write
\[
    c_I:T[N_I]\longrightarrow F
\]
for the restriction of \(\widetilde c_{[k]}\) to the \(I\)-face.

For \(E\in\mathcal E_S\), define
\[
    K_E^{(k)}(v)
    :=
    \operatorname{TotFib}_{I\subseteq[k]}
    \Gamma_{T[N_I]}^{\mathrm{st}}(E_{T[N_I]}).
\]
The coefficient cube is coherently split by restriction along the face inclusions and pullback along the retractions \(p_{J,I}\). Let
\[
    \overline D_{v,E}:
    \Gamma_{T[N_1,\ldots,N_k]}^{\mathrm{st}}(E_{T[N_1,\ldots,N_k]})
    \longrightarrow
    K_E^{(k)}(v)
\]
be the alternating projector of Lemma~\ref{lem:var-split-cube-projector}. For a map
\[
    \sigma:F\longrightarrow\Omega_S^\infty E[n],
\]
define
\[
    \boxed{
    \Delta_v^{(k)}\sigma
    :=
    \overline D_{v,E}\bigl(\widetilde c_{[k]}^*\sigma\bigr)
    \in
    \Omega^\infty K_E^{(k)}(v)[n].
    }
\]
Under the total-fibre inclusion, its full-vertex representative is
\[
    \sum_{I\subseteq[k]}
    (-1)^{k-|I|}\sigma_I,
    \qquad
    \sigma_I:=p_{[k],I}^*c_I^*\sigma.
\]
These points form a natural section on \((\mathcal V_F^{(k)})^{\mathrm{op}}\).
\end{definition}

\begin{definition}[Parameterized higher-difference coefficient of a field variation]
\label{def:var-parameterized-k-field-difference}
Let \(M\in\mathcal E_Q\) and let \(v\in\mathcal V_F^{(k)}\) be a finite cubical field variation with affine centre \(c:T\to F\) and structural map \(t:T\to S\). Put
\[
    M_T:=t_Q^*M\in\mathcal E_{Q_T}.
\]
For every \(I\subseteq[k]\), write
\[
    Q_I:=Q_{T[N_I]}=Q\times_ST[N_I]
\]
and let
\[
    r_I:Q_I\longrightarrow Q_T
\]
be induced by the split retraction \(T[N_I]\to T\). Define a \(k\)-cube in \(\mathcal E_{Q_T}\) by
\[
    \mathscr G_{M,v}(I)
    :=
    \Pi_{r_I}r_I^*M_T.
\]
For \(I\subseteq J\), restriction along
\[
    Q_I\longrightarrow Q_J
\]
and composition of dependent products give the structure map
\[
    \mathscr G_{M,v}(J)
    \longrightarrow
    \mathscr G_{M,v}(I).
\]
Define
\[
    \boxed{
    \mathscr K^{(k)}_{M,v}
    :=
    \operatorname{TotFib}_{I\subseteq[k]}
    \mathscr G_{M,v}(I)
    \in
    \mathcal E_{Q_T}.
    }
\]
For \(k=1\), this is canonically \(\mathscr K^{(1)}_{M,v}\) of Definition~\ref{def:var-parameterized-first-field-difference}.
\end{definition}

\begin{lemma}[Higher-difference Beck--Chevalley and local evaluation]
\label{lem:var-k-difference-beck-chevalley}
For every finite \(k\ge1\), there is a canonical equivalence
\[
    \boxed{
    \Gamma_{Q_T}^{\mathrm{st}}
    \mathscr K^{(k)}_{M,v}
    \simeq
    K_{\Pi_qM}^{(k)}(v)
    \simeq
    \operatorname{TotFib}_{I\subseteq[k]}
    \Gamma_{Q_{T[N_I]}}^{\mathrm{st}}(M).
    }
\]
For every affine point \(u:U\to Q_T\), base change of the complete cube gives the induced local structured \(k\)-variation
\[
    v_u\in\mathcal V^{(k)}_{\bar A/Q}
\]
and a canonical equivalence
\[
    \boxed{
    \mathscr K^{(k)}_{M,v}(U,u)
    \simeq
    K_M^{(k)}(v_u).
    }
\]
Writing \(\phi_T:Q_T\to\bar A\) for the centre field composed with \(\bar A_T\to\bar A\), the right-Kan evaluation morphisms assemble to a canonical morphism
\[
    \boxed{
    \operatorname{ev}^{(k)}_v:
    \phi_T^*\mathbb V_M^{(k)}(A)
    \longrightarrow
    \mathscr K^{(k)}_{M,v}.
    }
\]
The construction is natural in the complete cubical variation and in \(M\). For \(k=1\) it recovers Lemma~\ref{lem:var-first-difference-beck-chevalley}.
\end{lemma}

\begin{proof}
Stable sections preserve finite limits. For every \(I\subseteq[k]\), adjunction gives
\[
    \Gamma_{Q_T}^{\mathrm{st}}
    \bigl(\Pi_{r_I}r_I^*M_T\bigr)
    \simeq
    \Gamma_{Q_{T[N_I]}}^{\mathrm{st}}(M).
\]
Taking the total fibre over \(I\subseteq[k]\) gives
\[
    \Gamma_{Q_T}^{\mathrm{st}}
    \mathscr K^{(k)}_{M,v}
    \simeq
    \operatorname{TotFib}_{I\subseteq[k]}
    \Gamma_{Q_{T[N_I]}}^{\mathrm{st}}(M).
\]
Right Beck--Chevalley at every vertex \(T[N_I]\) identifies this vertexwise cube with the cube defining \(K_{\Pi_qM}^{(k)}(v)\), compatibly with all face restrictions. Since total fibre is a finite limit, this gives the first boxed equivalence.

Now let \(u:U\to Q_T\) be affine. For every \(I\subseteq[k]\),
\[
    U[N_{I,u}]
    :=
    U\times_{Q_T}Q_{T[N_I]}
    \simeq
    U\times_TT[N_I]
\]
is affine. Right Beck--Chevalley for the pullback of \(r_I\) along \(u\) gives
\[
    u^*\mathscr G_{M,v}(I)
    \simeq
    \Pi_{r_{I,u}}r_{I,u}^*M_U.
\]
Pullback preserves the finite total fibre, so after taking stable sections over \(U\) one obtains
\[
    \mathscr K^{(k)}_{M,v}(U,u)
    \simeq
    \operatorname{TotFib}_{I\subseteq[k]}
    \Gamma_{U[N_{I,u}]}^{\mathrm{st}}(M)
    =
    K_M^{(k)}(v_u).
\]

Finally,
\[
    \bigl(\phi_T^*\mathbb V_M^{(k)}(A)\bigr)(U,u)
    \simeq
    \mathbb V_M^{(k)}(A)(U,\phi_Tu).
\]
The induced local cube \(v_u\) is an object in the right-Kan indexing category over the centre \(\phi_Tu\). Theorem~\ref{thm:var-k-kan-sheaf} therefore gives a canonical evaluation map
\[
    \mathbb V_M^{(k)}(A)(U,\phi_Tu)
    \longrightarrow
    K_M^{(k)}(v_u)
    \simeq
    \mathscr K^{(k)}_{M,v}(U,u).
\]
Base change of the complete field cube is functorial under affine restriction, and the right-Kan projections form the universal cone. Hence the affine-point maps are natural. The affine-point presentation of \(\mathcal E_{Q_T}\) assembles them uniquely to \(\operatorname{ev}^{(k)}_v\).
\end{proof}

\begin{theorem}[Higher cubical variation of Lagrangian field evaluation]
\label{thm:var-field-k-variation-evaluation}
Let \(L\) be any Lagrangian with coefficient \(M\), and let \(v\in\mathcal V_F^{(k)}\) be a finite cubical field variation. Under the equivalence
\[
    K_{\Pi_qM}^{(k)}(v)
    \simeq
    \Gamma_{Q_T}^{\mathrm{st}}
    \mathscr K^{(k)}_{M,v}
\]
of Lemma~\ref{lem:var-k-difference-beck-chevalley}, the global cubical difference
\[
    \Delta_v^{(k)}\mathbf L_L
\]
is the global section of \(\mathscr K^{(k)}_{M,v}[n]\) represented by
\[
    \boxed{
    \one_{Q_T}
    \xrightarrow{\;\phi_T^*\Delta_L^{(k)}\;}
    \phi_T^*\mathbb V_M^{(k)}(A)[n]
    \xrightarrow{\;\operatorname{ev}^{(k)}_v[n]\;}
    \mathscr K^{(k)}_{M,v}[n].
    }
\]
For every affine point \(u:U\to Q_T\), base change of the complete cube produces the induced local structured \(k\)-variation
\[
    v_u\in\mathcal V^{(k)}_{\bar A/Q},
\]
and the affine-point evaluation of the global identity is canonically
\[
    \boxed{
    \left.\Delta_v^{(k)}\mathbf L_L\right|_{v_u}
    \simeq
    \left\langle\Delta_L^{(k)},v_u\right\rangle.
    }
\]
For \(k=1\), Proposition~\ref{prop:var-k1-euler} identifies the right side with \(\langle\mathsf{EL}_L,v_u\rangle\), recovering Theorem~\ref{thm:var-euler-represents-first-variation}. For \(k=2\), it is evaluation of the intrinsic symmetric mixed second variation.
\end{theorem}

\begin{proof}
At every vertex \(I\subseteq[k]\), the universal field evaluation pulls the local Lagrangian section \(\ell_L\) back along the corresponding face field. Right Beck--Chevalley identifies the resulting split coefficient cube with the cube of coefficient sections on \(Q_{T[N_I]}\). Its alternating projector therefore represents \(\Delta_v^{(k)}\mathbf L_L\) as a global section of \(\mathscr K^{(k)}_{M,v}[n]\).

For an affine point \(u:U\to Q_T\), pullback of the complete cube gives the local structured variation \(v_u\). Naturality of the alternating projector of Lemma~\ref{lem:var-split-cube-projector} under this pullback identifies the affine value of \(\Delta_v^{(k)}\mathbf L_L\) with
\[
    \Delta_{v_u}^{(k)}L
    =
    \left\langle\Delta_L^{(k)},v_u\right\rangle
\]
by Construction~\ref{con:var-k-variation-lagrangian} and the right-Kan evaluation map.

The composite
\[
    \one_{Q_T}
    \xrightarrow{\phi_T^*\Delta_L^{(k)}}
    \phi_T^*\mathbb V_M^{(k)}(A)[n]
    \xrightarrow{\operatorname{ev}^{(k)}_v[n]}
    \mathscr K^{(k)}_{M,v}[n]
\]
has the same affine value by the defining affine-point formula for \(\operatorname{ev}^{(k)}_v\). Under the Chapter 10 affine-point presentation, two morphisms of parameterized spectra agree if their natural affine-point components agree. Hence this composite is exactly the global cubical-difference section. The displayed local identity is its affine-point evaluation.
\end{proof}

\begin{theorem}[Integration commutes with every finite cubical variation]
\label{thm:var-action-k-variation}
Let \(L\) be a differential Lagrangian. For every finite \(k\ge1\), the coefficient equivalence
\[
    \mathcal I^\partial_{x,D}:
    \Pi_qM^\partial_{x,D}
    \xrightarrow{\sim}
    \mathbb A^\partial_{x,D}
\]
induces an equivalence of cubical-difference functors on \((\mathcal V_F^{(k)})^{\mathrm{op}}\), and under that equivalence the natural cubical-difference section of \(\mathbf L_L\) is carried to the natural cubical-difference section of \(\mathbf S_L\). Equivalently, for every \(v\in\mathcal V_F^{(k)}\),
\[
    \boxed{
    \Delta_v^{(k)}\mathbf S_L
    \simeq
    K^{(k)}(\mathcal I^\partial_{x,D})(v)
    \bigl(\Delta_v^{(k)}\mathbf L_L\bigr).
    }
\]
The map \(K^{(k)}(\mathcal I^\partial_{x,D})(v)\) is an equivalence. On every induced affine local \(k\)-variation, the coefficient-valued term is the evaluation of \(\Delta_L^{(k)}\). Thus every finite cubical variation of the action is literally differential exceptional integration on the actual represented cubical square-zero parameter.
\end{theorem}

\begin{proof}
A coefficient equivalence induces an equivalence of the entire coherently split vertexwise cube and hence of its total fibre, naturally in structured cubical variations. Since
\[
    \mathbf S_L
    =
    \Omega_S^\infty(\mathcal I^\partial_{x,D}[n])\circ\mathbf L_L,
\]
naturality of the alternating projector of Lemma~\ref{lem:var-split-cube-projector} gives the displayed natural equivalence of cubical-difference sections. The affine-local identification is Theorem~\ref{thm:var-field-k-variation-evaluation}. Represented-parameter recovery identifies every vertexwise coefficient equivalence with the differential exceptional adjunction on that represented square-zero test.
\end{proof}

\begin{corollary}[Integrated symmetric mixed second variation]
\label{prop:var-action-mixed-second-variation}
For \(k=2\), Theorem~\ref{thm:var-action-k-variation} gives
\[
    \boxed{
    \Delta_{12}\mathbf S_L
    \simeq
    K^{(2)}(\mathcal I^\partial_{x,D})(v_{12})
    \bigl(\Delta_{12}\mathbf L_L\bigr),
    }
\]
and the affine-local coefficients are evaluations of
\[
    \Delta_L^{(2),\mathrm{mix}}:
    \one_{\bar A}\to\Vsecond_{M^\partial_{x,D}}(A)[n].
\]
The transposition symmetry is preserved by exceptional integration.
\end{corollary}

\begin{remark}[No symmetry or Noether input]
\label{rem:var-action-no-noether}
The action construction uses only the big differential mapping-spectrum coefficients, the dualizing coefficient obtained from them, the universal field-moduli evaluation, and the already constructed differential exceptional adjunction.  The relative Lepage and universal presymplectic-transgression objects of Section~\ref{sec:var-lepage-presymplectic} may be specialized to the differential dualizing-density coefficient and hence coexist with the action, but they require no Lie symmetry object.  No infinitesimal symmetry action, Noether current, moment map, BRST differential, or gauge-reduction datum is constructed here.  Those structures begin only after a separate brave-new Lie and symmetry theory has been developed.
\end{remark}

\section{The brave-new variational package}
\label{sec:var-final-package}

\begin{theorem}[Higher-categorical brave-new variational geometry and action package]
\label{thm:var-final-package}
Let \(x:X\to S\) be a morphism in the native big spectral Nisnevich \(\infty\)-topos, let \((A,\rho_A)\) be a Cartan object for the Chapter 11 jet comonad, let
\[
    a:\bar A\to Q:=Q_{X/S}
\]
be its constructed descent object, and let \(M\in\mathcal E_Q\). Then the constructions of this chapter canonically provide the following higher-categorical variational package.

\begin{enumerate}
\item The Cartan descent equivalence
\[
    A\simeq X\times_Q\bar A
\]
and the complete Cartan--vertical bisimplicial relation
\[
    \Rvar^{p,q}(A)
    =
    X^{\times_Q(q+1)}\times_QR_{\bar A/Q,p}.
\]

\item The bicosimplicial stable coefficient object
\[
    \Cvar^{p,q}(A;M)
    =
    \Pi_{r_{p,q}}r_{p,q}^*M,
\]
horizontal effective descent
\[
    \Tot_H\Cvar^{\bullet,\bullet}(A;M)
    \simeq
    C_{\dR}^{\bullet}(\bar A/Q;M),
\]
and the total equivalence
\[
    \DRvar(A;M)
    \simeq
    \mathrm{DR}_{\mathrm{eff}}(\bar A/Q;M).
\]

\item The primary complete bifiltered variational object
\[
    \boxed{
    \mathfrak{Var}_{X/S}(A;M)
    =
    F^{\bullet,\bullet}\DRvar(A;M)
    \in
    \BiFil(\mathcal E_Q),
    }
\]
with its canonically equivalent reverse-order construction and iterated associated graded
\[
    \operatorname{gr}_V^p\operatorname{gr}_H^q\mathfrak{Var}_{X/S}(A;M)
    \simeq
    \operatorname{gr}_H^q\operatorname{gr}_V^p\mathfrak{Var}_{X/S}(A;M)
    \simeq
    \Omegavar^{p,q}(A;M)[-p-q].
\]

\item The \(\infty\)-functor
\[
    \mathfrak{Var}_{X/S}:
    ((\XSp)_{/Q})^{\mathrm{op}}\times\mathcal E_Q
    \longrightarrow
    \BiFil(\mathcal E_Q),
\]
its coherent double-cochain presentation, the unsigned coherent boundaries \(\partial_V,\partial_H\), and the signed bicomplex shadow
\[
    d_V=\partial_V,
    \qquad
    d_H|_{(p,q)}=(-1)^p\partial_H
\]
with the specified coherent relations
\[
    d_V^2\simeq0,
    \qquad
    d_H^2\simeq0,
    \qquad
    d_Hd_V+d_Vd_H\simeq0
\]
and all higher compatible cells.

\item The Cartan--vertical filtration
\[
    F_{\mathcal C}^p\DRvar:=F^{p,0}\DRvar,
\]
its associated graded
\[
    \operatorname{gr}_{\mathcal C}^p\DRvar
    \simeq
    \Sec^p(A;M)[-p],
\]
and the explicit statement that no identification with powers of a contact ideal is made.

\item The secondary form spectra
\[
    \Sec^p(A;M)
    =
    \Tot_HN_V^p\Cvar
    \simeq
    \Omegainf^p(\bar A/Q;M),
\]
with horizontal degree interpreted as a complete Cartan descent resolution, and the coherent full variational operator
\[
    \delta^\infty_{\mathrm{var}}:
    \Sec^p(A;M)\to\Sec^{p+1}(A;M)
\]
identified with the Chapter 10 coherent infinitesimal differential of \(\bar A/Q\).

\item Coherently closed secondary forms, exact classifiers, and the universal exact-form moduli map; primitive spaces and parameterized primitive objects are its homotopy fibres rather than separately postulated structures.

\item The local full one-variation coefficient
\[
    \Vfull_M(A)
    =
    \fib(\Pi_ce^*a^*M\to a^*M),
    \qquad
    \Pi_a\Vfull_M(A)\simeq\Sec^1(A;M),
\]
and the right-Kan one-stage coefficient \(\Vfirst_M(A)\), constructed from the complete \(\infty\)-category of split structured square-zero variation diagrams. The raw right Kan extension is already a spectral Nisnevich sheaf.

\item The canonical full-to-first restriction
\[
    \lambda^1_{\mathrm{loc}}:\Vfull_M(A)\to\Vfirst_M(A)
\]
and its pushed-forward map
\[
    \lambda^1_{\mathrm{var}}:
    \Sec^1(A;M)\to\Sec^{1,(1)}(A;M).
\]
If \(a\) is representable, universal affine split variations constructed from the cotangent complex are initial in the correct comma categories
\[
    (c^{\mathrm{op}}\downarrow\pi_1^{\mathrm{op}}),
\]
giving the recognition
\[
    \Vfirst_M(A)
    \simeq
    L^{\mathrm{st}}_{\bar A/Q}(M).
\]
For scalar coefficients this recovers the Chapter 10 first-order cotangent linearization after represented affine base change.

\item The universal variational moduli diagram
\[
    \boxed{
    \mathbf{Lag}^n_{A/Q}(M)
    \longrightarrow
    \mathbf{ClVar}^n_{A/Q}(M)
    \longrightarrow
    \mathbf{Eul}^n_{A/Q}(M),
    }
\]
whose composite is the universal Euler map. Parameterized Helmholtz lifts, first-order Lagrangian realizations, and primitive objects are its derived fibres, and the previously defined global spaces are their global-point shadows.

\item The full and first local Euler sections
\[
    \mathsf E_L^\infty:\one_{\bar A}\to\Vfull_M(A)[n],
    \qquad
    \mathsf{EL}_L:\one_{\bar A}\to\Vfirst_M(A)[n],
\]
their homotopy zero fibres \(\Crit^\infty(L)\) and \(\Crit(L)\), and their intrinsic formal integrability as objects over \(Q\).

\item The universal derived critical families
\[
    \mathbf{Crit}^{\infty}_{A/Q}(M;n),
    \qquad
    \mathbf{Crit}^{(1)}_{A/Q}(M;n)
\]
over
\[
    \bar A\times_Q\mathbf{Lag}^n_{A/Q}(M),
\]
with each individual critical locus recovered by base change along the corresponding Lagrangian section.

\item For a cofree field object \(K(Y)\), the ambient Euler jet morphism obtained by pulling back the descended critical object, automatically as a Cartan morphism; no separate integrability or prolongation datum is required.

\item The complete horizontal realization tower and its universal morphism
\[
    \Omega_Q^\infty F_H^q\Sec^p[n]
    \longrightarrow
    \Omega_Q^\infty\Sec^p[n],
\]
with individual horizontal realization spaces as fibres, and the complete transgression coefficient
\[
    \TransH^{p,q}
    \simeq
    \Omega\Sec_H^{p,<q}.
\]

\item For every Lagrangian and every \(q\ge1\), the derived relative Lepage-realization object
\[
    \boxed{
    \mathbf{Lep}^{\,q}(L)
    \longrightarrow Q,
    }
\]
its tautological complete boundary realization
\[
    \Theta_L^{(q)}:
    \one_{\mathbf{Lep}^{\,q}(L)}
    \longrightarrow
    (\pi_L^q)^*F_H^q\Sec^1(A;M)[n],
\]
and the universal presymplectic transgression
\[
    \boxed{
    \omega_L^{(q)}:
    \one_{\mathbf{Lep}^{\,q}(L)}
    \longrightarrow
    (\pi_L^q)^*\TransH^{2,q}(A;M)[n].
    }
\]
The defining pullback gives the relative integration-by-parts factorization of Theorem~\ref{thm:var-relative-integration-by-parts}; the resulting current is coherently vertically closed, satisfies the complete presymplectic recurrence of Theorem~\ref{thm:var-presymplectic-recurrence}, and, within each fibre over a fixed coherent Euler lift, changes by \(d_V^{\mathrm{tr}}\) of the complete transgression when the horizontal realization changes.  No preferred Lepage or Poincar\'e--Cartan representative and no chosen integration-by-parts operator is inserted.

\item For every finite \(k\ge1\), the complete \(\infty\)-category \(\mathcal V^{(k)}_{\bar A/Q}\) of structured split \(k\)-cubical variations, its centre-pullback Nisnevich topology, the total-fibre functor
\[
    K_M^{(k)}(v)
    =
    \operatorname{TotFib}_{I\subseteq[k]}
    \Gamma_{T[N_I]}^{\mathrm{st}}(M),
\]
and the right-Kan coefficient
\[
    \boxed{
    \mathbb V_M^{(k)}(A)
    =
    \operatorname{Ran}_{\pi_k^{\mathrm{op}}}K_M^{(k)},
    }
\]
where the displayed raw right Kan extension is already a spectral Nisnevich sheaf. These coefficients satisfy arbitrary base change and carry coherent \(\Sigma_k\)-actions.

\item For every Lagrangian, the alternating higher variation
\[
    \boxed{
    \Delta_L^{(k)}:
    \one_{\bar A}\to\mathbb V_M^{(k)}(A)[n],
    }
\]
with \(\Delta_L^{(1)}=\mathsf{EL}_L\) and \(\Delta_L^{(2)}=\Delta_L^{(2),\mathrm{mix}}\). The second variation is intrinsically cubical: the explicit affine example proves that it can detect the mixed \(N_1\otimes N_2\) lift even with fixed side variations.

\item The side-restriction functor
\[
    \partial_2:\mathcal V^{(2)}_{\bar A/Q}
    \to
    \mathcal V^{(1)}_{\bar A/Q}
    \times_{\mathcal C^{\mathrm{aff}}_{/\bar A}}
    \mathcal V^{(1)}_{\bar A/Q}
\]
and the unrestricted derived Hessian-descent space \(\HessDesc(L)\). Thus off-shell side descent remains additional data and can fail. Pullback to \(\Crit(L)\) gives the on-shell mixed second variation. In the represented smooth scalar sector, the chapter identifies the affine-point linear category \(\mathcal Q_Z\) with the Chapter 1 category \(\operatorname{QCoh}(Z)\), proves finite affine spectral Nisnevich module descent and that its underlying affine-point spectrum presheaf is already Nisnevich-local, and constructs
\[
    U_Z:\mathcal Q_Z\to\mathcal E_Z,
    \qquad
    U_Z(\mathcal O_Z^{\mathrm{lin}})\simeq\Oadd_Z.
\]
The represented cotangent coefficient refines to a finite-locally-free object \(L^{\mathrm{lin}}_{\bar A/Q}\in\mathcal Q_{\bar A}\) with underlying coefficient \(L^{\mathrm{st}}_{\bar A/Q}\), and its dual is \(T^{\mathrm{lin}}_{\bar A/Q}\).  The chapter further constructs the mixed boundary as a structured square-zero extension by \(N_1\otimes_BN_2\), proves smooth fillability and the mixed-fill torsor, proves that translation in that torsor changes the mixed variation by Euler contraction, and uses criticality to descend canonically to side data. The scalar total-fibre coefficient is
\[
    K_{\Oadd}^{(k)}(v)
    \simeq
    N_1\otimes_B\cdots\otimes_BN_k.
\]
The universal descended value gives the symmetric bilinear critical Hessian
\[
    \Hess_L^{\mathrm{crit}}:
    j_L^*T^{\mathrm{lin}}_{\bar A/Q}
    \otimes
    j_L^*T^{\mathrm{lin}}_{\bar A/Q}
    \longrightarrow
    \mathcal O_{C_L}^{\mathrm{lin}}[n],
\]
its linear adjoint \(\Jac_L^{\mathrm{lin}}\), the underlying parameterized Jacobi operator
\[
    \Jac_L:
    j_L^*T^{\mathrm{st}}_{\bar A/Q}
    \longrightarrow
    j_L^*L^{\mathrm{st}}_{\bar A/Q}[n],
\]
and the derived Jacobi-variation object \(\mathbf{Jac}(L)\to\Crit(L)\).  The Hessian is the canonical represented critical side descent of the intrinsic on-shell mixed second variation.  No fillability, multilinearity, symmetry, Hessian, or Jacobi axiom is imposed.

\item Arbitrary base change, formal-\'etale locality, and spectral Nisnevich descent for all preceding internal coefficient and moduli objects, together with the corresponding restriction/descent statements for their global point spaces. The derived Hessian-descent space carries the functorial restriction maps induced by pullback of the side-data diagram, but no arbitrary base-change equivalence or independent Čech-reconstruction theorem for that global space is asserted. The chapter also gives the represented smooth local formal-translation and tubular models and the explicitly limited reduced-classical and classical secondary-calculus comparisons.
\end{enumerate}

Assume in addition that \(S\) is a qcqs spectral scheme in the standing differential-motivic base range of Chapters 9--10, that \(x:X\to S\) is represented, smooth, separated, and of finite presentation, and let \(D\in\mathcal R_S^0\). Then the chapter further provides:

\begin{enumerate}
\setcounter{enumi}{19}
\item For every represented qcqs differential base \(Z\), the big differential mapping-spectrum coefficient
\[
    \underline{\operatorname{map}}_Z^\partial(E,F)\in\mathcal E_Z,
\]
its finite spectral Nisnevich descent, arbitrary represented base change, and global-section equivalence
\[
    \Gamma_Z^{\mathrm{st}}
    \underline{\operatorname{map}}_Z^\partial(E,F)
    \simeq
    \operatorname{map}_{\mathcal R_Z^0}(E,F).
\]

\item The big dualizing-density coefficient
\[
    \mathscr D_{x,D}^{\mathrm{big}}
    :=
    \underline{\operatorname{map}}_X^\partial
    (\one_X^\partial,x_e^!D),
    \qquad
    M^\partial_{x,D}
    :=
    \Pi_\epsilon\mathscr D_{x,D}^{\mathrm{big}},
\]
with density realization
\[
    \Gamma_Q^{\mathrm{st}}M^\partial_{x,D}
    \simeq
    \operatorname{map}_{\mathcal R_X^0}
    (\one_X^\partial,x_e^!D)
\]
and arbitrary represented base change.

\item The big action-value coefficient
\[
    \mathbb A^\partial_{x,D}
    :=
    \underline{\operatorname{map}}_S^\partial
    (x_{e!}\one_X^\partial,D)
\]
and the big exceptional-integration equivalence
\[
    \boxed{
    \mathcal I^\partial_{x,D}:
    \Pi_qM^\partial_{x,D}
    \xrightarrow{\sim}
    \mathbb A^\partial_{x,D},
    }
\]
whose affine values are the Chapter 10 differential Umkehr equivalences.

\item The higher field-moduli object
\[
    \mathbf{Fld}_S(A)=\Pi_q\bar A
\]
and, for every differential Lagrangian, the universal coefficient-valued field evaluation and the orientation-free action morphism
\[
    \mathbf L_L:
    \mathbf{Fld}_S(A)
    \to
    \Omega_S^\infty(\Pi_qM^\partial_{x,D}[n]),
\]
\[
    \boxed{
    \mathbf S_L:
    \mathbf{Fld}_S(A)
    \to
    \Omega_S^\infty(\mathbb A^\partial_{x,D}[n]).
    }
\]
The ordinary action on global fields is the global-points shadow and has value
\[
    x_{e!}\one_X^\partial
    \xrightarrow{x_{e!}\mathcal L_L(\phi)}
    x_{e!}x_e^!D[n]
    \xrightarrow{\operatorname{tr}^\partial}
    D[n].
\]

\item The \(\infty\)-category \(\mathcal V_F^{(1)}\) of first field variations and, more generally, \(\mathcal V_F^{(k)}\) for every finite \(k\). For each field \(k\)-variation \(v\), the parameterized total-fibre coefficient
\[
    \mathscr K^{(k)}_{M,v}
    \in
    \mathcal E_{Q_T}
\]
has global sections
\[
    \Gamma_{Q_T}^{\mathrm{st}}\mathscr K^{(k)}_{M,v}
    \simeq
    K_{\Pi_qM}^{(k)}(v)
\]
and receives the canonical right-Kan evaluation morphism
\[
    \operatorname{ev}^{(k)}_v:
    \phi_T^*\mathbb V_M^{(k)}(A)
    \longrightarrow
    \mathscr K^{(k)}_{M,v}.
\]
The first-difference and higher cubical-difference sections are natural sections of these corresponding total-fibre coefficients, rather than merely pointwise formulas.

\item The global Euler representation theorem
\[
    \Delta_v\mathbf L_L
    \simeq
    \operatorname{ev}^{(1)}_v
    \bigl(\phi_T^*\mathsf{EL}_L\bigr)
\]
under the canonical identification with global sections of \(\mathscr K^{(1)}_{M,v}\), and its higher cubical extension
\[
    \Delta_v^{(k)}\mathbf L_L
    \simeq
    \operatorname{ev}^{(k)}_v
    \bigl(\phi_T^*\Delta_L^{(k)}\bigr).
\]
Their affine-point evaluations are respectively
\[
    \left.\Delta_v\mathbf L_L\right|_{v_u}
    \simeq
    \langle\mathsf{EL}_L,v_u\rangle
\]
and
\[
    \left.\Delta_v^{(k)}\mathbf L_L\right|_{v_u}
    \simeq
    \langle\Delta_L^{(k)},v_u\rangle.
\]

\item For every finite \(k\ge1\), the natural equivalence of cubical-difference functors
\[
    \boxed{
    \Delta^{(k)}\mathbf S_L
    \simeq
    K^{(k)}(\mathcal I^\partial_{x,D})
    \bigl(\Delta^{(k)}\mathbf L_L\bigr),
    }
\]
whose evaluation on each structured field \(k\)-variation is differential exceptional integration on the actual represented cubical square-zero parameter. The first-variation and symmetric mixed second-variation formulas are the cases \(k=1\) and \(k=2\).

\item The critical-field moduli object
\[
    \mathbf{CritFld}_S(L)
    :=
    \Pi_q\Crit(L)
    \longrightarrow
    \mathbf{Fld}_S(A)
\]
and the stationarity theorem: a chosen critical lift canonically nullhomotopes the first variation of the big action. No converse is asserted because no generation theorem identifies all affine-local right-Kan variations with global field variations.
\end{enumerate}

Every construction above is made from machinery already available in the manuscript: effective relation groupoids, finite limits in the native \(\infty\)-topos, stable limits in parameterized spectrum categories, dependent adjunction and Beck--Chevalley, complete totalization and coherent normalization, structured split square-zero extensions and their classifying derivations, finite tensor products of those extensions, cotangent obstruction theory, pointwise right Kan extension, finite Nisnevich descent, relative zeroth-space moduli, homotopy fibres, the affine-point linear coefficient categories identified with Chapter 1 quasi-coherent modules together with finite affine spectral Nisnevich module descent, Thom-stable differential realization, and differential exceptional adjunction. No finite-order jet tower, tangent-distribution axiom, exterior-cotangent identification of the full normalized cochains, multiplicative contact-ideal axiom, chosen integration-by-parts operator, orientation, admissibility condition, fillability axiom, multilinearity axiom, Hessian axiom, Jacobi axiom, action-functional axiom, or Noether datum is introduced.
\end{theorem}

\begin{proof}
Cartan descent, vertical infinitesimal geometry, the double relation, and horizontal effective descent are Sections~\ref{sec:var-cartan-descent}--\ref{sec:var-double-cochains}. The primary bifiltered object, its associated graded, the coherent double-cochain object, its signed shadow, and the \(\infty\)-functorial packaging are Section~\ref{sec:var-normalization} and Subsection~\ref{subsec:var-infinity-functor}. Secondary forms, closed forms, primitive moduli, horizontal realization, relative Lepage realizations, and universal presymplectic transgression are Sections~\ref{sec:var-secondary-forms}, \ref{sec:var-closed-primitive}, \ref{sec:var-moduli}, \ref{sec:var-horizontal-transgression}, and \ref{sec:var-lepage-presymplectic}.

The full and first variation coefficients, their right-Kan descent, represented cotangent recognition, Euler sections, individual and universal critical loci, and Helmholtz/inverse-problem fibres are Sections~\ref{sec:var-square-zero-sector}, \ref{sec:var-first-represented}, \ref{sec:var-euler-lagrange}, \ref{sec:var-helmholtz}, and \ref{sec:var-inverse-problem}; the universal variational moduli diagram is Subsection~\ref{sec:var-universal-variational-moduli}. The higher cubical variation system, its descent, base change, symmetric-group action, mixed second variation, unrestricted Hessian-descent space, on-shell second variation, and the represented critical Hessian and Jacobi constructions are Section~\ref{sec:var-higher-variations} and Subsection~\ref{subsec:var-represented-jacobi}. Base change, formal-\'etale locality, Nisnevich descent, local models, and the classical comparisons are the subsequent comparison sections.

The differential mapping coefficient, density coefficient, action-value coefficient, field moduli, action morphism, parameterized field-difference coefficients, natural first-difference statement, Euler representation theorem, higher cubical field variations, and the natural equivalence expressing commutation of exceptional integration with every finite cubical variation are Section~\ref{sec:var-action-functional}. Each statement is obtained from the previously constructed adjunctions, Beck--Chevalley equivalences, finite limits, and descent results; no additional axiom is used.
\end{proof}

The chapter has therefore constructed a higher-categorical brave-new variational theory in which the complete bifiltered object and its coherent double-cochain presentation are primary, classical bicomplex formulas are signed shadows, inverse problems and realizations are derived fibres of universal moduli maps, Euler equations form a universal derived family over Lagrangian moduli, and higher variations are intrinsically cubical before any possible multilinear descent.  The horizontal realization tower now carries derived relative Lepage moduli and a universal coherently closed presymplectic transgression without selecting a Poincar\'e--Cartan representative.  In the represented smooth scalar sector the critical Euler equation has a constructed Jacobi linearization and symmetric bilinear critical Hessian whose side evaluation is the intrinsic on-shell mixed second variation.  In the represented smooth differential sector, the action remains a morphism of higher field moduli and differential exceptional integration commutes naturally with every finite cubical variation.  A later brave-new Lie and symmetry theory may therefore add infinitesimal symmetries, Noether maps, Hamiltonian reduction, BRST data, and gauge-theoretic BV--BFV structures without reopening the variational, presymplectic-transgression, or represented Jacobi foundations constructed here.
